% concatenated from iclr2026_conference.tex, math_commands.tex, sections/abstract.tex, sections/introduction.tex, sections/related_work.tex, sections/background.tex, sections/draft_charlie_method.tex, figures/chatbot_arena_figure.tex, figures/processbench_figure.tex, figures/factuality_figure.tex, sections/appendix/generalization.tex, sections/appendix/algorithm.tex, sections/appendix/additional_experiments.tex, figures/chatbotarena_figure_10.tex, figures/processbench_figure_10.tex, figures/factuality_figure_10.tex, figures/chatbot_arena_figure1000.tex, figures/processbench_figure_1000.tex, figures/factuality_figure_1000.tex, figures/chatbotarena_shrinkage.tex, figures/processbench_shrinkage.tex, figures/factuality_shrinkages.tex, sections/appendix/chatbotarena_restricted.tex, sections/appendix/processbench_restricted.tex, sections/appendix/factuality_restricted.tex, sections/appendix/additional_theory.tex, sections/appendix/proofs.tex, sections/appendix/computational.tex, sections/appendix/dual.tex, sections/appendix/experimental_details.tex

\documentclass{article} % For LaTeX2e
\usepackage{iclr2026_conference,times}

% Optional math commands from https://github.com/goodfeli/dlbook_notation.
\newcommand{\p}{\mathbb{P}}
\newcommand{\q}{\mathbb{Q}}
\newcommand{\R}{\mathbb{R}}
\newcommand{\C}{\mathcal{C}}
\newcommand{\Ex}{\mathbb{E}}
\usepackage{outlines}

\usepackage{amsmath}
\usepackage{amsfonts}
\usepackage{amssymb}
\usepackage{amsthm}
\usepackage{bbm}
\usepackage{natbib}
\usepackage{enumitem}
\usepackage[utf8]{inputenc} % allow utf-8 input
\usepackage[T1]{fontenc}    % use 8-bit T1 fonts
\usepackage{hyperref}       % hyperlinks
\usepackage{cleveref}
\usepackage{url}            % simple URL typesetting
\usepackage{booktabs}       % professional-quality tables
\usepackage{kantlipsum, lipsum}
\usepackage{xcolor}
\definecolor{darkblue}{rgb}{0, 0, 0.5}
\hypersetup{colorlinks=true, citecolor=darkblue, linkcolor=darkblue, urlcolor=darkblue}
\usepackage{etoc}
\usepackage{dsfont}
\usepackage{epigraph} 
\usepackage{etoc}
\usepackage{tikz}
\usepackage{pgfplots}
\usepackage{subcaption}
\renewcommand{\cite}{\citep}
\renewcommand{\sectionautorefname}{\S}
\renewcommand{\subsectionautorefname}{\S}
%\renewcommand{\autoref}{\Cref}
\usepackage{aliascnt}


\newcommand{\charlie}[1]{ {\color{red}[CCB: #1]}}
\newcommand{\adam}[1]{ {\color{magenta}[AF: #1]}}
\newcommand{\jb}[1]{ {\color{blue}[JB: #1]}}
\newcommand{\jacob}[1]{ {\color{blue}[JE: #1]}}
\newcommand{\alekh}[1]{ {\color{brown}[AA: #1]}}
\newcommand{\amirg}[1]{ {\color{green}[AG: #1]}}
\usepackage[most,skins,theorems]{tcolorbox}

\DeclareMathOperator{\Var}{Var}
\DeclareMathOperator{\argmax}{argmax}
\DeclareMathOperator{\Cov}{Cov}
\DeclareMathOperator{\multiPPI}{MultiPPI}
\DeclareMathOperator{\Corr}{Corr}
\DeclareMathOperator{\supp}{supp}
\DeclareMathOperator{\range}{range}
\DeclareMathOperator{\tr}{tr}

\newtheorem{theorem}{Theorem}
%\newtheorem{theorem}[theorem]{Theorem}
\newaliascnt{lemma}{theorem}
\newaliascnt{proposition}{theorem}
\newaliascnt{corollary}{theorem}
\newtheorem{lemma}[lemma]{Lemma}
\aliascntresetthe{lemma}

\newtheorem{corollary}[corollary]{Corollary}
\newtheorem{proposition}[proposition]{Proposition}
%\theoremstyle{definition}
\newtheorem{remark}[theorem]{Remark}
\newtheorem{definition}[theorem]{Definition}
\newtheorem{assumption}[theorem]{Assumption}
\newtheorem{example}[theorem]{Example}
\newtheorem{fact}[theorem]{Fact}
\newtheorem{question}{Question}
\newcommand{\I}{\mathbbm{1}}
\newcommand{\wh}[1]{\widehat{#1}}
\newcommand{\ul}[1]{\underline{#1}}
\DeclareMathOperator{\argmin}{argmin}
\DeclareMathOperator{\classical}{classical}

\aliascntresetthe{proposition}
\aliascntresetthe{corollary}

\crefname{lemma}{lemma}{lemmas}
\Crefname{lemma}{Lemma}{Lemmas}

\tcbset{
  aibox/.style={
    width=\linewidth,
    top=6pt,
    bottom=0pt,
    colback=blue!6!white,
    colframe=black,
    colbacktitle=black,
    enhanced,
    center,
    attach boxed title to top left={yshift=-0.1in,xshift=0.15in},
    boxed title style={boxrule=0pt,colframe=white,},
  }
}
\newtcolorbox{AIbox}[2][]{aibox,title=#2,#1}

\tcbuselibrary{listings, breakable}
\tcbset{
    graybox/.style={
        colback=gray!20,  % Background color (light gray)
        colframe=black,   % Border color (black)
        boxrule=1pt,      % Border thickness
        arc=4mm,          % Box corner rounding
        width=\textwidth, % Box width to match the text width
        before=\vspace{10pt}, % Space before the box
        after=\vspace{10pt},  % Space after the box
        center             % Center the box
    }
}

% \usepackage{mathtools}
% \mathtoolsset{showonlyrefs}
% \usepackage{afterpage}
\iclrfinalcopy

\usepackage{hyperref}
\usepackage{url}


% Paper Title
\title{Multiple-Prediction-Powered Inference}


\author{Charlie Cowen-Breen$^{1\dagger}$ \hfill \quad Alekh Agarwal$^2$ \quad \hfill Stephen Bates$^1$ \AND  William W. Cohen$^3$ \qquad Jacob Eisenstein$^3$ \qquad Amir Globerson$^2$ \qquad  Adam Fisch$^3$ \\\\ $~^{1}$Massachusetts Institute of Technology, $~^2$Google Research, $~^3$Google DeepMind\\ $~^\dagger$Work done during internship at Google DeepMind.\\
\href{mailto:ccbreen@mit.edu}{\texttt{ccbreen@mit.edu}}, \href{mailto:fisch@google.com}{\texttt{fisch@google.com}}. 
}

% \author{%
%   \begin{tabular}[t]{ccc}
%     Charlie Cowen-Breen$~^{1\dagger}$ & Alekh Agarwal & Stephen Bates \\
%     William W. Cohen & Jacob Eisenstein & Amir Globerson \\
%     \multicolumn{3}{c}{Adam Fisch}
%   \end{tabular}\\[1ex] % Adds a little vertical space between authors and affiliations
%   $~^{1}$Massachusetts Institute of Technology \\
%   $^\dagger$Work done during internship at Google DeepMind
% }
% \author{%\\
%   \begin{tabular*}{\linewidth}{@{\extracolsep{\fill}}ccc}
%     Charlie Cowen-Breen$~^{1\dagger}$ & Alekh Agarwal & Stephen Bates \\\\\\
%     William W. Cohen & Jacob Eisenstein & Amir Globerson \\\\\\
%     \multicolumn{3}{c}{Adam Fisch}
%   \end{tabular*}\\[1ex] 
%   $~^{1}$Massachusetts Institute of Technology \\
%   $^\dagger$Work done during internship at Google DeepMind
% }

% The \author macro works with any number of authors. There are two commands
% used to separate the names and addresses of multiple authors: \And and \AND.
%
% Using \And between authors leaves it to \LaTeX{} to determine where to break
% the lines. Using \AND forces a linebreak at that point. So, if \LaTeX{}
% puts 3 of 4 authors names on the first line, and the last on the second
% line, try using \AND instead of \And before the third author name.

\newcommand{\fix}{\marginpar{FIX}}
\newcommand{\new}{\marginpar{NEW}}



\begin{document}

\maketitle


\begin{abstract}

% Statistical estimation
% %, such as the mean accuracy of generative AI model responses to random queries,
% often introduces tradeoffs between expensive, high-quality measurements and a variety of cheaper, lower-quality proxies. In this work, we introduce Multiple-Prediction-Powered Inference (MultiPPI), a general framework for constructing statistically efficient estimates of a ground-truth parameter by leveraging cheaper proxies.
% %While optimizing over combinations of subsets of available proxies may seem daunting \jacob{it's not yet clear why one would want to optimize over subsets of proxies. maybe we could borrow a sentence from the intro to explain this, e.g. via the link to PPI and the extension to more complex estimators}, we show that the problem of finding the optimal sampling strategy---which variables to query, how many times, and how to weight them---can be formulated and solved efficiently as a second-order cone program (SOCP) using off-the-shelf tools.
% We provide theoretical guarantees about the minimax optimality, finite-sample performance, and asymptotic normality of the MultiPPI estimator. Through experiments across three diverse large language model (LLM) evaluation scenarios, we demonstrate that MultiPPI consistently achieves lower estimation error for a given budget than existing baselines. In particular, we show  that it accomplishes this with an allocation strategy that strategically varies with the available budget: it relies on cheap proxies in low-budget regimes, and incorporates more expensive, high-quality predictors as the budgetincreases.\looseness=-1


%\begin{center}\textbf{Version 1:}\end{center}
Statistical estimation often involves tradeoffs between expensive, high-quality measurements and a variety of lower-quality proxies. We introduce Multiple-Prediction-Powered Inference (MultiPPI): a general framework for constructing statistically efficient estimates by optimally allocating resources across these diverse data sources. This work provides theoretical guarantees about the minimax optimality, finite-sample performance, and asymptotic normality of the MultiPPI estimator. Through experiments across three diverse large language model (LLM) evaluation scenarios, we show that MultiPPI consistently achieves lower estimation error than existing baselines. This advantage stems from its budget-adaptive allocation strategy, which strategically combines subsets of models by learning their complex cost and correlation structures.\looseness=-1


% \begin{center}\textbf{Version 2:}\end{center}
% A core challenge in modern AI model development is obtaining high-quality evaluation metrics in a cost-effective way. Such evaluation often involves tradeoffs between expensive, high-quality measurements and a variety of lower-quality proxies. We introduce Multiple-Prediction-Powered Inference (MultiPPI), a general framework for constructing statistically efficient estimates by optimally allocating resources across these diverse data sources. We provide theoretical guarantees about the minimax optimality, finite-sample performance, and asymptotic normality of the MultiPPI estimator. Through experiments across three diverse large language model (LLM) evaluation scenarios, we show that MultiPPI consistently achieves lower estimation error than existing baselines. This advantage stems from its budget-adaptive allocation strategy, which strategically combines subsets of models by learning their complex cost and correlation structures.

% and strategically incorporates more expensive, high-quality predictors as the budget increases.
\end{abstract}

\section{Introduction}
\label{sec:intro}

Efficiently estimating expectations of random variables under a fixed budget is a fundamental problem in many scientific settings. This paper addresses the general problem of optimally estimating such expectations when each random variable has an associated sampling cost, subject to a fixed total budget constraint. We are specifically motivated by the challenge presented by AI model evaluation, which is a critical, but often resource-intensive, step in model development and maintenance. %This paper focuses on the common scenario of choosing between a high-quality, but expensive, measurement process and \textbf{various}  cheaper, but lower-quality, proxies. 


More concretely, in the AI model evaluation setting, a variable $X_1$ might represent a high-quality but expensive metric computed for every model response to an input query, such as a score from a human annotator or a powerful proprietary model used as an "autorater". The remaining variables, $X_2, \ldots, X_k$,  might represent cheaper evaluation options (e.g., scores from smaller autoraters or rule-based systems), which can be viewed as covariates or proxies for the true score. Given the option to obtain samples of $X_1, \ldots, X_k$ (either jointly or independently), the primary objective is  often to then estimate the mean of the high-quality score, $\Ex[X_1]$. In other cases, we may be interested in the mean difference between two scores, say, $\Ex[X_1-X_2]$. The core difficulty in each case is in determining \emph{which} of these variables to query, how \emph{many times} to query them, and then finally how to \emph{combine} them together to produce a statistically efficient, consistent estimate of the ground truth.\looseness=-1
%Other objectives are also possible, such as estimating the difference between two models, $\Ex[X_1 - X_2]$\alekh{By two models here, we mean two autoraters? It is not two LLMs based on our notation here, if that i what we wanted to say.}

To formalize this, let $X := (X_1, \ldots, X_k)$ be a set of random variables with finite variance. We then consider the general problem of efficiently estimating any linear function of the mean of $X$ subject to a total observation budget $B$. That is, for some $a \in \mathbb{R}^k$, we want to estimate $\theta^* = a^\top \Ex [X]$ while spending no more than a total budget $B$ on collecting subsets of joint random variables $X_I = \{X_i\}_{i\in I}$ at cost $c_I$ for index subsets $I \subseteq \{1, \ldots, k\}$. More precisely, if $n_I$ is the number of times the subset $X_I$ is observed, we require that the $n_I$ satisfy a system of linear budget constraints of the form $\sum_{I} c_I n_I \leq B$, where the sum is over all such collected subsets $I$.\looseness=-1 %In fact, we further generalize this constraint to the case of satisfying multiple budget inequalities simultaneously, $\{B_1, \ldots, B_b\}$, where $B_i$ might represent an additional constraint such as a limit on the total number of observations of set $I$ (e.g., if we have a limit on the number of available prompts).\looseness=-1

Estimating linear functions of $\Ex[X]$ allows for flexibility in how $\theta^*$ is defined. Given the AI evaluation setting above, for example, measuring $\Ex[X_1]$ corresponds to  $a=(1,0,\ldots,0)$, while measuring  $\Ex[X_1 - X_2]$ corresponds to  $a=(1,{-}1,0,\ldots,0)$. 
The flexibility to observe subsets of $X$ also introduces a key trade-off that is unique with respect to previous related approaches to estimation. As we will show, observing variables jointly can be advantageous by reducing overall estimation variance. This benefit, however, must be weighed against the data acquisition costs, $c_I$. We make no assumptions about the structure of these costs (e.g., they may be non-additive over the components in $I$). For instance, in our AI evaluation setting, obtaining predictions from multiple autoraters can often be parallelized, so the cost of multiple predictions (in latency) is not significantly more than that of the single slowest one. This is not always true; in medical diagnostics, for example, ordering many tests may become too taxing for a patient, and therefore undesireable or impossible to do jointly.\looseness=-1


% Being able to observe general subsets of $X$ also adds flexibility.
% As we will show, observing variables jointly can be statistically advantageous, as it can reduce overall estimation variance. This benefit, however, must be weighed against the  cost of data acquisition. Furthermore, with respect to the data acquisition costs $c_I$: we do not make any assumptions on their structure (e.g., they may be non-additive over components indexed in the subset $I$). For example, in our AI model evaluation setting, obtaining predictions from multiple autoraters is often something that can be parallelized---so that the cost of obtaining multiple predictions is not significantly more expensive (at least in terms of latency) than obtaining the single slowest prediction. That said, this is not always the case. In medical diagnostics, for example, ordering too many tests may be taxing for a patient after a certain point and undesireable or impossible to do jointly.\looseness=-1

To solve this cost-optimal, multi-variate estimation problem,  we introduce the Multiple-Prediction-Powered Inference (MultiPPI) estimator, which is a cost-aware generalization of the Efficient Prediction-Powered Inference (PPI{++}) estimator of \citet{angelopoulos2023ppi++}, and extends it to \textbf{optimally leverage multiple types of predictions to power inference}. The MultiPPI estimator constructs a low-variance, consistent estimate of $\theta^\star$ by combining observations from judiciously chosen subsets of $X$. The core of our method is an optimization procedure that jointly determines the number of samples $n_I$ to draw from each subset $I$ and the corresponding linear weights $\lambda_I$ used to form the final estimate. We demonstrate that this allocation problem can be formulated as a second-order cone program (SOCP) for a single budget constraint, and a semidefinite program (SDP) for multiple budget constraints, and thus solved efficiently using standard techniques.\looseness=-1


Theoretically, we show that the MultiPPI estimator is minimax optimal when the joint covariance matrix, ${\Sigma = \Cov(X)}$, is known. For the typical case where it is unknown, however, we provide a  framework for  integrating an initial estimation phase where an approximation of the required covariance matrix, $\widehat{\Sigma}$, can be derived from either a small "burn-in" sample or a pre-existing labeled "transfer" dataset (a common scenario in applied settings)---and  provide finite-sample bounds on the performance degradation that is incurred by substituting $\widehat{\Sigma}$ for  $\Sigma$. Finally,  we empirically demonstrate the effectiveness of this approach across three diverse LLM evaluation settings, including choosing between autoraters of different sizes, autoraters with different test-time reasoning configurations, and complex multi-autorater-debate scenarios. In all cases, our method achieves lower mean-squared error and tighter confidence intervals for a given annotation budget than existing baselines. We demonstrate that MultiPPI achieves this by automatically tailoring its strategy to the available budget $B$: that is,  it learns to rely primarily on the cheaper autoraters when the budget is small, and naturally begins to incorporate more expensive, better autoraters as the budget increases. Taken together, our work provides a principled and computationally tractable framework for cost-effective, model-aided statistical inference, in settings with complex cost-versus-performance tradeoffs.\looseness=-1


In summary, our main contributions are as follows:
\begin{itemize}[leftmargin=*]
\item We introduce the \textbf{MultiPPI estimator} and frame the problem of finding the optimal subset sampling strategy and estimator weights as an efficient second-order cone program (SOCP).
\item We prove that the MultiPPI estimator is \textbf{minimax optimal} when the covariance matrix $\Sigma$ of $X_1, \ldots, X_k$ is known, and provide finite-sample performance guarantees for the practical setting where the covariance matrix must first be estimated as a part of the overall inference problem.\looseness=-1
\item We demonstrate MultiPPI’s applicability across multiple  LLM evaluation settings, and show how it can effectively combine signals from different model sizes, reasoning configurations, and multi-agent debates to achieve \textbf{lower error and tighter confidence intervals} for a given budget.\looseness=-1
\end{itemize}

%\alekh{Should we say somewhere in here that we do not expect $k$ to be very large, because otherwise enumerating over subsets starts to look daunting?}
%\adam{TODO: add comment for reasonable $k$}


% \begin{itemize}
%     \item We introduce \textit{Multi-PPI}, a strategy for statistically-efficient semi-supervised estimation with many proxies on a budget.
%     \item We prove that our strategy is minimax in the setting in which the covariance matrix is known, and yields natural confidence intervals which are asymptotically normal.
%     \item Our strategy guarantees that budget constraints will always be met, through easy-to-implement deterministic sampling strategies, in contrast to prior work.
%     \item We provide a detailed analysis of the estimator in each budget-regime, and compute in closed-form its performance.
%     \item We provide a finite-sample analysis of the performance of our estimator.
%     \item We demonstrate the effectiveness of our estimator empirically on three datasets, and showcase its versatility through implementation for (i) win-ratio estimation for model ranking, (ii) process evaluation through test-time scaling with a variable thinking budget, and (iii) multi-agent debate for factuality evaluation.
% \end{itemize}

% \charlie{Would like to make the following point: All existing methods involve estimating the covariance matrix $\Sigma$ and then computing the optimal parametrized estimator for the covariance matrix. Here we ask: is that even optimal? In many cases, especially the case of multiple models, the answer is no. We derive the minimax optimal strategy and demonstrate it in several experiments. The optimal strategy can involve nuanced behavior such as cascading.}
% \adam{I don't think this should go here. Put it later in the paper.}
% \begin{center}
% \fbox{%
%   \begin{minipage}{1.0\textwidth}
%     \paragraph{Multi-PPI:} 
%     \paragraph{Input:} the covariance matrix $\Sigma$ between the label $Y$ and the predictive models $(\hat{Y}_1,\dots,\hat{Y}_k)$, and a set of linear constraints on the number of times each model may be queried.
%     \paragraph{Output:} The minimax-optimal strategy for estimating $\Ex Y$ which respects the hard linear constraints, in the family of distributions of covariance $\Sigma$.
%   \end{minipage}
% }\label{fact}
% \end{center}
\section{Related work}
% \amirg{there are works on regression under a budget on measurements, such as \url{https://arxiv.org/pdf/1108.4559} and \url{https://www.cs.huji.ac.il/~shais/papers/2011_JMLR_CesShalSham.pdf}. Would be good to mention, and to also discuss active learning in our setting.}
%\amirg{At some point we should emphasize that our setting is not regression from $X_2,\ldots,X_n$ to $X_1$. Presumably one could view our problem via that lens and use standard regression (with partial observation). Maybe need to say what that would give.}

Our work builds upon  Prediction-Powered Inference~\cite[PPI;][]{angelopoulos2023prediction}, a statistical framework for efficiently estimating population-level quantities by augmenting a small set of labeled data with predictions from a machine learning (ML) model. We specifically build on PPI{++}, the efficient extension of PPI introduced in \citet{angelopoulos2023ppi++}, which also further improves variance by optimally reweighting these predictions. We describe PPI in greater depth in \Cref{sec:preliminaries}.

%that also allows for learning an optimal reweighting of the machine learning model's predictions to improve estimator variance.

PPI is part of a broader class of statistical methods that leverage ML predictions for estimation. Its principles connect to classical control variates and difference estimators~\cite{Ripley87, model1992sarndal, chaganty-etal-2018-price}, which reduce  variance by subtracting a correlated random variable with a known mean; the correlated variable in PPI is the ML prediction, whose mean can be (cheaply) estimated on unlabeled data. This approach also shares theoretical foundations with modern semi-parametric  inference, particularly methods from the causal inference literature like Augmented Inverse Propensity Weighting~\cite[AIPW;][]{aipw_robins}, Targeted Maximum Likelihood Estimation~\cite[TMLE;][]{vanderLaan2006tmle}, and double  machine learning~\cite[DML;][]{chernozhukov2018dml}. Recently, PPI has been applied to  Generative AI evaluation, where human annotations (or more generally, annotations from some trusted source) are combined with cheaper "autorater" outputs for efficient, unbiased estimates of model performance~\cite{boyeau2024autoeval, chatzi2024ppirank, fisch2024stratified,angelopoulos2025costoptimalactiveaimodel, saad-falcon-etal-2024-ares, demirel2024prediction}.\looseness=-1


Existing PPI frameworks, however, assume either a single predictor~\cite{angelopoulos2023prediction, angelopoulos2023ppi++} or a fixed set of predictors queried together~\cite{miao2024assumptionleandataadaptivepostpredictioninference}. We address the common scenario where multiple predictors (e.g., autoraters) with different cost-performance profiles are available. This introduces a complex budget allocation problem: determining which  predictors to query (individually, jointly, or in any joint subset), how often to query them, and how to combine the measurements they provide for a minimum-variance estimate under a fixed budget. Our work partially generalizes \citet{angelopoulos2025costoptimalactiveaimodel}, which optimizes a sampling policy for a single predictor. Unlike that work, however, which focuses on input-conditional policies and expected budget constraints, we find a fixed allocation policy that always satisfies a hard budget constraint for every run.\looseness=-1


% Existing PPI frameworks, however, typically operate in a setting with a single predictor~\cite{angelopoulos2023prediction, angelopoulos2023ppi++} or a fixed, vector-valued set of predictors~\cite{miao2024assumptionleandataadaptivepostpredictioninference} that are always queried together. This leaves open the problem of optimal budget allocation in the common scenario where a practitioner has access to a wide variety of potential autoraters, each with a different cost-performance profile---which is the problem we study here. This introduces a complex resource allocation problem: one must determine not only the sampling ratios, but also which predictors to query, how often to query them (both individually and jointly), and how to combine the resulting measurements to form a minimum-variance estimate, all subject to a total budget constraint. In particular, this problem is also a partial generalization of the one studied in the related, recent work of \citet{angelopoulos2025costoptimalactiveaimodel}, which considers optimizing the sampling policy for a single predictor under cost constraints (though in this current work, we leave the derivation of sampling policies that are conditional on individual inputs to future study). Furthermore, \citet{angelopoulos2025costoptimalactiveaimodel} optimize for constraints that are required to hold only in expectation, whereas the budget constraints that we consider hold for every run.\looseness=-1

Our allocation problem is also related to budgeted regression with partially observed features~\cite{JMLR:v12:cesa-bianchi11a, hazan2012linear} and active learning or testing~\cite{settles2009active, activetesting, zhang2015active}. We emphasize, however, that our goal is estimation of a linear function of a population mean (i.e., $a^\top\Ex[X]$), and not regression (e.g., predicting $X_1$ from $X_2, \ldots, X_k$). While related, standard approaches to regression, including with partial observations, optimize for sample-wise predictive accuracy rather than for predictive accuracy of a population-level quantity. Our problem also connects to multi-armed bandit allocation for adaptive Monte Carlo estimation~\cite{pmlr-v32-neufeld14}. A key difference is that these frameworks often use sequential, input-dependent policies to minimize regret, making it difficult to derive valid confidence intervals (CIs). Our framework, in contrast, computes a fixed allocation policy over predictive models (not individual inputs as in active learning or testing) and guarantees unbiased estimates with valid CIs. Even more broadly, our work shares similar high-level goals with transfer learning and domain adaptation~\cite[][\emph{inter alia}]{Pan2010ASO, bendavid2010theory}---i.e., leveraging signals of varying quality and potential bias---though the statistical techniques are distinct.\looseness=-1

 
% Lastly, our allocation problem is broadly related to work on budgeted regression with partially observed features~\cite{JMLR:v12:cesa-bianchi11a, hazan2012linear}, as well as  active learning and active testing~\cite{settles2009active, activetesting}, where the goal is to select the most informative samples or measurements to maximize predictive performance  subject to some cost. It is also conceptually related to resource allocation problems studied in multi-armed bandits~\cite{pmlr-v32-neufeld14, perchet2016batched}. However, these frameworks often focus on sequential policies to minimize regret, and it is less clear how one would derive valid confidence intervals from their resulting estimates. Our framework, in contrast, computes a  fixed allocation policy over proxy predictive \emph{models} versus individual inputs, and  is guaranteed to produce unbiased, consistent estimates with valid confidence intervals.\looseness=-1

% \alekh{Maybe some more refs: https://arxiv.org/abs/1510.02847 looks relevant for active learning. Do we need to comment on the broader literature on transfer learning and domain adaptation as well, which is qualitatively quite different in the techniques, but shares the same basic flavor of leveraging signals with different levels of utility, even in the presence of bias?}



%Our work introduces Multi-Prediction-Powered Inference (MultiPPI), an estimator that directly addresses this multi-predictor allocation problem. Concretely, MultiPPI generalizes the PPI{$++$} estimator by establishing a principled and computationally tractable framework that jointly optimizes the sampling counts ($n_I$) for arbitrary subsets of predictors ($I$) and the linear weights ($\lambda_I$) used to construct the final estimate from subset predictions. We demonstrate that this optimal allocation problem can be formulated as a second-order cone program (SOCP) and thus solved efficiently. 

%While conceptually related to resource allocation problems studied in the context of multi-armed bandits~\cite{pmlr-v32-neufeld14, pmlr-v40-Perchet15}, our framework is designed to provide confidence intervals, and is best used in the setting where covariance information can be pre-estimated, which permits the direct computation of an allocation policy. \charlie{I made a slight change here to emphasize that our goal is confidence intervals, whereas it's not clear to me how one would obtain confidence intervals from the bandit policies. But maybe I'm being naive.}

\section{Preliminaries}
\label{sec:preliminaries}

In the following section, we introduce the general estimation problem of interest and summarize existing approaches. Suppose that we are interested in the mean of a random variable $X_1$, which is dependent upon another random variable $X_2$ (corresponding to estimating $a^\top \Ex[X]$ for $a = (1, 0)$ as described in \S\ref{sec:intro}). For example, in the AI model evaluation setting, $X_2$ may be an autorater's score for a model output to a user's query, and $X_1$ may be the ground truth quality of the response as measured by an expert human annotator. Suppose we have access to a small number ($n$) of i.i.d. samples that contain labels from both the target rater ($X_1$) and autorater ($X_2$), and a large number ($N$) of i.i.d. samples that contain only the autorater predictions ($\tilde{X}_2$). A na\"ive approach to estimating the mean is to simply take the sample average of $X_1$ and ignore  $X_2$ entirely, which we denote by $\hat{\theta}_{\text{classic}} = \frac{1}{n}\sum_{j=1}^n X_1^{(j)}$. When the prediction $X_2$ is correlated with $X_1$ and easy to query, however, it is natural to consider the "prediction-powered" PPI estimator~\cite{angelopoulos2023prediction, angelopoulos2023ppi++}:\looseness=-1
\begin{equation}
\hat{\theta}_{\text{PPI}} = \frac{1}{n}\sum_{j=1}^n \left(X_1^{(j)}-X_2^{(j)}\right)
%_{\substack{\text{Variance }\frac{1}{n}\sigma^2_{X_1-X_2}\ll \frac{1}{n}\sigma^2_{X_1} \\ \text{if $X_1\approx X_2$ in $L^2$}}} + \underbrace{
+\frac{1}{N}\sum_{j=1}^{N}\tilde{X}_2^{(j)}
%}_{\text{Variance }\mathcal{O}(\frac{1}{N})}.
\end{equation}
%
%\jacob{the underbrace states that $\text{Variance }\frac{1}{n}\sigma^2_{X_1-X_2}\ll \frac{1}{n}\sigma^2_{X_1}$; we should acknowledge that this depends on $X_1$ and $X_2$ rather than imply that we think it is generally true.}\charlie{Good point, I suppose this is done in the subsequent text, but let me amend the equation}
When we can afford to take $N$ to be very large, it is clear that the variance of $\hat{\theta}_{\text{PPI}}$ is much smaller than that of $\hat{\theta}_{\text{classic}}$ provided that our model predictions $X_2$ are close to $X_1$ in mean-squared error. When that fails, \citet{angelopoulos2023ppi++} propose adding a linear fit of the form:
%
\begin{equation}\label{eqn:ppi++}
    \hat{\theta}_{\text{PPI{++}}} = \frac{1}{n}\sum_{j=1}^n \left(X_1^{(j)}-\lambda X_2^{(j)}\right) + \frac{1}{N}\sum_{j=1}^N\lambda\tilde{X}_2^{(j)}.
\end{equation}
The parameter $\lambda$ may be chosen to minimize the variance of $\hat{\theta}_{\text{PPI$++$}}$ based on the observed labeled data. This strategy yields an estimator which asymptotically improves on $\hat{\theta}_{\text{classic}}$ and $\hat{\theta}_{\text{PPI}}$ in the limit that $n\to\infty$ and $N\gg n$. Toward the setting where $n$ and $N$ may be comparable in size, if one is able to choose to or not to request a label $X_1^{(j)}$ for every observed unlabeled point $X_2^{(j)},$ a modification of $\hat{\theta}_{\text{PPI{++}}}$ allows one to do so in a cost-optimal way \cite{angelopoulos2025costoptimalactiveaimodel}. %In each case, knowledge of the \textit{covariance} matrix of $(X_1, X_2)$ is enough to determine the optimal strategy.

\subsection{Multiple predictive models}
\label{sec:multiple_models}
How should one adapt the preceding setting when one has access to many predictions, rather than just $X_2$? One option is to \textit{stack} all predictions into a vector $X_{2:k}:=(X_2,\dots,X_k)$ and choose $\lambda\in\R^{k-1}$ to be a vector in $\hat{\theta}_{\text{PPI++}}$; this is the estimator proposed by \citet{miao2024assumptionleandataadaptivepostpredictioninference}, and can be written\looseness=-1
%
\begin{equation}\label{eqn:ppi++vector}
    \hat{\theta}_{\text{PPI++ vector}} = \frac{1}{n}\sum_{j=1}^n \left(X_1^{(j)} - \lambda^\top X_{2:k}\right) + \frac{1}{N}\sum_{j=1}^N \lambda^\top \tilde{X}_{2:k}^{(j)}
\end{equation}
%
But this approach is suboptimal when (as is  becoming standard) the best models may be available only for the highest prices: if any of $X_2,\dots,X_k$ is expensive to obtain, our ability to sample $X_{2:k}$ will be limited. This yields suboptimal results, as we show in \S\ref{sec:results}. One may instead decide to perform PPI with just one model $X_{i}$, for whichever $i \neq 1$ has the best cost/accuracy tradeoff---but it is not clear \textit{a priori} which one this is, or how much worse it may be compared to some combination of a cost-effective subset of $X$.
%\amirg{If you have the covariance matrix this is probably easy to do, so not sure it should be mentioned as the weakness. Isn't it more that one model is suboptimal to multiple?}\charlie{Good point. I suppose I meant that, to a practitioner familiar with PPI++, it still might not be clear which model is best. But this point is not very important.}
Alternatively, perhaps it is  possible for cheaper models be used to recursively estimate the means of more expensive models, thus creating a \textit{PPI{++} cascade}: for instance, if $k=3$ and $(X_1,X_2,X_3)$ are in decreasing order of cost, we might consider\looseness=-1
\begin{equation}\label{eqn:cascade}
    \hat{\theta}_{\text{PPI++ cascade}} = \frac{1}{n}\sum_{j=1}^n \left(X_1^{(j)} - \lambda X_2^{(j)}\right) + \frac{1}{N}\sum_{j=1}^N\left(\lambda \tilde{X}_2^{(j)}-\lambda'\tilde{X}_3^{(j)}\right) + \frac{1}{M}\sum_{j=1}^M\lambda'\tilde{\tilde{X}}_3^{(j)}
\end{equation}
%
% \amirg{not clear if you are presenting this as how it related to what we do.}\jacob{not crucial for this version, but it might be nice for motivation to actually write this out as you've done in the slides} \looseness=-1
%
Each of these strategies can be realized as possible instances of the MultiPPI estimator we propose in the next section. Rather than coarsely limiting ourselves to sampling $X_{2:k}=(X_2,\dots,X_k)$ together, we allow the flexibility of sampling $X_I$ for generic index subsets $I\subseteq\{1,\dots,k\}$. %\textcolor{red}{We allow sampling of $X_I$ for $I$ lying in a predetermined family of index subsets $\mathcal{I}\subseteq 2^k$; this family is chosen to be large enough to afford the strategies described above, but not so large that our computational procedure (see \autoref{sec:computation}) is intractable or our finite-sample convergence is too slow (see \Cref{cor:finite}).
%For example, \Cref{eqn:ppi++} can be understood as the case $\mathcal{I}=\{\{1,2\},\{2\}\}$, while \Cref{eqn:ppi++vector} corresponds to the case $\mathcal{I}=\{\{1,\dots,k\},\{2,\dots,k\}\}$.} \adam{TODO: see where this shakes out with the experiments} 
% In the experiments for this paper \jacob{the analysis doesn't make this restriction?}, we consider the minimal modification to this setup in which we take $\mathcal{I}=\{\{1,\dots,k\},\{2,\dots,k\},\{2\},\dots,\{k\}\}\}$. The resulting estimator, to be specified in the next section, then allows for flexible adjudication between (1) returning $\hat{\theta}_{\text{PPI++}}$ with the best (in terms of cost/accuracy tradeoff) model among $X_2,\dots,X_k$, (2) returning $\hat{\theta}_{\text{PPI++ vector}}$, and (3) cascading a more expensive model by a cheaper one. We achieve this while maintaining that $|\mathcal{I}|=\mathcal{O}(k)$, which ensures that our computational procedure is tractable and that our finite-sample convergence is not too slow.

% The following questions are left open when one has access to many models:
% \begin{enumerate}
%     \item Which predictive model, in terms of cost/accuracy tradeoff, is the best one to use in PPI?
%     \item Can cheaper models be used to estimate the means of more expensive models, thus creating a \textit{PPI cascade}?
% \end{enumerate}

% \subsection{Multiple-prediction-powered inference \charlie{[WIP]}}
% \adam{Lets move this into the next section.}

% Each of the preceding estimators proposed is of the following flavor: First, some amount of \textit{burn-in} or \textit{transfer} data is used to estimate the relevant information needed to determine the optimal parameters $\lambda$; this is the covariance matrix. Second, $\lambda$ is estimated according to this covariance. Third, the resulting estimator is deployed.

% It is not clear that this strategy is optimal, even if one estimates the covariance matrix perfectly. In this paper, we answer the following question:
% \begin{center}
%     \textit{What is the minimax-optimal strategy for estimating the mean of $Y$ on if one has access to many predictive models $f_1(X),\cdots,f_k(X)$, each for a price, and has knowledge of their joint covariance?}
% \end{center}
% We describe the minimax estimator in this setting in the following section. We analyze this estimator theoretically, and show that it is asymptotically normal, among other things. Lastly, we demonstrate our estimator in several experimental settings described in \autoref{sec:exp-setup}.

\section{Multiple-prediction-powered inference (MultiPPI)}\label{sec:methods}
As \Cref{sec:multiple_models} highlights, it is not obvious how to best allocate a budget across a diverse suite of predictive models, where each model has its own cost and performance tradeoffs. We begin by defining the class of permissible estimators: We require that the number of times, $n_I$, that $X_I$ is sampled satisfies a linear budget constraint, specified by a set of non-negative costs $c_I\geq0$ and total budget $B\geq0$, for each index subset $I\subseteq\{1,\dots,k\}$.\footnote{In \Cref{app:generalized}, we extend the methodology to multiple budget constraints.}
\begin{definition}
\label{def:budget}
    An estimator $\hat{\theta}$ is \textbf{budget-satisfying} if it a measurable function of $n_I$ i.i.d. samples of $X_I$, for each $I\subseteq\{1,\dots,k\}$, such that $\sum_{I} n_Ic_I \leq B$.
\end{definition}
To develop a principled search for the best budget-satisfying estimator, we begin by asking a simple question under idealized conditions:
\begin{question}
\label{eqn:minimax}
    If the  covariance matrix, $\Sigma = \Cov(X)$, is exactly  known, what is the minimax optimal, budget-satisfying estimator of $\theta^* = a^\top\mu$ with respect to the mean-squared error, $\Ex[(\hat{\theta} - \theta^*)^2]$?
\end{question}
The answer to \Cref{eqn:minimax} will provide us with a set of allocations $n_I$ and a corresponding budget-satisying estimator $\hat{\theta}_{\multiPPI}$ which we will evaluate on the $n_I$ samples of $X_I$, for each $I$. Once we have addressed this question, we address the case of unknown $\Sigma$ by describing strategies depending on the empirical covariance matrix $\wh{\Sigma}$, which may be estimated from data.

It turns out, perhaps surprisingly, that \Cref{eqn:minimax} reduces to the following tractable alternative:
\begin{question}
\label{eqn:linear}
If the  covariance matrix, $\Sigma = \Cov(X)$, is exactly  known, what is the minimum variance, \underline{linear}, \underline{unbiased} budget-satisfying estimator of $\theta^*= a^\top\mu$?
\end{question}
We demonstrate the equivalence of \Cref{eqn:minimax} and \Cref{eqn:linear} in \Cref{thm:minimax}. For now, the "oracle" assumption on knowing the covariance matrix $\Sigma$  allows us to isolate the resource allocation problem from the separate challenge of estimating how closely related  $(X_1, \ldots, X_k)$ are to begin with, and to analyze what a good procedure for leveraging multiple predictive models under cost constraints should look like in theory. All proofs of our theoretical results are deferred to Appendix~\ref{app:proofs}.

%Furthermore, while the current restriction on $\hat{\theta}$ to the class of linear unbiased estimators may seem strong, we will later show that the resulting best such estimator of this kind is in fact minimax optimal among \emph{all} possible estimators that satisfy our observation budget constraint $B$. 


% Under this assumption, we derive a minimax-optimal, unbiased, and consistent estimator for our general task: estimating a linear function $a^\top\mu$ of the mean $\mu = \Ex[X]$ for a collection of random variables $X = (X_1, \ldots, X_k)$ subject to a budget on querying subsets of these variables.

\subsection{MultiPPI($\Sigma$): A minimax optimal algorithm}
Recalling notation from \Cref{sec:intro}, let $X\in\R^k$ denote a random vector of finite second moment with distribution $\p$. Let $\mathcal{I} \subseteq 2^{\{1, \ldots, k\}}$ denote a collection of index subsets which may be queried, and for any $I \in \mathcal{I}$, let  $X_I = \{X_i\}_{i \in I}$ be the corresponding subset of $X$. Next, let $\underline{n} = \{n_I\}_{I \in \mathcal{I}}$, $n_I \in \mathbb{N}$ be an allocation of sample sizes, where $n_I$ i.i.d. samples are drawn for each subset $I$, and let $\underline{\lambda} = \{\lambda_I\}_{I \in \mathcal{I}}$, $\lambda_I \in \mathbb{R}^{|I|}$ define a corresponding set of weighting vectors for each subset  $I$. Finally, let $\hat{\theta}(\underline{n}, \underline{\lambda})$ denote the weighted sum of sample means from each non-empty subset, i.e.,\looseness=-1
\begin{equation}\label{eq:theta-n-lambda}
    \hat{\theta}(\underline{n}, \underline{\lambda}) = \sum_{I:n_I>0} \frac{1}{n_I}\sum_{j=1}^{n_I} \lambda_I^\top X_I^{(j)}.
\end{equation}
The MultiPPI estimator, $\hat{\theta}_{\text{MultiPPI}}$, is then defined as the optimal estimator in this class that minimizes the MSE subject to our unbiasedness (\textbf{U}) and budget (\textbf{B}) constraints:
\begin{equation}
\label{eq:multippi_opt}
    \hat{\theta}_{\text{MultiPPI}} = \underset{\hat{\theta}(\underline{n}, \underline{\lambda})}{\argmin} ~\Ex\left[\Big(\hat{\theta}(\underline{n}, \underline{\lambda}) - \theta^*\Big)^2\right] \quad \text{s.t.} \quad \text{
    \textbf{U} and \textbf{B} hold}, 
\end{equation}
where the constraints \textbf{U}  and \textbf{B} are
\begin{align*}
    \textbf{U}\iff \Ex[ \hat{\theta}(\ul{n},\ul{\lambda})] = \theta^*\text{ for all $\p$ of finite second moment} \quad\quad\text{and} \quad\quad \textbf{B}\iff\sum_{I}n_Ic_I\leq B.
\end{align*}
It can be shown that \textbf{U} reduces to a linear constraint on $\ul{\lambda}$, which makes our optimization convenient.

As previously discussed, the estimators of \Cref{eqn:ppi++vector} and \Cref{eqn:cascade} can be viewed as special cases of this setup. For instance, \Cref{eqn:ppi++vector} corresponds to imposing the additional restriction that $\lambda_I=0$ for all $I\in 2^{\{1,\dots,k\}}$ except for $I=\{1,\dots,k\}$ and $I=\{2,\dots,k\}$; \Cref{eqn:cascade} corresponds to the additional restriction that $\lambda_I=0$ for all $I$ except for $\{1,2\}$, $\{2,3\}$ and $\{3\}$.\looseness=-1

\subsubsection{Optimization}
Solving \Cref{eq:multippi_opt} is, in general, non-trivial. Since $\hat{\theta}(\ul{n},\ul{\lambda})$ is linear in $X$, it can be shown that the optimal $(\ul{n},\ul{\lambda})$ depend only on the covariance matrix $\Sigma$ of $X$, and so we will denote by $\hat{\theta}_{\multiPPI(\Sigma)}$ the solution to \Cref{eq:multippi_opt} given any distribution such that $\Sigma=\Cov(X)$. Then, it can be further shown (this follows from \Cref{thm:minimax}, presented next) that the MSE of $\hat{\theta}_{\multiPPI(\Sigma)}$ is
\begin{equation}
    \label{eq:variance}
    \mathcal{V}_B=\underset{\substack{\ul{n}\,:\,\text{\textbf{B} holds} \\ \supp(a)\subseteq\bigcup\{I:n_I>0\}}}{\min} a^\top S(\ul{n}) a,\quad\quad S(\ul{n})=\left(\sum_{I\in\mathcal{I}} n_I \Sigma_I^\dagger\right)^\dagger
\end{equation}
where $\Sigma_I$ denotes the principle submatrix of $\Sigma$ on $I$, embedded back into $\R^{k\times k}$, and $\dagger$ denotes the Moore-Penrose pseudo-inverse.\footnote{More formally, if $P_I\in\R^{k\times k}$ denotes the orthogonal projection onto $\text{span}(I)\subseteq\R^k$, we define $\Sigma_I=P_I\Sigma P_I^\top$, and so $\Sigma_I^\dagger := (P_I\Sigma P_I^\top)^\dagger$.}
%
The minimizing $\ul{n}$ of the above expression then also determines the optimal $\lambda_I$ to be the restriction of $n_I\Sigma_I^\dagger S(\ul{n})a$ to the coordinates $I$. If the integrality constraints on $n_I$ are relaxed, we show in the appendix that this reduces to a second-order cone problem in the case of a single budget constraint, and a semi-definite program in the case of multiple budget constraints. This allows for  \Cref{eq:variance} to be solved efficiently using standard techniques (\Cref{sec:computation}).\looseness=-1 %, as we show in the appendix, and provide justification for the integral relaxation.

\subsubsection{Minimax optimality}
The minimal MSE $\mathcal{V}_B$ shown in \Cref{eq:variance} has a more fundamental characterization. Here we show that it is in fact the minimax optimal MSE achievable by \textbf{any} budget-satisfying estimator, taken over the set of distributions $\mathcal{P}$ of covariance $\Sigma$. Consequently, the estimator defined by $\hat{\theta}_{\text{MultiPPI}(\Sigma)}$ is minimax optimal over the set of distributions $$\mathcal{P}_\Sigma = \{\text{distribution }P\text{ on }\mathbb{R}^k \colon \Cov(X) = \Sigma\text{ for }X \sim P\}.$$ 
%
Specifically, given costs $(c_I)_I$ and a budget $B$, let $\Theta_B$ denote the set of budget-satisfying estimators $\hat{\theta}$ per \Cref{def:budget}. We emphasize that we make \textbf{no}  restriction on $\Theta_B$ to include only linear or unbiased estimators. 
 Then the following result holds:
 
\begin{theorem}[Minimax optimality of MultiPPI for known $\Sigma$]
\label{thm:minimax}
For all $\Sigma\succ0$, we have
    \[
    \inf_{\hat{\theta}\in\Theta_B}\sup_{P\in\mathcal{P}_\Sigma} \Ex\left[(\hat{\theta}-\theta^*)^2\right] = \Var 
    \left(\hat{\theta}_{\multiPPI(\Sigma)}\right) = \mathcal{V}_B,
    \]
    where the variance is with respect to any distribution $P\in\mathcal{P}_\Sigma$.
\end{theorem}


% Our MultiPPI estimator is a linear function of the queried samples from subsets of $(X_1, \ldots, X_k)$. Specifically, let
% \begin{equation}\label{eqn:multippi}
%     \hat{\theta}_{\text{MultiPPI}} = \sum_{I:n_I>0} \frac{1}{n_I}\sum_{j=1}^{n_I} \lambda_I^\top X_I^{(I,j)},
% \end{equation}
% where $n_I\geq0$ is a non-negative integer to be specified for each index subset $I$,  $\lambda_I\in\R^{I}$ is a parameter learned from the observed data to be specified for each subset $I$, $X_I$ denotes $(X_i)_{i\in I}$, and the sets of random variables $\{X_I^{(I,j)}\}_{I\in \mathcal{I},j\leq n_I}$ are drawn independently.

%  Suppose the covariance matrix $\Sigma$ is known. Then the parameters $n_I,\lambda_I$ are chosen to minimize the variance of \Cref{eqn:multippi} under $\Sigma$ subject to the following two restrictions:\amirg{nit: this condition implies unbiasedness but not the other way around, right? The naming is a bit confusing in this sense. Can we provide an argument that this is wlog?} \charlie{Hmm, I might be missing something, but it seems to me that it's equivalent to unbiasedness, unless I've made a mistake. Maybe the wording is unclear in any case, though.} \adam{It's not a necessary condition for unbiasedness, no? Consider X with mean 0.} \charlie{I see, I mean my impression is that ``unbiased'' generically means ``unbiased'' for all distributions for which means exist? Would it be appropriate if I specify that?}
% \begin{align*}
%     \textbf{Unbiasedness:}&\quad\quad\textbf{U}\iff  \sum_{I \colon j \in I, n_I>0} \lambda_{I(j)}=a_j\,\quad\forall 1\leq j\leq k \\
%     \textbf{Budget satisfaction:}&
%     \quad\quad \textbf{B}\iff\sum_{I}n_Ic_I\leq B.
% \end{align*}
% Note that although here we consider a budget constraint described by only a single inequality, we extend our theory to the case of many inequalities in \Cref{sec:generalized}. Also note that the unbiasedness constraint ensures that $\Ex[ \hat{\theta}_{\text{MultiPPI}}] = \theta^*$ when the $\lambda_I$ are independent from $X_I$.
% We refer to this optimization process via the $\textrm{MultiAllocate}(\Sigma)$ subroutine, shown below. 
% \begin{center}
% \fbox{%
%   \begin{minipage}{1.0\textwidth}
%     \paragraph{Subroutine 1:} 
%     MultiAllocate$(\Sigma)$
%     \paragraph{Input:} a symmetric positive-define matrix $\Sigma\in\R^{k\times k}$. Implicitly given $a$, $(c_I)_I$, $B$.
%     \paragraph{Output:} the solution $(n_I,\lambda_I)_{I}$ to the problem
%     \begin{align*}
%         \min_{(n_I,\lambda_I)_{I\in\mathcal{I}}} \sum_{I:n_I>0} \frac{1}{n_I}\lambda_I^\top \Sigma_I \lambda_I \quad\text{s.t.} \quad \text{
%     \textbf{U} and \textbf{B} hold}, 
%     \end{align*}
%     where $\Sigma_I=(\Sigma_{ij})_{i,j\in I}$.
%   \end{minipage}
% }\label{fact}
% \end{center}
%  $\textrm{MultiAllocate}(\Sigma)$ reduces to solving a second-order cone program (SOCP), which can be solved efficiently using off the shelf tools. See \Cref{sec:computation} for details.\amirg{need to explain why this is an SOCP.}
%  \jacob{Maybe not the highest priority, but I read the appendix and I still don't understand how we are solving the optimization over the integers $\{n_I\}$. Can we give some more intuition and/or a citation to a similar derivation?}\alekh{It is also not clear at this point how MultiAllocate arises as the objective for estimation of $\theta^*$. In fact, this is ill-defined as we have not even set up MSE as the criterion of interest at this point, I think. Defining the performance metric and explaining that this is equivalent due to unbiasedness seems important?}\amirg{I agree. Would be good to explicitly start from MSE, then variance then MultiAllocate.}

 


% We consider the task of estimating a linear function $a^\top \mu$ of the mean $\mu=\Ex X$ of a collection of random variables $X=(X_1,\dots,X_k)$. These random variables, and subsets thereof, are associated with certain costs, and our objective is to construct the most accurate estimator $\hat{\theta}$ of $\theta^*=a^\top \mu$ subject to the corresponding budget restriction. 

% We next introduce the minimax optimal estimator in the setting in which $\Sigma=\Cov X\succ0$ is known. Let $c_I$ denote the (vector-valued) cost of querying the random variables associated to an index subset $I\subseteq\{1,\dots,k\}$; the reason for discussing costs at the level of sets rather than singletons is to incorporate the frequently non-additive cost structure exhibited by many tasks in this domain (c.f. \autoref{sec:exp-setup}). We allow our costs to be vector-valued to incorporate multiple types of budget constraints, e.g. monetary and temporal, hard limits on the number of samples available from each model, etc.

% The MultiPPI estimator is
% \begin{equation}\label{eqn:multippi}
%     \hat{\theta}_{\text{MultiPPI}} = \sum_{I:n_I>0} \frac{1}{n_I}\sum_{j=1}^{n_I} \lambda_I^\top X_I^{(I,j)}
% \end{equation}
% where $n_I\geq0$ is a non-negative integer to be specified for each subset $I$, $\lambda_I\in\R^{|I|}$ is a tuning parameter\amirg{not sure this is the best name here. This is more similar to a regression coefficient, so maybe call it parameter/weight learned from data.} to be specified for each subset $I$, $X_I$ denotes $(X_i)_{i\in I}$, and the set of random variables $\{X_I^{(I,j)}\}_{I\in I,j\leq n_I}$ is drawn independently.




% Here, $B$ is a vector-valued quantity representing the maximum budget available across (potentially) multiple types of budget. The budget satisfaction inequality should be interpreted in the vector-valued sense. The first condition is equivalent to the statement that $\Ex \hat{\theta}_{\text{MultiPPI}(\Sigma)}=\theta^*$ for all distributions of finite second moment \jacob{is the first condition sometimes stronger the expectation-matching condition?}, and the second condition is equivalent to budget satisfaction under the model in which we query $n_I$ times the model subset $X_I$, in the vector sense.

% %We prove shortly that solving for the parameters $n_I,\lambda_I$ reduces to a second-order cone program (SOCP) and can be solved efficiently with off-the-shelf tools. First, we address the consideration that $\Sigma$ is, in general, unknown.

% We describe efficient methods for solving for the parameters $n_I,\lambda_I$ subsequently. First, we address the consideration that $\Sigma$ is, in general, unknown.

% \charlie{Just say we make the simplifying and realistic consideration that we have some amount of \textit{fully-labeled} data. We use this data to approximate the covariance. This preserves asymptotic normality.}


\subsection{MultiPPI($\widehat{\Sigma}$): A practical algorithm}\label{sec:practical}


%\subsubsection{Finite-sample bounds}
In practice, $\Sigma$ is rarely known and must be approximated by an estimated covariance matrix $\wh{\Sigma}$. In general, there are many methods for constructing an estimate $\wh{\Sigma}$ of $\Sigma$, and many of the theoretical properties of MultiPPI are agnostic to the particular choice made. The following theorem shows that, for any $\wh{\Sigma}$ which converges in probability to $\Sigma$ as our budget tends to infinity, the MultiPPI estimator is asymptotically normal and achieves the optimal variance of \Cref{thm:minimax}. For this result, we need a technical condition which amounts to \Cref{eq:multippi_opt} having a unique minimizer $\ul{n}$ as $B\to\infty$; we state it formally in \Cref{sec:asymptotic_proof}.\looseness=-1

\begin{theorem}\label{thm:asymptotic_revised}
Suppose $X\in\R^k$ has finite second moment, and suppose that $\Sigma=\Cov(X)$ satisfies condition \ref{eqn:condition}. Suppose that $\wh{\Sigma}\overset{p}{\to}\Sigma$ in the operator norm as $B\to\infty$. Then for $\hat{\theta}_{\multiPPI(\widehat{\Sigma})}$ arbitrarily dependent on any potential samples used to estimate $\widehat{\Sigma}$, we have
    \[
    \sqrt{B}\left(\hat{\theta}_{\multiPPI(\widehat{\Sigma})}-\theta^*\right)\overset{d}{\to}\mathcal{N}(0,\mathcal{V}^*)
    \]
    as $B\to\infty$, where $\mathcal{V}^*=\lim_{B\to\infty} B\mathcal{V}_{B}$, and $\mathcal{V}_{B}$ is defined in \Cref{eq:variance}.
\end{theorem}
It is important to note that the estimator $\hat{\theta}_{\multiPPI(\wh{\Sigma})}$ continues to enjoy unbiasedness, budget satisfaction, and asymptotic normality regardless of mis-specification in $\wh{\Sigma}$. 


%$\hat{\theta}_{\multiPPI(\wh{\Sigma})}$ remains unbiased, budget-satisfying, and asymptotically normal (as we see next) in a broad range of scenarios, regardless of mis-specification in $\wh{\Sigma}$. However, better estimates $\wh{\Sigma}$ of $\Sigma$ will lead to better performance of the estimator.

A natural question concerns the level of suboptimality of $\hat{\theta}_{\multiPPI(\wh{\Sigma})}$ as a function of the degree of mis-specification of $\wh{\Sigma}$ in finite samples. Below, we present a meta-result which serves to quantify the sensitivity of our procedure to errors in the specification of $\wh{\Sigma}$.

\begin{theorem}[Stability of MultiPPI]\label{thm:stab}
    Let $P$ be a distribution of covariance $\Sigma$, and suppose that $\Sigma$ has minimum eigenvalue $\gamma_{\min}$. Let $\sigma^2_{\classical}$ denote the least MSE of any budget-satisfying sample mean of $\theta^*$. Let $\widehat{\Sigma}$ denote any non-random symmetric positive-definite matrix. Then we have
    \[
    \Ex\left[\left(\hat{\theta}_{\multiPPI(\widehat{\Sigma})}-\theta^*\right)^2\right] \leq \mathcal{V}_B + \frac{4\sigma^2_{\classical}}{\gamma_{\min}}\|\wh{\Sigma}-\Sigma\|_{F}.
    \]
    whenever $\|\wh{\Sigma}-\Sigma\|_{F}\leq \gamma_{\min}/2$, where $\|\cdot\|_F$ denotes the Frobenius norm.
\end{theorem}

In general, there are many methods for constructing an estimate $\wh{\Sigma}$ of $\Sigma$, and \Cref{thm:stab} is agnostic to the particular choice made. In \Cref{app:finite-sample-bounds}, we show how to apply the meta-result above to derive a family of finite-sample bounds in a variety of distributional settings and for a variety of estimates $\wh{\Sigma}$.
In practice, we estimate ${\Sigma}$ from data, and find the Ledoit-Wolf estimator $\wh{\Sigma}$ of the covariance matrix $\Sigma$ to perform best in our experiments. This is consistent with the fact that the Ledoit-Wolf estimate is designed to minimize $\Ex \|\wh{\Sigma}-\Sigma\|_F$, and \Cref{thm:stab} shows that the error of $\hat{\theta}_{\multiPPI(\widehat{\Sigma})}$ is controlled by $\|\wh{\Sigma}-\Sigma\|_{F}$. In \Cref{thm:ledoit}, we apply \Cref{thm:stab} to provide finite-sample performance guarantees on MultiPPI when the Ledoit-Wolf estimator is used to estimate covariance.\looseness=-1

In our experiments, we evaluate $\hat{\theta}_{\ul{n},\ul{\lambda}}$ on the same samples we used to estimate $\wh{\Sigma}$. A similar approach was taken by \citet{angelopoulos2023ppi++} for PPI++ when tuning the value of $\lambda$, and we find that it is easy to implement and yields strong empirical results in practice.

Note that while the data reuse introduces bias in finite samples---due in part to the additional dependency of $\lambda_I$ on $X_I$ in \Cref{eq:theta-n-lambda}---it preserves consistency and asymptotic normality in the limit as our budget $B$ and the number of (reused) burn-in samples tend to infinity. For an analysis of the bias introduced in finite samples, see \Cref{sec:coverage-decay}.\looseness=-1
%\adam{we say something further about this in the appendix, point to it?}
% In practice, we evaluate $\hat{\theta}_{\ul{n},\ul{\lambda}}$ on the same samples we used to estimate $\wh{\Sigma}$. A similar approach was taken by \citet{angelopoulos2023ppi++} for PPI++, and we find that it is easy to implement and yields strong empirical results in practice. While doing so introduces bias in finite samples---due in part to the additional dependency of $\lambda_I$ on $X_I$ in \Cref{eq:theta-n-lambda}---it preserves consistency and asymptotic normality in the limit as our budget $B$ and the number of (reused) burn-in  samples tend to infinity, as we see next. For this result, we need a technical condition which amounts to \Cref{eq:multippi_opt} having a unique minimizer $\ul{n}$ as $B\to\infty$; we state it formally in \Cref{sec:asymptotic_proof}. \looseness=-1

% \begin{theorem}\label{thm:asymptotic_revised}
% Suppose $X\in\R^k$ has finite second moment, and suppose that $\Sigma=\Cov(X)$ satisfies condition \ref{eqn:condition}. Suppose that $\wh{\Sigma}\overset{p}{\to}\Sigma$ in the operator norm as $B\to\infty$. Then for $\hat{\theta}_{\multiPPI(\widehat{\Sigma})}$ arbitrarily dependent on any potential samples used to estimate $\widehat{\Sigma}$, we have
%     \[
%     \sqrt{B}\left(\hat{\theta}_{\multiPPI(\widehat{\Sigma})}-\theta^*\right)\overset{d}{\to}\mathcal{N}(0,\mathcal{V}^*)
%     \]
%     as $B\to\infty$, where $\mathcal{V}^*=\lim_{B\to\infty} B\mathcal{V}_{B}$, and $\mathcal{V}_{B}$ is defined in \Cref{eq:variance}.
% \end{theorem}
%It is important to note that the estimator $\hat{\theta}_{\MultiPPI(\wh{\Sigma})}$ continues to enjoy unbiasedness, budget satisfaction, and asymptotic normality regardless of mis-specification in $\wh{\Sigma}$. 

%\charlie{TODO: Emphasize that this holds for any estimate $\wh{\Sigma}_N$ of $\Sigma$.}

%It is important to note that the estimator $\hat{\theta}_{\MultiPPI(\wh{\Sigma})}$ continues to enjoy unbiasedness, budget satisfaction, and asymptotic normality (shown below) regardless of mis-specification in $\wh{\Sigma}$. However, better estimates $\wh{\Sigma}$ of $\Sigma$ will lead to better performance of the estimator.

% The following theorem shows that, for any $\wh{\Sigma}$ which converges in probability to $\Sigma$ as our budget tends to infinity, the MultiPPI estimator is asymptotically normal and achieves the optimal variance of \Cref{thm:minimax}. 

%$\wh{\Sigma}$ may be estimated from data---this is what we do in our experiments---or inherited 

%Here, we quantify the stability of our procedure to mis-specification in $\wh{\Sigma}$.


\looseness=-1
% \begin{theorem}\label{thm:finite}
%     Let $P$ be a distribution of covariance $\Sigma$, and suppose $|X_i|\leq1$ almost surely, for each $i$. Suppose that $\Sigma$ has minimum and maximum eigenvalues $\lambda_{\min}$ and $\lambda_{\max}$, respectively. Let $\sigma^2_{\classical}$ denote the least MSE of any budget-satisfying sample mean of $a^\top X$. Let $\widehat{\Sigma}$ denote the empirical covariance matrix determined by $N$ i.i.d. copies of $X$. Then for an absolute constant $c$,
%     \[
%     \Ex\left[\left(\hat{\theta}_{\multiPPI(\widehat{\Sigma})}-\theta^*\right)^2\right] \leq \mathcal{V}_B + c\left(\frac{\lambda_{\max}^{1/2}}{\lambda_{\min}} \sqrt{\frac{k}{N}} + \frac{1}{\lambda_{\min}}\frac{k}{N}\right)\sigma^2_{\classical}
%     \]
%     where the randomness is taken over both the evaluation of $\hat{\theta}_{\multiPPI(\wh{\Sigma})}$ and the $N$ burn-in samples.
% \end{theorem}

% \begin{theorem}[Stability of MultiPPI]\label{thm:stab}
%     Let $P$ be a distribution of covariance $\Sigma$, and suppose that $\Sigma$ has minimum eigenvalue $\gamma_{\min}$. Let $\sigma^2_{\classical}$ denote the least MSE of any budget-satisfying sample mean of $a^\top X$. Let $\widehat{\Sigma}$ denote any symmetric positive-definite matrix. Then we have
%     \[
%     \Ex\left[\left(\hat{\theta}_{\multiPPI(\widehat{\Sigma})}-\theta^*\right)^2\right] \leq \mathcal{V}_B + \frac{4\sigma^2_{\classical}}{\gamma_{\min}}\|\wh{\Sigma}-\Sigma\|_{op}.
%     \]
%     whenever $\|\wh{\Sigma}-\Sigma\|_{op}\leq \gamma_{\min}/2$.
% \end{theorem}

% A natural approach is to estimate $\Sigma$ from data; this is what we do in our experiments. Specifically, we use an empirical estimate $\widehat{\Sigma}$ of the covariance matrix based on some number of \textit{fully-labeled samples} for which the entire joint vector $X=(X_1,\dots,X_k)$ is observed. Given $N$ such samples, we provide a finite-sample bound on the MSE of our estimator in this setting.

% In the appendix, we apply \Cref{thm:stab} in a variety of realistic settings to derive finite-sammple bounds on the performance  

% \Cref{thm:finite} holds for bounded $X$ (and subgaussian $X$; see \Cref{app:finite-sample-bounds}). In \Cref{app:finite-sample-bounds} we also provide even stronger results for settings in which $X$ has certain distributional characteristics. \charlie{TODO: Emphasize here that this is a meta theorem that holds for any estimate $\wh{\Sigma}_N$ of $\Sigma$, say that we prove a similar result for shrinkage in the appendix, for reviewer bTBc. Emphasize that, in particular, this makes a direct statement about covariance mis-specification.}\looseness=-1



% \Cref{thm:finite} holds for general bounded $X$ (and for subgaussian $X$, as we see in \Cref{sec:finite}) but it becomes significantly more powerful in the distributional settings relevant to our study here. For example, test-time compute and model distillation represent settings in which it is natural to assume that each of the models $X_1,\dots,X_k$ is a noisy approximation of its predecessor. Such settings motivate the statistical \textit{autoregressive model} of order 1, in which we model $X_{t+1}=\rho X_t+\epsilon_t$, where $\epsilon_t$ is a white noise process with zero mean and constant variance. In this model, the $\lambda_{\min}$ and $\lambda_{\max}$ defined in \Cref{thm:finite} are of constant order (see \Cref{cor:gersch}), and we may conclude that
% \begin{equation}\label{eq:lim}
% \Ex\left[\left(\hat{\theta}_{\text{MultiPPI}}-\theta^*\right)^2\right] = \mathcal{V}+o(1)
% \end{equation}
% in any limit $N,k\to\infty$ such that $N\gg k$. This result provides a natural heuristic for the burn-in size needed in the case of an autoregressive model: the burn-in size should be large compared to the number of models. While this may seem like a relatively specific model, \Cref{cor:gersch} shows that its conclusion holds in relative generality provided sufficiently weak long-range dependencies between our models.

%\subsection{MultiPPI($\widehat{\Sigma}$): A practical algorithm}\label{sec:practical}

% The optimization procedure described above (and to which the finite-sample bounds of \Cref{thm:finite} apply) relies on an estimate $\widehat{\Sigma}$ that involves data splitting: the best parameters ${\ul{n},\ul{\lambda}}$ for $\hat{\theta}_{\ul{n},\ul{\lambda}}$ are \textit{estimated} on a held-out dataset, and then $\hat{\theta}_{\ul{n},\ul{\lambda}}$ is \textit{evaluated} on a distinct dataset.
% In principle, however, we can also evaluate $\hat{\theta}_{\ul{n},\ul{\lambda}}$ on the same samples we used to estimate $\wh{\Sigma}$. A similar approach was taken by \citet{angelopoulos2023ppi++} for PPI++, and we find that it is easy to implement and yields strong empirical results in practice. While doing so introduces bias in finite samples---due in part to the additional dependency of $\lambda_I$ on $X_I$ in \Cref{eq:theta-n-lambda}---it preserves consistency and asymptotic normality in the limit as our budget $B$ and the number of (reused) burn-in  samples tend to infinity, as we see next. For this result, we need a technical condition which amounts to \Cref{eq:multippi_opt} having a unique minimizer $\ul{n}$ as $B\to\infty$; we state it formally in \Cref{sec:asymptotic_proof}. \looseness=-1





% % 

% %We consider a limit in which our budget $B$ tends to infinity, and  $\widehat{\Sigma}\overset{p}{\to} \Sigma$ in the operator norm; under mild moment conditions, this may be achieved by letting $\widehat{\Sigma}$ be the empirical covariance estimated on a subset of the data~\cite{wainwright2019high}, as we do here. Importantly, if $\hat{\theta}_{\multiPPI(\widehat{\Sigma})}$ is evaluated on the same samples used to construct $\widehat{\Sigma}$, we still have asymptotic normality:\looseness=-1
% \begin{theorem}\label{thm:asymptotic}
% Suppose $X$ has finite second moment, and let $\widehat{\Sigma}$ denote the empirical covariance of $X$ estimated from $N$ i.i.d. samples. Assume that $\Sigma$ satisfies condition \ref{eqn:condition}. Then for $\hat{\theta}_{\multiPPI(\widehat{\Sigma})}$ arbitrarily dependent on the $N$ samples used to estimate $\widehat{\Sigma}$, we have
%     \[
%     \sqrt{B}\left(\hat{\theta}_{\multiPPI(\widehat{\Sigma})}-\theta^*\right)\overset{d}{\to}\mathcal{N}(0,\mathcal{V}^*)
%     \]
%     in any limit in which $B,N\to\infty$, where $\mathcal{V}^*=\lim_{B\to\infty} B\mathcal{V}_{B}$, and $\mathcal{V}_{B}$ is defined in \Cref{eq:variance}.
% \end{theorem}
% \charlie{TODO: Emphasize that this holds for any estimate $\wh{\Sigma}_N$ of $\Sigma$.}

\subsubsection{Procedure}
\label{sec:procedure}
We now specify an easy-to-implement procedure that makes use of a burn-in of $N$ fully labeled samples to estimate $\widehat{\Sigma}$, and then also reuses the $N$ samples when estimating $\hat{\theta}_{\text{MultiPPI}(\widehat{\Sigma})}$. Specifically, we target the practical setting where we are given $N$ fully-labeled samples  \textit{a priori}, and have no ability to obtain more. This is typical of real-world settings in which we are given, or have already collected, a fixed dataset of "gold" labels that we are then trying to augment with PPI related techniques---and may be encapsulated by the budget constraint $n_{\{1, \ldots, k\}}\leq N$. While we may not be able to obtain more fully-labeled samples, we may be afforded a separate computational budget for querying model predictions that then augment the $N$ fully-labeled samples; taken together, this setting is represented by a system of budget constraints.\footnote{We explain how to solve the optimization problem posed by such systems in \autoref{app:generalized}.} In summary, we propose the following:\looseness=-1
% In this setting, we propose the following procedure:
\begin{enumerate}[leftmargin=*]
    \item Estimate the covariance matrix $\widehat{\Sigma}$ on the $N$ fully-labeled samples, which we will reuse in Step 3.
    \item Solve for the $n_I,\lambda_I$ which minimize \Cref{eq:multippi_opt} given $\wh{\Sigma}$. We refer to this as MultiAllocate$(\widehat{\Sigma})$. %\footnote{We note that we leverage the ability to reuse the $N$ fully-labeled samples by including it as the set $I = \{1, \ldots, k\}$, with the additional budget constraint that $n_{I} \leq N$. We explain how to solve the optimization problem posed by such systems with \emph{multiple} types of budget constraints in \Cref{app:generalized}.}
    \item Sample the $n_I$, $\forall I \in \mathcal{I}$ additional data points accordingly, and return $\hat{\theta}_{\text{MultiPPI}(\widehat{\Sigma})}$, evaluated on both the $n_I$ additional data points for each $I\in\mathcal{I}$, and the $N$ (reused) initial samples.
\end{enumerate}
%The form of $\mathcal{V}$ specified in \autoref{eq:variance} provides an efficient recipe for solving the MultiAllocate optimization problem, as we explore in \Cref{sec:computation}. \charlie{TODO: comparison to semiparametric efficient estimator.}
%\adam{Remove and put in appendix, or respond in rebuttal}

% \alekh{A natural question at this point would be if $V$ is the optimal variance, making this a semiparametrically efficient estimator? Even if we do not know, we should at least say that as this would definitely draw a question from any theory reviewer.} \charlie{This is a great point! If I'm understanding correctly, I believe the answer is that $V$ is not optimal. I believe that the semiparametric efficient estimator will also involve higher moments of $X_i$. However, I think, but I'm not sure, that it could be constructed (for fixed $k$) from the MultiPPI estimator by expanding the covariance, sufficiently slowly (my intuition is at rate $o(B^{1/2})$) with a bounded basis for the continuous functions on $X$, provided regularity on the regression function. This is my intuition based on e.g. https://projecteuclid.org/journals/annals-of-statistics/volume-19/issue-3/Efficient-Estimation-of-Linear-Functionals-of-a-Probability-Measure-P/10.1214/aos/1176348251.full and https://arxiv.org/pdf/1606.07268, but I'm not sure.}


%The strength of the theorem is in handling the case in which $\widehat{\Sigma}$ is not independent of some of the sampled variables, as in the approach described in \autoref{sec:practical}.



%. To make concrete how this may occur in practice, consider the situation explored in \autoref{sec:exp-setup}: $\widehat{\Sigma}$ depends on some number $\texttt{num\_labels}$ of labeled samples $X^{(1)},\dots,X^{(\texttt{num\_labels})}$ which we have access to, and we may query no more. However, we use $\widehat{\Sigma}$ to determine our allocation for our budget $B$. This formally corresponds to a problem in which the budget constraints are $n_{[k]}\leq \texttt{num\_labels}$ and $\sum_I c_In_I\leq B$.


% \begin{center}
% \fbox{%
%   \begin{minipage}{1.0\textwidth}
%     \paragraph{Subroutine 1:} 
%     Multi-allocate$(\Sigma)$
%     \paragraph{Input:} a symmetric positive-define matrix $\Sigma\in\R^{k\times k}$. Implicitly given $\mu$, $\mathcal{I}$, $(c_I)_I$, $C$.
%     \paragraph{Output:} the solution $(n_I,\lambda_I)_{I\in\mathcal{I}}$ to the problem
%     \begin{align*}
%         \min_{(n_I,\lambda_I)_{I\in\mathcal{I}}} &\sum_{I\in\mathcal{I}} \frac{1}{n_I}\lambda_I^\top \Sigma_I \lambda_I \\
%         \text{s.t.} & \text{
%     \textbf{U} and \textbf{B} hold} 
%     \end{align*}
%     where $\Sigma_I=(\Sigma_{ij})_{i,j\in I}$.
%   \end{minipage}
% }\label{fact}
% \end{center}

% In what follows, we propose a simple prescriptive approach for a common setting. Suppose our goal is to estimate the mean of $X_1$, and we have access to a finite number of samples $\texttt{num\_labels}$ from $X_1$ but can query no more. We suppose that the remaining indices $\{2,\dots,k\}:=I_\text{models}$ index our set of autoraters. We may query each autorater $X_i$ individually for cost $c_i<\infty$, or all autoraters jointly for cost $c_{I_\text{models}}<\infty$.

% In this setting, we propose the following procedure:
% \begin{outline}
%     \1 Autorate each of the $\texttt{num\_labels}$ labeled examples with each autorater, $X_2,\dots,X_k$ and form the empirical covariance matrix $\widehat{\Sigma}$.
%     \1 Construct the estimator $\hat{\theta}_{\text{MultiPPI}(\widehat{\Sigma})}$ subject to the (system of) budget conditions that
%     \2 $n_{\{1,\dots,k\}}\leq \texttt{num\_labels}$, and
%     \2 $c_{I_\text{models}}n_{I_\text{models}} +\sum_{i\in I_\text{models}} n_ic_i\leq B$.
% \end{outline}
% This procedure is not optimal, but we believe that it suffices as a rule of thumb. Note that the limitation to the prescribed set of index collection means that this set grows linearly in $k$, rather than exponentially. In particular, this choice of cost profile is minimal such that $\hat{\theta}_\text{MultiPPI}$ improves on both 

% Instead, we make a simple prescriptive approach here: Suppose that we are given some number of labeled samples \textit{a priori}, and have no ability to query any more. 

% This is quite common, and corresponds to a linear constraint of the form $\sum_{I\ni 0} n_I\leq \texttt{num\_labels}$. We propose to label all labeled samples with all autoraters. Once this has been done, we obtain an estimate $\widehat{\Sigma}$, and solve the corresponding problem. This procedure is likely not optimal, but we believe that it suffices as a rule of thumb.

% Note that this yields asymptotic normality. It is this procedure that we study in all experiments. The budget profile we set up is therefore
% \[
% n_{[k]}\leq \texttt{num\_labels},\quad \quad \sum_{I}
% \]
% This is the minimal comparison to existing methods.

% \subsection{Relationship to previous estimators}
% Several previous estimators proposed can be understood as special cases of $\hat{\theta}_{\text{MultiPPI}(\Sigma)}$. For example, the PPI++ estimator (Equation \ref{eqn:ppi++}) can be cast as a MultiPPI estimator. Optimizing the sample sizes $n$ and $N$ in PPI++ is equivalent to solving the MultiPPI problem with $k=2$, target vector $a=(1,0)$, and finite costs given only to the subsets $I_1=\{1,2\}$ and $I_2=\{2\}$.\amirg{do you mean equal costs?}

% A natural generalization of PPI++ to multiple models $(X_2, \dots, X_k)$ involves stacking these into a vector and using a vector-valued coefficient $\lambda \in \R^{k-1}$. This corresponds to MultiPPI with $a=(1,0,\dots,0)$, and finite costs given only to the subsets $I_1=\{1,\dots,k\}$ and $I_2 = \{2,\dots,k\}$. Note that this choice of $I_2$ is inherently suboptimal in the generic case in which $X_1,\dots,X_k$ are positively correlated and costs are additive: it would reduce variance to instead replace $I_2$ by the collection of singletons $I'_2=\{2\},\dots,I'_k=\{k\}$.\amirg{maybe explain this more. Do you mean if the cost of $I_2$ is the sum of the individual costs? Also, do you mean the sets are \{1,i\}?
% Would be good to provide some more details here.}

% For our experiments, we enrich the collection of finite-cost subsets beyond this basic vector generalization. We include $I_1=\{1, \dots, k\}$ (fully labeled data), $I_2=\{2, \dots, k\}$ (all models), and also each individual model $I'_i=\{i\}$ for $i \in \{2, \dots, k\}$. This configuration ensures that MultiPPI can perform at least as well as using PPI++ with any single model -- often optimal in low-budget scenarios -- and at least as well as the vectorized PPI++ in high-budget regimes. Furthermore, this flexible setup allows MultiPPI to automatically discover more complex strategies, such as \textit{cascaded PPI}, where cheaper models are used to reduce the need for querying more expensive ones, which in turn are strong predictors of the true label $X_1$.






\section{Experimental setup}\label{sec:exp-setup}
% In each experiment, our goal is to estimate the mean $\theta^*=\mathbb{E}[X_1]$ of some random random variable $X_1$ to be specified. We will refer to $\theta^*$ as our \textit{target}. We will also specify a family of \textit{models} $(X_2,\dots,X_k)$, together with their \textit{cost structure}.


% For each experiment, we will specify a \textit{target}, and a family of \textit{models}, chosen together with their \textit{cost structure}.\amirg{need clearer mapping to our setting. What is $X_1$, what's $a$ etc.} Broadly speaking, a \textit{model} is any map $f:\mathcal{X}\to\mathcal{Y}$. In our first experiment, the models are chosen conventionally: they are simply two off-the-shelf language models, Gemini 2.5 Pro and Gemini 2.5 Flash. In subsequent experiments, we explore how they may be chosen more creatively, e.g. with $f$ being an ensemble of simpler models. In this section, we detail the model constructions and their associated cost structure.\alekh{Do we ever spell out the exact strategy used for estimating the covariance in the experiments, and do we account for any burn-in samples used in the plots?}\charlie{Good point, will do, thanks.}



% In each experiment, we consider a setting in which we are given some number of samples for which the entire vector $X=(X_1,\dots,X_k)$ is visible; we refer to such samples as \textit{fully-labeled}. 


% \begin{enumerate}
%     \item Observe 100 \textit{fully-labeled} samples, and estimate $\widehat{\Sigma}$ accordingly.
%     \item Solve for the parameters $n_I,\lambda_I$ which solve MultiAllocate$(\widehat{\Sigma})$.
%     \item Sample $X_I$ accordingly, and evaluate $\hat{\theta}_{\text{MultiPPI}}$ on both the fresh samples and the 100 fully-labeled samples. 
% \end{enumerate}
% \charlie{TODO: complete}

In each experiment, our goal is to estimate the mean $\theta^*=\mathbb{E}[X_1]$ of some random variable $X_1$ to be specified, which we will refer to as the \textit{target}. This corresponds to the choice $a=(1,0,\dots,0)$ in our notation. We will also specify a \textit{model family} $(X_2,\dots,X_k)$, together with a \textit{cost structure} $(c_I)_{I\in\mathcal{I}}$.
%
In each experiment, we are given some number of samples for which the entire vector $X=(X_1,\dots,X_k)$ is visible; we refer to such samples as \textit{fully-labeled}. Given these samples, we perform the procedure outlined in \Cref{sec:procedure}: we estimate $\widehat{\Sigma}$ using these samples, sample from the auxiliary models $(X_2, \ldots, X_k)$ according to the allocation  specified by MultiAllocate$(\widehat{\Sigma})$, and return $\hat{\theta}_{\multiPPI(\widehat{\Sigma})}$, evaluated on both the $N$ fully-labeled samples and the additional auxiliary data.\looseness=-1

\textbf{Baselines:} In each experiment, we compare to several baselines. First, we compare to classical sampling. Second, we compare to PPI++ with each model included in the family (specified in \Cref{eqn:ppi++}), and to vector PPI++ with every model in the family (specified in \Cref{eqn:ppi++vector}). %TODO: additional comparisons to (i) higher order features, (ii) resampling randomly, and (iii) active procedures.

\textbf{Experiment 1: Estimating Arena win-rates by autorater ensembles.}
We focus on the Chatbot Arena dataset~\citep{chiang2024chatbot}, where of interest is the \emph{win-rate} between a pair of models, which is the probability that a given user prefers the response of one model to that of the other. The randomness is taken over the prompt, the user, and the model responses. Here, we aim to estimate the win-rate between Claude-2.1 and GPT-4-1106-Preview; this is our \emph{target}. Our \emph{model family} consists of autoraters built on Gemini 2.5 Pro (without thinking) and Gemini 2.5 Flash. In our notation, we have $(X_1,X_2,X_3)=(\textrm{human label}, \textrm{Gemini 2.5 Pro label}, \textrm{Gemini 2.5 Flash label})$. We draw model \emph{costs} from the Gemini developer API pricing guide \citep{gemini}, see \Cref{sec:experimental-details}. In this case, the cost of querying both models is simply the sum of the costs of querying each model independently. %, but we emphasize that this will not be the case in the subsequent experiments.\looseness=-1
% \begin{description}
%     \item[Target:] Win-rate between Claude-2.1 and GPT-4-1106-Preview on ChatBot Arena
%     \item[Model family:] $\{\text{Gemini 2.5 Pro, Gemini2.5 Flash}\}$
%     \item[Cost structure:] $c_\text{Pro}=1.25,\,c_\text{Flash}=0.30,\,c_{\{\text{Pro, Flash}\}}=1.55$
% \end{description}

% \textbf{Target:} Win-rate between Claude-2.1 and GPT-4-1106-Preview on ChatBot Arena. \textbf{Model family:} $\{\text{Gemini 2.5 Pro, Gemini2.5 Flash}\}$. \textbf{Costs:} $c_\text{Pro}=1.25,\,c_\text{Flash}=0.30,\,c_{\{\text{Pro, Flash}\}}=1.55$.

% \begin{table}[]
%     \centering
%     \begin{tabular}{c|c}
%         Model collection & Cost \\
%         \hline
%         Gemini 2.5 Pro & \$1.25 \\
%         Gemini 2.5 Flash & \$0.30 \\
%         Both & \$1.55
%     \end{tabular}
%     \caption{Cost structure for experiment 1 \jacob{we don't need a table for this}}
%     \label{tab:cost-exp1}
% \end{table}
\textbf{Experiment 2: Optimal test-time autorater scaling on ProcessBench.}
In this experiment, we aim to estimate the fraction of correct solutions in the ProcessBench dataset \citep{zheng2024processbench}, given a small number of labeled examples. The task is simplified from its original form to a binary classification problem: determining whether a given math proof solution contains a process error, without identifying the specific step. We employ Gemini 2.5 Pro with a variable thinking budget as our autorater. Its accuracy correlates with the number of words expended in the thought, with performance gains saturating after approximately 500 words (see \Cref{fig:scaling-budget} in the appendix). We create a family of four autoraters by checkpointing the model's thought process at 125, 250, 375, and 500 words. A key aspect of this setup is the non-additive, cascading cost structure. Generating a response from a model with a larger thinking budget makes the outputs of all smaller-budget models available at a marginal cost. Consequently, the total cost for a subset of models S is modeled with two components: an input cost proportional to the sum of the word budgets in S, and an output cost proportional to the maximum word budget in S. Explicitly, for $S\subseteq\{125,250,375,500\},$ we set
    \begin{equation}\label{eq:cost}
    c_S=\texttt{output\_cost\_per\_word}\cdot \max S+\texttt{input\_cost\_per\_word}\cdot\sum S    
    \end{equation} 
For concision, in our results we refer to the model assessments after 125, 250, 375, and 500 words by ``tiny,'' ``small,'' ``medium,'' and ``large,'' respectively.

% In the ProcessBench dataset, the task is to determine whether or not a process error occurred in an attempted solution to a math proof, and to identify the step at which the process error occurred if there was one \citep{zheng2024processbench}. The dataset consists of attempted solutions, paired with labels (process error or not, and if yes, in which step). To make this into an evaluation problem, we disregard the step at which the process error may have occured, and consider only a binary label (process error or not) for every attempted solution.

% Here, we consider the problem of estimating the fraction of correct solutions---our \textit{target}---given only a small number of binary labels $n$. We employ Gemini 2.5 Pro with a variable thinking budget as our autorater. An increased thinking budget leads to higher accuracy at label prediction~\citep{zheng2024processbench}, but at higher cost. We find that most of the improvement occurs within the first 500 thought tokens (\autoref{fig:scaling-budget}). We therefore build a family of autoraters by checkpointing the thought process at 4 points before this saturation point: we interrupt the thought process after 125 (``tiny''), 250 (``small''), 375 (``medium''), and 500 (``large'') words of thought, respectively. %While this choice is convenient for the purpose of exposition, it is not essential to specify \textit{a priori}: in further experiments in the appendix \jacob{ref}\charlie{TODO}, we demonstrate that by including models corresponding to numbers of words of thought past the saturation point, our procedure learns to ignore such models.

% Crucially, the cost structure is not additive: if Gemini 2.5 Pro has already thought for 500 words, then we can obtain its answer after 250 words of thought for a significantly reduced price. We capture this by modeling total cost of a subset $S$ as involving an \textit{input cost}, proportional to the sum of all word budgets in $S$, and an \textit{output cost}, proportional to the maximum word budget in $S$. This enables a cascade in which we may easily rectify the bias incurred by under- or over-thinking for the same price as the maximum amount of thinking performed. The resulting cost structure is:

% \begin{description}
%     \item[Target:] Fraction of correct solutions in dataset
%     \item[Model family:] $\{\text{Demand answer from Gemini 2.5 Pro after $T$ words}:T\in\{125,250,375,500\}\}$
%     \item[Cost structure:] For $S\subseteq\{125,250,375,500\},$ 
%     \begin{equation}\label{eq:cost}
%     c_S=\texttt{output\_cost\_per\_word}\cdot \max S+\texttt{input\_cost\_per\_word}\cdot\sum S    
%     \end{equation}
% \end{description}

% \textbf{Target:} Fraction of correct solutions in dataset. \textbf{Model family:} $\{\text{Demand answer from Gemini 2.5 Pro after $T$ words}:T\in\{125,250,375,500\}\}$. \textbf{Costs:} For $S\subseteq\{125,250,375,500\},$ 
%     \begin{equation}\label{eq:cost}
%     c_S=\texttt{output\_cost\_per\_word}\cdot \max S+\texttt{input\_cost\_per\_word}\cdot\sum S    
%     \end{equation} 
%Note for later about method: Multi-PPI is always as good as just doing PPI several times. So in particular, it's always as good as you can do by modeling each model using the next one. And this is best when they're sequentially correlated.
\textbf{Experiment 3: Hybrid factuality evaluation through multi-autorater debate.}
%\charlie{Would be great if we could condense this subsection, in particular.}
Following \citet{du2023improving}, we evaluate the factual consistency of biographies for 524 computer scientists generated by Gemini 2.5 Pro. For each person $p\in\mathcal{P}$, we compare their Gemini-generated biography $b^p$ against a set of known grounding facts $\mathcal{F}^p=\{f_1^p,\dots,f_{m_p}^p\}$ about the person. 
%
Our target metric is the proportion of \textit{factually consistent pairs} $(b,f)$ within the total set $\mathcal{S}=\{(b^p,f^p): p\in\mathcal{P},f^p\in\mathcal{F}^p\}$. Concretely, we \textit{target} the proportion
$
|\{(b,f)\in\mathcal{S}:(b,f)\text{ is factually consistent}\}| ~/~ |\mathcal{S}|
$.

Ground-truth consistency of a pair $(b,f)$ is established by majority voting over five independent judgments from Gemini 2.5 Pro with thinking, a method validated by \citet{du2023improving} to have over 95\% agreement with human annotators. Our experiment, illustrated in \Cref{fig:scaling-accuracy-biographies}, assesses the performance of a more cost-effective model, Gemini 2.0 Flash Lite, as an autorater.
%
To elicit better autoratings from queries to Gemini 2.0 Flash Lite, we bootstrap performance via multi-round debate. For a fixed number of agents $A\in\{1,2,3\}$, and a fixed number of maximum rounds $R\in\{1,2\}$, we perform the following procedure: In each round, $A$ instances of Flash Lite are independently prompted to provide a reasoned judgment on the consistency of a pair $(b,f)\in\mathcal{S}$. A "pooler" instance of Flash Lite then consolidates these responses into a single \emph{yes}, \emph{no}, or \emph{uncertain} output. A definitive \emph{yes} or \emph{no} concludes the process. If the pooler outputs \emph{uncertain}, and the number of maximum rounds $R$ has not yet been reached, the $A$ agents review all prior responses and continue their debate in a new round. If the output remains \emph{uncertain} after the final round, either \emph{yes} or \emph{no} is reported with equal probability---since the dataset is balanced, this outcome is fair insofar as it is as good as random guessing. We impose the maximum round restriction to encapsulate our budget constraint. For a given $(A,R)$, the cost is $A\cdot R$; for collections, the cost follows \Cref{eq:cost}.\looseness=-1

% \begin{description}
%     \item[Target:] Proportion of factually-consistent pairs, $\#\{(b,f)\in\mathcal{S}:(b,f)\text{ is factually consistent}\}/\#\mathcal{S}$
%     \item[Model family:] $\{\text{The output of the above procedure given $(A,R)$}: A\in\{1,2,3\},R\in\{1,2\}\}$
%     \item[Cost structure:] For a given $(A,R)$, the cost is $
% A\cdot R$; for collections, the cost follows \autoref{eq:cost}.
% \end{description}

% \begin{figure}
%     \centering
%     \includegraphics[width=0.9\linewidth]{figures/Factuality task2.png}
%     \caption{Depiction of biography-fact pairs $(b,f)$ as in Experiment 3. Judgements about factual consistency of $(b,f)$ are made by a language model.}
%     \label{fig:biography-prompt}
% \end{figure}

\section{Empirical results}\label{sec:results}
%Please see appendix for additional plots. \textbf{TODO:} change color scheme, add additional comparisons for ChatBot Arena, run more trials to reduce variance, bootstrap CIs. \textbf{Other:} please suggest!
We plot MultiPPI, and the baselines described in \Cref{sec:exp-setup}, for a budgets between 0 and 2$k$ units of cost. We normalize model costs so that one unit of cost always represents exactly one query to our most expensive model. 
For each fixed budget and each method, we estimate the target, and construct asymptotic 95\% confidence intervals $\mathcal{C}$ based on \Cref{thm:asymptotic_revised}. We plot (i) coverage, $\mathbb{P}(\theta^*\in\mathcal{C})$; (ii) confidence interval width, $|\mathcal{C}|$; and (iii) mean-squared error $\mathbb{E}[(\hat{\theta}-\theta^*)^2]$. We report both the confidence interval width and the mean-squared error as a fraction of what classical sampling achieves (lower is better). In each case, the target is $\theta^*=\mathbb{E}[X_1]$, and $\mathbb{P}$ and $\mathbb{E}$ are computed with respect to the empirical distribution over the dataset observed (we perform 500$k$ random trials with 250 given labels). Note that these 250 labeled points are evidently enough for all estimators considered to achieve good coverage (in \Cref{sec:experiments-num-labeled} we also include results with 1$k$ labeled points). We implement the optimization scheme in \texttt{cvxpy}, and use CVXOPT as our choice of optimizer.\looseness=-1



\textbf{Experiment 1: Chatbot Arena.}  Results are shown in \Cref{fig:results-all} (top). Observe that different baselines dominate in different budget regimes. In the low-budget regime, scalar PPI++ with Gemini 2.5 Flash is the best baseline, while in the large-budget regime, vector PPI++ with both Gemini 2.5 Pro and Gemini 2.5 Flash is the best baseline.
However, we see that MultiPPI improves on all baselines in all regimes. In the appendix, \Cref{fig:lambdas-chatbotarena} and \Cref{fig:allocations-chatbotarena} plot the $\lambda_I$ and $n_I$ values learned by MultiPPI across budget regimes. Note that the learned values tend to the specifications for PPI++ with Gemini 2.5 Flash in the low-budget regime, and to the specifications for vector PPI++ in the large-budget regime, a finding that we rigorously prove happens in broader generality in \Cref{sec:limiting}. Lastly, note that PPI++ with Gemini 2.5 Pro is suboptimal in all regimes. In other words, PPI++ with Gemini 2.5 Pro is not included in the Pareto frontier. This is because, for this task, its correlation with the label is no more than that of PPI++ with Gemini 2.5 Flash, yet it is strictly more expensive.\looseness=-1 %Once again, we prove that this behavior occurs in greater generality in \Cref{sec:pareto}.



% \subsection{Experiment 2: ProcessBench}
% \begin{figure}
%     \centering
%     % \includegraphics[width=\linewidth]{figures/processbench_results2.png}
%     \input{figures/processbench_figure}
%     \caption{Results by budget, Experiment 2: ProcessBench.}
%     \label{fig:results-processbench}
% \end{figure}

\textbf{Experiment 2: ProcessBench.} Results are shown in \Cref{fig:results-all} (middle). Again, we see that each baseline has  a range of budgets for which it outperforms all other baselines. In particular, the cheaper models yield better performance when used in PPI++ in the smaller-budget regimes, while the more-expensive models yield better performance in the higher-budget regimes. In particular, vector PPI++, which uses all $k - 1$ models, steadily improves as the budget increases, but only outperforms the other baselines at the highest budgets. This behavior is explained by \Cref{fig:scaling-budget} in the appendix, which shows that predictive performance improves for larger thinking budgets. Thus the more expensive models yield higher correlation with the label and thus yield low-variance rectifiers; on the other hand, their high cost means that this decrease in rectifier variance is outweighed by our inability to draw an adequate number of samples from them in the low-budget regimes.\looseness=-1

Of note is the fact that \textit{the models which think for longer are not in general less biased}. This phenomenon is shown in \Cref{fig:need-to-debias}, which shows that thinking for longer is not enough to resolve the systematic bias present in the autorater. However, the figure also shows that simple debiasing schemes like PPI resolve this issue. Note that this trend is not reflected in the correlations between these models and the label, because correlation is invariant to addition of constants.

Finally, MultiPPI improves on all baselines methods in all regimes. Interestingly, \Cref{fig:allocations-processbench} and \Cref{fig:lambdas-processbench} show that the parameters $\lambda_I$ and $n_I$ learned by MultiPPI transition from emulating PPI++ with the tiny model (which is the best baseline in the low-budget regime) to emulating a \textit{cascaded version of PPI} (see \Cref{eqn:cascade}), in which the medium model is used to debias the larger model.

% TODO: Separate sections in experimental setup for
% 1. Model construction
% 2. Cost structure (can include here number of labeled samples)
% 3. Baselines. Maybe baselines can be more broad



\textbf{Experiment 3: Biography factuality evaluation.}
% \begin{figure}
%     \centering
%     %\includegraphics[width=1.0\linewidth]{figures/factuality_results2.png}
%     \input{figures/factuality_figure}
%     \caption{Result by budget, Experiment 3: biography factuality.}
%     \label{fig:results-biographies}
% \end{figure}
Results are shown in \Cref{fig:results-all} (bottom). Once again, each baseline is dominant over the others in certain regimes; MultiPPI improves on all baselines in all regimes. Of note, however, is the fact that the coverage of all estimators considered, but MultiPPI and vector PPI++ in particular, degrades slightly in the large-budget regime (i.e., the 95\% CI under-covers by $\approx 1\%)$. We discuss this interesting phenomenon in  \Cref{sec:coverage-decay}, and find that it does not occur when the number of labeled samples grows in constant proportion with the budget (see, for example, our additional results with $N = 1000$ fully-labeled samples in \Cref{sec:experiments-num-labeled}).

In terms of the performance-vs-cost profile that MultiPPI leverages: \Cref{fig:scaling-accuracy-biographies} shows that predictive performance increases, across many metrics, as the number of agents and number of rounds increases. Note, however, that a marginal increase in number of agents yields a greater increase in accuracy than a marginal increase in number of rounds (this is largely because the pooler is more likely to report "uncertain" after the end of the first round than after the end of the second; see \Cref{fig:scaling-uncertainty-biographies}).\looseness=-1


\begin{figure}[!t]
    \centering
    %\includegraphics[width=\linewidth]{figures/chatbot_arena_results2.png}
    \textrm{\bf \scriptsize Experiment 1: Chatbot Arena}
\scalebox{.8}{\centering
\begin{tikzpicture}
\begin{axis}[
    title={Coverage},
    xlabel={Budget},
    xlabel style={yshift=5pt},   % Moves xlabel closer to the axis
    grid=both,
    outer sep=0pt,
    enlarge x limits=false,    minor grid style={dashed, gray!50},
    major grid style={solid, gray!80},
    width=\textwidth,
    yticklabel style={/pgf/number format/fixed}, % No scientific notation
    yticklabel={\pgfmathparse{\tick*100}\pgfmathprintnumber{\pgfmathresult}\%},
    xmin=0, xmax=5000,                          % Set x-axis range
    font=\scriptsize,             % Sets the base size for everything
    title style={yshift=-5pt, font=\small},    % Optional: Make title slightly bigger than the rest
    label style={font=\scriptsize}, % Font for xlabel and ylabel
    tick label style={font=\scriptsize}, % Font for the numbers on the axes
      xshift=0.0cm,
      width=6.5cm,
      height=5.5cm,
            ymin=0.9, ymax=1,
]

\addplot[color=black, dashed, very thick] coordinates {(0,.95) (5000,.95)};
\addplot[color=blue, thick] coordinates {(250,0.946442) (500,0.946224) (750,0.946858) (1000,0.946918) (1250,0.946958) (1500,0.947248) (1750,0.94647) (2000,0.947084) (2250,0.946102) (2500,0.946536) (2750,0.946618) (3000,0.946862) (3250,0.945972) (3500,0.946684) (3750,0.947166) (4000,0.94704) (4250,0.946958) (4500,0.947234) (4750,0.94648) (5000,0.947006)};
\addlegendentry{MultiPPI}

\addplot[color=green!50!black, thick, dashdotted] coordinates {(250,0.947788) (500,0.947136) (750,0.946664) (1000,0.946666) (1250,0.946944) (1500,0.947214) (1750,0.94811) (2000,0.948416) (2250,0.946338) (2500,0.946002) (2750,0.947484) (3000,0.947148) (3250,0.947352) (3500,0.947386) (3750,0.94667) (4000,0.946496) (4250,0.946034) (4500,0.946354) (4750,0.946236) (5000,0.946988)};
\addlegendentry{PPI{\textnormal{+}}{\textnormal{+}} Vector}

\addplot[color=brown!90, thick, densely dashed] coordinates {(250,0.948656) (500,0.948826) (750,0.948184) (1000,0.948902) (1250,0.947474) (1500,0.948362) (1750,0.94773) (2000,0.94752) (2250,0.947648) (2500,0.947904) (2750,0.947262) (3000,0.9471) (3250,0.948134) (3500,0.948122) (3750,0.948418) (4000,0.947844) (4250,0.947346) (4500,0.947644) (4750,0.947386) (5000,0.947828)};
\addlegendentry{PPI{\textnormal{+}}{\textnormal{+}} Scalar (Pro)}

\addplot[color=brown!60, thick, densely dashed] coordinates {(250,0.947434) (500,0.94749) (750,0.946594) (1000,0.946776) (1250,0.94738) (1500,0.947356) (1750,0.946844) (2000,0.94737) (2250,0.947606) (2500,0.947558) (2750,0.94729) (3000,0.94675) (3250,0.94756) (3500,0.94778) (3750,0.947296) (4000,0.948016) (4250,0.947618) (4500,0.947942) (4750,0.947426) (5000,0.947662)};
\addlegendentry{PPI{\textnormal{+}}{\textnormal{+}} Scalar (Flash)}

\legend{};
\end{axis}
\begin{axis}[
    title={Mean-squared CI width (frac. of classical)},
    xlabel={Budget},
    xlabel style={yshift=5pt},   % Moves xlabel closer to the axis
    grid=both,
    outer sep=0pt,
    enlarge x limits=false,    minor grid style={dashed, gray!50},
    major grid style={solid, gray!80},
    width=\textwidth,
    yticklabel style={/pgf/number format/fixed}, % No scientific notation
    yticklabel={\pgfmathparse{\tick*100}\pgfmathprintnumber{\pgfmathresult}\%},
    xmin=0, xmax=5000,                          % Set x-axis range
    font=\scriptsize,             % Sets the base size for everything
    title style={yshift=-5pt, font=\small},    % Optional: Make title slightly bigger than the rest
    label style={font=\scriptsize}, % Font for xlabel and ylabel
    tick label style={font=\scriptsize}, % Font for the numbers on the axes
      xshift=5.9cm,
      width=6.5cm,
      height=5.5cm,
]

\addplot[color=blue, thick] coordinates {(250,0.9332786678486313) (500,0.924066156392822) (750,0.9192880015348308) (1000,0.9164218331407227) (1250,0.914601019679555) (1500,0.9131166282187997) (1750,0.912249302954298) (2000,0.9113239831110261) (2250,0.9106880371095362) (2500,0.9101785809163747) (2750,0.9096224237680951) (3000,0.9092925870311739) (3250,0.9087647933760736) (3500,0.9085761483028764) (3750,0.9084362376577875) (4000,0.9081541808688721) (4250,0.9079687484563952) (4500,0.9078216379165787) (4750,0.9075877498474475) (5000,0.9075404573839976)};
\addlegendentry{MultiPPI}

\addplot[color=green!50!black, thick, dashdotted] coordinates {(250,0.9613838679182574) (500,0.9454363865919806) (750,0.9362332874552723) (1000,0.9305577133393179) (1250,0.9266377240749157) (1500,0.9238766142893752) (1750,0.9214411775011296) (2000,0.9197619186225694) (2250,0.9181242974218763) (2500,0.916846235028651) (2750,0.9160383757256494) (3000,0.9151459675135678) (3250,0.9144354940201401) (3500,0.9138526103275946) (3750,0.9133015425620248) (4000,0.9127981888753653) (4250,0.9122010786052722) (4500,0.9118150936344519) (4750,0.9115086527972619) (5000,0.9111842343093346)};
\addlegendentry{PPI{\textnormal{+}}{\textnormal{+}} Vector}

\addplot[color=brown!90, thick, densely dashed] coordinates {(250,0.9573944397530033) (500,0.9434304732227339) (750,0.9364306548465037) (1000,0.9321265164222806) (1250,0.9294588471327857) (1500,0.9274056121603059) (1750,0.9258267064328568) (2000,0.924619693358391) (2250,0.9237958183526989) (2500,0.922998664506488) (2750,0.9223486424707397) (3000,0.9218202406910652) (3250,0.921222227181355) (3500,0.9207769408494638) (3750,0.9205208168930091) (4000,0.9201239954108568) (4250,0.9197861360748052) (4500,0.9195919125242314) (4750,0.919444726384083) (5000,0.9192736086468563)};
\addlegendentry{PPI{\textnormal{+}}{\textnormal{+}} Scalar (Pro)}

\addplot[color=brown!60, thick, densely dashed] coordinates {(250,0.9337415127010026) (500,0.9266667916217228) (750,0.9238693613479424) (1000,0.9223725341372523) (1250,0.9215488667035905) (1500,0.9209427437520867) (1750,0.9205305068075634) (2000,0.9201192900681335) (2250,0.9197809680406205) (2500,0.9195720887957706) (2750,0.9194090888329987) (3000,0.9192947838439127) (3250,0.9192021803427878) (3500,0.9191101836242109) (3750,0.9190368604739912) (4000,0.918928832183079) (4250,0.9188502386927782) (4500,0.9187508322621617) (4750,0.9186953536668976) (5000,0.9186612952049926)};
\addlegendentry{PPI{\textnormal{+}}{\textnormal{+}} Scalar (Flash)}

\legend{};
\end{axis}
\begin{axis}[
    title={Mean-squared Error (frac. of classical)},
    xlabel={Budget},
    xlabel style={yshift=5pt},   % Moves xlabel closer to the axis
    grid=both,
    outer sep=0pt,
    enlarge x limits=false,    minor grid style={dashed, gray!50},
    major grid style={solid, gray!80},
    width=\textwidth,
    yticklabel style={/pgf/number format/fixed}, % No scientific notation
    yticklabel={\pgfmathparse{\tick*100}\pgfmathprintnumber{\pgfmathresult}\%},
    xmin=0, xmax=5000,                          % Set x-axis range
    font=\scriptsize,             % Sets the base size for everything
    title style={yshift=-5pt, font=\small},    % Optional: Make title slightly bigger than the rest
    label style={font=\scriptsize}, % Font for xlabel and ylabel
    tick label style={font=\scriptsize}, % Font for the numbers on the axes
      xshift=11.7cm,
      width=6.5cm,
      height=5.5cm,
            legend pos=north east,
            legend cell align={left},
            legend style={font=\tiny,inner sep=1pt,row sep=-2pt},
]

\addplot[color=blue, thick] coordinates {(250,0.9415337353745189) (500,0.9333510753539588) (750,0.9264635446215488) (1000,0.9238424932400724) (1250,0.9220326287689277) (1500,0.9208392833276946) (1750,0.920491404604133) (2000,0.9188830769953523) (2250,0.9210484713710463) (2500,0.9191161171542401) (2750,0.9184462921865898) (3000,0.9169569949811881) (3250,0.9182015744844775) (3500,0.9174623933814354) (3750,0.9170755075974752) (4000,0.9166220389980524) (4250,0.9157667401591891) (4500,0.9152934257141438) (4750,0.9174215990020387) (5000,0.9149228282238807)};
\addlegendentry{MultiPPI}

\addplot[color=green!50!black, thick, dashdotted] coordinates {(250,0.9691687291795068) (500,0.9552789449233803) (750,0.943862393907119) (1000,0.9397525235825022) (1250,0.9354782735333677) (1500,0.9338712094583039) (1750,0.9255704544331537) (2000,0.9224674840833711) (2250,0.9288635591756605) (2500,0.9288094733365517) (2750,0.9227653524808638) (3000,0.922419880924207) (3250,0.9224392275507964) (3500,0.9209008794343542) (3750,0.9225570018955682) (4000,0.9230057788500835) (4250,0.9229316611332712) (4500,0.9217589723804286) (4750,0.9210063140616493) (5000,0.9189422080379199)};
\addlegendentry{PPI{\textnormal{+}}{\textnormal{+}} Vector}

\addplot[color=brown!90, thick, densely dashed] coordinates {(250,0.9585365772303855) (500,0.9458291383636018) (750,0.941045123275949) (1000,0.9348668408508402) (1250,0.9337263339346455) (1500,0.9311420174276922) (1750,0.9317970488372972) (2000,0.9299622949878438) (2250,0.9290298442338046) (2500,0.9271085733516884) (2750,0.9272410167978503) (3000,0.9272536243779175) (3250,0.9247041607592478) (3500,0.9251084135438157) (3750,0.922553717762973) (4000,0.9237983508193314) (4250,0.926352478369087) (4500,0.9248338187615982) (4750,0.9248807336448018) (5000,0.924189657060446)};
\addlegendentry{PPI{\textnormal{+}}{\textnormal{+}} Scalar (Pro)}

\addplot[color=brown!60, thick, densely dashed] coordinates {(250,0.9412839119307392) (500,0.931640999354131) (750,0.9323027901039446) (1000,0.9324674320012433) (1250,0.9285871623973225) (1500,0.9270601781126871) (1750,0.926462642377259) (2000,0.9265295059215264) (2250,0.9247907914744056) (2500,0.9260196418958097) (2750,0.9259669101831602) (3000,0.9265479982815203) (3250,0.9242408766193834) (3500,0.9233601310251223) (3750,0.9251584903180471) (4000,0.9226598066242832) (4250,0.925256553770054) (4500,0.9238577197920242) (4750,0.924160138384959) (5000,0.9250580754211756)};
\addlegendentry{PPI{\textnormal{+}}{\textnormal{+}} Scalar (Flash)}

\end{axis}
\end{tikzpicture}
}
    \textrm{\bf \scriptsize Experiment 2: ProcessBench} 
% https://colab.corp.google.com/drive/1MYZmW9xoa4OXEF8RJJrsfsjL8ONU6Vdl?resourcekey=0-ZPSTOPKAT7JjfECQPDklxg#scrollTo=GKQBldUIRK99&line=4&uniqifier=1
\centering
\scalebox{.8}{\centering
\begin{tikzpicture}
\begin{axis}[
    title={Coverage},
    xlabel={Budget},
    xlabel style={yshift=5pt},   % Moves xlabel closer to the axis
    grid=both,
    outer sep=0pt,
    enlarge x limits=false,    minor grid style={dashed, gray!50},
    major grid style={solid, gray!80},
    width=\textwidth,
    yticklabel style={/pgf/number format/fixed}, % No scientific notation
    yticklabel={\pgfmathparse{\tick*100}\pgfmathprintnumber{\pgfmathresult}\%},
    xmin=0, xmax=5000,                          % Set x-axis range
    font=\scriptsize,             % Sets the base size for everything
    title style={yshift=-5pt, font=\small},    % Optional: Make title slightly bigger than the rest
    label style={font=\scriptsize}, % Font for xlabel and ylabel
    tick label style={font=\scriptsize}, % Font for the numbers on the axes
      xshift=0.0cm,
      width=6.5cm,
      height=5.5cm,
            ymin=0.9, ymax=1,
]

\addplot[color=black, dashed, very thick] coordinates {(0,.95) (5000,.95)};
\addplot[color=gray!10!black!100, thick, densely dashed] coordinates {(250,0.946932) (500,0.94719) (750,0.948278) (1000,0.948212) (1250,0.94842) (1500,0.948454) (1750,0.948098) (2000,0.947658) (2250,0.948316) (2500,0.949274) (2750,0.947802) (3000,0.947512) (3250,0.948522) (3500,0.94842) (3750,0.948552) (4000,0.948464) (4250,0.94753) (4500,0.947918) (4750,0.948018) (5000,0.947966)};
\addlegendentry{PPI{\textnormal{+}}{\textnormal{+}} Scalar Large}

\addplot[color=gray!40!black!70, thick, densely dashed] coordinates {(250,0.947842) (500,0.947696) (750,0.948782) (1000,0.948396) (1250,0.948398) (1500,0.948914) (1750,0.947948) (2000,0.948254) (2250,0.94825) (2500,0.947864) (2750,0.94802) (3000,0.947312) (3250,0.948716) (3500,0.947864) (3750,0.947686) (4000,0.947378) (4250,0.947272) (4500,0.947488) (4750,0.94753) (5000,0.946726)};
\addlegendentry{PPI{\textnormal{+}}{\textnormal{+}} Scalar Medium}

\addplot[color=gray!70!black!40, thick, densely dashed] coordinates {(250,0.949338) (500,0.949142) (750,0.947208) (1000,0.947714) (1250,0.947932) (1500,0.948044) (1750,0.947794) (2000,0.948454) (2250,0.946878) (2500,0.9469) (2750,0.947308) (3000,0.9471) (3250,0.94776) (3500,0.94773) (3750,0.947474) (4000,0.94689) (4250,0.947452) (4500,0.946964) (4750,0.946734) (5000,0.94653)};
\addlegendentry{PPI{\textnormal{+}}{\textnormal{+}} Scalar Small}

\addplot[color=gray!100!black!10, thick, densely dashed] coordinates {(250,0.947956) (500,0.948202) (750,0.948162) (1000,0.947022) (1250,0.946802) (1500,0.946858) (1750,0.947146) (2000,0.946986) (2250,0.94754) (2500,0.94778) (2750,0.947386) (3000,0.947374) (3250,0.946948) (3500,0.947096) (3750,0.947318) (4000,0.947422) (4250,0.947068) (4500,0.947304) (4750,0.947406) (5000,0.947212)};
\addlegendentry{PPI{\textnormal{+}}{\textnormal{+}} Scalar Tiny}

\addplot[color=blue, thick] coordinates {(250,0.945794) (500,0.945372) (750,0.946556) (1000,0.946582) (1250,0.945202) (1500,0.944918) (1750,0.945556) (2000,0.9454) (2250,0.945552) (2500,0.94498) (2750,0.945684) (3000,0.945368) (3250,0.94543) (3500,0.944558) (3750,0.944898) (4000,0.945242) (4250,0.944564) (4500,0.945068) (4750,0.945104) (5000,0.944734)};
\addlegendentry{MultiPPI}

\addplot[color=green!50!black, thick, dashdotted] coordinates {(250,0.949342) (500,0.949676) (750,0.948178) (1000,0.948368) (1250,0.947804) (1500,0.947608) (1750,0.947634) (2000,0.947968) (2250,0.945678) (2500,0.94571) (2750,0.947188) (3000,0.947288) (3250,0.946878) (3500,0.946222) (3750,0.945672) (4000,0.946096) (4250,0.945106) (4500,0.945758) (4750,0.94499) (5000,0.9455)};
\addlegendentry{PPI{\textnormal{+}}{\textnormal{+}} Vector}

\legend{};
\end{axis}
\begin{axis}[
    title={Mean-squared CI width (frac. of classical)},
    xlabel={Budget},
    xlabel style={yshift=5pt},   % Moves xlabel closer to the axis
    grid=both,
    outer sep=0pt,
    enlarge x limits=false,    minor grid style={dashed, gray!50},
    major grid style={solid, gray!80},
    width=\textwidth,
    yticklabel style={/pgf/number format/fixed}, % No scientific notation
    yticklabel={\pgfmathparse{\tick*100}\pgfmathprintnumber{\pgfmathresult}\%},
    xmin=0, xmax=5000,                          % Set x-axis range
    font=\scriptsize,             % Sets the base size for everything
    title style={yshift=-5pt, font=\small},    % Optional: Make title slightly bigger than the rest
    label style={font=\scriptsize}, % Font for xlabel and ylabel
    tick label style={font=\scriptsize}, % Font for the numbers on the axes
      xshift=5.9cm,
      width=6.5cm,
      height=5.5cm,
]

\addplot[color=gray!10!black!100, thick, densely dashed] coordinates {(250,0.7980297650159358) (500,0.7187469550967315) (750,0.6755218729759318) (1000,0.6481370401632676) (1250,0.6292909534968988) (1500,0.6153440150324012) (1750,0.6048330567074213) (2000,0.5964050640915959) (2250,0.5894116240873336) (2500,0.5835956995785824) (2750,0.5789486160221142) (3000,0.5749131325842569) (3250,0.5712367792981972) (3500,0.5681752270723763) (3750,0.5656556160884652) (4000,0.5634787128049807) (4250,0.561210126430431) (4500,0.5591082972304448) (4750,0.5576348274477432) (5000,0.5559577442621371)};
\addlegendentry{PPI{\textnormal{+}}{\textnormal{+}} Scalar Large}

\addplot[color=gray!40!black!70, thick, densely dashed] coordinates {(250,0.7963602976508679) (500,0.7316735701615502) (750,0.6990905830805777) (1000,0.6796850681017162) (1250,0.6662302277390959) (1500,0.6567788714404474) (1750,0.6499184780218964) (2000,0.6443351010075139) (2250,0.6400755553832811) (2500,0.6364959339775272) (2750,0.6334551376437744) (3000,0.6309574796664749) (3250,0.6287173797106966) (3500,0.626842233538315) (3750,0.6251799431571733) (4000,0.6237743290642711) (4250,0.6224306335401725) (4500,0.621303587508996) (4750,0.6201877356214713) (5000,0.619351992711427)};
\addlegendentry{PPI{\textnormal{+}}{\textnormal{+}} Scalar Medium}

\addplot[color=gray!70!black!40, thick, densely dashed] coordinates {(250,0.7923906988852455) (500,0.7441585532024705) (750,0.7219886616660748) (1000,0.7097514519347978) (1250,0.7016941656725879) (1500,0.6961683752140193) (1750,0.6920404901906884) (2000,0.6888481430669717) (2250,0.6864255745789448) (2500,0.6844398102627519) (2750,0.682661087849334) (3000,0.6811937529143596) (3250,0.6798432837371655) (3500,0.6788385116340983) (3750,0.6779854990340339) (4000,0.6770652984167209) (4250,0.6763498944678022) (4500,0.6757594489542277) (4750,0.6752436065268836) (5000,0.6745294664857189)};
\addlegendentry{PPI{\textnormal{+}}{\textnormal{+}} Scalar Small}

\addplot[color=gray!100!black!10, thick, densely dashed] coordinates {(250,0.7583625683043073) (500,0.7265679282749303) (750,0.7139481440135692) (1000,0.7070298790494929) (1250,0.7029648325781889) (1500,0.6999264851211573) (1750,0.6978227042041539) (2000,0.6962712723100042) (2250,0.6951734643056862) (2500,0.6942419766699929) (2750,0.6932510179974493) (3000,0.6924220146513287) (3250,0.6919293368535262) (3500,0.6913013635389091) (3750,0.6908067597112355) (4000,0.6904746396869361) (4250,0.6903383179303129) (4500,0.6901324758694185) (4750,0.6898117322374355) (5000,0.6895393523740894)};
\addlegendentry{PPI{\textnormal{+}}{\textnormal{+}} Scalar Tiny}

\addplot[color=blue, thick] coordinates {(250,0.7364485594588359) (500,0.6652491358255594) (750,0.6247753451854042) (1000,0.5984838484379978) (1250,0.5802005324242506) (1500,0.5662884217603211) (1750,0.555594095993414) (2000,0.5470787303967217) (2250,0.5403756456955711) (2500,0.5346494162878811) (2750,0.5296870866843192) (3000,0.5255543462630179) (3250,0.521851550014873) (3500,0.5187525821937171) (3750,0.5160298460248852) (4000,0.5134915825414207) (4250,0.5112914031725155) (4500,0.5092197292442158) (4750,0.5075191236737114) (5000,0.5058913429926757)};
\addlegendentry{MultiPPI}

\addplot[color=green!50!black, thick, dashdotted] coordinates {(250,0.85678954064154) (500,0.7780793603749747) (750,0.7272420193553959) (1000,0.69111205455506) (1250,0.6639482419331265) (1500,0.6430607889754084) (1750,0.6259984881052804) (2000,0.6122389685787858) (2250,0.6010931009349222) (2500,0.5913010556432227) (2750,0.582917791107265) (3000,0.5756673246034175) (3250,0.5692526046813898) (3500,0.5638234517156931) (3750,0.5585350569763351) (4000,0.5540033007975464) (4250,0.5499897150514994) (4500,0.5463411841089482) (4750,0.543008632193073) (5000,0.5399293274288467)};
\addlegendentry{PPI{\textnormal{+}}{\textnormal{+}} Vector}

\legend{};
\end{axis}
\begin{axis}[
    title={Mean-squared Error (frac. of classical)},
    xlabel={Budget},
    xlabel style={yshift=5pt},   % Moves xlabel closer to the axis
    grid=both,
    outer sep=0pt,
    enlarge x limits=false,    minor grid style={dashed, gray!50},
    major grid style={solid, gray!80},
    width=\textwidth,
    yticklabel style={/pgf/number format/fixed}, % No scientific notation
    yticklabel={\pgfmathparse{\tick*100}\pgfmathprintnumber{\pgfmathresult}\%},
    xmin=0, xmax=5000,                          % Set x-axis range
    font=\scriptsize,             % Sets the base size for everything
    title style={yshift=-5pt, font=\small},    % Optional: Make title slightly bigger than the rest
    label style={font=\scriptsize}, % Font for xlabel and ylabel
    tick label style={font=\scriptsize}, % Font for the numbers on the axes
      xshift=11.7cm,
      width=6.5cm,
      height=5.5cm,
            legend pos=north east,
            legend cell align={left},
            legend style={font=\tiny,inner sep=1pt,row sep=-2pt},
]

\addplot[color=gray!10!black!100, thick, densely dashed] coordinates {(250,0.8061141364618488) (500,0.7275332876486599) (750,0.679060185879669) (1000,0.6517272002953728) (1250,0.6335831692514948) (1500,0.6192817033460803) (1750,0.609336497750568) (2000,0.6016372119634451) (2250,0.5940363274869541) (2500,0.5871187145479723) (2750,0.5848482223531245) (3000,0.5811777332565368) (3250,0.5753609282716077) (3500,0.5728593406455009) (3750,0.5697631691443428) (4000,0.5672447103439366) (4250,0.5675054772423553) (4500,0.5657865303479424) (4750,0.5631455941702603) (5000,0.5608239633710729)};
\addlegendentry{PPI{\textnormal{+}}{\textnormal{+}} Scalar Large}

\addplot[color=gray!40!black!70, thick, densely dashed] coordinates {(250,0.8034027812418819) (500,0.7364615733618085) (750,0.7021288754974137) (1000,0.6843218581282777) (1250,0.6694094225823446) (1500,0.6591665198934142) (1750,0.6542816770058498) (2000,0.6488161048057528) (2250,0.6429414022691462) (2500,0.6408674128083341) (2750,0.638130461347414) (3000,0.6369257766430112) (3250,0.6323267631539511) (3500,0.631581581508514) (3750,0.630226990571085) (4000,0.6293185044001302) (4250,0.6279307305781405) (4500,0.6274555665270142) (4750,0.6264383021976611) (5000,0.6258414127227888)};
\addlegendentry{PPI{\textnormal{+}}{\textnormal{+}} Scalar Medium}

\addplot[color=gray!70!black!40, thick, densely dashed] coordinates {(250,0.7918504913569675) (500,0.7439558885239508) (750,0.7288925671437027) (1000,0.715543801574256) (1250,0.7068494247287058) (1500,0.7010260107234209) (1750,0.6960664104952957) (2000,0.6915616360757079) (2250,0.6928162480450434) (2500,0.6904420937827911) (2750,0.6878207320864346) (3000,0.6863619189655898) (3250,0.6841655996106406) (3500,0.68431140033583) (3750,0.6833163160722117) (4000,0.6834265094884598) (4250,0.6815018745818692) (4500,0.6820145009059648) (4750,0.6818716578266709) (5000,0.681499540574864)};
\addlegendentry{PPI{\textnormal{+}}{\textnormal{+}} Scalar Small}

\addplot[color=gray!100!black!10, thick, densely dashed] coordinates {(250,0.7628591321102924) (500,0.7299301634703218) (750,0.7180030230475876) (1000,0.7136607667986825) (1250,0.710143805763209) (1500,0.7071268402673393) (1750,0.7035490919506011) (2000,0.7030981801508427) (2250,0.7006784248887161) (2500,0.6980966596206589) (2750,0.6989068239688977) (3000,0.6979370735881584) (3250,0.6983299768281421) (3500,0.6971807220250954) (3750,0.6967984099520609) (4000,0.6965525429607541) (4250,0.696751068505863) (4500,0.6961244509425846) (4750,0.6960175856663144) (5000,0.6950498346909258)};
\addlegendentry{PPI{\textnormal{+}}{\textnormal{+}} Scalar Tiny}

\addplot[color=blue, thick] coordinates {(250,0.7453192708389018) (500,0.676869238998873) (750,0.6331965129094288) (1000,0.6074481903280788) (1250,0.5918713823219656) (1500,0.5775090193359138) (1750,0.5671710533294675) (2000,0.5583028623031042) (2250,0.551236275520762) (2500,0.5461491048826467) (2750,0.5403229193198908) (3000,0.536702030961984) (3250,0.5319401212163277) (3500,0.531041060732247) (3750,0.5281111089295817) (4000,0.5245224957080034) (4250,0.5237557886810801) (4500,0.520877944890352) (4750,0.5189489485894658) (5000,0.5174176639688769)};
\addlegendentry{MultiPPI}

\addplot[color=green!50!black, thick, dashdotted] coordinates {(250,0.8548499694606531) (500,0.7790853722176825) (750,0.7313181690886361) (1000,0.6951430362330507) (1250,0.668285275853806) (1500,0.6483192190984365) (1750,0.6315733574661345) (2000,0.6178969947958857) (2250,0.6119929902466942) (2500,0.601444908696594) (2750,0.5901337580293808) (3000,0.582873596212884) (3250,0.5778029412131999) (3500,0.5728810965533042) (3750,0.5681340865239867) (4000,0.5631551591003348) (4250,0.5606555165437129) (4500,0.5564752701658192) (4750,0.5547710694737332) (5000,0.5514524911241405)};
\addlegendentry{PPI{\textnormal{+}}{\textnormal{+}} Vector}

\end{axis}
\end{tikzpicture}
}

    \textrm{\bf \scriptsize Experiment 3: Factuality}
\scalebox{.8}{
\centering
\begin{tikzpicture}
\begin{axis}[
    title={Coverage},
    xlabel={Budget},
    xlabel style={yshift=5pt},   % Moves xlabel closer to the axis
    grid=both,
    outer sep=0pt,
    enlarge x limits=false,    minor grid style={dashed, gray!50},
    major grid style={solid, gray!80},
    width=\textwidth,
    yticklabel style={/pgf/number format/fixed}, % No scientific notation
    yticklabel={\pgfmathparse{\tick*100}\pgfmathprintnumber{\pgfmathresult}\%},
    xmin=0, xmax=5000,                          % Set x-axis range
    font=\scriptsize,             % Sets the base size for everything
    title style={yshift=-5pt, font=\small},    % Optional: Make title slightly bigger than the rest
    label style={font=\scriptsize}, % Font for xlabel and ylabel
    tick label style={font=\scriptsize}, % Font for the numbers on the axes
      xshift=0.0cm,
      width=6.5cm,
      height=5.5cm,
            ymin=0.9, ymax=1,
]

\addplot[color=black, dashed, very thick] coordinates {(0,.95) (5000,.95)};
\addplot[color=blue, thick] coordinates {(250,0.945538) (500,0.945476) (750,0.945278) (1000,0.943838) (1250,0.94372) (1500,0.943052) (1750,0.942456) (2000,0.942378) (2250,0.94217) (2500,0.940856) (2750,0.941818) (3000,0.94139) (3250,0.941416) (3500,0.941236) (3750,0.940634) (4000,0.940704) (4250,0.941094) (4500,0.940598) (4750,0.940252) (5000,0.940664)};
\addlegendentry{MultiPPI}

\addplot[color=green!50!black, thick, dashdotted] coordinates {(250,0.946518) (500,0.946872) (750,0.944198) (1000,0.943856) (1250,0.944146) (1500,0.943878) (1750,0.943688) (2000,0.942932) (2250,0.94235) (2500,0.941556) (2750,0.94261) (3000,0.943076) (3250,0.941458) (3500,0.941888) (3750,0.941724) (4000,0.94199) (4250,0.941248) (4500,0.9416478832957665) (4750,0.941526) (5000,0.94074)};
\addlegendentry{PPI{\textnormal{+}}{\textnormal{+}} Vector}

\addplot[color=red!40,  thick, densely dotted] coordinates {(250,0.947736) (500,0.94722) (750,0.94721) (1000,0.947314) (1250,0.947246) (1500,0.947676) (1750,0.946948) (2000,0.947388) (2250,0.947488) (2500,0.946952) (2750,0.946948) (3000,0.947706) (3250,0.947418) (3500,0.947368) (3750,0.947194) (4000,0.947054) (4250,0.947504) (4500,0.947782) (4750,0.947214) (5000,0.947684)};
\addlegendentry{PPI{\textnormal{+}}{\textnormal{+}} Scalar 1 agent, 1 round}

\addplot[color=red!70,  thick, densely dotted] coordinates {(250,0.948592) (500,0.94818) (750,0.94739) (1000,0.947828) (1250,0.947194) (1500,0.947054) (1750,0.94763) (2000,0.947332) (2250,0.946914) (2500,0.94729) (2750,0.947038) (3000,0.947666) (3250,0.947092) (3500,0.947048) (3750,0.947562) (4000,0.947882) (4250,0.9477) (4500,0.94748) (4750,0.94695) (5000,0.946868)};
\addlegendentry{PPI{\textnormal{+}}{\textnormal{+}} Scalar 2 agents, 1 round}

\addplot[color=red,  thick, densely dotted] coordinates {(250,0.949082) (500,0.947858) (750,0.947972) (1000,0.947402) (1250,0.94711) (1500,0.946394) (1750,0.94729) (2000,0.947906) (2250,0.94729) (2500,0.94756) (2750,0.947282) (3000,0.947546) (3250,0.948076) (3500,0.947432) (3750,0.946998) (4000,0.947156) (4250,0.94686) (4500,0.947584) (4750,0.94704) (5000,0.94712)};
\addlegendentry{PPI{\textnormal{+}}{\textnormal{+}} Scalar 3 agents, 1 round}

\addplot[color=red!40,  thick, densely dashed] coordinates {(250,0.948506) (500,0.947204) (750,0.947912) (1000,0.947434) (1250,0.947222) (1500,0.947024) (1750,0.94723) (2000,0.947564) (2250,0.947432) (2500,0.947048) (2750,0.947496) (3000,0.947808) (3250,0.947014) (3500,0.947362) (3750,0.9468) (4000,0.947326) (4250,0.948028) (4500,0.94763) (4750,0.947938) (5000,0.94757)};
\addlegendentry{PPI{\textnormal{+}}{\textnormal{+}} Scalar 1 agent, 2 rounds}

\addplot[color=red!70,  thick, densely dashed] coordinates {(250,0.947954) (500,0.947922) (750,0.94791) (1000,0.947296) (1250,0.947542) (1500,0.947132) (1750,0.947938) (2000,0.947884) (2250,0.94794) (2500,0.94772) (2750,0.947672) (3000,0.94771) (3250,0.947114) (3500,0.947054) (3750,0.94752) (4000,0.94779) (4250,0.947448) (4500,0.947178) (4750,0.94703) (5000,0.947374)};
\addlegendentry{PPI{\textnormal{+}}{\textnormal{+}} Scalar 2 agents, 2 rounds}

\addplot[color=red,  thick, densely dashed] coordinates {(250,0.947816) (500,0.94743) (750,0.94672) (1000,0.946968) (1250,0.947368) (1500,0.94764) (1750,0.948064) (2000,0.9478) (2250,0.947508) (2500,0.946924) (2750,0.947716) (3000,0.947716) (3250,0.94788) (3500,0.947562) (3750,0.94702) (4000,0.947404) (4250,0.94725) (4500,0.947156) (4750,0.947378) (5000,0.947052)};
\addlegendentry{PPI{\textnormal{+}}{\textnormal{+}} Scalar 3 agents, 2 rounds}

\legend{};
\end{axis}
\begin{axis}[
    title={Mean-squared CI width (frac. of classical)},
    xlabel={Budget},
    xlabel style={yshift=5pt},   % Moves xlabel closer to the axis
    grid=both,
    outer sep=0pt,
    enlarge x limits=false,    minor grid style={dashed, gray!50},
    major grid style={solid, gray!80},
    width=\textwidth,
    yticklabel style={/pgf/number format/fixed}, % No scientific notation
    yticklabel={\pgfmathparse{\tick*100}\pgfmathprintnumber{\pgfmathresult}\%},
    xmin=0, xmax=5000,                          % Set x-axis range
    font=\scriptsize,             % Sets the base size for everything
    title style={yshift=-5pt, font=\small},    % Optional: Make title slightly bigger than the rest
    label style={font=\scriptsize}, % Font for xlabel and ylabel
    tick label style={font=\scriptsize}, % Font for the numbers on the axes
      xshift=5.9cm,
      width=6.5cm,
      height=5.5cm,
]

\addplot[color=blue, thick] coordinates {(250,0.7815891318787707) (500,0.7572098412105606) (750,0.7466481899524897) (1000,0.7407283809898375) (1250,0.7368031403344583) (1500,0.7339701782676044) (1750,0.7318499809275609) (2000,0.7302118240440247) (2250,0.7287291638020583) (2500,0.7276003028132875) (2750,0.7267786382532437) (3000,0.7259347830093046) (3250,0.7250132246869493) (3500,0.7244139395369087) (3750,0.72395190339402) (4000,0.7234819760514071) (4250,0.7231552776930326) (4500,0.7227769164645585) (4750,0.7223682918798461) (5000,0.7221515401134438)};
\addlegendentry{MultiPPI}

\addplot[color=green!50!black, thick, dashdotted] coordinates {(250,0.904242808363255) (500,0.8577838581849664) (750,0.8295478681759987) (1000,0.810725682526691) (1250,0.7976348877104296) (1500,0.7874093984279751) (1750,0.7795445804211422) (2000,0.7731819012153286) (2250,0.7679933504204509) (2500,0.7638187919604417) (2750,0.7599342645764056) (3000,0.7569573209221803) (3250,0.7541019409911524) (3500,0.7516759956705275) (3750,0.7495924544516358) (4000,0.7476750829161933) (4250,0.7460780969384618) (4500,0.7445910577762498) (4750,0.7433199942747893) (5000,0.7419010161520146)};
\addlegendentry{PPI{\textnormal{+}}{\textnormal{+}} Vector}

\addplot[color=red!40,  thick, densely dotted] coordinates {(250,0.8145599347045719) (500,0.8047088734042923) (750,0.8011036528215908) (1000,0.799378235627502) (1250,0.7982728586535643) (1500,0.7975034831739398) (1750,0.7969587209842053) (2000,0.7966603840852226) (2250,0.7963234832211699) (2500,0.7960159923878171) (2750,0.7958503735736842) (3000,0.795684369390539) (3250,0.795474245258313) (3500,0.7953278072100014) (3750,0.7952302887282232) (4000,0.7950959034738128) (4250,0.7950146048410192) (4500,0.7949870384800093) (4750,0.7948568062595484) (5000,0.794867567497817)};
\addlegendentry{PPI{\textnormal{+}}{\textnormal{+}} Scalar 1 agent, 1 round}

\addplot[color=red!70,  thick, densely dotted] coordinates {(250,0.7975003143703929) (500,0.7765465267082413) (750,0.7683616186450305) (1000,0.7641199619808533) (1250,0.7614000292932283) (1500,0.7596853735448843) (1750,0.7582928029347317) (2000,0.7572619553043435) (2250,0.7565520903113171) (2500,0.7560199428560952) (2750,0.755467846379558) (3000,0.7549687345100314) (3250,0.7547207326292005) (3500,0.7544172086926421) (3750,0.7541149444588526) (4000,0.7537338608395145) (4250,0.7536785628253279) (4500,0.75358615497684) (4750,0.7533001260859777) (5000,0.7531872108902771)};
\addlegendentry{PPI{\textnormal{+}}{\textnormal{+}} Scalar 2 agents, 1 round}

\addplot[color=red,  thick, densely dotted] coordinates {(250,0.8077397128802674) (500,0.7792440280291139) (750,0.767641689607561) (1000,0.7613914224499864) (1250,0.7575040626384333) (1500,0.7548458432322815) (1750,0.7527804964738845) (2000,0.7513023727189286) (2250,0.7502696112175546) (2500,0.7492454973032927) (2750,0.7485728817367753) (3000,0.7479590520049251) (3250,0.7473168411553578) (3500,0.7468406308943033) (3750,0.7464973194305706) (4000,0.7461549623335709) (4250,0.745681646387856) (4500,0.7454106976460372) (4750,0.7452580000654304) (5000,0.7450652200738855)};
\addlegendentry{PPI{\textnormal{+}}{\textnormal{+}} Scalar 3 agents, 1 round}

\addplot[color=red!40,  thick, densely dashed] coordinates {(250,0.8154673894918791) (500,0.7966250594791043) (750,0.7893637136179762) (1000,0.7856560560919792) (1250,0.7834230991875888) (1500,0.7817804065973237) (1750,0.7805828272939067) (2000,0.7797306157564002) (2250,0.779039920961176) (2500,0.7784503269690237) (2750,0.7781038705401989) (3000,0.7776788825918429) (3250,0.7773624795439881) (3500,0.7770735603019362) (3750,0.7768463204422404) (4000,0.7764752090937018) (4250,0.7764818857141428) (4500,0.7763976342074234) (4750,0.7760688115404396) (5000,0.7760064823984619)};
\addlegendentry{PPI{\textnormal{+}}{\textnormal{+}} Scalar 1 agent, 2 rounds}

\addplot[color=red!70,  thick, densely dashed] coordinates {(250,0.8311138163254586) (500,0.7998011835790856) (750,0.7862672580651878) (1000,0.7787619529895309) (1250,0.7742052858889021) (1500,0.7710243116447721) (1750,0.7685908039845144) (2000,0.7666750571136998) (2250,0.7653770752992336) (2500,0.7641404520660701) (2750,0.7632082301876095) (3000,0.762353410902316) (3250,0.7617930858759311) (3500,0.7611699754541712) (3750,0.7607216868918384) (4000,0.7603094995463453) (4250,0.7598264863595037) (4500,0.7594562057599211) (4750,0.7592165852451561) (5000,0.7588808000206959)};
\addlegendentry{PPI{\textnormal{+}}{\textnormal{+}} Scalar 2 agents, 2 rounds}

\addplot[color=red,  thick, densely dashed] coordinates {(250,0.8539714261471679) (500,0.8169431225895919) (750,0.799771912783853) (1000,0.7899213085547999) (1250,0.783454798512838) (1500,0.7789063899707016) (1750,0.7759932201878386) (2000,0.7733936200784055) (2250,0.7712958462508623) (2500,0.7694932070737023) (2750,0.7683082067443473) (3000,0.7671258510239723) (3250,0.7660477861028351) (3500,0.7652340245322122) (3750,0.7646836284075451) (4000,0.763976408802095) (4250,0.7632473093618016) (4500,0.7627147599937791) (4750,0.7622313686994081) (5000,0.76176560138543)};
\addlegendentry{PPI{\textnormal{+}}{\textnormal{+}} Scalar 3 agents, 2 rounds}

\legend{};
\end{axis}
\begin{axis}[
    title={Mean-squared Error (frac. of classical)},
    xlabel={Budget},
    xlabel style={yshift=5pt},   % Moves xlabel closer to the axis
    grid=both,
    outer sep=0pt,
    enlarge x limits=false,    minor grid style={dashed, gray!50},
    major grid style={solid, gray!80},
    width=\textwidth,
    yticklabel style={/pgf/number format/fixed}, % No scientific notation
    yticklabel={\pgfmathparse{\tick*100}\pgfmathprintnumber{\pgfmathresult}\%},
    xmin=0, xmax=5000,                          % Set x-axis range
    font=\scriptsize,             % Sets the base size for everything
    title style={yshift=-5pt, font=\small},    % Optional: Make title slightly bigger than the rest
    label style={font=\scriptsize}, % Font for xlabel and ylabel
    tick label style={font=\scriptsize}, % Font for the numbers on the axes
      xshift=11.7cm,
      width=6.5cm,
      height=5.5cm,
            legend pos=north east,
            legend cell align={left},
            legend style={font=\tiny,inner sep=1pt,row sep=-2pt},
]

\addplot[color=blue, thick] coordinates {(250,0.7932126143326675) (500,0.7684026174991527) (750,0.7593893835503579) (1000,0.755272338372958) (1250,0.752031407151611) (1500,0.7506657074847831) (1750,0.7508367297182216) (2000,0.7489675870313275) (2250,0.7489625801202315) (2500,0.7497372718511056) (2750,0.7462883593261183) (3000,0.7463243789791116) (3250,0.746167889440732) (3500,0.7461831609224707) (3750,0.7468721202621297) (4000,0.7466511685300513) (4250,0.7447057253784731) (4500,0.7449749407953106) (4750,0.7466985536539664) (5000,0.7453568185356301)};
\addlegendentry{MultiPPI}

\addplot[color=green!50!black, thick, dashdotted] coordinates {(250,0.9129688726142113) (500,0.8668421301066459) (750,0.8453967070936118) (1000,0.8277390609664009) (1250,0.8130518125147695) (1500,0.8039930034809354) (1750,0.7966795320559148) (2000,0.7917922277806836) (2250,0.7877841589518322) (2500,0.7850836634301456) (2750,0.7788393602955023) (3000,0.7744496930851945) (3250,0.7762419371806527) (3500,0.7730500064424409) (3750,0.7705486607235917) (4000,0.7669562632572285) (4250,0.7686379615534535) (4500,0.766601416796857) (4750,0.7665106926889446) (5000,0.7657581440642486)};
\addlegendentry{PPI{\textnormal{+}}{\textnormal{+}} Vector}

\addplot[color=red!40,  thick, densely dotted] coordinates {(250,0.8193295659118004) (500,0.8103040621457978) (750,0.8072231113332462) (1000,0.8043572907470224) (1250,0.8034136341463657) (1500,0.8026937693775978) (1750,0.8043956602301641) (2000,0.8019609880921327) (2250,0.8024786163912986) (2500,0.8021481893393759) (2750,0.8018199335683782) (3000,0.800561595864214) (3250,0.8007836018375716) (3500,0.8007303170447515) (3750,0.8008199480278352) (4000,0.8003479480379248) (4250,0.8004802962866358) (4500,0.7989755736372824) (4750,0.8012374885722333) (5000,0.7988094911991698)};
\addlegendentry{PPI{\textnormal{+}}{\textnormal{+}} Scalar 1 agent, 1 round}

\addplot[color=red!70,  thick, densely dotted] coordinates {(250,0.7989838784516625) (500,0.7804503207289006) (750,0.7737510879832402) (1000,0.7683457836458774) (1250,0.7668183644683522) (1500,0.7660337613210183) (1750,0.7629729476868574) (2000,0.7625845013652953) (2250,0.7634026559270034) (2500,0.7616217281100441) (2750,0.7606385303639466) (3000,0.7594333669907456) (3250,0.7607460970095683) (3500,0.759776136617649) (3750,0.7582643471005144) (4000,0.7577118607389522) (4250,0.7580077114354602) (4500,0.7582810466212503) (4750,0.7594099229969855) (5000,0.7594378254627586)};
\addlegendentry{PPI{\textnormal{+}}{\textnormal{+}} Scalar 2 agents, 1 round}

\addplot[color=red,  thick, densely dotted] coordinates {(250,0.808438699854797) (500,0.7820585698927633) (750,0.7714216284383852) (1000,0.7666523951122317) (1250,0.763783686345367) (1500,0.7627102496172011) (1750,0.7577077050234847) (2000,0.7560302158698625) (2250,0.75567113416502) (2500,0.7537768923362747) (2750,0.7537672995679462) (3000,0.752315320217536) (3250,0.7509347717401323) (3500,0.7523143074561516) (3750,0.7523984366180293) (4000,0.7510081813218378) (4250,0.7514152160307179) (4500,0.7503833725713546) (4750,0.7515147438559914) (5000,0.7519011435549817)};
\addlegendentry{PPI{\textnormal{+}}{\textnormal{+}} Scalar 3 agents, 1 round}

\addplot[color=red!40,  thick, densely dashed] coordinates {(250,0.8178789686128312) (500,0.8016236471186714) (750,0.7937985258943859) (1000,0.7913215508475595) (1250,0.7889999216593798) (1500,0.7873263077734705) (1750,0.7863042279839697) (2000,0.7846604812626384) (2250,0.784187412165729) (2500,0.7846960501570198) (2750,0.7830136142523306) (3000,0.7815842900848065) (3250,0.7823926993253906) (3500,0.7818226168480067) (3750,0.78262415255693) (4000,0.7832507887623079) (4250,0.7802754203865566) (4500,0.7808749952919724) (4750,0.7799705110862153) (5000,0.7814259673727549)};
\addlegendentry{PPI{\textnormal{+}}{\textnormal{+}} Scalar 1 agent, 2 rounds}

\addplot[color=red!70,  thick, densely dashed] coordinates {(250,0.8350410069350743) (500,0.8029438529093469) (750,0.7905939755499624) (1000,0.7841911756033078) (1250,0.7768426708901177) (1500,0.7761441222492257) (1750,0.7718423545073578) (2000,0.7704473973165701) (2250,0.7682812195198443) (2500,0.7685642257604614) (2750,0.7670402485575956) (3000,0.7668312939162978) (3250,0.7669913688473404) (3500,0.7661006477147392) (3750,0.7647912747874197) (4000,0.7642871926210213) (4250,0.7651100414330595) (4500,0.7649002801829996) (4750,0.7640746422386885) (5000,0.7640383112869098)};
\addlegendentry{PPI{\textnormal{+}}{\textnormal{+}} Scalar 2 agents, 2 rounds}

\addplot[color=red,  thick, densely dashed] coordinates {(250,0.857386370586729) (500,0.8201690333175979) (750,0.8063371036620116) (1000,0.7968173302248817) (1250,0.787124409107033) (1500,0.7834030490287359) (1750,0.7805052755607268) (2000,0.7781655929004995) (2250,0.7769942194602811) (2500,0.7763348181773418) (2750,0.77344781230839) (3000,0.7722918386682269) (3250,0.7707914505644414) (3500,0.7697815237095043) (3750,0.7707992340709693) (4000,0.7686922058837109) (4250,0.7674015304954125) (4500,0.7674946955696178) (4750,0.7676974552141823) (5000,0.7675303322441959)};
\addlegendentry{PPI{\textnormal{+}}{\textnormal{+}} Scalar 3 agents, 2 rounds}

\end{axis}
\end{tikzpicture}
}

    \vspace{-10pt}
    \caption{Results by budget for the experiments on Chatbot Arena (a), ProcessBench (b), and Factuality (c). For each estimator (all baselines and MultiPPI), the left column plots the empirical coverage of the 95\% CI, the middle column plots the width of the 95\% CI, and the right column plots the empirical mean-squared error of the point estimate. The fully-labeled sample size $N$ is 250.\looseness=-1}
    \label{fig:results-all}
    \vspace{-10pt}
\end{figure}

% \newpage
% \begin{figure}[!t]
%     \centering
%     %\includegraphics[width=\linewidth]{figures/chatbot_arena_results2.png}
%     \textrm{\bf \scriptsize Experiment 1: Chatbot Arena}
%     \input{figures/chatbot_arena_figure_100}  
%     \textrm{\bf \scriptsize Experiment 2: ProcessBench} 
%     \input{figures/processbench_figure_100}
%     \textrm{\bf \scriptsize Experiment 3: Factuality}
%     \input{figures/factuality_figure_100}
%     \vspace{-10pt}
%     \caption{Results by budget for the experiments on Chatbot Arena (a), ProcessBench (b), and Factuality (c). For each estimator (all baselines and MultiPPI), the left column plots the empirical coverage of the 95\% CI, the middle column plots the width of the 95\% CI, and the right column plots the empirical mean-squared error of the point estimate. The fully-labeled sample size $N$ is 100.\looseness=-1}
%     \label{fig:results-all}
%     \vspace{-10pt}
% \end{figure}
% \newpage
% \begin{figure}[!t]
%     \centering
%     %\includegraphics[width=\linewidth]{figures/chatbot_arena_results2.png}
%     \textrm{\bf \scriptsize Experiment 1: Chatbot Arena}
%     \input{figures/chatbot_arena_figure_250}
%     \textrm{\bf \scriptsize Experiment 2: ProcessBench} 
%     \input{figures/processbench_figure_250}
%     \textrm{\bf \scriptsize Experiment 3: Factuality}
%     \input{figures/factuality_figure_250}
%     \vspace{-10pt}
%     \caption{Results by budget for the experiments on Chatbot Arena (a), ProcessBench (b), and Factuality (c). For each estimator (all baselines and MultiPPI), the left column plots the empirical coverage of the 95\% CI, the middle column plots the width of the 95\% CI, and the right column plots the empirical mean-squared error of the point estimate. The fully-labeled sample size $N$ is 250.\looseness=-1}
%     \label{fig:results-all}
%     \vspace{-10pt}
% \end{figure}
% \newpage
%\charlie{TODO: At a minimum here, comment on our results for very small sample size, with applications to low-resource etc. Show that it can be robustly estimated with a small sample. For reviewer YZU1.}


% \section{Theoretical results}\label{sec:theory}

% \subsection{A setting in which the estimator is minimax optimal}
% Let $\mathcal{P}_\Sigma$ denote the set of distributions on $X$ with covariance matrix $\Sigma$. Let $\Theta_B$ denote the set of budget-satisfying estimators $\hat{\theta}$: formally,
% \[
% \Theta_B = \left\{\hat{\theta}:\hat{\theta}\text{ is a measurable function of $n_I$ independent copies of $X_I$, for each $I\in\mathcal{I}$, and $\mathbf{B}(\underline{n},B)$.}\right\}
% \]
% Our estimator is minimax optimal in this class.
% \begin{theorem}
%     We have
%     \[
%     \inf_{\hat{\theta}\in \Theta_B} \sup_{ P\in\mathcal{P}_\Sigma} \Ex\left[(\hat{\theta}-\Ex Y)^2\right] = \Var \hat{\theta}_{\text{MultiPPI}(\Sigma)}
%     \]
% \end{theorem}

% \subsection{Asymptotic normality}
% Fix a vector $B_0$. We consider the limit in which our budget is $B=t\cdot B_0$ and let $t\to\infty$. Suppose that $\widehat{\Sigma}\overset{p}{\to} \Sigma$ in the operator norm, potentially dependent on the variables sampled $X_I$. Then we have
% \begin{theorem}\label{thm:asymptotic}
%     \[
%     \sqrt{t}\left(\hat{\theta}_{\text{MultiPPI}(\widehat{\Sigma})}-\theta^*\right)\overset{d}{\to}\mathcal{N}(0,V)
%     \]
%     where
%     \[
%     V = \min_{\underline{n}:\mathbf{B}(\underline{n},B_0)}a^\top\left(\sum_I n_IP_I^\top (P_I\Sigma P_I^\top)^{-1}P_I\right)^\dagger a
%     \]
% \end{theorem}
% The point of the theorem is to handle the case in which $\widehat{\Sigma}$ depends strongly on some of the sampled variables. To make concrete how this may occur in practice, consider the situation explored in \autoref{sec:exp-setup}: $\widehat{\Sigma}$ depends on some number $\texttt{num\_labels}$ of labeled samples $X^{(1)},\dots,X^{(\texttt{num\_labels})}$ which we have access to, and we may query no more. However, we use $\widehat{\Sigma}$ to determine our allocation for our budget $B$. This formally corresponds to a problem in which the budget constraints are $n_{[k]}\leq \texttt{num\_labels}$ and $\sum_I c_In_I\leq B$.


% \subsection{Finite-sample bounds}

% %The remaining budget is thus $B-mc_{[k]}$. For ease of notation, let us denote by $\hat{\theta}_{\multiPPI(\Sigma,B)}$ the Multi-PPI estimator with budget $B$.
% \subsection{When are two model collections comparable? \charlie{[WIP]}}
% We introduce a strict partial order on the set of model collections $(X_I)_{I\in\mathcal{I}}$: Let $\rho_I$ denote the correlation coefficient of $X_I$ with $Y$, and let $c_I$ denote the cost of sampling $X_I$. Let $\hat{\theta}^*_{\text{PPI++}(I)}$ denote the PPI++ estimator with $X_I$, using the optimal $\lambda$-tuning parameter. This estimator depends on $n$, the number of labeled examples seen, and $B$ the budget available to query independent samples of $X_I$. Note crucially that we have oracle access to the optimal $\lambda$-tuning parameter here.

% We introduce the following partial order: say that
% \[
% I\succ J \quad\iff\quad \rho_I>\rho_J,\quad c_I<c_J
% \]
% If and only if $I\succ J$, then we should never sample $X_J$. This is made precise in the following result:
% \begin{theorem} We have that
% \[
% \Var(\hat{\theta}^*_{\text{PPI++}(I)}) < \Var(\hat{\theta}^*_{\text{PPI++}(J)})
% \]
% for all $n,B$ if and only if $I\succ J$.
% \end{theorem}
% or, equivalently,
% \begin{theorem} Fix $n$. There exist budget regimes in which $\hat{\theta}^*_{\text{PPI++}(I)}$ and $\hat{\theta}^*_{\text{PPI++}(J)}$ each dominate if and only if $I$ and $J$ are incomparable.
% \end{theorem}
% \[
% \textbf{Note:}\quad I\succ J \quad\implies \frac{\rho_I}{c_I} > \frac{\rho_J}{c_J}
% \]
% \[
% \text{Consider}\quad I^* = \argmax_{I\in\mathcal{I}} \frac{\rho_I}{c_I}
% \]
% \[
% \text{As}\quad B\to0,\quad \hat{\theta}_{\text{MultiPPI}}\to \hat{\theta}^*_{\text{PPI++($I^*$)}}
% \]

% \subsection{Behavior in limiting regimes \charlie{[WIP]}}
% We next describe a certain limiting behavior. Slightly changing notation, we will target $\Ex Y$ given $n$ samples from $(Y,X_1,\dots,X_k)$, and a budget of $C$ to expend querying subsets $X_I$ of the model vector $X=(X_1,\dots,X_k)$, with associated cost $c_I$. Let $\rho_I$ denote the correlation coefficient of $X_I$ with $Y$.

% \begin{theorem}
%     Fix $B$ and consider the limit as $n\to\infty$. For each $I\in\mathcal{I}$, let $a_I:=\rho_I/c_I$. Suppose that $I^*$ uniquely minimizes $a_I$. Then 
%     \[
%     %n^*_\longrightarrow\texttt{unit\_vector}(I^*)\cdot \frac{C}{c_{I^*}}
%     n^*_{I^*}\to B/c_{I^*},\quad\quad n_I^*\to0\,\forall I\neq I^*
%     \]
% \end{theorem}

% \subsection{Performance in closed-form \charlie{[WIP]}}

\section{Conclusion}
In this work, we introduce Multiple-Prediction-Powered Inference (MultiPPI), a framework for efficiently estimating expectations under budget constraints by optimally leveraging multiple information sources of varying costs. MultiPPI formulates the optimal allocation of queries across subsets of variables as a second-order cone program in the case of a single budget constraints, or a semi-definite program in the case of multiple---both can be efficiency solved using off-the-shelf tools. We provide theoretical guarantees, including minimax optimality when covariances are known, and demonstrate empirically across diverse LLM evaluation tasks that MultiPPI outperforms existing methods. By adaptively balancing cost and information, MultiPPI achieves lower error for a given budget, automatically shifting its strategy from cheaper proxies to more expensive, accurate predictors as the budget increases, thus offering a principled and practical solution for cost-effective inference.\looseness=-1
\newpage

\section*{Acknowledgements}
We are very grateful for insightful comments and suggestions by the anonymous reviewers. We are also grateful for helpful discussions with Jonathan Berant, Michael Beukman, Frederic J\o rgensen, Petros Karypis, Chris Dyer, and Martin Wainwright.

\section*{Reproducibility statement}
To ensure reproducibility, we provide a detailed specification of the algorithm in \Cref{sec:algorithm}. We also include implementation details in \Cref{sec:experimental-details}, and address computational considerations in \Cref{sec:computation}. Finally, all experiments shown in \S\ref{sec:results} were averaged over 500k trials.



\bibliographystyle{iclr2026_conference}
\bibliography{iclr2026_conference}
\input
\clearpage
\appendix
\addcontentsline{toc}{chapter}{Appendices} % For main TOC
\etocsetnexttocdepth{3}
\localtableofcontents
\clearpage

\section{Ethics statement}
\label{sec:impacts}
This paper describes fundamental research on techniques for constructing statistically efficient estimates of a target metric by optimally allocating resources across multiple types of proxy measurements. The primary intended use case which is analyzed in this work is the evaluation of generative AI systems, for which reliable evaluation is a core technical challenge. Efficient and precise estimates of model performance  can help make AI systems easier to build, deploy, and monitor. We do not speculate about broader impacts that may follow from this technical contribution. Gemini was used for light copy-editing during the writing of this work.  

%\input{iclr2026/macros}




\section{Generalization to multiple budget inequalities}\label{app:generalized}
We recall some notation. Fix a set $\mathcal{I}$ of index subsets $I\subseteq[k]$. For each $I\in\mathcal{I}$, let $c_I=(c_I^{(1)},\dots,c_I^{(m)})\in\R_{\geq0}^m$ denote the vector-valued cost of querying the collection of models indexed by $I$. Similarly, for each $I\in\mathcal{I}$, we let $n_I\geq0$ be an integer denoting the number of times that the collection of models indexed by $I$ is queried. We let $\underline{n}=(n_I)_{I\in\mathcal{I}}$ refer to the associated allocation.

For a vector-valued budget $B\in\R_{\geq0}^m$, we say that the allocation $\underline{n}$ satisfies the budget $B$, and write $\mathbf{B}(\underline{n},B)$, if
\[
\sum_{I\in\mathcal{I}} n_Ic_I^{(1)}\leq B^{(1)},\dots,\sum_{I\in\mathcal{I}}n_Ic_I^{(m)}\leq B^{(m)},
\]
or more succinctly, 
\[
\sum_{I\in\mathcal{I}}n_Ic_I\leq B.
\]
Similarly, for each $I\in\mathcal{I}$, we let $\lambda_I\in\R^{|I|}$, and denote by $\underline{\lambda}=(\lambda_I)_{I\in\mathcal{I}}$ their collection. Let
\[
\hat{\theta}_{\underline{n},\underline{\lambda}} = \sum_{I\in\mathcal{I}:n_I>0}\frac{1}{n_I}\sum_{j=1}^{n_I}\lambda_I^\top X_I^{(I,j)}
\]
where $X^{(I,j)}$ denote independent copies of $X$ for every $I\in\mathcal{I}$ and $1\leq j\leq n_I$. We say that the unbiased condition holds for $\underline{n},\underline{\lambda}$, and write $\mathbf{U}$, if $\Ex \hat{\theta}_{\underline{n},\underline{\lambda}}=a^\top\Ex X$ for every distribution of finite second-moment on $X$.

Note that the variance of $\hat{\theta}_{\underline{n},\underline{\lambda}}$ depends only upon $\Sigma=\Cov(X)$. Thus we let
\[
\hat{\theta}_{\text{MultiPPI}(\Sigma)} := \hat{\theta}_{\underline{n},\underline{\lambda}}
\quad\begin{array}[t]{@{}l}
\text{ where $\underline{n},\underline{\lambda}$ are chosen so that the resulting estimator} \\ 
\text{ has minimal variance under $\Sigma$ such that $\mathbf{B}$ and $\mathbf{U}$ hold.} 
\end{array}
\]
\section{Detailed specification of the algorithm}\label{sec:algorithm}

In this section, we outline the procedure used in all experiments in greater detail. First, we describe the algorithm for the case of a single budget inequality, for which a more-efficient procedure exists; second, we describe the general case, in which the procedure reduces to a semi-definite program (SDP).
We first suppose that $\Sigma$ is known, and later explain the procedure in the case that it must be estimated from data.\looseness=-1

\subsection{The case of a single budget inequality, known $\Sigma$}
We suppose that there is a random vector $X\in\R^k$ with known covariance $\Sigma$, and our goal is to estimate $\theta^*=a^\top \Ex X$ for some fixed $a\in\R^k$. There is some fixed collection $\mathcal{I}$ of index subsets $I\subset\{1,\dots,k\}$ such that we may sample $X_I:=(X_i)_{i\in I}$. We may sample $X_I$ a maximum of $n_I$ times, subject to the constraint that $\sum_{I\in\mathcal{I}} c_In_I\leq B$ for some $c_I\geq0$ and $B>0$.

\paragraph{Step 1:} Solve the SOCP
\[
\sup_{y\in \R^k} a^\top y \quad\quad\text{s.t.}\quad \bigwedge_{I\in\mathcal{I}} \big\{y_I^\top {\Sigma_I} y_I\leq c_I^{-1}\big\}
\]
and obtain the solution $y_I^\star$ and the multipliers $\alpha^\star_I\geq0$ for each $I\in\mathcal{I}$.

\paragraph{Step 2:} Set
\begin{align*}
    \lambda^\star_I &= 2\alpha^\star_I \Sigma_I^{-1} y_I^\star \\
    n^\star_I &= \left\lfloor \left(\frac{B}{c_I}\right) \frac{\sqrt{c_I (\lambda^\star_I)^\top \Sigma_I \lambda^\star_I}}{\sum_{J\in\mathcal{I}}\sqrt{c_J (\lambda^\star_I)^\top \Sigma_J \lambda^\star_J} }\right\rfloor
\end{align*}
for each $I\in\mathcal{I}$.

\paragraph{Step 3:} For each $I\in\mathcal{I}$, independently sample $X_I$ $n^\star_I$ times, and compute the sample mean $\overline{\lambda^\star_I\cdot X_I}$. Return
\[
\hat{\theta}_{\text{MultiPPI}(\Sigma)} = \sum_{I\in\mathcal{I}} \overline{\lambda_I^\star \cdot X_I}
\]
with $(1-\alpha)$-confidence intervals given by
\[
\mathcal{C} = \hat{\theta}_{\text{MultiPPI}(\Sigma)} \pm z_{1-\alpha/2} \sqrt{\sum_{I\in\mathcal{I}} \frac{1}{n^\star_I} \widehat{\sigma^2}_{\lambda_I^\star \cdot X_I}}
\]
where $\widehat{\sigma^2}_{\lambda_I^\star \cdot X_I}$ denotes the sample variance of $\lambda_I^\star \cdot X_I$, and $z_p$ denotes the $p$\textsuperscript{th} quantile of the standard normal distribution.

\subsection{The case of multiple budget inequalities, known $\Sigma$}\label{sec:multi-budget-alg}
We again suppose that there is a random vector $X\in\R^k$ with known covariance $\Sigma$, and our goal is to estimate $\theta^*=a^\top \Ex X$ for some fixed $a\in\R^k$. We may now sample $X_I$ a maximum of $n_I$ times, subject to the constraints that $\sum_{I\in\mathcal{I}} c_I^{(\ell)}n_I\leq B^{(\ell)}$ for some $c_I^{(\ell)}\geq0$ and $B^{(\ell)}>0$, with $1\leq\ell\leq m$.

\paragraph{Step 1:} Solve the SDP
\begin{align*}
    \sup_{t\in\R} t \quad\text{s.t.}\quad \begin{pmatrix}
    \sum_{I\in\mathcal{I}} n_I P_I^\top \Sigma_I^{-1}P_I & a \\
    a^\top & t
\end{pmatrix} &\succeq0, \\
n_I &\geq0\quad&\forall I\in\mathcal{I} \\
\sum_{I\in\mathcal{I}}c_I^{(\ell)}n_I&\leq B^{(\ell)}\quad &\forall \ell\leq m
\end{align*}
for real valued $n_I$, and obtain solutions $n^\star_{I,\text{frac}}$.

\paragraph{Step 2:} Set
\begin{align*}
n^\star_I &= \left\lfloor n^\star_{I,\text{frac}}\right\rfloor \\    
\lambda^\star_I &= n^\star_I\Sigma_I^{-1} P_I \left(\sum_{I\in\mathcal{I}} n^\star_I P_I^\top \Sigma_I^{-1} P_I\right)^\dagger a
\end{align*}
for all $I\in\mathcal{I}$.

\paragraph{Step 3:} 
As in the previous section, for each $I\in\mathcal{I}$, independently sample $X_I$ $n^\star_I$ times, and compute the sample mean $\overline{\lambda^\star_I\cdot X_I}$. Return
\[
\hat{\theta}_{\text{MultiPPI}(\Sigma)} = \sum_{I\in\mathcal{I}} \overline{\lambda_I^\star \cdot X_I}
\]
with $(1-\alpha)$-confidence intervals given by
\[
\mathcal{C} = \hat{\theta}_{\text{MultiPPI}(\Sigma)} \pm z_{1-\alpha/2} \sqrt{\sum_{I\in\mathcal{I}} \frac{1}{n^\star_I} \widehat{\sigma^2}_{\lambda_I^\star \cdot X_I}}
\]
where $\widehat{\sigma^2}_{\lambda_I^\star \cdot X_I}$ denotes the sample variance of $\lambda_I^\star \cdot X_I$, and $z_p$ denotes the $p$\textsuperscript{th} quantile of the standard normal distribution.

\subsection{The case of unknown $\Sigma$}\label{sec:algorithm_unknown}
In general, the approach is to construct an estimate $\widehat{\Sigma}$ of $\Sigma$ from data, and use this estimate for $\Sigma$ in the steps outlined above. In principle, it is possible to recycle the data used to construct $\wh{\Sigma}$ in step 3 of the above procedures; this preserves asymptotic normality as a consequence of \Cref{thm:asymptotic_revised}. Below, we detail one approach to doing this---the approach used in our experiments, and the approach outlined in \Cref{sec:practical}.

Suppose that $a=(1,0,\dots,0)$, and we have some hard limit $N$ on the number of samples available from $X_1$. This typically represents a ``gold'' label which is invaluable in some sense. We also suppose that these labeled samples are fully labeled---that is, that the entire vector $X=(X_1,\dots,X_k)$ is visible in each case---or alternatively, that $N$ is small enough that they are relatively inexpensive to obtain model predictions for. 

\paragraph{Step 1:} Construct the empirical covariance matrix $\wh{\Sigma}$ from the $N$ fully-labeled samples.

\paragraph{Step 2:} Take $\mathcal{I}$ to be all subsets of models---that is, all subsets of $\{2,\dots,k\}$---together with the set of all indices $\{1,\dots,k\}$. Formally, $\mathcal{I}=\{\{1,\dots,k\}\}\cup 2^{\{2,\dots,k\}}$.

\paragraph{Step 3:} Run \Cref{sec:multi-budget-alg} with any existing budget constraints, together with the constraint that $n_{\{1,\dots,k\}}\leq N$, and obtain allocations $n^\star_I,\lambda^*_I$.

\paragraph{Step 4:} Sample accordingly, with the guarantee that the number of fully labeled samples $X_{\{1,\dots,k\}}$ queried won't exceed the number available, $N$. These samples from step 1 may be reused for this.

\paragraph{Step 5:} Return the resulting estimator, as described in \Cref{sec:multi-budget-alg}.

% In this section, we will (1) demonstrate the reduction to an SOCP, and (2) explain the procedure for reconstructing the optimal allocation from the solution to the SOCP.

% We first recall the problem of interest. Let $\Sigma\in\R^{k\times k}$ be SPD, let $\mathcal{I}$ denote a collection of index subsets $I\subseteq \{1,\dots,k\}$, and let $c_I$ be a positive scalar defined for every $I\in\mathcal{I}$. Denoting $\underline{c}=(c_I)_{I\in\mathcal{I}}\in\R_{>0}^{|\mathcal{I}|}$, our budget constraint may be written $\underline{c}^\top\underline{n}\leq B$. Lastly, let $P_I\in\R^{|I|\times k}$ denote the orthogonal projection onto the coordinates indexed by $I\in\mathcal{I}$. With all of this said, recall that Multi-Allocate$(\Sigma)$ solves
% \begin{equation}\label{eqn:og}
%     V(a) = \min_{\underline{n},\underline{\lambda}} \sum_{I\in\mathcal{I}:n_I>0} \frac{1}{n_I}\lambda_I^\top\Sigma_I\lambda_I\quad\quad\text{s.t.}\quad 
%     \sum_{I\in\mathcal{I}:n_I>0}P_I^\top \lambda_I=a,\quad     \underline{c}^\top\underline{n}\leq B
% \end{equation}
% We begin by deriving tractable methods to solve \autoref{eqn:og}. Let us assume for the moment that $\underline{n}\in\R_{\geq0}^{|\mathcal{I}|}$; we will later construct the final budget allocation by rounding. Our first step is to remove the dependence on $\underline{n}$: we show that the above problem is equivalent to the following:
% \begin{equation}\label{eqn:primal}
%     U(a)=\min_{\underline{\lambda}\in\Lambda} \sum_{I\in\mathcal{I}} \sqrt{c_I}\|\lambda_I\|_{\Sigma_I}\quad\quad\text{s.t.}\quad A\underline{\lambda} = a
% \end{equation}
% We next show that this is equivalent to the dual problem
% \begin{equation}\label{eqn:dual}
%     U(a) = \sup_{y\in \R^k} a^\top y \quad\quad\text{s.t.}\quad \bigwedge_{I\in\mathcal{I}} \big\{\|y\|_{\Sigma_I}^2\leq c_I\big\}
% \end{equation}
% Finally, this is a second order cone program, and can be solved with off-the-shelf tools. After we have shown these things, we describe how to convert solutions to \autoref{eqn:dual} into solutions to \autoref{eqn:og}.

% \begin{proposition}\label{prop:og-primal}
%     The problems described in \autoref{eqn:og} and \autoref{eqn:primal} yield the same optimum $V=U^2/B$.
% \end{proposition}
% \begin{proposition}\label{prop:primal-dual}
%     The problems described in \autoref{eqn:primal} and \autoref{eqn:dual} yield the same optimum $U$.
% \end{proposition}

\clearpage
\section{Additional experiments}

\subsection{Learned allocations and linear parameters}
\begin{figure}[!h]
    \centering
    \includegraphics[width=0.5\linewidth]{figures/allocations_by_budget_chatbot_arena.png}
    \caption{Proportion of budget allocated to different models in Experiment 1: ChatBot Arena. Gemini 2.5 Flash, the cheapest model, is most sampled in the low-budget regime, while the proportion of budget allocated to the joint (both models combined) increases monotonically with budget.}
    \label{fig:allocations-chatbotarena}
\end{figure}
\begin{figure}[!h]
    \centering
    \includegraphics[width=0.5\linewidth]{figures/allocations_by_budget_processbench.png}
    \caption{Proportion of budget allocated to different models in Experiment 2: ProcessBench. Tiny (125 word thinking budget) is most sampled in the low-budget regime, while the proportion of budget allocated to the joint (all models combined) increases monotonically with budget.}
    \label{fig:allocations-processbench}
\end{figure}


\begin{figure}[!h]
    \centering
    \includegraphics[width=\linewidth]{figures/lambdas_by_budget_processbench.png}
    \caption{Linear parameters $\lambda_I$ learned across budget regimes in Experiment 2: ProcessBench. While only the tiny model (125 word thinking budget) has a nonzero linear parameter in the low-budget regime, a \textit{cascading} behavior is learned in the large-budget regime: the cheaper models are prescribed the opposite sign from the more-expensive models in the joint term.}
    \label{fig:lambdas-processbench}
\end{figure}
\clearpage

\begin{figure}[!h]
    \centering
    \includegraphics[width=\linewidth]{figures/lambdas_by_budget_Chatbot_arena.png}
    \caption{Linear parameters $\lambda_I$ learned across budget regimes in Experiment 1: ChatBot Arena. While only Gemini 2.5 Pro has a nonzero linear parameter in the low-budget regime, a \textit{cascading} behavior is learned in the large-budget regime: the cheaper model (Gemini 2.5 Flash) is prescribed the opposite sign from the more-expensive model (Gemini 2.5 Pro) in the joint term.}
    \label{fig:lambdas-chatbotarena}
\end{figure}
    
\subsection{MultiPPI by varying number of labeled samples}
\label{sec:experiments-num-labeled}
In this section, we compare results of MultiPPI for a variable number of fully-labeled samples. MultiPPI continues to achieve smaller MSE than all baselines in all settings considered. In \Cref{fig:varying}, we plot the performance of MultiPPI with number of fully-labeled sample varying between $N=10$ and $N=200$. The extreme case $N=10$ is shown in \Cref{fig:10}, while  $N=1000$ is shown in \Cref{fig:1000}. Even for $N=10$, we find that MultiPPI improves on all baselines in MSE. It is important to note that, for all methods, including the baselines, the coverage is significantly below 95\% due to the small sample size. Nevertheless, even in this extreme setting, MultiPPI performs best in MSE.

% \begin{figure}[!h]
%     \centering
%   \input{figures/factuality_figure_variable.tex}
%   \vspace{-10pt}
%     \caption{Results by number of labeled examples, with fixed 1:1 ratio between budget and number of labeled examples.}
%     \label{fig:chatbotarena1000}
%     \vspace{-10pt}
% \end{figure}

% \begin{figure}[!h]
%     \centering
%   \input{iclr2026/figures/factuality_figure_variable_5}
%   \vspace{-10pt}
%     \caption{Results by number of labeled examples, with fixed 5:1 ratio between budget and number of labeled examples.}
%     \label{fig:chatbotarena1000}
%     \vspace{-10pt}
% \end{figure}

\begin{figure}[!t]
    \centering
    %\includegraphics[width=\linewidth]{figures/chatbot_arena_results2.png}
    \textrm{\bf \scriptsize Experiment 1: Chatbot Arena}
\scalebox{0.8}{
\centering
\begin{tikzpicture}
\begin{axis}[
    title={Coverage},
    xlabel={Budget},
    xlabel style={yshift=5pt},   % Moves xlabel closer to the axis
    grid=both,
    outer sep=0pt,
    enlarge x limits=false,    minor grid style={dashed, gray!50},
    major grid style={solid, gray!80},
    width=\textwidth,
    yticklabel style={/pgf/number format/fixed}, % No scientific notation
    yticklabel={\pgfmathparse{\tick*100}\pgfmathprintnumber{\pgfmathresult}\%},
    xmin=0, xmax=200,                          % Set x-axis range
    font=\scriptsize,             % Sets the base size for everything
    title style={yshift=-5pt, font=\small},    % Optional: Make title slightly bigger than the rest
    label style={font=\scriptsize}, % Font for xlabel and ylabel
    tick label style={font=\scriptsize}, % Font for the numbers on the axes
      xshift=0.0cm,
      width=6.5cm,
      height=5.5cm,
            ymin=0.8, ymax=1,
]

\addplot[color=black, dashed, very thick] coordinates {(0,.95) (5000,.95)};
\addplot[color=blue, thick] coordinates {(10,0.87695) (20,0.870882) (30,0.871124) (40,0.86913) (50,0.869456) (60,0.867726) (70,0.866332) (80,0.865782) (90,0.864686) (100,0.864118) (110,0.864576) (120,0.86331) (130,0.861808) (140,0.862414) (150,0.862882) (160,0.86173) (170,0.86073) (180,0.861652) (190,0.861682) (200,0.860516)};
\addlegendentry{MultiPPI}

\addplot[color=green!50!black, thick, dashdotted] coordinates {(10,0.891188) (20,0.885738) (30,0.87878) (40,0.875896) (50,0.876562) (60,0.87369) (70,0.871824) (80,0.87047) (90,0.868588) (100,0.867354) (110,0.867862) (120,0.866862) (130,0.867452) (140,0.866906) (150,0.865702) (160,0.866384) (170,0.86516) (180,0.863872) (190,0.864886) (200,0.863698)};
\addlegendentry{PPI{\textnormal{+}}{\textnormal{+}} Vector}

\addplot[color=brown!90, thick, densely dashed] coordinates {(10,0.891156) (20,0.887226) (30,0.886968) (40,0.884166) (50,0.881202) (60,0.881096) (70,0.878404) (80,0.87904) (90,0.878442) (100,0.877934) (110,0.878918) (120,0.877908) (130,0.875996) (140,0.877212) (150,0.876336) (160,0.876842) (170,0.877266) (180,0.87707) (190,0.876522) (200,0.87717)};
\addlegendentry{PPI{\textnormal{+}}{\textnormal{+}} Scalar (Pro)}

\addplot[color=brown!60, thick, densely dashed] coordinates {(10,0.880468) (20,0.87735) (30,0.879208) (40,0.878134) (50,0.876884) (60,0.876182) (70,0.876796) (80,0.876848) (90,0.875544) (100,0.87505) (110,0.875136) (120,0.874976) (130,0.874454) (140,0.875396) (150,0.874258) (160,0.874242) (170,0.874802) (180,0.8751) (190,0.873392) (200,0.874372)};
\addlegendentry{PPI{\textnormal{+}}{\textnormal{+}} Scalar (Flash)}

\legend{};
\end{axis}
\begin{axis}[
    title={Mean-squared CI width (frac. of classical)},
    xlabel={Budget},
    xlabel style={yshift=5pt},   % Moves xlabel closer to the axis
    grid=both,
    outer sep=0pt,
    enlarge x limits=false,    minor grid style={dashed, gray!50},
    major grid style={solid, gray!80},
    width=\textwidth,
    yticklabel style={/pgf/number format/fixed}, % No scientific notation
    yticklabel={\pgfmathparse{\tick*100}\pgfmathprintnumber{\pgfmathresult}\%},
    xmin=0, xmax=200,                          % Set x-axis range
    font=\scriptsize,             % Sets the base size for everything
    title style={yshift=-5pt, font=\small},    % Optional: Make title slightly bigger than the rest
    label style={font=\scriptsize}, % Font for xlabel and ylabel
    tick label style={font=\scriptsize}, % Font for the numbers on the axes
      xshift=5.9cm,
      width=6.5cm,
      height=5.5cm,
]

\addplot[color=blue, thick] coordinates {(10,0.8393697175424515) (20,0.8204322158482938) (30,0.8111769132190574) (40,0.8051062974122828) (50,0.8008182398939092) (60,0.7979173995302055) (70,0.7961997066525859) (80,0.7941863797524963) (90,0.7914922315057951) (100,0.7904646327912951) (110,0.7908260088196442) (120,0.7896023406859873) (130,0.7871118698583753) (140,0.7866320981005582) (150,0.7871743976839608) (160,0.7862906276083721) (170,0.7850717483185908) (180,0.7851825711865148) (190,0.7850028694775115) (200,0.7845046376061751)};
\addlegendentry{MultiPPI}

\addplot[color=green!50!black, thick, dashdotted] coordinates {(10,0.8834397617228924) (20,0.8558891804153179) (30,0.8398228471705688) (40,0.8292001113927214) (50,0.8248010008462136) (60,0.8186236645797176) (70,0.811686256495771) (80,0.8080079123648414) (90,0.8051968295621642) (100,0.8022687201519043) (110,0.8028975517785619) (120,0.8008412205948183) (130,0.7981046160383922) (140,0.7969028853233376) (150,0.7968657791608081) (160,0.7955830938710845) (170,0.7939316881731148) (180,0.7928099664814627) (190,0.7928675236506435) (200,0.7914193806342619)};
\addlegendentry{PPI{\textnormal{+}}{\textnormal{+}} Vector}

\addplot[color=brown!90, thick, densely dashed] coordinates {(10,0.8937688019065857) (20,0.8755820169817747) (30,0.8655121795697914) (40,0.8594951521695591) (50,0.8536337895234599) (60,0.850693616461009) (70,0.8497578675901312) (80,0.8477438292237578) (90,0.844709658642407) (100,0.8439665243975306) (110,0.8438904863654465) (120,0.843200932903508) (130,0.8418490781597662) (140,0.841393432158215) (150,0.8393645985800108) (160,0.8394308823911337) (170,0.8402557517894477) (180,0.8398177278693004) (190,0.839170583132366) (200,0.8390010564311886)};
\addlegendentry{PPI{\textnormal{+}}{\textnormal{+}} Scalar (Pro)}

\addplot[color=brown!60, thick, densely dashed] coordinates {(10,0.8604351496071132) (20,0.8498434647940092) (30,0.848428312965607) (40,0.845943023490771) (50,0.8424448415576357) (60,0.84194856464545) (70,0.8413149862998526) (80,0.8409753981011674) (90,0.840421841487838) (100,0.8401966730458358) (110,0.839181053540726) (120,0.8396164707746125) (130,0.8382583861239175) (140,0.8387871855446078) (150,0.8390075387690442) (160,0.8382604179128181) (170,0.8384797631973508) (180,0.838528264308914) (190,0.838619322456463) (200,0.8386660213207286)};
\addlegendentry{PPI{\textnormal{+}}{\textnormal{+}} Scalar (Flash)}

\legend{};
\end{axis}
\begin{axis}[
    title={Mean-squared Error (frac. of classical)},
    xlabel={Budget},
    xlabel style={yshift=5pt},   % Moves xlabel closer to the axis
    grid=both,
    outer sep=0pt,
    enlarge x limits=false,    minor grid style={dashed, gray!50},
    major grid style={solid, gray!80},
    width=\textwidth,
    yticklabel style={/pgf/number format/fixed}, % No scientific notation
    yticklabel={\pgfmathparse{\tick*100}\pgfmathprintnumber{\pgfmathresult}\%},
    xmin=0, xmax=200,                          % Set x-axis range
    font=\scriptsize,             % Sets the base size for everything
    title style={yshift=-5pt, font=\small},    % Optional: Make title slightly bigger than the rest
    label style={font=\scriptsize}, % Font for xlabel and ylabel
    tick label style={font=\scriptsize}, % Font for the numbers on the axes
      xshift=11.7cm,
      width=6.5cm,
      height=5.5cm,
            legend pos=north east,
            legend cell align={left},
            legend style={font=\tiny,inner sep=1pt,row sep=-2pt},
]

\addplot[color=blue, thick] coordinates {(10,0.9855233237477353) (20,0.9800532412030322) (30,0.9716025805386603) (40,0.9682306999194581) (50,0.9671833287906229) (60,0.9674806583863179) (70,0.9674910888723504) (80,0.9653720524127563) (90,0.9655341164810459) (100,0.9647034323069652) (110,0.9663598935042933) (120,0.9675463877223046) (130,0.9686743781260481) (140,0.9660033333779143) (150,0.9640527678637557) (160,0.9659707591641421) (170,0.9656213021260939) (180,0.9650723448751212) (190,0.9650298167337702) (200,0.9652229018265256)};
\addlegendentry{MultiPPI}

\addplot[color=green!50!black, thick, dashdotted] coordinates {(10,0.9907194664793624) (20,0.9829503902141552) (30,0.9845996343251454) (40,0.9797696013451132) (50,0.9746103064375347) (60,0.9754516243642646) (70,0.9714149960743317) (80,0.9701967201034377) (90,0.9742374178801364) (100,0.9747526843468401) (110,0.9693636568770007) (120,0.9702475238015471) (130,0.9688773030121346) (140,0.9682430045490189) (150,0.9696849908248601) (160,0.9687526690549582) (170,0.9660575711684225) (180,0.9660975137592577) (190,0.9670412355772597) (200,0.9688302457363824)};
\addlegendentry{PPI{\textnormal{+}}{\textnormal{+}} Vector}

\addplot[color=brown!90, thick, densely dashed] coordinates {(10,0.9910978370307952) (20,0.9866786274420013) (30,0.9813061290778713) (40,0.9805107703310276) (50,0.9822856242238708) (60,0.9796663737442669) (70,0.9845703701401789) (80,0.9831511550444273) (90,0.9829263045644364) (100,0.982465515016734) (110,0.98095357292486) (120,0.9824336176383369) (130,0.9846037950630674) (140,0.9807708924412083) (150,0.9822683903579161) (160,0.9826207831515937) (170,0.9820177248636136) (180,0.9824811558659589) (190,0.9794791196576359) (200,0.9778352369905925)};
\addlegendentry{PPI{\textnormal{+}}{\textnormal{+}} Scalar (Pro)}

\addplot[color=brown!60, thick, densely dashed] coordinates {(10,0.9907475707940188) (20,0.9882060823219496) (30,0.9838207134353095) (40,0.9825688642570305) (50,0.9797548111203452) (60,0.9801463686570374) (70,0.9790859565287556) (80,0.978231528940719) (90,0.9824938594750814) (100,0.9826049049156026) (110,0.9829501530458838) (120,0.9823130167863096) (130,0.9837958722705573) (140,0.9822752760602019) (150,0.9830867352459651) (160,0.9818789148571772) (170,0.9822650467976909) (180,0.9796811555207197) (190,0.9843210702872266) (200,0.9827099274811516)};
\addlegendentry{PPI{\textnormal{+}}{\textnormal{+}} Scalar (Flash)}

\end{axis}
\end{tikzpicture}
}
    \textrm{\bf \scriptsize Experiment 2: ProcessBench} 
\scalebox{0.8}{
\centering
\begin{tikzpicture}
\begin{axis}[
    title={Coverage},
    xlabel={Budget},
    xlabel style={yshift=5pt},   % Moves xlabel closer to the axis
    grid=both,
    outer sep=0pt,
    enlarge x limits=false,    minor grid style={dashed, gray!50},
    major grid style={solid, gray!80},
    width=\textwidth,
    yticklabel style={/pgf/number format/fixed}, % No scientific notation
    yticklabel={\pgfmathparse{\tick*100}\pgfmathprintnumber{\pgfmathresult}\%},
    xmin=0, xmax=200,                          % Set x-axis range
    font=\scriptsize,             % Sets the base size for everything
    title style={yshift=-5pt, font=\small},    % Optional: Make title slightly bigger than the rest
    label style={font=\scriptsize}, % Font for xlabel and ylabel
    tick label style={font=\scriptsize}, % Font for the numbers on the axes
      xshift=0.0cm,
      width=6.5cm,
      height=5.5cm,
            ymin=0.8, ymax=1,
]

\addplot[color=black, dashed, very thick] coordinates {(0,.95) (5000,.95)};
\addplot[color=gray!10!black!100, thick, densely dashed] coordinates {(10,0.899574) (20,0.893248) (30,0.892498) (40,0.893264) (50,0.892754) (60,0.89441) (70,0.89518) (80,0.895) (90,0.893236) (100,0.894826) (110,0.899468) (120,0.897112) (130,0.897938) (140,0.898114) (150,0.897872) (160,0.898996) (170,0.897114) (180,0.898154) (190,0.898636) (200,0.897956)};
\addlegendentry{PPI{\textnormal{+}}{\textnormal{+}} Scalar Large}

\addplot[color=gray!40!black!70, thick, densely dashed] coordinates {(10,0.892556) (20,0.888528) (30,0.892634) (40,0.892796) (50,0.88822) (60,0.888756) (70,0.893634) (80,0.892092) (90,0.891974) (100,0.890734) (110,0.889896) (120,0.890796) (130,0.888478) (140,0.888334) (150,0.888986) (160,0.887274) (170,0.886534) (180,0.887408) (190,0.88584) (200,0.885064)};
\addlegendentry{PPI{\textnormal{+}}{\textnormal{+}} Scalar Medium}

\addplot[color=gray!70!black!40, thick, densely dashed] coordinates {(10,0.892302) (20,0.885658) (30,0.883002) (40,0.87833) (50,0.87638) (60,0.873638) (70,0.86693) (80,0.865634) (90,0.86137) (100,0.85949) (110,0.856656) (120,0.853766) (130,0.848282) (140,0.845542) (150,0.843048) (160,0.839782) (170,0.83538) (180,0.834422) (190,0.828948) (200,0.826046)};
\addlegendentry{PPI{\textnormal{+}}{\textnormal{+}} Scalar Small}

\addplot[color=gray!100!black!10, thick, densely dashed] coordinates {(10,0.888274) (20,0.884712) (30,0.881596) (40,0.878812) (50,0.872928) (60,0.87027) (70,0.864766) (80,0.86072) (90,0.856894) (100,0.853452) (110,0.8491) (120,0.848256) (130,0.841118) (140,0.83849) (150,0.834802) (160,0.832106) (170,0.827642) (180,0.826128) (190,0.823434) (200,0.820482)};
\addlegendentry{PPI{\textnormal{+}}{\textnormal{+}} Scalar Tiny}

\addplot[color=blue, thick] coordinates {(10,0.87482) (20,0.867944) (30,0.864986) (40,0.857936) (50,0.851732) (60,0.847976) (70,0.845002) (80,0.842446) (90,0.842418) (100,0.839638) (110,0.839062) (120,0.837342) (130,0.836672) (140,0.835694) (150,0.830906) (160,0.828954) (170,0.82926) (180,0.827018) (190,0.828566) (200,0.827634)};
\addlegendentry{MultiPPI}

\addplot[color=green!50!black, thick, dashdotted] coordinates {(10,0.894754) (20,0.886364) (30,0.884224) (40,0.882416) (50,0.879758) (60,0.877798) (70,0.881104) (80,0.878314) (90,0.87751) (100,0.876274) (110,0.873096) (120,0.872804) (130,0.872412) (140,0.871984) (150,0.869156) (160,0.868642) (170,0.868692) (180,0.868028) (190,0.866748) (200,0.865522)};
\addlegendentry{PPI{\textnormal{+}}{\textnormal{+}} Vector}

\legend{};
\end{axis}
\begin{axis}[
    title={Mean-squared CI width (frac. of classical)},
    xlabel={Budget},
    xlabel style={yshift=5pt},   % Moves xlabel closer to the axis
    grid=both,
    outer sep=0pt,
    enlarge x limits=false,    minor grid style={dashed, gray!50},
    major grid style={solid, gray!80},
    width=\textwidth,
    yticklabel style={/pgf/number format/fixed}, % No scientific notation
    yticklabel={\pgfmathparse{\tick*100}\pgfmathprintnumber{\pgfmathresult}\%},
    xmin=0, xmax=200,                          % Set x-axis range
    font=\scriptsize,             % Sets the base size for everything
    title style={yshift=-5pt, font=\small},    % Optional: Make title slightly bigger than the rest
    label style={font=\scriptsize}, % Font for xlabel and ylabel
    tick label style={font=\scriptsize}, % Font for the numbers on the axes
      xshift=5.9cm,
      width=6.5cm,
      height=5.5cm,
]

\addplot[color=gray!10!black!100, thick, densely dashed] coordinates {(10,0.7604037701150433) (20,0.6834775645750324) (30,0.6368174532475046) (40,0.6076747560020481) (50,0.5855986329222285) (60,0.5696095362403844) (70,0.5582953270783699) (80,0.5475686157257844) (90,0.5387817653075705) (100,0.5315942277267944) (110,0.5247122499707297) (120,0.5192989030090932) (130,0.5141004366335763) (140,0.5098828886758607) (150,0.5066102090022324) (160,0.5027781170015786) (170,0.4999442951125094) (180,0.4972768424482151) (190,0.4940495237332589) (200,0.4916047509924428)};
\addlegendentry{PPI{\textnormal{+}}{\textnormal{+}} Scalar Large}

\addplot[color=gray!40!black!70, thick, densely dashed] coordinates {(10,0.7460089706686795) (20,0.6825320010105506) (30,0.649568233567738) (40,0.6253190331877653) (50,0.6090827874039473) (60,0.597051411898955) (70,0.5878244830923177) (80,0.5801777898443033) (90,0.5740583885103815) (100,0.5676449247957128) (110,0.563273186239766) (120,0.5598678051271274) (130,0.5551901523814461) (140,0.552119388285241) (150,0.5511467547232968) (160,0.5477353101359859) (170,0.5465191250321427) (180,0.5436442191924445) (190,0.5409611395313437) (200,0.5385970710946318)};
\addlegendentry{PPI{\textnormal{+}}{\textnormal{+}} Scalar Medium}

\addplot[color=gray!70!black!40, thick, densely dashed] coordinates {(10,0.741479026450834) (20,0.6860585805917642) (30,0.6578731812837224) (40,0.641527879271075) (50,0.6309138887091591) (60,0.6230168104987097) (70,0.6164402557094593) (80,0.6115064048239831) (90,0.6057439944812233) (100,0.6025574087777384) (110,0.600097733222985) (120,0.5975760255119518) (130,0.5959992916153244) (140,0.5935258727884744) (150,0.5921615233905756) (160,0.5903194596051635) (170,0.5893095552849221) (180,0.58733957238333) (190,0.5856436818910817) (200,0.5853157020532103)};
\addlegendentry{PPI{\textnormal{+}}{\textnormal{+}} Scalar Small}

\addplot[color=gray!100!black!10, thick, densely dashed] coordinates {(10,0.6989762936180757) (20,0.6585247194245739) (30,0.6398374886181891) (40,0.6285134086349774) (50,0.6212483131008404) (60,0.6162493767763944) (70,0.612597110389704) (80,0.6102337546334241) (90,0.6072316152277067) (100,0.6053229507497871) (110,0.6038796185435296) (120,0.6021414980564077) (130,0.6002783285543326) (140,0.5988830226658992) (150,0.5979007301969708) (160,0.5974038099852663) (170,0.5957224234921296) (180,0.5953641012004142) (190,0.594242425432665) (200,0.5945593229041849)};
\addlegendentry{PPI{\textnormal{+}}{\textnormal{+}} Scalar Tiny}

\addplot[color=blue, thick] coordinates {(10,0.620200467364889) (20,0.5347513693434114) (30,0.4883530007872466) (40,0.4589218928146549) (50,0.43893396985941135) (60,0.4235797363702762) (70,0.4108934496825846) (80,0.4011453536971576) (90,0.39195015243713194) (100,0.38401300783543296) (110,0.37917519317799775) (120,0.37361428185607365) (130,0.3686158450588804) (140,0.3641511059227025) (150,0.3602271779787554) (160,0.35621874538864473) (170,0.353701467782074) (180,0.3501868676778684) (190,0.3470437141107698) (200,0.34414462074299396)};
\addlegendentry{MultiPPI}

\addplot[color=green!50!black, thick, dashdotted] coordinates {(10,0.7769101889283802) (20,0.6989894285958996) (30,0.6446760490925153) (40,0.6045369932798864) (50,0.5728021518545097) (60,0.5482066599399847) (70,0.5284159402617903) (80,0.5082944683545394) (90,0.4940211718910014) (100,0.48099233951835624) (110,0.46976394403380345) (120,0.45985334591958976) (130,0.4528213831633346) (140,0.4445028889126892) (150,0.4366303190484803) (160,0.42923462427959685) (170,0.4225607123432423) (180,0.41750680868741347) (190,0.4125284126633244) (200,0.40755108847309246)};
\addlegendentry{PPI{\textnormal{+}}{\textnormal{+}} Vector}

\legend{};
\end{axis}
\begin{axis}[
    title={Mean-squared Error (frac. of classical)},
    xlabel={Budget},
    xlabel style={yshift=5pt},   % Moves xlabel closer to the axis
    grid=both,
    outer sep=0pt,
    enlarge x limits=false,    minor grid style={dashed, gray!50},
    major grid style={solid, gray!80},
    width=\textwidth,
    yticklabel style={/pgf/number format/fixed}, % No scientific notation
    yticklabel={\pgfmathparse{\tick*100}\pgfmathprintnumber{\pgfmathresult}\%},
    xmin=0, xmax=200,                          % Set x-axis range
    font=\scriptsize,             % Sets the base size for everything
    title style={yshift=-5pt, font=\small},    % Optional: Make title slightly bigger than the rest
    label style={font=\scriptsize}, % Font for xlabel and ylabel
    tick label style={font=\scriptsize}, % Font for the numbers on the axes
      xshift=11.7cm,
      width=6.5cm,
      height=5.5cm,
            legend pos=north east,
            legend cell align={left},
            legend style={font=\tiny,inner sep=1pt,row sep=-2pt},
]

\addplot[color=gray!10!black!100, thick, densely dashed] coordinates {(10,0.9017477078613161) (20,0.8523306567847153) (30,0.8311860508549418) (40,0.8117760770474213) (50,0.8026619249001475) (60,0.7929300357814594) (70,0.7828405573798843) (80,0.7762912696301082) (90,0.7719614221390926) (100,0.7690849079699311) (110,0.7629128531090615) (120,0.7623209555430103) (130,0.7596063700566271) (140,0.7577390187813069) (150,0.7582554305067342) (160,0.7572847119451765) (170,0.7540621370425076) (180,0.7523548455711452) (190,0.7518931996725428) (200,0.7519480132197011)};
\addlegendentry{PPI{\textnormal{+}}{\textnormal{+}} Scalar Large}

\addplot[color=gray!40!black!70, thick, densely dashed] coordinates {(10,0.9056601672608945) (20,0.8676218860169075) (30,0.8434945864839168) (40,0.8296505600228087) (50,0.8266811958092154) (60,0.8216851898610502) (70,0.8097580008277128) (80,0.807714585083752) (90,0.8005898292078357) (100,0.8004029825251561) (110,0.7975687824545952) (120,0.797002943508103) (130,0.7945247225487396) (140,0.7919482132251456) (150,0.7930221876074475) (160,0.7923696950120531) (170,0.7894542828436474) (180,0.7878605310607619) (190,0.7887309078537649) (200,0.787809554061553)};
\addlegendentry{PPI{\textnormal{+}}{\textnormal{+}} Scalar Medium}

\addplot[color=gray!70!black!40, thick, densely dashed] coordinates {(10,0.9033410360831637) (20,0.871543829307358) (30,0.8566095874711541) (40,0.8483638545542765) (50,0.841767544298797) (60,0.8379668417501583) (70,0.8353351844360142) (80,0.8322380119743988) (90,0.8320089787696717) (100,0.8292856827937366) (110,0.8252446794535441) (120,0.8253183673210267) (130,0.8280095024588456) (140,0.827662526533966) (150,0.8265399036102744) (160,0.825478624557103) (170,0.8264570478449931) (180,0.8225166929897811) (190,0.8283035810722282) (200,0.8283816259658607)};
\addlegendentry{PPI{\textnormal{+}}{\textnormal{+}} Scalar Small}

\addplot[color=gray!100!black!10, thick, densely dashed] coordinates {(10,0.8855890781132787) (20,0.8640184134223184) (30,0.8526004475802046) (40,0.8465385920837519) (50,0.8439730136162029) (60,0.8416527882933189) (70,0.8402243290433781) (80,0.8400345227593548) (90,0.839545123777051) (100,0.8379224533381257) (110,0.8371922545743491) (120,0.8348973720736304) (130,0.8381296606123672) (140,0.8364306421378538) (150,0.8374093656409336) (160,0.8375439662220576) (170,0.8359416412217776) (180,0.8343006587260656) (190,0.8348147072358469) (200,0.8341307605259535)};
\addlegendentry{PPI{\textnormal{+}}{\textnormal{+}} Scalar Tiny}

\addplot[color=blue, thick] coordinates {(10,0.7974717593341353) (20,0.7324391365562989) (30,0.6863866988249859) (40,0.6646764640587609) (50,0.6575773693765966) (60,0.6470585621379623) (70,0.6362635084335677) (80,0.6282894999127835) (90,0.6225640092562202) (100,0.6180210274627588) (110,0.6134651337582745) (120,0.6117928498427966) (130,0.6079203171026024) (140,0.6041574070362893) (150,0.6032595017290512) (160,0.6015980664675182) (170,0.6001858788544742) (180,0.5982019410761285) (190,0.5946630140457768) (200,0.5936647427555224)};
\addlegendentry{MultiPPI}

\addplot[color=green!50!black, thick, dashdotted] coordinates {(10,0.8895079985306005) (20,0.8260324866577812) (30,0.7855108873854014) (40,0.7552809312949279) (50,0.7385110230435312) (60,0.721503374271584) (70,0.6996375969578933) (80,0.6882010117138563) (90,0.6740660290874722) (100,0.6651949482151669) (110,0.6611670101780066) (120,0.6531128405150349) (130,0.6492907615023116) (140,0.6437046327252014) (150,0.6410063416243277) (160,0.6355451252180763) (170,0.6327850129454623) (180,0.6286502398981799) (190,0.6234719277344478) (200,0.6202823334456746)};
\addlegendentry{PPI{\textnormal{+}}{\textnormal{+}} Vector}

\end{axis}
\end{tikzpicture}
}
    \textrm{\bf \scriptsize Experiment 3: Factuality}
\scalebox{0.8}{
\centering
\begin{tikzpicture}
\begin{axis}[
    title={Coverage},
    xlabel={Budget},
    xlabel style={yshift=5pt},   % Moves xlabel closer to the axis
    grid=both,
    outer sep=0pt,
    enlarge x limits=false,    minor grid style={dashed, gray!50},
    major grid style={solid, gray!80},
    width=\textwidth,
    yticklabel style={/pgf/number format/fixed}, % No scientific notation
    yticklabel={\pgfmathparse{\tick*100}\pgfmathprintnumber{\pgfmathresult}\%},
    xmin=0, xmax=200,                          % Set x-axis range
    font=\scriptsize,             % Sets the base size for everything
    title style={yshift=-5pt, font=\small},    % Optional: Make title slightly bigger than the rest
    label style={font=\scriptsize}, % Font for xlabel and ylabel
    tick label style={font=\scriptsize}, % Font for the numbers on the axes
      xshift=0.0cm,
      width=6.5cm,
      height=5.5cm,
            ymin=0.7, ymax=1,
]

\addplot[color=black, dashed, very thick] coordinates {(0,.95) (5000,.95)};
\addplot[color=blue, thick] coordinates {(10,0.8483) (20,0.819094) (30,0.797788) (40,0.78003) (50,0.764996) (60,0.754926) (70,0.746254) (80,0.740546) (90,0.734262) (100,0.729058) (110,0.724816) (120,0.719848) (130,0.716736) (140,0.713824) (150,0.71249) (160,0.71056) (170,0.707674) (180,0.70743) (190,0.705998) (200,0.703272)};
\addlegendentry{MultiPPI}

\addplot[color=green!50!black, thick, dashdotted] coordinates {(10,0.88565) (20,0.882406) (30,0.87538) (40,0.866976) (50,0.861842) (60,0.853644) (70,0.846172) (80,0.840448) (90,0.83648) (100,0.830052) (110,0.823302) (120,0.818444) (130,0.81595) (140,0.810288) (150,0.807698) (160,0.801876) (170,0.798934) (180,0.793962) (190,0.788216) (200,0.783166)};
\addlegendentry{PPI{\textnormal{+}}{\textnormal{+}} Vector}

\addplot[color=red!40,  thick, densely dotted] coordinates {(10,0.853362) (20,0.84435) (30,0.838634) (40,0.8351) (50,0.832286) (60,0.830664) (70,0.828802) (80,0.828292) (90,0.824666) (100,0.825966) (110,0.8238) (120,0.823354) (130,0.82294) (140,0.822748) (150,0.82217) (160,0.822252) (170,0.821894) (180,0.821066) (190,0.821136) (200,0.819572)};
\addlegendentry{PPI{\textnormal{+}}{\textnormal{+}} Scalar 1 agent, 1 round}

\addplot[color=red!70,  thick, densely dotted] coordinates {(10,0.855776) (20,0.838458) (30,0.827474) (40,0.816926) (50,0.809636) (60,0.80644) (70,0.80434) (80,0.801706) (90,0.799934) (100,0.799404) (110,0.799842) (120,0.798646) (130,0.797446) (140,0.79737) (150,0.796858) (160,0.796818) (170,0.796892) (180,0.796754) (190,0.797328) (200,0.79626)};
\addlegendentry{PPI{\textnormal{+}}{\textnormal{+}} Scalar 2 agents, 1 round}

\addplot[color=red,  thick, densely dotted] coordinates {(10,0.862038) (20,0.845396) (30,0.833776) (40,0.82544) (50,0.818896) (60,0.810712) (70,0.806574) (80,0.80176) (90,0.798966) (100,0.797024) (110,0.79698) (120,0.794968) (130,0.793656) (140,0.792858) (150,0.792396) (160,0.791406) (170,0.792232) (180,0.790818) (190,0.79108) (200,0.791002)};
\addlegendentry{PPI{\textnormal{+}}{\textnormal{+}} Scalar 3 agents, 1 round}

\addplot[color=red!40,  thick, densely dashed] coordinates {(10,0.851326) (20,0.82934) (30,0.813106) (40,0.804892) (50,0.799064) (60,0.79493) (70,0.793772) (80,0.792194) (90,0.79131) (100,0.792184) (110,0.79084) (120,0.791604) (130,0.7904) (140,0.790168) (150,0.789942) (160,0.790324) (170,0.790418) (180,0.79046) (190,0.78844) (200,0.788676)};
\addlegendentry{PPI{\textnormal{+}}{\textnormal{+}} Scalar 1 agent, 2 rounds}

\addplot[color=red!70,  thick, densely dashed] coordinates {(10,0.865972) (20,0.848584) (30,0.837552) (40,0.825792) (50,0.813638) (60,0.806778) (70,0.802464) (80,0.79792) (90,0.79238) (100,0.789024) (110,0.785536) (120,0.783216) (130,0.783844) (140,0.78183) (150,0.78137) (160,0.77938) (170,0.77849) (180,0.77949) (190,0.779426) (200,0.779392)};
\addlegendentry{PPI{\textnormal{+}}{\textnormal{+}} Scalar 2 agents, 2 rounds}

\addplot[color=red,  thick, densely dashed] coordinates {(10,0.87004) (20,0.860526) (30,0.851572) (40,0.841502) (50,0.830742) (60,0.821916) (70,0.816058) (80,0.810188) (90,0.806898) (100,0.800914) (110,0.795338) (120,0.79163) (130,0.790562) (140,0.787544) (150,0.784616) (160,0.783556) (170,0.781886) (180,0.781494) (190,0.778336) (200,0.776654)};
\addlegendentry{PPI{\textnormal{+}}{\textnormal{+}} Scalar 3 agents, 2 rounds}

\legend{};
\end{axis}
\begin{axis}[
    title={Mean-squared CI width (frac. of classical)},
    xlabel={Budget},
    xlabel style={yshift=5pt},   % Moves xlabel closer to the axis
    grid=both,
    outer sep=0pt,
    enlarge x limits=false,    minor grid style={dashed, gray!50},
    major grid style={solid, gray!80},
    width=\textwidth,
    yticklabel style={/pgf/number format/fixed}, % No scientific notation
    yticklabel={\pgfmathparse{\tick*100}\pgfmathprintnumber{\pgfmathresult}\%},
    xmin=0, xmax=200,                          % Set x-axis range
    font=\scriptsize,             % Sets the base size for everything
    title style={yshift=-5pt, font=\small},    % Optional: Make title slightly bigger than the rest
    label style={font=\scriptsize}, % Font for xlabel and ylabel
    tick label style={font=\scriptsize}, % Font for the numbers on the axes
      xshift=5.9cm,
      width=6.5cm,
      height=5.5cm,
]

\addplot[color=blue, thick] coordinates {(10,0.6530159350195583) (20,0.6093148512569959) (30,0.5885407990484566) (40,0.5745968557398743) (50,0.5642047740040436) (60,0.5574790037273248) (70,0.5522686261739675) (80,0.5482241709227709) (90,0.5443816369081783) (100,0.5412661504488766) (110,0.5389758986382105) (120,0.535697166418024) (130,0.5339206802871236) (140,0.5323968783195645) (150,0.5305663452358521) (160,0.5298873036420335) (170,0.5290229772797255) (180,0.5277008139941646) (190,0.5257848851761715) (200,0.5245958310803884)};
\addlegendentry{MultiPPI}

\addplot[color=green!50!black, thick, dashdotted] coordinates {(10,0.8256209101839945) (20,0.7615094493633643) (30,0.7195883540086702) (40,0.6909038213510774) (50,0.6711352666582738) (60,0.6542193625213788) (70,0.6401729571561857) (80,0.6293462426382511) (90,0.6215127815249903) (100,0.6136552930589911) (110,0.606649435172469) (120,0.6003812563598772) (130,0.5946551007179078) (140,0.5902917681723284) (150,0.587306599000811) (160,0.5832341761014774) (170,0.5811013691180461) (180,0.5772843321272988) (190,0.5733742329311398) (200,0.5708667674312328)};
\addlegendentry{PPI{\textnormal{+}}{\textnormal{+}} Vector}

\addplot[color=red!40,  thick, densely dotted] coordinates {(10,0.7498682441089992) (20,0.7367795755814678) (30,0.7317136952967208) (40,0.7293646343588702) (50,0.7281954889620114) (60,0.727499220590288) (70,0.7261866937511618) (80,0.7252854248596179) (90,0.7239935538808693) (100,0.724800409502794) (110,0.7243346280553937) (120,0.72418664143499) (130,0.7236550114992666) (140,0.7235451366679188) (150,0.7234326523418364) (160,0.722955888398972) (170,0.7226739663413101) (180,0.7229484082799995) (190,0.7223678950583043) (200,0.7221028105488899)};
\addlegendentry{PPI{\textnormal{+}}{\textnormal{+}} Scalar 1 agent, 1 round}

\addplot[color=red!70,  thick, densely dotted] coordinates {(10,0.7384280921519144) (20,0.7112510556142422) (30,0.6999797471628534) (40,0.6935852821474642) (50,0.689927039000659) (60,0.6876916453268395) (70,0.6852673962776346) (80,0.6834881251278062) (90,0.6842730412987149) (100,0.6830486819348985) (110,0.6811768321356471) (120,0.6814005245178553) (130,0.6803994213076349) (140,0.6802046040472625) (150,0.6785376967933031) (160,0.677815964831827) (170,0.6769526533075223) (180,0.6777737198187849) (190,0.6776617956834842) (200,0.6774280754155405)};
\addlegendentry{PPI{\textnormal{+}}{\textnormal{+}} Scalar 2 agents, 1 round}

\addplot[color=red,  thick, densely dotted] coordinates {(10,0.7529282422477855) (20,0.7187478184797984) (30,0.7029954854154825) (40,0.694456759637327) (50,0.6882640811331814) (60,0.6848094978413092) (70,0.6810314258203293) (80,0.678991958785667) (90,0.6758179000946297) (100,0.6744759651871448) (110,0.6750108673374731) (120,0.6734319719115834) (130,0.6731348787422213) (140,0.6725798272186692) (150,0.6707631579405362) (160,0.6701645117837156) (170,0.6698440618250299) (180,0.6689678983423462) (190,0.6694488925343928) (200,0.6682165086179443)};
\addlegendentry{PPI{\textnormal{+}}{\textnormal{+}} Scalar 3 agents, 1 round}

\addplot[color=red!40,  thick, densely dashed] coordinates {(10,0.7533806865392823) (20,0.7283298698865082) (30,0.7167308805803659) (40,0.7108940686401307) (50,0.7079688010023516) (60,0.7056515366863989) (70,0.7023269606183908) (80,0.7011964141326789) (90,0.7007913382786946) (100,0.7003699763435677) (110,0.6986285411228609) (120,0.6983002420916173) (130,0.6972744231943896) (140,0.6969102786855115) (150,0.6967797561579103) (160,0.6964948893875593) (170,0.6956394224976912) (180,0.6952010930408576) (190,0.6944885336152418) (200,0.6944729409100514)};
\addlegendentry{PPI{\textnormal{+}}{\textnormal{+}} Scalar 1 agent, 2 rounds}

\addplot[color=red!70,  thick, densely dashed] coordinates {(10,0.7769389861210382) (20,0.7397208471006786) (30,0.7217786322799141) (40,0.7111518878463372) (50,0.705071132526155) (60,0.7006943571328351) (70,0.6962434858579297) (80,0.6931499488603214) (90,0.692579932591043) (100,0.6907738459447323) (110,0.6894502195477149) (120,0.6879385201616055) (130,0.6867995389253889) (140,0.6857455109332334) (150,0.6842548969458355) (160,0.6835399726101318) (170,0.6828221916921523) (180,0.6826530862838183) (190,0.682027473090026) (200,0.682106624761564)};
\addlegendentry{PPI{\textnormal{+}}{\textnormal{+}} Scalar 2 agents, 2 rounds}

\addplot[color=red,  thick, densely dashed] coordinates {(10,0.8001999923715719) (20,0.7582804170075924) (30,0.7389194317464699) (40,0.7258144823749508) (50,0.7168819304093401) (60,0.7101554879223969) (70,0.7050226061876348) (80,0.701365428851182) (90,0.7006692966417221) (100,0.6973481615349448) (110,0.6947015671793896) (120,0.6933769145546895) (130,0.6915332588652331) (140,0.6899027718030242) (150,0.6885246953984662) (160,0.6875568969845044) (170,0.6869556511225504) (180,0.6862806751657384) (190,0.6851773623600161) (200,0.6839136830128844)};
\addlegendentry{PPI{\textnormal{+}}{\textnormal{+}} Scalar 3 agents, 2 rounds}

\legend{};
\end{axis}
\begin{axis}[
    title={Mean-squared Error (frac. of classical)},
    xlabel={Budget},
    xlabel style={yshift=5pt},   % Moves xlabel closer to the axis
    grid=both,
    outer sep=0pt,
    enlarge x limits=false,    minor grid style={dashed, gray!50},
    major grid style={solid, gray!80},
    width=\textwidth,
    yticklabel style={/pgf/number format/fixed}, % No scientific notation
    yticklabel={\pgfmathparse{\tick*100}\pgfmathprintnumber{\pgfmathresult}\%},
    xmin=0, xmax=200,                          % Set x-axis range
    font=\scriptsize,             % Sets the base size for everything
    title style={yshift=-5pt, font=\small},    % Optional: Make title slightly bigger than the rest
    label style={font=\scriptsize}, % Font for xlabel and ylabel
    tick label style={font=\scriptsize}, % Font for the numbers on the axes
      xshift=11.7cm,
      width=6.5cm,
      height=5.5cm,
            legend pos=north east,
            legend cell align={left},
            legend style={font=\tiny,inner sep=1pt,row sep=-2pt},
]

\addplot[color=blue, thick] coordinates {(10,0.8402818260306603) (20,0.8197756276185428) (30,0.8112039350361567) (40,0.8087442317747134) (50,0.8036881038309551) (60,0.8018578561459611) (70,0.8021920726354884) (80,0.800378302033071) (90,0.7990428015881574) (100,0.7996393318935692) (110,0.8015037418318643) (120,0.800557004095866) (130,0.799445681820447) (140,0.7986025681533975) (150,0.7969353790014108) (160,0.7966284121838013) (170,0.7990665347567754) (180,0.7959962630399583) (190,0.7939700802160584) (200,0.7956883703063046)};
\addlegendentry{MultiPPI}

\addplot[color=green!50!black, thick, dashdotted] coordinates {(10,0.9372965808407719) (20,0.8983824281444172) (30,0.874236602395182) (40,0.8597115699897763) (50,0.8508585346519404) (60,0.8446473029901482) (70,0.841497265767906) (80,0.8365600615415592) (90,0.8284746017670033) (100,0.8266411968901303) (110,0.8272151439685005) (120,0.8244361651670751) (130,0.8177598527387161) (140,0.8176492783715785) (150,0.8171510746916238) (160,0.8171944673520487) (170,0.8126065964018458) (180,0.812059311754402) (190,0.811468362849966) (200,0.812446817874578)};
\addlegendentry{PPI{\textnormal{+}}{\textnormal{+}} Vector}

\addplot[color=red!40,  thick, densely dotted] coordinates {(10,0.9444537645307376) (20,0.940333545007819) (30,0.9392010379029183) (40,0.9386951406882821) (50,0.9374794823836593) (60,0.9368218636857824) (70,0.934907690567548) (80,0.9334643667775314) (90,0.9370609873656306) (100,0.9336039470442354) (110,0.9352922500058619) (120,0.9361873187511447) (130,0.9329264537250468) (140,0.9343622539954655) (150,0.9353666870197483) (160,0.9345065954324725) (170,0.9358864473975481) (180,0.9352999629411106) (190,0.9351283438471252) (200,0.935428577068748)};
\addlegendentry{PPI{\textnormal{+}}{\textnormal{+}} Scalar 1 agent, 1 round}

\addplot[color=red!70,  thick, densely dotted] coordinates {(10,0.92315464126594) (20,0.9104362080450542) (30,0.9049406728730934) (40,0.9057264243721646) (50,0.9043200583833868) (60,0.9027287101607805) (70,0.9005610942098446) (80,0.901933365366832) (90,0.9043738304711848) (100,0.9018782096996542) (110,0.900137982740625) (120,0.9004851881540998) (130,0.9021891133117521) (140,0.9026994325894527) (150,0.8999836851097667) (160,0.8984952587968555) (170,0.8980313399189371) (180,0.9007289164152231) (190,0.899630684124246) (200,0.9014845509628511)};
\addlegendentry{PPI{\textnormal{+}}{\textnormal{+}} Scalar 2 agents, 1 round}

\addplot[color=red,  thick, densely dotted] coordinates {(10,0.9289515532363204) (20,0.9129424476147066) (30,0.905608705373494) (40,0.9010081079220001) (50,0.8973373475354665) (60,0.8987429225692323) (70,0.8965538206995348) (80,0.898159386752027) (90,0.8945813294697242) (100,0.8952717125221087) (110,0.8930632551592376) (120,0.8936493492012115) (130,0.8951163315481313) (140,0.8954433576552722) (150,0.8918548917301523) (160,0.8931765479899142) (170,0.8938755441785964) (180,0.8920039226163952) (190,0.8938870331357404) (200,0.8916700936821903)};
\addlegendentry{PPI{\textnormal{+}}{\textnormal{+}} Scalar 3 agents, 1 round}

\addplot[color=red!40,  thick, densely dashed] coordinates {(10,0.9344049886395556) (20,0.9226875191007586) (30,0.9163216478391047) (40,0.9126764808982272) (50,0.91425018194041) (60,0.9140505841208225) (70,0.9104326080891904) (80,0.9112248537650972) (90,0.9124728572722204) (100,0.9102142328619839) (110,0.9108647299910934) (120,0.909521744768859) (130,0.9105314913159996) (140,0.9108417662523121) (150,0.9106087066327003) (160,0.9109106641095426) (170,0.9093345135263595) (180,0.9092802161286864) (190,0.9112300529202804) (200,0.9109462823487524)};
\addlegendentry{PPI{\textnormal{+}}{\textnormal{+}} Scalar 1 agent, 2 rounds}

\addplot[color=red!70,  thick, densely dashed] coordinates {(10,0.9391129115632306) (20,0.9234440230727463) (30,0.9125706882626938) (40,0.9078542024241256) (50,0.9090647372520264) (60,0.9079323763890091) (70,0.9037463158201677) (80,0.903745147930805) (90,0.903743604442422) (100,0.9041514710738933) (110,0.9046868874075771) (120,0.9043713536915134) (130,0.9014020117443058) (140,0.9017162788014701) (150,0.9005933657712972) (160,0.9016592154198587) (170,0.8998831170603869) (180,0.8990251404946047) (190,0.8989180278263007) (200,0.9009003307958694)};
\addlegendentry{PPI{\textnormal{+}}{\textnormal{+}} Scalar 2 agents, 2 rounds}

\addplot[color=red,  thick, densely dashed] coordinates {(10,0.9517495182411896) (20,0.9306055700633229) (30,0.9184550295904318) (40,0.9112817966527086) (50,0.9123019039737074) (60,0.9103593993241759) (70,0.9069189051251753) (80,0.9061846788787609) (90,0.9070774159142678) (100,0.9066750561326186) (110,0.9038159719988187) (120,0.9047617556732512) (130,0.903572347886444) (140,0.9047134531099894) (150,0.9023566783870449) (160,0.9008735783058744) (170,0.9019102071957558) (180,0.902142024101357) (190,0.8991448228207706) (200,0.9027798086152858)};
\addlegendentry{PPI{\textnormal{+}}{\textnormal{+}} Scalar 3 agents, 2 rounds}

\end{axis}
\end{tikzpicture}
}
    \vspace{-10pt}
    \caption{Results given only $N=10$ labeled examples. Results are shown by budget for the experiments on Chatbot Arena (top), ProcessBench (middle), and Factuality (bottom). For each estimator (all baselines and MultiPPI), the left column plots the empirical coverage of the 95\% CI, the middle column plots the width of the 95\% CI, and the right column plots the empirical mean-squared error of the point estimate. \looseness=-1}
    \label{fig:10}
    \vspace{-10pt}
\end{figure}

\begin{figure}[!t]
    \centering
    %\includegraphics[width=\linewidth]{figures/chatbot_arena_results2.png}
    \textrm{\bf \scriptsize Experiment 1: Chatbot Arena}
\scalebox{.8}{\centering
\begin{tikzpicture}
\begin{axis}[
    title={Coverage},
    xlabel={Budget},
    xlabel style={yshift=5pt},   % Moves xlabel closer to the axis
    grid=both,
    outer sep=0pt,
    enlarge x limits=false,    minor grid style={dashed, gray!50},
    major grid style={solid, gray!80},
    width=\textwidth,
    yticklabel style={/pgf/number format/fixed}, % No scientific notation
    yticklabel={\pgfmathparse{\tick*100}\pgfmathprintnumber{\pgfmathresult}\%},
    xmin=0, xmax=20000,                          % Set x-axis range
    font=\scriptsize,             % Sets the base size for everything
    title style={yshift=-5pt, font=\small},    % Optional: Make title slightly bigger than the rest
    label style={font=\scriptsize}, % Font for xlabel and ylabel
    tick label style={font=\scriptsize}, % Font for the numbers on the axes
      xshift=0.0cm,
      width=6.5cm,
      height=5.5cm,
            ymin=0.9, ymax=1,
]

\addplot[color=black, dashed, very thick] coordinates {(0,.95) (5000,.95)};
\addplot[color=blue, thick] coordinates {(1000,0.949274) (2000,0.949068) (3000,0.949544) (4000,0.949628) (5000,0.949098) (6000,0.94912) (7000,0.949168) (8000,0.949224) (9000,0.949262) (10000,0.94917) (11000,0.9495) (12000,0.949082) (13000,0.948256) (14000,0.948504) (15000,0.94937) (16000,0.948826) (17000,0.948922) (18000,0.948998) (19000,0.94862) (20000,0.949084)};
\addlegendentry{MultiPPI}

\addplot[color=green!50!black, thick, dashdotted] coordinates {(1000,0.949586) (2000,0.949714) (3000,0.947598) (4000,0.947668) (5000,0.949256) (6000,0.949566) (7000,0.949148) (8000,0.948438) (9000,0.949978) (10000,0.94927) (11000,0.948962) (12000,0.948836) (13000,0.949516) (14000,0.949496) (15000,0.948638) (16000,0.949436) (17000,0.949308) (18000,0.948924) (19000,0.949052) (20000,0.94919)};
\addlegendentry{PPI{\textnormal{+}}{\textnormal{+}} Vector}

\addplot[color=brown!90, thick, densely dashed] coordinates {(1000,0.950554) (2000,0.950276) (3000,0.950218) (4000,0.95046) (5000,0.94905) (6000,0.949046) (7000,0.948746) (8000,0.949476) (9000,0.949514) (10000,0.949248) (11000,0.948976) (12000,0.94894) (13000,0.949134) (14000,0.949586) (15000,0.949906) (16000,0.949334) (17000,0.949342) (18000,0.949378) (19000,0.94924) (20000,0.94944)};
\addlegendentry{PPI{\textnormal{+}}{\textnormal{+}} Scalar (Pro)}

\addplot[color=brown!60, thick, densely dashed] coordinates {(1000,0.94958) (2000,0.949078) (3000,0.94932) (4000,0.949528) (5000,0.949268) (6000,0.949678) (7000,0.949804) (8000,0.948732) (9000,0.94917) (10000,0.949392) (11000,0.949014) (12000,0.949008) (13000,0.94928) (14000,0.949488) (15000,0.94961) (16000,0.94867) (17000,0.950424) (18000,0.949852) (19000,0.9498) (20000,0.9489)};
\addlegendentry{PPI{\textnormal{+}}{\textnormal{+}} Scalar (Flash)}

\legend{};
\end{axis}
\begin{axis}[
    title={Mean-squared CI width (frac. of classical)},
    xlabel={Budget},
    xlabel style={yshift=5pt},   % Moves xlabel closer to the axis
    grid=both,
    outer sep=0pt,
    enlarge x limits=false,    minor grid style={dashed, gray!50},
    major grid style={solid, gray!80},
    width=\textwidth,
    yticklabel style={/pgf/number format/fixed}, % No scientific notation
    yticklabel={\pgfmathparse{\tick*100}\pgfmathprintnumber{\pgfmathresult}\%},
    xmin=0, xmax=20000,                          % Set x-axis range
    font=\scriptsize,             % Sets the base size for everything
    title style={yshift=-5pt, font=\small},    % Optional: Make title slightly bigger than the rest
    label style={font=\scriptsize}, % Font for xlabel and ylabel
    tick label style={font=\scriptsize}, % Font for the numbers on the axes
      xshift=5.9cm,
      width=6.5cm,
      height=5.5cm,
]

\addplot[color=blue, thick] coordinates {(1000,0.9365540172347733) (2000,0.9279940885903003) (3000,0.923518110627097) (4000,0.9208088633342846) (5000,0.9189570792325394) (6000,0.9175921832790981) (7000,0.9164910606169235) (8000,0.9156748506793676) (9000,0.9150781006109757) (10000,0.9145540661240397) (11000,0.9141006460793741) (12000,0.9137310845770253) (13000,0.9134594032113424) (14000,0.9131799481667346) (15000,0.9128952893824723) (16000,0.912657871522569) (17000,0.9124753940024685) (18000,0.912324172787189) (19000,0.9121151782585596) (20000,0.9119629639458621)};
\addlegendentry{MultiPPI}

\addplot[color=green!50!black, thick, dashdotted] coordinates {(1000,0.9639657044565335) (2000,0.9484556175677382) (3000,0.9398404248930575) (4000,0.9342783597182173) (5000,0.9304986332528292) (6000,0.927630510500333) (7000,0.9254801094919953) (8000,0.9237780902903421) (9000,0.9223576063447055) (10000,0.9212179955030245) (11000,0.9202306240530287) (12000,0.919398244873218) (13000,0.9186878019002828) (14000,0.9180912703905452) (15000,0.9175073593788518) (16000,0.9170270755620303) (17000,0.916639811612487) (18000,0.9162503238381591) (19000,0.9158526476543661) (20000,0.9155650273802454)};
\addlegendentry{PPI{\textnormal{+}}{\textnormal{+}} Vector}

\addplot[color=brown!90, thick, densely dashed] coordinates {(1000,0.9597935282311227) (2000,0.9461787935240409) (3000,0.9391989174472801) (4000,0.935062650116862) (5000,0.9322726836032441) (6000,0.9302708517772372) (7000,0.9287760122331542) (8000,0.9276465291701175) (9000,0.9267095031599436) (10000,0.9259700455431293) (11000,0.9252672736862207) (12000,0.9247044332813973) (13000,0.9242885819279071) (14000,0.9239083420304247) (15000,0.9235369394312855) (16000,0.923218542339616) (17000,0.9229730258621824) (18000,0.9227166176605403) (19000,0.9224446370488687) (20000,0.9222080248575197)};
\addlegendentry{PPI{\textnormal{+}}{\textnormal{+}} Scalar (Pro)}

\addplot[color=brown!60, thick, densely dashed] coordinates {(1000,0.9365769002447676) (2000,0.929569197730841) (3000,0.9268631179392386) (4000,0.9254108510119781) (5000,0.9245446208985812) (6000,0.9239996604145924) (7000,0.9235710638924928) (8000,0.9232375507141142) (9000,0.9229116262846602) (10000,0.922751222594485) (11000,0.9225922220084404) (12000,0.9224273004809259) (13000,0.9223118599236867) (14000,0.9221943468283949) (15000,0.922073921732708) (16000,0.9220673190371799) (17000,0.9220005414898103) (18000,0.9219040613941991) (19000,0.9218194508250789) (20000,0.92173896587307)};
\addlegendentry{PPI{\textnormal{+}}{\textnormal{+}} Scalar (Flash)}

\legend{};
\end{axis}
\begin{axis}[
    title={Mean-squared Error (frac. of classical)},
    xlabel={Budget},
    xlabel style={yshift=5pt},   % Moves xlabel closer to the axis
    grid=both,
    outer sep=0pt,
    enlarge x limits=false,    minor grid style={dashed, gray!50},
    major grid style={solid, gray!80},
    width=\textwidth,
    yticklabel style={/pgf/number format/fixed}, % No scientific notation
    yticklabel={\pgfmathparse{\tick*100}\pgfmathprintnumber{\pgfmathresult}\%},
    xmin=0, xmax=20000,                          % Set x-axis range
    font=\scriptsize,             % Sets the base size for everything
    title style={yshift=-5pt, font=\small},    % Optional: Make title slightly bigger than the rest
    label style={font=\scriptsize}, % Font for xlabel and ylabel
    tick label style={font=\scriptsize}, % Font for the numbers on the axes
      xshift=11.7cm,
      width=6.5cm,
      height=5.5cm,
            legend pos=north east,
            legend cell align={left},
            legend style={font=\tiny,inner sep=1pt,row sep=-2pt},
]

\addplot[color=blue, thick] coordinates {(1000,0.9384498810371166) (2000,0.9313113294917891) (3000,0.9242329734308448) (4000,0.9226711796567072) (5000,0.9226758407203406) (6000,0.9207988662423808) (7000,0.9194148756502066) (8000,0.91894157986197) (9000,0.9172800163160336) (10000,0.9171373619178241) (11000,0.9166653334989701) (12000,0.9160451525371209) (13000,0.9173553451668081) (14000,0.917143839336111) (15000,0.9158584252305314) (16000,0.9150898854591314) (17000,0.9161098047416942) (18000,0.9161982471464877) (19000,0.9173728379559655) (20000,0.9150856808585291)};
\addlegendentry{MultiPPI}

\addplot[color=green!50!black, thick, dashdotted] coordinates {(1000,0.9666514694736841) (2000,0.9507786271722175) (3000,0.9458251048289251) (4000,0.940485854401145) (5000,0.9324614086405651) (6000,0.9289962410137284) (7000,0.9267412898253855) (8000,0.9271386153574581) (9000,0.923623095786627) (10000,0.923975208123204) (11000,0.9231448949832828) (12000,0.9225715264675385) (13000,0.9191996062498826) (14000,0.9201477136864146) (15000,0.9213456699466588) (16000,0.9189904094619791) (17000,0.9184862578829008) (18000,0.921003823341632) (19000,0.9192943308712991) (20000,0.919147650886344)};
\addlegendentry{PPI{\textnormal{+}}{\textnormal{+}} Vector}

\addplot[color=brown!90, thick, densely dashed] coordinates {(1000,0.9575793011001906) (2000,0.9450526691261466) (3000,0.9393146526181683) (4000,0.9344926966308179) (5000,0.9331603008938709) (6000,0.9321926247363359) (7000,0.9316087290243962) (8000,0.9300077317181406) (9000,0.9282055883223109) (10000,0.928359816952962) (11000,0.9277688141970771) (12000,0.927564152350475) (13000,0.9270350763787834) (14000,0.9256058316398537) (15000,0.9220084915586295) (16000,0.924681740023388) (17000,0.9261377834754526) (18000,0.9253119263663484) (19000,0.9243482754663763) (20000,0.9247081312328499)};
\addlegendentry{PPI{\textnormal{+}}{\textnormal{+}} Scalar (Pro)}

\addplot[color=brown!60, thick, densely dashed] coordinates {(1000,0.9395527659908117) (2000,0.9314545581633918) (3000,0.9284324614355017) (4000,0.9252810376040651) (5000,0.9265215102173489) (6000,0.9238010873553013) (7000,0.9243901480966175) (8000,0.9262675607884072) (9000,0.9253685309615148) (10000,0.9246033829229643) (11000,0.925446353107072) (12000,0.9255972310075582) (13000,0.9224224035289161) (14000,0.9237850592327386) (15000,0.9229810042636704) (16000,0.9245512819674873) (17000,0.9209141785702917) (18000,0.9217949335005752) (19000,0.9212295044842476) (20000,0.9239553395180174)};
\addlegendentry{PPI{\textnormal{+}}{\textnormal{+}} Scalar (Flash)}

\end{axis}
\end{tikzpicture}}
    \textrm{\bf \scriptsize Experiment 2: ProcessBench} 
\scalebox{.8}{\centering
\begin{tikzpicture}
\begin{axis}[
    title={Coverage},
    xlabel={Budget},
    xlabel style={yshift=5pt},   % Moves xlabel closer to the axis
    grid=both,
    outer sep=0pt,
    enlarge x limits=false,    minor grid style={dashed, gray!50},
    major grid style={solid, gray!80},
    width=\textwidth,
    yticklabel style={/pgf/number format/fixed}, % No scientific notation
    yticklabel={\pgfmathparse{\tick*100}\pgfmathprintnumber{\pgfmathresult}\%},
    xmin=0, xmax=20000,                          % Set x-axis range
    font=\scriptsize,             % Sets the base size for everything
    title style={yshift=-5pt, font=\small},    % Optional: Make title slightly bigger than the rest
    label style={font=\scriptsize}, % Font for xlabel and ylabel
    tick label style={font=\scriptsize}, % Font for the numbers on the axes
      xshift=0.0cm,
      width=6.5cm,
      height=5.5cm,
            ymin=0.9, ymax=1,
]

\addplot[color=black, dashed, very thick] coordinates {(0,.95) (5000,.95)};
\addplot[color=gray!10!black!100, thick, densely dashed] coordinates {(1000,0.949726) (2000,0.948916) (3000,0.950298) (4000,0.95051) (5000,0.948764) (6000,0.949434) (7000,0.949678) (8000,0.949752) (9000,0.949158) (10000,0.949344) (11000,0.94929) (12000,0.95005) (13000,0.94988) (14000,0.950352) (15000,0.949148) (16000,0.949564) (17000,0.949592) (18000,0.949254) (19000,0.949926) (20000,0.949486)};
\addlegendentry{PPI{\textnormal{+}}{\textnormal{+}} Scalar Large}

\addplot[color=gray!40!black!70, thick, densely dashed] coordinates {(1000,0.95073) (2000,0.94983) (3000,0.948972) (4000,0.949924) (5000,0.949242) (6000,0.949234) (7000,0.949524) (8000,0.94862) (9000,0.949012) (10000,0.948912) (11000,0.949494) (12000,0.949322) (13000,0.949178) (14000,0.949984) (15000,0.948794) (16000,0.949436) (17000,0.949506) (18000,0.94998) (19000,0.949028) (20000,0.949382)};
\addlegendentry{PPI{\textnormal{+}}{\textnormal{+}} Scalar Medium}

\addplot[color=gray!70!black!40, thick, densely dashed] coordinates {(1000,0.949232) (2000,0.949562) (3000,0.95158) (4000,0.950662) (5000,0.94911) (6000,0.949194) (7000,0.949176) (8000,0.949262) (9000,0.9497) (10000,0.948692) (11000,0.949436) (12000,0.949512) (13000,0.949482) (14000,0.9495) (15000,0.949502) (16000,0.95011) (17000,0.949732) (18000,0.949698) (19000,0.949034) (20000,0.94928)};
\addlegendentry{PPI{\textnormal{+}}{\textnormal{+}} Scalar Small}

\addplot[color=gray!100!black!10, thick, densely dashed] coordinates {(1000,0.94904) (2000,0.949022) (3000,0.949114) (4000,0.950068) (5000,0.949638) (6000,0.950062) (7000,0.949394) (8000,0.950052) (9000,0.949232) (10000,0.94938) (11000,0.948926) (12000,0.948594) (13000,0.949236) (14000,0.948856) (15000,0.949074) (16000,0.949192) (17000,0.9495) (18000,0.949012) (19000,0.94932) (20000,0.949124)};
\addlegendentry{PPI{\textnormal{+}}{\textnormal{+}} Scalar Tiny}

\addplot[color=blue, thick] coordinates {(1000,0.949778) (2000,0.950046) (3000,0.948412) (4000,0.949118) (5000,0.947954) (6000,0.948388) (7000,0.947736) (8000,0.948868) (9000,0.94882) (10000,0.94862) (11000,0.948006) (12000,0.948216) (13000,0.949182) (14000,0.94803) (15000,0.948738) (16000,0.948952) (17000,0.948216) (18000,0.947882) (19000,0.948684) (20000,0.949186)};
\addlegendentry{MultiPPI}

\addplot[color=green!50!black, thick, dashdotted] coordinates {(1000,0.950464) (2000,0.95092) (3000,0.94969) (4000,0.949894) (5000,0.948018) (6000,0.948406) (7000,0.948924) (8000,0.948596) (9000,0.949502) (10000,0.949286) (11000,0.949558) (12000,0.949466) (13000,0.949466) (14000,0.948896) (15000,0.947732) (16000,0.94812) (17000,0.949346) (18000,0.948576) (19000,0.949686) (20000,0.949344)};
\addlegendentry{PPI{\textnormal{+}}{\textnormal{+}} Vector}

\legend{};
\end{axis}
\begin{axis}[
    title={Mean-squared CI width (frac. of classical)},
    xlabel={Budget},
    xlabel style={yshift=5pt},   % Moves xlabel closer to the axis
    grid=both,
    outer sep=0pt,
    enlarge x limits=false,    minor grid style={dashed, gray!50},
    major grid style={solid, gray!80},
    width=\textwidth,
    yticklabel style={/pgf/number format/fixed}, % No scientific notation
    yticklabel={\pgfmathparse{\tick*100}\pgfmathprintnumber{\pgfmathresult}\%},
    xmin=0, xmax=20000,                          % Set x-axis range
    font=\scriptsize,             % Sets the base size for everything
    title style={yshift=-5pt, font=\small},    % Optional: Make title slightly bigger than the rest
    label style={font=\scriptsize}, % Font for xlabel and ylabel
    tick label style={font=\scriptsize}, % Font for the numbers on the axes
      xshift=5.9cm,
      width=6.5cm,
      height=5.5cm,
]

\addplot[color=gray!10!black!100, thick, densely dashed] coordinates {(1000,0.7993986681639637) (2000,0.7204024002423617) (3000,0.6773686849449404) (4000,0.6501894426739608) (5000,0.6314844764771151) (6000,0.6176589556775789) (7000,0.6070838632489277) (8000,0.5988008342495335) (9000,0.5920559550937897) (10000,0.5865738668146425) (11000,0.5817764486920511) (12000,0.5777116944203661) (13000,0.574311261536469) (14000,0.5712479590306425) (15000,0.5686288761162066) (16000,0.5662896862939432) (17000,0.5642606734078073) (18000,0.5624835235023419) (19000,0.5607647260234458) (20000,0.5592842547881676)};
\addlegendentry{PPI{\textnormal{+}}{\textnormal{+}} Scalar Large}

\addplot[color=gray!40!black!70, thick, densely dashed] coordinates {(1000,0.7980462551146105) (2000,0.7335142791606915) (3000,0.701105620678157) (4000,0.6817058312082624) (5000,0.6689182472652306) (6000,0.6596779807528715) (7000,0.6528818114981991) (8000,0.6474106239596273) (9000,0.6429661468841279) (10000,0.6394597756457072) (11000,0.6366369119516598) (12000,0.6341240836248286) (13000,0.631824810277656) (14000,0.6299661079612306) (15000,0.6284083879568589) (16000,0.6270318327349971) (17000,0.6255931526747479) (18000,0.6244809153235246) (19000,0.6234477641820944) (20000,0.6224919559762154)};
\addlegendentry{PPI{\textnormal{+}}{\textnormal{+}} Scalar Medium}

\addplot[color=gray!70!black!40, thick, densely dashed] coordinates {(1000,0.7942022195454868) (2000,0.7462299161979068) (3000,0.7247781962878307) (4000,0.7125387405260917) (5000,0.7046255524359053) (6000,0.6991328571838096) (7000,0.6949345541290121) (8000,0.6917510151152929) (9000,0.689315343486781) (10000,0.6873600442855361) (11000,0.685662434751897) (12000,0.6842296769676771) (13000,0.683109178469432) (14000,0.6820537598194459) (15000,0.6810848191878927) (16000,0.6802939677017565) (17000,0.6796559751333162) (18000,0.6789969214735978) (19000,0.6784153394495488) (20000,0.6779681385433327)};
\addlegendentry{PPI{\textnormal{+}}{\textnormal{+}} Scalar Small}

\addplot[color=gray!100!black!10, thick, densely dashed] coordinates {(1000,0.7607220320392857) (2000,0.7293194337803357) (3000,0.7168969346432841) (4000,0.7102233887887247) (5000,0.7060176409274405) (6000,0.7031098746670451) (7000,0.701143867919019) (8000,0.6994898601307323) (9000,0.6983607733254269) (10000,0.6973929751488377) (11000,0.6964888719225243) (12000,0.6957873987387294) (13000,0.6952061977916085) (14000,0.6946843539854054) (15000,0.6943152232279661) (16000,0.693903636146796) (17000,0.6935759229756202) (18000,0.6932714628290435) (19000,0.692999353120526) (20000,0.6927362229028315)};
\addlegendentry{PPI{\textnormal{+}}{\textnormal{+}} Scalar Tiny}

\addplot[color=blue, thick] coordinates {(1000,0.742423324195798) (2000,0.6715449846726853) (3000,0.6313658624049199) (4000,0.6053543432743653) (5000,0.5871528195472392) (6000,0.573583838214499) (7000,0.563034867302201) (8000,0.5547108491229069) (9000,0.54786218528985) (10000,0.5422293797376039) (11000,0.5375026717843455) (12000,0.5333998219419874) (13000,0.5300564453789105) (14000,0.5270244156120542) (15000,0.5242314098843377) (16000,0.5218575112459359) (17000,0.5196181441688988) (18000,0.5176677386558809) (19000,0.5158423971030864) (20000,0.514284092408562)};
\addlegendentry{MultiPPI}

\addplot[color=green!50!black, thick, dashdotted] coordinates {(1000,0.8597190623851068) (2000,0.7815050649456956) (3000,0.7308603370238177) (4000,0.695191023146282) (5000,0.6686256020420445) (6000,0.6480744905066397) (7000,0.6315722483247952) (8000,0.618064814863785) (9000,0.606844864508784) (10000,0.5973500138660364) (11000,0.5892708217788137) (12000,0.5821847191086569) (13000,0.5759925598212412) (14000,0.5705354528749643) (15000,0.5655460606377604) (16000,0.5611684183590134) (17000,0.5573963794221144) (18000,0.5537434784410046) (19000,0.5505907322033217) (20000,0.5475913974519732)};
\addlegendentry{PPI{\textnormal{+}}{\textnormal{+}} Vector}

\legend{};
\end{axis}
\begin{axis}[
    title={Mean-squared Error (frac. of classical)},
    xlabel={Budget},
    xlabel style={yshift=5pt},   % Moves xlabel closer to the axis
    grid=both,
    outer sep=0pt,
    enlarge x limits=false,    minor grid style={dashed, gray!50},
    major grid style={solid, gray!80},
    width=\textwidth,
    yticklabel style={/pgf/number format/fixed}, % No scientific notation
    yticklabel={\pgfmathparse{\tick*100}\pgfmathprintnumber{\pgfmathresult}\%},
    xmin=0, xmax=20000,                          % Set x-axis range
    font=\scriptsize,             % Sets the base size for everything
    title style={yshift=-5pt, font=\small},    % Optional: Make title slightly bigger than the rest
    label style={font=\scriptsize}, % Font for xlabel and ylabel
    tick label style={font=\scriptsize}, % Font for the numbers on the axes
      xshift=11.7cm,
      width=6.5cm,
      height=5.5cm,
            legend pos=north east,
            legend cell align={left},
            legend style={font=\tiny,inner sep=1pt,row sep=-2pt},
]

\addplot[color=gray!10!black!100, thick, densely dashed] coordinates {(1000,0.801051500100294) (2000,0.7232042407980451) (3000,0.676425899065194) (4000,0.6485066754202288) (5000,0.6361733468664894) (6000,0.6206289163411193) (7000,0.6076366948678222) (8000,0.5980586367164196) (9000,0.5940551600288613) (10000,0.588427394798092) (11000,0.5820722108147971) (12000,0.5771641187226991) (13000,0.5750537344968617) (14000,0.5717131141583344) (15000,0.5708696351801444) (16000,0.5675877506267826) (17000,0.565709616398016) (18000,0.5643427436694836) (19000,0.5621022048259007) (20000,0.5614823937743342)};
\addlegendentry{PPI{\textnormal{+}}{\textnormal{+}} Scalar Large}

\addplot[color=gray!40!black!70, thick, densely dashed] coordinates {(1000,0.7957885389654996) (2000,0.7329912441981103) (3000,0.7041502773325703) (4000,0.6830831404013713) (5000,0.6708298187063596) (6000,0.6624941371127829) (7000,0.6542696737215914) (8000,0.6489991233068282) (9000,0.6448970309320533) (10000,0.6426304711767694) (11000,0.6381934040491014) (12000,0.6347695088474326) (13000,0.6339610216361619) (14000,0.6306425748501301) (15000,0.630546309555347) (16000,0.6283825342992819) (17000,0.6269141859240472) (18000,0.6245866733160189) (19000,0.6253796174817358) (20000,0.6238298070476496)};
\addlegendentry{PPI{\textnormal{+}}{\textnormal{+}} Scalar Medium}

\addplot[color=gray!70!black!40, thick, densely dashed] coordinates {(1000,0.7969645315602866) (2000,0.7473426516196) (3000,0.7217997622485525) (4000,0.7094304375878975) (5000,0.7052917262773435) (6000,0.6997175202893499) (7000,0.6970179478079139) (8000,0.6934025120491162) (9000,0.6904627644789738) (10000,0.6897673252147264) (11000,0.6866406886190122) (12000,0.6854748836171077) (13000,0.6842823525563732) (14000,0.6823702275484405) (15000,0.6822496838742339) (16000,0.6792976133940398) (17000,0.6808281555632786) (18000,0.6799785694018153) (19000,0.6811899780022606) (20000,0.6792345518457944)};
\addlegendentry{PPI{\textnormal{+}}{\textnormal{+}} Scalar Small}

\addplot[color=gray!100!black!10, thick, densely dashed] coordinates {(1000,0.763853531490764) (2000,0.7315859930665712) (3000,0.7191984099981242) (4000,0.7100452160594491) (5000,0.7061362305665546) (6000,0.702337577467954) (7000,0.7013019109405453) (8000,0.6987963319987393) (9000,0.6998693254826395) (10000,0.698456388365942) (11000,0.6999158672161083) (12000,0.6981755126011389) (13000,0.6973881036812757) (14000,0.6971917508371115) (15000,0.6971440027407387) (16000,0.6954628909662844) (17000,0.6946227933686822) (18000,0.6948261062439327) (19000,0.6940668120512336) (20000,0.6947420562238249)};
\addlegendentry{PPI{\textnormal{+}}{\textnormal{+}} Scalar Tiny}

\addplot[color=blue, thick] coordinates {(1000,0.743238453575548) (2000,0.6713966351505786) (3000,0.6358298263919373) (4000,0.6094697903661821) (5000,0.593358537478411) (6000,0.5784889381638559) (7000,0.5667988195899993) (8000,0.5573524398970974) (9000,0.5498211486034809) (10000,0.5444173240516947) (11000,0.5420145047283002) (12000,0.5372962492502489) (13000,0.5329942664505202) (14000,0.531187933654902) (15000,0.5275254084185564) (16000,0.5244115744626652) (17000,0.523943540088984) (18000,0.5217360224976609) (19000,0.5189776814478799) (20000,0.5169547968335159)};
\addlegendentry{MultiPPI}

\addplot[color=green!50!black, thick, dashdotted] coordinates {(1000,0.8596642675303847) (2000,0.7815001995091919) (3000,0.7332852559155222) (4000,0.6954824511604277) (5000,0.6732416244013755) (6000,0.6516929529214145) (7000,0.6333426519333174) (8000,0.6197957127861256) (9000,0.6080540357659716) (10000,0.5984667309090985) (11000,0.590903170880836) (12000,0.5835282944800904) (13000,0.5776261215292738) (14000,0.5731386867945923) (15000,0.5707094503400049) (16000,0.5656827528836897) (17000,0.5599732059050275) (18000,0.5568646618234584) (19000,0.5505760971260303) (20000,0.548657110806749)};
\addlegendentry{PPI{\textnormal{+}}{\textnormal{+}} Vector}

\end{axis}
\end{tikzpicture}}
    \textrm{\bf \scriptsize Experiment 3: Factuality}
\scalebox{.8}{\centering
\begin{tikzpicture}
\begin{axis}[
    title={Coverage},
    xlabel={Budget},
    xlabel style={yshift=5pt},   % Moves xlabel closer to the axis
    grid=both,
    outer sep=0pt,
    enlarge x limits=false,    minor grid style={dashed, gray!50},
    major grid style={solid, gray!80},
    width=\textwidth,
    yticklabel style={/pgf/number format/fixed}, % No scientific notation
    yticklabel={\pgfmathparse{\tick*100}\pgfmathprintnumber{\pgfmathresult}\%},
    xmin=0, xmax=20000,                          % Set x-axis range
    font=\scriptsize,             % Sets the base size for everything
    title style={yshift=-5pt, font=\small},    % Optional: Make title slightly bigger than the rest
    label style={font=\scriptsize}, % Font for xlabel and ylabel
    tick label style={font=\scriptsize}, % Font for the numbers on the axes
      xshift=0.0cm,
      width=6.5cm,
      height=5.5cm,
            ymin=0.9, ymax=1,
]

\addplot[color=black, dashed, very thick] coordinates {(0,.95) (5000,.95)};
\addplot[color=blue, thick] coordinates {(1000,0.948722) (2000,0.948448) (3000,0.948936) (4000,0.949024) (5000,0.948868) (6000,0.948958) (7000,0.947734) (8000,0.948922) (9000,0.948446) (10000,0.947852) (11000,0.947908) (12000,0.948364) (13000,0.948184) (14000,0.948326) (15000,0.948236) (16000,0.94798) (17000,0.948296) (18000,0.948022) (19000,0.947782) (20000,0.94776)};
\addlegendentry{MultiPPI}

\addplot[color=green!50!black, thick, dashdotted] coordinates {(1000,0.94811) (2000,0.947832) (3000,0.948938) (4000,0.94838) (5000,0.94834) (6000,0.948262) (7000,0.948814) (8000,0.948138) (9000,0.948524) (10000,0.948398) (11000,0.947332) (12000,0.947422) (13000,0.948494) (14000,0.948612) (15000,0.948738) (16000,0.948556) (17000,0.948028) (18000,0.947684) (19000,0.948276) (20000,0.948606)};
\addlegendentry{PPI{\textnormal{+}}{\textnormal{+}} Vector}

\addplot[color=red!40,  thick, densely dotted] coordinates {(1000,0.949378) (2000,0.949256) (3000,0.949596) (4000,0.949434) (5000,0.949006) (6000,0.949426) (7000,0.949092) (8000,0.949668) (9000,0.949476) (10000,0.949296) (11000,0.948782) (12000,0.948922) (13000,0.949368) (14000,0.949864) (15000,0.949178) (16000,0.949586) (17000,0.949044) (18000,0.949416) (19000,0.94921) (20000,0.949198)};
\addlegendentry{PPI{\textnormal{+}}{\textnormal{+}} Scalar 1 agent, 1 round}

\addplot[color=red!70,  thick, densely dotted] coordinates {(1000,0.949054) (2000,0.948666) (3000,0.949362) (4000,0.949574) (5000,0.949186) (6000,0.949176) (7000,0.949532) (8000,0.94937) (9000,0.949164) (10000,0.949532) (11000,0.949992) (12000,0.949944) (13000,0.949362) (14000,0.949438) (15000,0.949218) (16000,0.949266) (17000,0.949444) (18000,0.949242) (19000,0.949802) (20000,0.948692)};
\addlegendentry{PPI{\textnormal{+}}{\textnormal{+}} Scalar 2 agents, 1 round}

\addplot[color=red,  thick, densely dotted] coordinates {(1000,0.949984) (2000,0.949492) (3000,0.949334) (4000,0.949892) (5000,0.949132) (6000,0.949712) (7000,0.949142) (8000,0.948956) (9000,0.949646) (10000,0.949302) (11000,0.94957) (12000,0.949314) (13000,0.949728) (14000,0.949548) (15000,0.949352) (16000,0.948914) (17000,0.949494) (18000,0.94954) (19000,0.94871) (20000,0.948968)};
\addlegendentry{PPI{\textnormal{+}}{\textnormal{+}} Scalar 3 agents, 1 round}

\addplot[color=red!40,  thick, densely dashed] coordinates {(1000,0.948782) (2000,0.948446) (3000,0.949148) (4000,0.948776) (5000,0.94973) (6000,0.950544) (7000,0.949436) (8000,0.949268) (9000,0.949346) (10000,0.949418) (11000,0.949632) (12000,0.94934) (13000,0.949368) (14000,0.949726) (15000,0.94941) (16000,0.94951) (17000,0.950098) (18000,0.949812) (19000,0.949374) (20000,0.949172)};
\addlegendentry{PPI{\textnormal{+}}{\textnormal{+}} Scalar 1 agent, 2 rounds}

\addplot[color=red!70,  thick, densely dashed] coordinates {(1000,0.949472) (2000,0.949902) (3000,0.949624) (4000,0.949976) (5000,0.949224) (6000,0.949144) (7000,0.949782) (8000,0.949466) (9000,0.9497) (10000,0.94902) (11000,0.948352) (12000,0.94911) (13000,0.94953) (14000,0.94954) (15000,0.949588) (16000,0.949356) (17000,0.94925) (18000,0.949554) (19000,0.948914) (20000,0.94947)};
\addlegendentry{PPI{\textnormal{+}}{\textnormal{+}} Scalar 2 agents, 2 rounds}

\addplot[color=red,  thick, densely dashed] coordinates {(1000,0.948596) (2000,0.949268) (3000,0.949478) (4000,0.949882) (5000,0.949492) (6000,0.950234) (7000,0.950848) (8000,0.95068) (9000,0.94894) (10000,0.949436) (11000,0.94971) (12000,0.949362) (13000,0.948784) (14000,0.948564) (15000,0.949086) (16000,0.948924) (17000,0.94954) (18000,0.949338) (19000,0.949098) (20000,0.949114)};
\addlegendentry{PPI{\textnormal{+}}{\textnormal{+}} Scalar 3 agents, 2 rounds}

\legend{};
\end{axis}
\begin{axis}[
    title={Mean-squared CI width (frac. of classical)},
    xlabel={Budget},
    xlabel style={yshift=5pt},   % Moves xlabel closer to the axis
    grid=both,
    outer sep=0pt,
    enlarge x limits=false,    minor grid style={dashed, gray!50},
    major grid style={solid, gray!80},
    width=\textwidth,
    yticklabel style={/pgf/number format/fixed}, % No scientific notation
    yticklabel={\pgfmathparse{\tick*100}\pgfmathprintnumber{\pgfmathresult}\%},
    xmin=0, xmax=20000,                          % Set x-axis range
    font=\scriptsize,             % Sets the base size for everything
    title style={yshift=-5pt, font=\small},    % Optional: Make title slightly bigger than the rest
    label style={font=\scriptsize}, % Font for xlabel and ylabel
    tick label style={font=\scriptsize}, % Font for the numbers on the axes
      xshift=5.9cm,
      width=6.5cm,
      height=5.5cm,
]

\addplot[color=blue, thick] coordinates {(1000,0.7874537719748442) (2000,0.76366148045508) (3000,0.7536759216226399) (4000,0.7481230325816423) (5000,0.7445216189419008) (6000,0.7420866985000297) (7000,0.7401389845676584) (8000,0.7387592903599457) (9000,0.7377032556787488) (10000,0.7367726830129973) (11000,0.7359734037239933) (12000,0.7353056190595244) (13000,0.7347106869945155) (14000,0.734180240775452) (15000,0.7338383877632428) (16000,0.7334863985456879) (17000,0.7330876186664366) (18000,0.732790738930683) (19000,0.7324744777821536) (20000,0.7322394167746838)};
\addlegendentry{MultiPPI}

\addplot[color=green!50!black, thick, dashdotted] coordinates {(1000,0.9083021212709785) (2000,0.8632309402176137) (3000,0.8360883398529693) (4000,0.8180735063665057) (5000,0.8051799734276758) (6000,0.7954594967705594) (7000,0.7878925221430647) (8000,0.7818749738495607) (9000,0.7768962491065259) (10000,0.7727876818431703) (11000,0.7693013903723388) (12000,0.7663542719090215) (13000,0.7637426154902871) (14000,0.7614794629035301) (15000,0.7594802932283584) (16000,0.7576787673104192) (17000,0.7559777445569218) (18000,0.7545390665009282) (19000,0.7533112082452882) (20000,0.752184485643065)};
\addlegendentry{PPI{\textnormal{+}}{\textnormal{+}} Vector}

\addplot[color=red!40,  thick, densely dotted] coordinates {(1000,0.8170956948008645) (2000,0.8073145975516046) (3000,0.8038121139375091) (4000,0.8020551651039127) (5000,0.8009518070719034) (6000,0.8002330291864865) (7000,0.7996802638666743) (8000,0.79935211422614) (9000,0.7989957387724352) (10000,0.7987213781039374) (11000,0.7985059480582197) (12000,0.798388360560327) (13000,0.7982372368545994) (14000,0.7981462418372617) (15000,0.798018059206156) (16000,0.7979095117856929) (17000,0.7977917450434621) (18000,0.7978013425371566) (19000,0.7976495416701108) (20000,0.7975867295937934)};
\addlegendentry{PPI{\textnormal{+}}{\textnormal{+}} Scalar 1 agent, 1 round}

\addplot[color=red!70,  thick, densely dotted] coordinates {(1000,0.7996136563131179) (2000,0.7788875762036165) (3000,0.7709219416233579) (4000,0.7667452718904192) (5000,0.7639848423433906) (6000,0.7622881734422168) (7000,0.7610219854037389) (8000,0.7600917646827474) (9000,0.7594362064521808) (10000,0.758780160666984) (11000,0.758239659797614) (12000,0.7578242412981887) (13000,0.757534447502748) (14000,0.7572048839510253) (15000,0.7569314127443795) (16000,0.7567035463045713) (17000,0.7564873201592512) (18000,0.7563664379968426) (19000,0.7561426693467459) (20000,0.7559455917596553)};
\addlegendentry{PPI{\textnormal{+}}{\textnormal{+}} Scalar 2 agents, 1 round}

\addplot[color=red,  thick, densely dotted] coordinates {(1000,0.8097668879058788) (2000,0.781497026868552) (3000,0.7702254642281382) (4000,0.7640388751092588) (5000,0.7601096616325753) (6000,0.7574723056244828) (7000,0.7556600309181039) (8000,0.7541625870241078) (9000,0.7530340393752706) (10000,0.7521491097488033) (11000,0.7513899768972716) (12000,0.7506507885316961) (13000,0.7501119620650842) (14000,0.7496222824454822) (15000,0.7492292565067783) (16000,0.7489534023397071) (17000,0.7484949592305382) (18000,0.7482289921770957) (19000,0.7480032998043019) (20000,0.7478469762308034)};
\addlegendentry{PPI{\textnormal{+}}{\textnormal{+}} Scalar 3 agents, 1 round}

\addplot[color=red!40,  thick, densely dashed] coordinates {(1000,0.8178613978601144) (2000,0.7992670261966067) (3000,0.7921201282678774) (4000,0.7884061562161021) (5000,0.7860683045802801) (6000,0.7845180759008893) (7000,0.7833602815625905) (8000,0.7825093540973924) (9000,0.7818142154254886) (10000,0.7813170737054206) (11000,0.7808701190471492) (12000,0.7805255388246239) (13000,0.7802244512767151) (14000,0.7799590092612299) (15000,0.7796672504025698) (16000,0.7795294164540294) (17000,0.7793483055923163) (18000,0.7791767917720939) (19000,0.7789906964051951) (20000,0.7788548592579664)};
\addlegendentry{PPI{\textnormal{+}}{\textnormal{+}} Scalar 1 agent, 2 rounds}

\addplot[color=red!70,  thick, densely dashed] coordinates {(1000,0.8330117925348229) (2000,0.8018719276711593) (3000,0.788716307205017) (4000,0.7814447766747932) (5000,0.7767736632493689) (6000,0.7735672638376581) (7000,0.7712957743506521) (8000,0.7694849239077532) (9000,0.768080117542115) (10000,0.7669289674431881) (11000,0.7660295516185562) (12000,0.7652550164377168) (13000,0.7645579958120139) (14000,0.7639595949992751) (15000,0.7634616057053278) (16000,0.7630293430269606) (17000,0.7626192749671514) (18000,0.7622994297109948) (19000,0.7619289725850616) (20000,0.7616463846182467)};
\addlegendentry{PPI{\textnormal{+}}{\textnormal{+}} Scalar 2 agents, 2 rounds}

\addplot[color=red,  thick, densely dashed] coordinates {(1000,0.855751517286451) (2000,0.8190153413947049) (3000,0.8022162606638495) (4000,0.7925424816015979) (5000,0.7861330172163509) (6000,0.7816997060403084) (7000,0.7784830716251645) (8000,0.7759517290766368) (9000,0.773971763989829) (10000,0.7723677360056876) (11000,0.770957836124237) (12000,0.7698220514029038) (13000,0.768841458163209) (14000,0.7679477017419952) (15000,0.7672449395273337) (16000,0.7665736020125338) (17000,0.7660644365079428) (18000,0.7655137818713421) (19000,0.7650681219453208) (20000,0.7646384785088041)};
\addlegendentry{PPI{\textnormal{+}}{\textnormal{+}} Scalar 3 agents, 2 rounds}

\legend{};
\end{axis}
\begin{axis}[
    title={Mean-squared Error (frac. of classical)},
    xlabel={Budget},
    xlabel style={yshift=5pt},   % Moves xlabel closer to the axis
    grid=both,
    outer sep=0pt,
    enlarge x limits=false,    minor grid style={dashed, gray!50},
    major grid style={solid, gray!80},
    width=\textwidth,
    yticklabel style={/pgf/number format/fixed}, % No scientific notation
    yticklabel={\pgfmathparse{\tick*100}\pgfmathprintnumber{\pgfmathresult}\%},
    xmin=0, xmax=20000,                          % Set x-axis range
    font=\scriptsize,             % Sets the base size for everything
    title style={yshift=-5pt, font=\small},    % Optional: Make title slightly bigger than the rest
    label style={font=\scriptsize}, % Font for xlabel and ylabel
    tick label style={font=\scriptsize}, % Font for the numbers on the axes
      xshift=11.7cm,
      width=6.5cm,
      height=5.5cm,
            legend pos=north east,
            legend cell align={left},
            legend style={font=\tiny,inner sep=1pt,row sep=-2pt},
]

\addplot[color=blue, thick] coordinates {(1000,0.7896063123952093) (2000,0.767048705430098) (3000,0.7560932487937743) (4000,0.750465515513114) (5000,0.7468156173651137) (6000,0.7464041620705967) (7000,0.7460565733455519) (8000,0.7417517610896185) (9000,0.7421591528321526) (10000,0.7426867528064012) (11000,0.7405610288906155) (12000,0.7403588661401569) (13000,0.7390147990876944) (14000,0.7392991155309058) (15000,0.7384119944312285) (16000,0.7391996178837065) (17000,0.7377072151768297) (18000,0.7372056771219695) (19000,0.738347675405416) (20000,0.7375807866745597)};
\addlegendentry{MultiPPI}

\addplot[color=green!50!black, thick, dashdotted] coordinates {(1000,0.9122365111800839) (2000,0.8675433606021994) (3000,0.8403719590905124) (4000,0.8234541185055849) (5000,0.81013934045819) (6000,0.8007361501557815) (7000,0.7915127276280649) (8000,0.7866470388703565) (9000,0.7810987222051445) (10000,0.7774793857547548) (11000,0.7768127053380193) (12000,0.7739414044057532) (13000,0.7672093971372245) (14000,0.7650426915151568) (15000,0.7641073900662937) (16000,0.7613862000885676) (17000,0.7614259383073335) (18000,0.7605453690550232) (19000,0.7569890332490565) (20000,0.7556310074985588)};
\addlegendentry{PPI{\textnormal{+}}{\textnormal{+}} Vector}

\addplot[color=red!40,  thick, densely dotted] coordinates {(1000,0.8186581926295349) (2000,0.8092488200887128) (3000,0.8043642045850022) (4000,0.8033981761705563) (5000,0.8026848501976848) (6000,0.801276120816251) (7000,0.8018879369170352) (8000,0.799927985829889) (9000,0.8010453634855209) (10000,0.8004772246796321) (11000,0.8012009853658859) (12000,0.8006686616126288) (13000,0.7991155082954186) (14000,0.7985232637797788) (15000,0.8005857227190537) (16000,0.7989438973040054) (17000,0.7990689963137492) (18000,0.7998540643891806) (19000,0.7992058870065158) (20000,0.7989663766274161)};
\addlegendentry{PPI{\textnormal{+}}{\textnormal{+}} Scalar 1 agent, 1 round}

\addplot[color=red!70,  thick, densely dotted] coordinates {(1000,0.8013752103269364) (2000,0.7814728685872974) (3000,0.7718242249451124) (4000,0.7667277075277196) (5000,0.7652450452023996) (6000,0.7627925624262656) (7000,0.7614985191495337) (8000,0.7612074118826541) (9000,0.7612040214026741) (10000,0.758741084748182) (11000,0.7588718238007071) (12000,0.7576419552746106) (13000,0.7594786658119713) (14000,0.7582621180900274) (15000,0.7580746492797628) (16000,0.7577240713675909) (17000,0.7581828842524502) (18000,0.7585192031047971) (19000,0.7562125212082911) (20000,0.7586028620981932)};
\addlegendentry{PPI{\textnormal{+}}{\textnormal{+}} Scalar 2 agents, 1 round}

\addplot[color=red,  thick, densely dotted] coordinates {(1000,0.8091660736970486) (2000,0.7819602080004182) (3000,0.7727196693166701) (4000,0.7647703247950014) (5000,0.7619882676186783) (6000,0.7586803148386856) (7000,0.7567684646967152) (8000,0.7556856795971845) (9000,0.7531606427309213) (10000,0.7534616458669169) (11000,0.752212458900897) (12000,0.750934819544703) (13000,0.7513106682784622) (14000,0.7506842855608912) (15000,0.7513866603396195) (16000,0.7508044398896937) (17000,0.7490530005286211) (18000,0.7490902371604011) (19000,0.7514149408849863) (20000,0.7502055405753594)};
\addlegendentry{PPI{\textnormal{+}}{\textnormal{+}} Scalar 3 agents, 1 round}

\addplot[color=red!40,  thick, densely dashed] coordinates {(1000,0.8209784448964035) (2000,0.8029678421802054) (3000,0.7938753613872551) (4000,0.7912768079189572) (5000,0.7863039817388606) (6000,0.7819707199884486) (7000,0.7838672147404498) (8000,0.7830581159661929) (9000,0.7833131554107008) (10000,0.7831412921311398) (11000,0.7820237948603959) (12000,0.781668227599909) (13000,0.7818408723990392) (14000,0.7805465702375187) (15000,0.7812891291755002) (16000,0.7795696490407972) (17000,0.7792174329179398) (18000,0.7800821170249493) (19000,0.7798387910685449) (20000,0.7808682754598628)};
\addlegendentry{PPI{\textnormal{+}}{\textnormal{+}} Scalar 1 agent, 2 rounds}

\addplot[color=red!70,  thick, densely dashed] coordinates {(1000,0.8330458059577495) (2000,0.8007901113072569) (3000,0.7892027191967381) (4000,0.7807414372529561) (5000,0.7782757924715384) (6000,0.7749792872919339) (7000,0.7717502548616161) (8000,0.7703369899664338) (9000,0.7675749935958974) (10000,0.7682132047467934) (11000,0.7701039680346915) (12000,0.7673900401887919) (13000,0.7649668138791851) (14000,0.7639015651104458) (15000,0.7638880206948385) (16000,0.7636818006337445) (17000,0.7633590089321306) (18000,0.7632261245904897) (19000,0.7645408340836579) (20000,0.7624901363768035)};
\addlegendentry{PPI{\textnormal{+}}{\textnormal{+}} Scalar 2 agents, 2 rounds}

\addplot[color=red,  thick, densely dashed] coordinates {(1000,0.8574993939082323) (2000,0.8213790039710331) (3000,0.8032071801202839) (4000,0.7926309827475057) (5000,0.7866461150472068) (6000,0.7795629462937927) (7000,0.7770641793306843) (8000,0.7748983231658703) (9000,0.7763544184424629) (10000,0.7723943470728811) (11000,0.7717957945035724) (12000,0.7712575230274231) (13000,0.7716687961370242) (14000,0.7701864983048651) (15000,0.7694230078125368) (16000,0.7692004112253884) (17000,0.76639224267556) (18000,0.7664286778804754) (19000,0.7667208419600317) (20000,0.7658556761812225)};
\addlegendentry{PPI{\textnormal{+}}{\textnormal{+}} Scalar 3 agents, 2 rounds}

\end{axis}
\end{tikzpicture}}
    \vspace{-10pt}
    \caption{Results given $N=1000$ labeled examples. Results are shown by budget for the experiments on Chatbot Arena (top), ProcessBench (middle), and Factuality (bottom). For each estimator (all baselines and MultiPPI), the left column plots the empirical coverage of the 95\% CI, the middle column plots the width of the 95\% CI, and the right column plots the empirical mean-squared error of the point estimate. \looseness=-1}
    \label{fig:1000}
    \vspace{-10pt}
\end{figure}



\begin{figure}[!h]
    \centering
  \textrm{\bf \scriptsize Experiment 1: Chatbot Arena}
  \includegraphics[width=1.0\textwidth]{figures/comparison_chatbotarena3.png}
  %\vspace{-10pt}
  \textrm{\bf \scriptsize Experiment 2: ProcessBench}
  \includegraphics[width=1.0\textwidth]{figures/comparison_processbench3.png}
  \textrm{\bf \scriptsize Experiment 3: Factuality}
  
\includegraphics[width=1.0\textwidth]{figures/comparison_factuality3.png}
  %\vspace{-10pt}
    \caption{MultiPPI with varying number of fully labeled examples.}
    \label{fig:varying}
    \vspace{-10pt}
\end{figure}

\begin{figure}
    \centering
    \textrm{\bf \scriptsize Experiment 1: ChatBot Arena}
\scalebox{0.8}{
\centering
\begin{tikzpicture}
\begin{axis}[
    title={Coverage},
    xlabel={Budget},
    xlabel style={yshift=5pt},   % Moves xlabel closer to the axis
    grid=both,
    outer sep=0pt,
    enlarge x limits=false,    minor grid style={dashed, gray!50},
    major grid style={solid, gray!80},
    width=\textwidth,
    yticklabel style={/pgf/number format/fixed}, % No scientific notation
    yticklabel={\pgfmathparse{\tick*100}\pgfmathprintnumber{\pgfmathresult}\%},
    xmin=0, xmax=1000,                          % Set x-axis range
    font=\scriptsize,             % Sets the base size for everything
    title style={yshift=-5pt, font=\small},    % Optional: Make title slightly bigger than the rest
    label style={font=\scriptsize}, % Font for xlabel and ylabel
    tick label style={font=\scriptsize}, % Font for the numbers on the axes
      xshift=0.0cm,
      width=6.5cm,
      height=5.5cm,
            ymin=0.85, ymax=1,
]

\addplot[color=black, dashed, very thick] coordinates {(0,.95) (5000,.95)};
\addplot[color=blue, thick] coordinates {(50,0.93678) (100,0.935164) (150,0.934636) (200,0.934248) (250,0.934596) (300,0.934542) (350,0.934052) (400,0.933864) (450,0.933336) (500,0.934078) (550,0.933158) (600,0.933138) (650,0.933054) (700,0.93252) (750,0.932956) (800,0.932508) (850,0.933692) (900,0.933256) (950,0.933172) (1000,0.93293)};
\addlegendentry{MultiPPI (Ledoit-Wolf)}

\addplot[color=green, thick, densely dashed] coordinates {(50,0.936586) (100,0.934652) (150,0.93371) (200,0.932556) (250,0.93396) (300,0.933984) (350,0.934012) (400,0.933768) (450,0.932148) (500,0.931784) (550,0.931752) (600,0.931314) (650,0.93122) (700,0.93139) (750,0.931762) (800,0.931034) (850,0.932174) (900,0.931256) (950,0.931896) (1000,0.93131)};
\addlegendentry{MultiPPI (Empirical)}

\legend{};
\end{axis}
\begin{axis}[
    title={Mean-squared CI width (frac. of classical)},
    xlabel={Budget},
    xlabel style={yshift=5pt},   % Moves xlabel closer to the axis
    grid=both,
    outer sep=0pt,
    enlarge x limits=false,    minor grid style={dashed, gray!50},
    major grid style={solid, gray!80},
    width=\textwidth,
    yticklabel style={/pgf/number format/fixed}, % No scientific notation
    yticklabel={\pgfmathparse{\tick*100}\pgfmathprintnumber{\pgfmathresult}\%},
    xmin=0, xmax=1000,                          % Set x-axis range
    font=\scriptsize,             % Sets the base size for everything
    title style={yshift=-5pt, font=\small},    % Optional: Make title slightly bigger than the rest
    label style={font=\scriptsize}, % Font for xlabel and ylabel
    tick label style={font=\scriptsize}, % Font for the numbers on the axes
      xshift=5.9cm,
      width=6.5cm,
      height=5.5cm,
]

\addplot[color=blue, thick] coordinates {(50,0.9164090567391647) (100,0.9049872637176533) (150,0.8985714007169173) (200,0.8950913169506546) (250,0.8930025583082056) (300,0.8913318872342815) (350,0.8898057784642975) (400,0.8888455117949398) (450,0.8882116513110412) (500,0.8872140193872418) (550,0.8864905387046085) (600,0.8860129898831816) (650,0.8859909852961145) (700,0.8855633168885012) (750,0.8849541299493545) (800,0.8848428321495893) (850,0.8847501257046745) (900,0.8844771325699236) (950,0.8841572752710469) (1000,0.8838202537781186)};
\addlegendentry{MultiPPI (Ledoit-Wolf)}

\addplot[color=green, thick, densely dashed] coordinates {(50,0.9164029710465458) (100,0.9047662091539396) (150,0.8989154284434913) (200,0.8950622554627548) (250,0.8929169138791052) (300,0.8909064918500701) (350,0.8894354278999306) (400,0.8881723370314002) (450,0.887696819710712) (500,0.8872038010490143) (550,0.8861289726336172) (600,0.8854424788329114) (650,0.8853771044336499) (700,0.8847831248440161) (750,0.8846527373313299) (800,0.8843505298766667) (850,0.8838887206099857) (900,0.8836686375682391) (950,0.8832629261198237) (1000,0.8830848544201532)};
\addlegendentry{MultiPPI (Empirical)}

\legend{};
\end{axis}
\begin{axis}[
    title={Mean-squared Error (frac. of classical)},
    xlabel={Budget},
    xlabel style={yshift=5pt},   % Moves xlabel closer to the axis
    grid=both,
    outer sep=0pt,
    enlarge x limits=false,    minor grid style={dashed, gray!50},
    major grid style={solid, gray!80},
    width=\textwidth,
    yticklabel style={/pgf/number format/fixed}, % No scientific notation
    yticklabel={\pgfmathparse{\tick*100}\pgfmathprintnumber{\pgfmathresult}\%},
    xmin=0, xmax=1000,                          % Set x-axis range
    font=\scriptsize,             % Sets the base size for everything
    title style={yshift=-5pt, font=\small},    % Optional: Make title slightly bigger than the rest
    label style={font=\scriptsize}, % Font for xlabel and ylabel
    tick label style={font=\scriptsize}, % Font for the numbers on the axes
      xshift=11.7cm,
      width=6.5cm,
      height=5.5cm,
            legend pos=north east,
            legend cell align={left},
            legend style={font=\tiny,inner sep=1pt,row sep=-2pt},
]

\addplot[color=blue, thick] coordinates {(50,0.9496007015558312) (100,0.9433701489854492) (150,0.935205341841942) (200,0.9333489943887195) (250,0.9303813985404417) (300,0.9288850774062749) (350,0.9286613601463324) (400,0.929061804216448) (450,0.9285255334520865) (500,0.9271125030967432) (550,0.9298041075363036) (600,0.9287077367063746) (650,0.9282378376854579) (700,0.9281181416259152) (750,0.9267185063687005) (800,0.9286897879317082) (850,0.9251499124939538) (900,0.9259727394002606) (950,0.9268771123599007) (1000,0.9280628073961095)};
\addlegendentry{MultiPPI (Ledoit-Wolf)}

\addplot[color=green, thick, densely dashed] coordinates {(50,0.9477604301113974) (100,0.9423172663752447) (150,0.9408311341002188) (200,0.9402707618912466) (250,0.9331995540470072) (300,0.9304126940619714) (350,0.9311661356375273) (400,0.9298385124955051) (450,0.9324201433070022) (500,0.9332272788407633) (550,0.9331771339775407) (600,0.9330576213012906) (650,0.933213186601883) (700,0.9304565037865629) (750,0.9315459620824166) (800,0.9326117505678091) (850,0.9310526444364026) (900,0.9321344956674121) (950,0.9302626538538268) (1000,0.9321676634249987)};
\addlegendentry{MultiPPI (Empirical)}

\end{axis}
\end{tikzpicture}
}
    \textrm{\bf \scriptsize Experiment 2: ProcessBench}
\scalebox{0.8}{
\centering
\begin{tikzpicture}
\begin{axis}[
    title={Coverage},
    xlabel={Budget},
    xlabel style={yshift=5pt},   % Moves xlabel closer to the axis
    grid=both,
    outer sep=0pt,
    enlarge x limits=false,    minor grid style={dashed, gray!50},
    major grid style={solid, gray!80},
    width=\textwidth,
    yticklabel style={/pgf/number format/fixed}, % No scientific notation
    yticklabel={\pgfmathparse{\tick*100}\pgfmathprintnumber{\pgfmathresult}\%},
    xmin=0, xmax=1000,                          % Set x-axis range
    font=\scriptsize,             % Sets the base size for everything
    title style={yshift=-5pt, font=\small},    % Optional: Make title slightly bigger than the rest
    label style={font=\scriptsize}, % Font for xlabel and ylabel
    tick label style={font=\scriptsize}, % Font for the numbers on the axes
      xshift=0.0cm,
      width=6.5cm,
      height=5.5cm,
            ymin=0.85, ymax=1,
]

\addplot[color=black, dashed, very thick] coordinates {(0,.95) (5000,.95)};
\addplot[color=blue, thick] coordinates {(50,0.93701) (100,0.933986) (150,0.933592) (200,0.931134) (250,0.929048) (300,0.92696) (350,0.926958) (400,0.926142) (450,0.92481) (500,0.92493) (550,0.923816) (600,0.922892) (650,0.92134) (700,0.919996) (750,0.919808) (800,0.919018) (850,0.918258) (900,0.91756) (950,0.917824) (1000,0.917406)};
\addlegendentry{MultiPPI (Empirical)}

\addplot[color=green, thick, densely dashed] coordinates {(50,0.936016) (100,0.933684) (150,0.93258) (200,0.93024) (250,0.927264) (300,0.927026) (350,0.928326) (400,0.927124) (450,0.924768) (500,0.924354) (550,0.9237) (600,0.923138) (650,0.921582) (700,0.921322) (750,0.9212) (800,0.92029) (850,0.920528) (900,0.920106) (950,0.921034) (1000,0.920898)};
\addlegendentry{MultiPPI (Ledoit-Wolf)}

\legend{};
\end{axis}
\begin{axis}[
    title={Mean-squared CI width (frac. of classical)},
    xlabel={Budget},
    xlabel style={yshift=5pt},   % Moves xlabel closer to the axis
    grid=both,
    outer sep=0pt,
    enlarge x limits=false,    minor grid style={dashed, gray!50},
    major grid style={solid, gray!80},
    width=\textwidth,
    yticklabel style={/pgf/number format/fixed}, % No scientific notation
    yticklabel={\pgfmathparse{\tick*100}\pgfmathprintnumber{\pgfmathresult}\%},
    xmin=0, xmax=1000,                          % Set x-axis range
    font=\scriptsize,             % Sets the base size for everything
    title style={yshift=-5pt, font=\small},    % Optional: Make title slightly bigger than the rest
    label style={font=\scriptsize}, % Font for xlabel and ylabel
    tick label style={font=\scriptsize}, % Font for the numbers on the axes
      xshift=5.9cm,
      width=6.5cm,
      height=5.5cm,
]

\addplot[color=blue, thick] coordinates {(50,0.7102507256517969) (100,0.63674577184455) (150,0.5955944924444864) (200,0.5681884225793133) (250,0.5484604545677358) (300,0.5334150894353978) (350,0.5220121368478097) (400,0.5123486686220909) (450,0.5043375167882358) (500,0.497590007884571) (550,0.49211114135887196) (600,0.4867996682241361) (650,0.48227301513358023) (700,0.47818276731063464) (750,0.4742931805136345) (800,0.4709444817307608) (850,0.468286633009942) (900,0.46534583873440616) (950,0.4634117190756791) (1000,0.46106954686359936)};
\addlegendentry{MultiPPI (Empirical)}

\addplot[color=green, thick, densely dashed] coordinates {(50,0.7104284432326666) (100,0.6361804462741223) (150,0.5940047059860405) (200,0.5661303592061417) (250,0.5467412687554667) (300,0.53187574786712) (350,0.5205079192547275) (400,0.5108925702552023) (450,0.5026650395712908) (500,0.4958221120760036) (550,0.4903616984599653) (600,0.48518346433162696) (650,0.480221652270261) (700,0.47659221682635233) (750,0.47322376929977994) (800,0.4702380765212926) (850,0.4669210231427845) (900,0.46449234999354055) (950,0.4624221010837662) (1000,0.4603441459609113)};
\addlegendentry{MultiPPI (Ledoit-Wolf)}

\legend{};
\end{axis}
\begin{axis}[
    title={Mean-squared Error (frac. of classical)},
    xlabel={Budget},
    xlabel style={yshift=5pt},   % Moves xlabel closer to the axis
    grid=both,
    outer sep=0pt,
    enlarge x limits=false,    minor grid style={dashed, gray!50},
    major grid style={solid, gray!80},
    width=\textwidth,
    yticklabel style={/pgf/number format/fixed}, % No scientific notation
    yticklabel={\pgfmathparse{\tick*100}\pgfmathprintnumber{\pgfmathresult}\%},
    xmin=0, xmax=1000,                          % Set x-axis range
    font=\scriptsize,             % Sets the base size for everything
    title style={yshift=-5pt, font=\small},    % Optional: Make title slightly bigger than the rest
    label style={font=\scriptsize}, % Font for xlabel and ylabel
    tick label style={font=\scriptsize}, % Font for the numbers on the axes
      xshift=11.7cm,
      width=6.5cm,
      height=5.5cm,
            legend pos=north east,
            legend cell align={left},
            legend style={font=\tiny,inner sep=1pt,row sep=-2pt},
]

\addplot[color=blue, thick] coordinates {(50,0.7415286970695942) (100,0.6737538510757545) (150,0.6324481711200541) (200,0.609449288231811) (250,0.5937982800230911) (300,0.5819009858444025) (350,0.5707004864461146) (400,0.5624693247741782) (450,0.5549002287591317) (500,0.5485429902953898) (550,0.5447669752806734) (600,0.5415916054305423) (650,0.5385901159003871) (700,0.5366599178143323) (750,0.5315263112459652) (800,0.5285338348235763) (850,0.5283562644628885) (900,0.5260687078176699) (950,0.5236186783193508) (1000,0.5225354285857369)};
\addlegendentry{MultiPPI (Empirical)}

\addplot[color=green, thick, densely dashed] coordinates {(50,0.7416967483228185) (100,0.6718947012152297) (150,0.634458521221496) (200,0.6108025173290396) (250,0.5937414057867306) (300,0.5810104924542603) (350,0.5660112713132612) (400,0.5581287489166212) (450,0.5533334511453084) (500,0.5470577605999188) (550,0.5422202868635849) (600,0.5368751052839603) (650,0.5345443236510281) (700,0.5315798752273319) (750,0.5284327821128794) (800,0.5258934507824748) (850,0.5238361690821306) (900,0.5216083150432199) (950,0.5178458620961298) (1000,0.5167790894161537)};
\addlegendentry{MultiPPI (Ledoit-Wolf)}

\end{axis}
\end{tikzpicture}
}
    \textrm{\bf \scriptsize Experiment 3: Factuality}
\scalebox{0.8}{
\centering
\begin{tikzpicture}
\begin{axis}[
    title={Coverage},
    xlabel={Budget},
    xlabel style={yshift=5pt},   % Moves xlabel closer to the axis
    grid=both,
    outer sep=0pt,
    enlarge x limits=false,    minor grid style={dashed, gray!50},
    major grid style={solid, gray!80},
    width=\textwidth,
    yticklabel style={/pgf/number format/fixed}, % No scientific notation
    yticklabel={\pgfmathparse{\tick*100}\pgfmathprintnumber{\pgfmathresult}\%},
    xmin=0, xmax=1000,                          % Set x-axis range
    font=\scriptsize,             % Sets the base size for everything
    title style={yshift=-5pt, font=\small},    % Optional: Make title slightly bigger than the rest
    label style={font=\scriptsize}, % Font for xlabel and ylabel
    tick label style={font=\scriptsize}, % Font for the numbers on the axes
      xshift=0.0cm,
      width=6.5cm,
      height=5.5cm,
            ymin=0.85, ymax=1,
]

\addplot[color=black, dashed, very thick] coordinates {(0,.95) (5000,.95)};
\addplot[color=blue, thick] coordinates {(50,0.931684) (100,0.925516) (150,0.921914) (200,0.918916) (250,0.917062) (300,0.915444) (350,0.912292) (400,0.91245) (450,0.910884) (500,0.910956) (550,0.909946) (600,0.909544) (650,0.90814) (700,0.908614) (750,0.90794) (800,0.907174) (850,0.906936) (900,0.906322) (950,0.906064) (1000,0.906346)};
\addlegendentry{MultiPPI (Ledoit-Wolf)}

\addplot[color=green, thick, densely dashed] coordinates {(50,0.93021) (100,0.922508) (150,0.917794) (200,0.913154) (250,0.910076) (300,0.906918) (350,0.905752) (400,0.90341) (450,0.901504) (500,0.89958) (550,0.901462) (600,0.899956) (650,0.898336) (700,0.897342) (850,0.89585) (900,0.895328) (950,0.894384) (1000,0.894756)};
\addlegendentry{MultiPPI (Empirical)}

\legend{};
\end{axis}
\begin{axis}[
    title={Mean-squared CI width (frac. of classical)},
    xlabel={Budget},
    xlabel style={yshift=5pt},   % Moves xlabel closer to the axis
    grid=both,
    outer sep=0pt,
    enlarge x limits=false,    minor grid style={dashed, gray!50},
    major grid style={solid, gray!80},
    width=\textwidth,
    yticklabel style={/pgf/number format/fixed}, % No scientific notation
    yticklabel={\pgfmathparse{\tick*100}\pgfmathprintnumber{\pgfmathresult}\%},
    xmin=0, xmax=1000,                          % Set x-axis range
    font=\scriptsize,             % Sets the base size for everything
    title style={yshift=-5pt, font=\small},    % Optional: Make title slightly bigger than the rest
    label style={font=\scriptsize}, % Font for xlabel and ylabel
    tick label style={font=\scriptsize}, % Font for the numbers on the axes
      xshift=5.9cm,
      width=6.5cm,
      height=5.5cm,
]

\addplot[color=blue, thick] coordinates {(50,0.7545364668479843) (100,0.7255458162063846) (150,0.7119201995335014) (200,0.7033429812767781) (250,0.696524071161093) (300,0.6916367464120456) (350,0.6884894463427089) (400,0.6855218289916332) (450,0.6834237878735445) (500,0.6813894001504142) (550,0.67951864487514) (600,0.6782205725951692) (650,0.6766564862096361) (700,0.6755847039022301) (750,0.6745882831724183) (800,0.6740555022443864) (850,0.6732813897423021) (900,0.6727000496701429) (950,0.6722167921964767) (1000,0.6714457175976168)};
\addlegendentry{MultiPPI (Ledoit-Wolf)}

\addplot[color=green, thick, densely dashed] coordinates {(50,0.7566034640949048) (100,0.7320025753376791) (150,0.7192117416496288) (200,0.7106197783539396) (250,0.7039413156688763) (300,0.6986149833866808) (350,0.6946833525705519) (400,0.6909818862011394) (450,0.6873755873510305) (500,0.6851209171934288) (550,0.6828992114450426) (600,0.6811096255461176) (650,0.6791931202266553) (700,0.6778154505171291) (850,0.6738732937145907) (900,0.6729310667357478) (950,0.6720623691217819) (1000,0.6713383487546645)};
\addlegendentry{MultiPPI (Empirical)}

\legend{};
\end{axis}
\begin{axis}[
    title={Mean-squared Error (frac. of classical)},
    xlabel={Budget},
    xlabel style={yshift=5pt},   % Moves xlabel closer to the axis
    grid=both,
    outer sep=0pt,
    enlarge x limits=false,    minor grid style={dashed, gray!50},
    major grid style={solid, gray!80},
    width=\textwidth,
    yticklabel style={/pgf/number format/fixed}, % No scientific notation
    yticklabel={\pgfmathparse{\tick*100}\pgfmathprintnumber{\pgfmathresult}\%},
    xmin=0, xmax=1000,                          % Set x-axis range
    font=\scriptsize,             % Sets the base size for everything
    title style={yshift=-5pt, font=\small},    % Optional: Make title slightly bigger than the rest
    label style={font=\scriptsize}, % Font for xlabel and ylabel
    tick label style={font=\scriptsize}, % Font for the numbers on the axes
      xshift=11.7cm,
      width=6.5cm,
      height=5.5cm,
            legend pos=north east,
            legend cell align={left},
            legend style={font=\tiny,inner sep=1pt,row sep=-2pt},
]

\addplot[color=blue, thick] coordinates {(50,0.7949081400250686) (100,0.7786140849376558) (150,0.7727034268397065) (200,0.7688334978778444) (250,0.7656795057417095) (300,0.7631636433328535) (350,0.7660535033967734) (400,0.7622304664272552) (450,0.7635189532789504) (500,0.7603842564882731) (550,0.7598959203567006) (600,0.7586542968122403) (650,0.7606197542972088) (700,0.7583829198571969) (750,0.7593437537641989) (800,0.7594086697945458) (850,0.7576758823212617) (900,0.7585728076381071) (950,0.7587632228894517) (1000,0.7582761252461898)};
\addlegendentry{MultiPPI (Ledoit-Wolf)}

\addplot[color=green, thick, densely dashed] coordinates {(50,0.8023032442083847) (100,0.7964836333926696) (150,0.7914962378638787) (200,0.7928243127532993) (250,0.7934967123564726) (300,0.7948219009149913) (350,0.792902406025407) (400,0.7933934550562645) (450,0.7920635997248212) (500,0.7917695854520442) (550,0.7875103310661308) (600,0.7882057823869585) (650,0.7895671701337454) (700,0.7888430856556551) (850,0.7871563503970929) (900,0.7876793837210879) (950,0.7876552652257569) (1000,0.7864314600139845)};
\addlegendentry{MultiPPI (Empirical)}

\end{axis}
\end{tikzpicture}
}
    \caption{Comparison of results with different techniques for covariance estimation, for $N=50$. We find that Ledoit-Wolf shrinkage covariance estimation yields best performance in all regimes.}
    \label{fig:shrinkage}
\end{figure}

\subsection{The impact of shrinkage covariance estimation}

In this section, we discuss the impact of shrinkage covariance estimation on MultiPPI. We provide finite-sample bounds on the induced performance, and empirical results.

In \autoref{fig:shrinkage}, we compare performance of MultiPPI with covariance estimation via (a) the empirical covariance matrix, and (b) the Ledoit-Wolf estimated covariance matrix.

For a finite-sample bound on our estimator with shrinkage estimation, see \Cref{thm:ledoit}.
For more general results on {sensitivity to mis-specification,} please refer to \Cref{thm:stab}. 

\subsection{Scalability and computational tractability of the estimator}\label{sec:restricted}
SOCPs and SDPs are known to run in polynomial time in the number of contraints, which is, in our formulation, $|\mathcal{I}|$. In the preceding sections we have made the choice $\mathcal{I}=2^{\{1,\dots,k\}}$, but we show in this section that we may recover much of the same performance with a choice of $\mathcal{I}\subseteq 2^{\{1,\dots,k\}}$ which grows only linearly in $k$. Specifically, we take $\mathcal{I}=\{\{1,\dots,k\},\{2,\dots,k\},\{2\},\dots,\{k\}\}$, which corresponds to including terms for each model individually, as well as for their joint. We label the version of MultiPPI induced by this choice ``MultiPPI (Restricted).'' \Cref{fig:restricted} shows that the results of this method are very comparable to those of standard MultiPPI, in which we take $\mathcal{I}$ to be the collection of \textit{all} subsets of $\{1,\dots,k\}$. For more information on the computational tractability of the procedure, please see \Cref{sec:computation}. While our restricted formulation empirically recovers the performance of the full exponential search space by growing only linearly in $k$, a formal theoretical guarantee for this heuristic remains an open question. Proving that the restricted estimator incurs a small variance penalty in comparison to the full $O(2^k)$ formulation would validate its use for larger autorater ensembles.

\begin{figure}
    \centering

    \textrm{\bf \scriptsize Experiment 1: ChatBot Arena}
\scalebox{0.8}{
\centering
\begin{tikzpicture}
\begin{axis}[
    title={Coverage},
    xlabel={Budget},
    xlabel style={yshift=5pt},   % Moves xlabel closer to the axis
    grid=both,
    outer sep=0pt,
    enlarge x limits=false,    minor grid style={dashed, gray!50},
    major grid style={solid, gray!80},
    width=\textwidth,
    yticklabel style={/pgf/number format/fixed}, % No scientific notation
    yticklabel={\pgfmathparse{\tick*100}\pgfmathprintnumber{\pgfmathresult}\%},
    xmin=0, xmax=1000,                          % Set x-axis range
    font=\scriptsize,             % Sets the base size for everything
    title style={yshift=-5pt, font=\small},    % Optional: Make title slightly bigger than the rest
    label style={font=\scriptsize}, % Font for xlabel and ylabel
    tick label style={font=\scriptsize}, % Font for the numbers on the axes
      xshift=0.0cm,
      width=6.5cm,
      height=5.5cm,
            ymin=0.85, ymax=1,
]

\addplot[color=black, dashed, very thick] coordinates {(0,.95) (5000,.95)};
\addplot[color=blue, thick] coordinates {(50,0.93678) (100,0.935164) (150,0.934636) (200,0.934248) (250,0.934596) (300,0.934542) (350,0.934052) (400,0.933864) (450,0.933336) (500,0.934078) (550,0.933158) (600,0.933138) (650,0.933054) (700,0.93252) (750,0.932956) (800,0.932508) (850,0.933692) (900,0.933256) (950,0.933172) (1000,0.93293)};
\addlegendentry{MultiPPI}

\addplot[color=red, thick, densely dashed] coordinates {(50,0.938018) (100,0.935344) (150,0.935494) (200,0.935086) (250,0.933386) (300,0.933802) (350,0.933838) (400,0.933184) (450,0.933878) (500,0.933122) (550,0.9329) (600,0.933118) (650,0.933506) (700,0.933468) (750,0.932696) (800,0.932768) (850,0.933362) (900,0.932846) (950,0.93331) (1000,0.933302)};
\addlegendentry{MultiPPI (Restricted)}

\legend{};
\end{axis}
\begin{axis}[
    title={Mean-squared CI width (frac. of classical)},
    xlabel={Budget},
    xlabel style={yshift=5pt},   % Moves xlabel closer to the axis
    grid=both,
    outer sep=0pt,
    enlarge x limits=false,    minor grid style={dashed, gray!50},
    major grid style={solid, gray!80},
    width=\textwidth,
    yticklabel style={/pgf/number format/fixed}, % No scientific notation
    yticklabel={\pgfmathparse{\tick*100}\pgfmathprintnumber{\pgfmathresult}\%},
    xmin=0, xmax=1000,                          % Set x-axis range
    font=\scriptsize,             % Sets the base size for everything
    title style={yshift=-5pt, font=\small},    % Optional: Make title slightly bigger than the rest
    label style={font=\scriptsize}, % Font for xlabel and ylabel
    tick label style={font=\scriptsize}, % Font for the numbers on the axes
      xshift=5.9cm,
      width=6.5cm,
      height=5.5cm,
]

\addplot[color=blue, thick] coordinates {(50,0.9164090567391647) (100,0.9049872637176533) (150,0.8985714007169173) (200,0.8950913169506546) (250,0.8930025583082056) (300,0.8913318872342815) (350,0.8898057784642975) (400,0.8888455117949398) (450,0.8882116513110412) (500,0.8872140193872418) (550,0.8864905387046085) (600,0.8860129898831816) (650,0.8859909852961145) (700,0.8855633168885012) (750,0.8849541299493545) (800,0.8848428321495893) (850,0.8847501257046745) (900,0.8844771325699236) (950,0.8841572752710469) (1000,0.8838202537781186)};
\addlegendentry{MultiPPI}

\addplot[color=red, thick, densely dashed] coordinates {(50,0.9163886372616877) (100,0.9044658002309899) (150,0.89908123822742) (200,0.8955149156342862) (250,0.8928428412781502) (300,0.8911642728378455) (350,0.8897300000768683) (400,0.8887156283964417) (450,0.8878069985685989) (500,0.8870060766255824) (550,0.8863330039471056) (600,0.8861074112773726) (650,0.8857462622050291) (700,0.8853649617148853) (750,0.8846886779986036) (800,0.8844806315326778) (850,0.8843125174909024) (900,0.8842326446976574) (950,0.8838120435921081) (1000,0.8836551330629692)};
\addlegendentry{MultiPPI (Restricted)}

\legend{};
\end{axis}
\begin{axis}[
    title={Mean-squared Error (frac. of classical)},
    xlabel={Budget},
    xlabel style={yshift=5pt},   % Moves xlabel closer to the axis
    grid=both,
    outer sep=0pt,
    enlarge x limits=false,    minor grid style={dashed, gray!50},
    major grid style={solid, gray!80},
    width=\textwidth,
    yticklabel style={/pgf/number format/fixed}, % No scientific notation
    yticklabel={\pgfmathparse{\tick*100}\pgfmathprintnumber{\pgfmathresult}\%},
    xmin=0, xmax=1000,                          % Set x-axis range
    font=\scriptsize,             % Sets the base size for everything
    title style={yshift=-5pt, font=\small},    % Optional: Make title slightly bigger than the rest
    label style={font=\scriptsize}, % Font for xlabel and ylabel
    tick label style={font=\scriptsize}, % Font for the numbers on the axes
      xshift=11.7cm,
      width=6.5cm,
      height=5.5cm,
            legend pos=north east,
            legend cell align={left},
            legend style={font=\tiny,inner sep=1pt,row sep=-2pt},
]

\addplot[color=blue, thick] coordinates {(50,0.9496007015558312) (100,0.9433701489854492) (150,0.935205341841942) (200,0.9333489943887195) (250,0.9303813985404417) (300,0.9288850774062749) (350,0.9286613601463324) (400,0.929061804216448) (450,0.9285255334520865) (500,0.9271125030967432) (550,0.9298041075363036) (600,0.9287077367063746) (650,0.9282378376854579) (700,0.9281181416259152) (750,0.9267185063687005) (800,0.9286897879317082) (850,0.9251499124939538) (900,0.9259727394002606) (950,0.9268771123599007) (1000,0.9280628073961095)};
\addlegendentry{MultiPPI}

\addplot[color=red, thick, densely dashed] coordinates {(50,0.944409601227813) (100,0.9385569764187278) (150,0.9342760713931344) (200,0.9323295544745394) (250,0.9344328055133783) (300,0.929775106353744) (350,0.9296865067456744) (400,0.931480279927532) (450,0.9287010319765756) (500,0.9291913440631457) (550,0.9290507498888383) (600,0.929153195213068) (650,0.9273878580012023) (700,0.9275339395822341) (750,0.9294422219180304) (800,0.9267493637694326) (850,0.9258932418576218) (900,0.9267631287256513) (950,0.9260098509426862) (1000,0.9261048393939854)};
\addlegendentry{MultiPPI (Restricted)}

\end{axis}
\end{tikzpicture}
}
    
    \textrm{\bf \scriptsize Experiment 2: ProcessBench}
\scalebox{0.8}{
\centering
\begin{tikzpicture}
\begin{axis}[
    title={Coverage},
    xlabel={Budget},
    xlabel style={yshift=5pt},   % Moves xlabel closer to the axis
    grid=both,
    outer sep=0pt,
    enlarge x limits=false,    minor grid style={dashed, gray!50},
    major grid style={solid, gray!80},
    width=\textwidth,
    yticklabel style={/pgf/number format/fixed}, % No scientific notation
    yticklabel={\pgfmathparse{\tick*100}\pgfmathprintnumber{\pgfmathresult}\%},
    xmin=0, xmax=1000,                          % Set x-axis range
    font=\scriptsize,             % Sets the base size for everything
    title style={yshift=-5pt, font=\small},    % Optional: Make title slightly bigger than the rest
    label style={font=\scriptsize}, % Font for xlabel and ylabel
    tick label style={font=\scriptsize}, % Font for the numbers on the axes
      xshift=0.0cm,
      width=6.5cm,
      height=5.5cm,
            ymin=0.85, ymax=1,
]

\addplot[color=black, dashed, very thick] coordinates {(0,.95) (5000,.95)};
\addplot[color=blue, thick] coordinates {(50,0.936016) (100,0.933684) (150,0.93258) (200,0.93024) (250,0.927264) (300,0.927026) (350,0.928326) (400,0.927124) (450,0.924768) (500,0.924354) (550,0.9237) (600,0.923138) (650,0.921582) (700,0.921322) (750,0.9212) (800,0.92029) (850,0.920528) (900,0.920106) (950,0.921034) (1000,0.920898)};
\addlegendentry{MultiPPI}

\addplot[color=red, thick, densely dashed] coordinates {(50,0.938018) (100,0.934368) (150,0.930828) (200,0.92988) (250,0.929184) (300,0.929022) (350,0.926872) (400,0.926476) (450,0.92551) (500,0.923524) (550,0.923768) (600,0.923436) (650,0.923254) (700,0.92325) (750,0.921804) (800,0.92087) (850,0.920632) (900,0.919598) (950,0.919906) (1000,0.919544)};
\addlegendentry{MultiPPI (Restricted)}

\legend{};
\end{axis}
\begin{axis}[
    title={Mean-squared CI width (frac. of classical)},
    xlabel={Budget},
    xlabel style={yshift=5pt},   % Moves xlabel closer to the axis
    grid=both,
    outer sep=0pt,
    enlarge x limits=false,    minor grid style={dashed, gray!50},
    major grid style={solid, gray!80},
    width=\textwidth,
    yticklabel style={/pgf/number format/fixed}, % No scientific notation
    yticklabel={\pgfmathparse{\tick*100}\pgfmathprintnumber{\pgfmathresult}\%},
    xmin=0, xmax=1000,                          % Set x-axis range
    font=\scriptsize,             % Sets the base size for everything
    title style={yshift=-5pt, font=\small},    % Optional: Make title slightly bigger than the rest
    label style={font=\scriptsize}, % Font for xlabel and ylabel
    tick label style={font=\scriptsize}, % Font for the numbers on the axes
      xshift=5.9cm,
      width=6.5cm,
      height=5.5cm,
]

\addplot[color=blue, thick] coordinates {(50,0.7104284432326666) (100,0.6361804462741223) (150,0.5940047059860405) (200,0.5661303592061417) (250,0.5467412687554667) (300,0.53187574786712) (350,0.5205079192547275) (400,0.5108925702552023) (450,0.5026650395712908) (500,0.4958221120760036) (550,0.4903616984599653) (600,0.48518346433162696) (650,0.480221652270261) (700,0.47659221682635233) (750,0.47322376929977994) (800,0.4702380765212926) (850,0.4669210231427845) (900,0.46449234999354055) (950,0.4624221010837662) (1000,0.4603441459609113)};
\addlegendentry{MultiPPI}

\addplot[color=red, thick, densely dashed] coordinates {(50,0.7103746351985277) (100,0.6353667103229861) (150,0.5930623028207371) (200,0.5655594921813271) (250,0.5466380308208925) (300,0.5317741803673091) (350,0.5193398117109714) (400,0.5100757028603162) (450,0.5020793202965687) (500,0.4955333871227922) (550,0.48971347354611766) (600,0.4847425943890366) (650,0.4808520638209038) (700,0.477073501214465) (750,0.47325741735640775) (800,0.4699571518255548) (850,0.46747143402388447) (900,0.46489478489925323) (950,0.4623208294779253) (1000,0.4602580838886658)};
\addlegendentry{MultiPPI (Restricted)}

\legend{};
\end{axis}
\begin{axis}[
    title={Mean-squared Error (frac. of classical)},
    xlabel={Budget},
    xlabel style={yshift=5pt},   % Moves xlabel closer to the axis
    grid=both,
    outer sep=0pt,
    enlarge x limits=false,    minor grid style={dashed, gray!50},
    major grid style={solid, gray!80},
    width=\textwidth,
    yticklabel style={/pgf/number format/fixed}, % No scientific notation
    yticklabel={\pgfmathparse{\tick*100}\pgfmathprintnumber{\pgfmathresult}\%},
    xmin=0, xmax=1000,                          % Set x-axis range
    font=\scriptsize,             % Sets the base size for everything
    title style={yshift=-5pt, font=\small},    % Optional: Make title slightly bigger than the rest
    label style={font=\scriptsize}, % Font for xlabel and ylabel
    tick label style={font=\scriptsize}, % Font for the numbers on the axes
      xshift=11.7cm,
      width=6.5cm,
      height=5.5cm,
            legend pos=north east,
            legend cell align={left},
            legend style={font=\tiny,inner sep=1pt,row sep=-2pt},
]

\addplot[color=blue, thick] coordinates {(50,0.7416967483228185) (100,0.6718947012152297) (150,0.634458521221496) (200,0.6108025173290396) (250,0.5937414057867306) (300,0.5810104924542603) (350,0.5660112713132612) (400,0.5581287489166212) (450,0.5533334511453084) (500,0.5470577605999188) (550,0.5422202868635849) (600,0.5368751052839603) (650,0.5345443236510281) (700,0.5315798752273319) (750,0.5284327821128794) (800,0.5258934507824748) (850,0.5238361690821306) (900,0.5216083150432199) (950,0.5178458620961298) (1000,0.5167790894161537)};
\addlegendentry{MultiPPI}

\addplot[color=red, thick, densely dashed] coordinates {(50,0.7375200693029724) (100,0.6699749165571814) (150,0.6366488045293939) (200,0.6099611093948141) (250,0.590014579123059) (300,0.5757540022226781) (350,0.5672717043225987) (400,0.5593309996398323) (450,0.5503492663348937) (500,0.5465708543474609) (550,0.5417653690168481) (600,0.5366346354003534) (650,0.5319919037207459) (700,0.5281181825316223) (750,0.5279104855264986) (800,0.5252051916333937) (850,0.5228088675905042) (900,0.5215331564885237) (950,0.5189633470154591) (1000,0.5174524723413267)};
\addlegendentry{MultiPPI (Restricted)}

\end{axis}
\end{tikzpicture}
}
    
    \textrm{\bf \scriptsize Experiment 3: Factuality}
\scalebox{0.8}{
\centering
\begin{tikzpicture}
\begin{axis}[
    title={Coverage},
    xlabel={Budget},
    xlabel style={yshift=5pt},   % Moves xlabel closer to the axis
    grid=both,
    outer sep=0pt,
    enlarge x limits=false,    minor grid style={dashed, gray!50},
    major grid style={solid, gray!80},
    width=\textwidth,
    yticklabel style={/pgf/number format/fixed}, % No scientific notation
    yticklabel={\pgfmathparse{\tick*100}\pgfmathprintnumber{\pgfmathresult}\%},
    xmin=0, xmax=1000,                          % Set x-axis range
    font=\scriptsize,             % Sets the base size for everything
    title style={yshift=-5pt, font=\small},    % Optional: Make title slightly bigger than the rest
    label style={font=\scriptsize}, % Font for xlabel and ylabel
    tick label style={font=\scriptsize}, % Font for the numbers on the axes
      xshift=0.0cm,
      width=6.5cm,
      height=5.5cm,
            ymin=0.85, ymax=1,
]

\addplot[color=black, dashed, very thick] coordinates {(0,.95) (5000,.95)};
\addplot[color=blue, thick] coordinates {(50,0.931684) (100,0.925516) (150,0.921914) (200,0.918916) (250,0.917062) (300,0.915444) (350,0.912292) (400,0.91245) (450,0.910884) (500,0.910956) (550,0.909946) (600,0.909544) (650,0.90814) (700,0.908614) (750,0.90794) (800,0.907174) (850,0.906936) (900,0.906322) (950,0.906064) (1000,0.906346)};
\addlegendentry{MultiPPI}

\addplot[color=red, thick, densely dashed] coordinates {(50,0.93072) (100,0.926848) (150,0.924302) (200,0.92173) (250,0.91794) (300,0.915694) (350,0.915026) (400,0.91378) (450,0.913466) (500,0.912188) (550,0.91237) (600,0.910498) (650,0.91021) (700,0.90968) (750,0.908588) (800,0.908268) (850,0.908248) (900,0.907662) (950,0.907436) (1000,0.906722)};
\addlegendentry{MultiPPI (Restricted)}

\legend{};
\end{axis}
\begin{axis}[
    title={Mean-squared CI width (frac. of classical)},
    xlabel={Budget},
    xlabel style={yshift=5pt},   % Moves xlabel closer to the axis
    grid=both,
    outer sep=0pt,
    enlarge x limits=false,    minor grid style={dashed, gray!50},
    major grid style={solid, gray!80},
    width=\textwidth,
    yticklabel style={/pgf/number format/fixed}, % No scientific notation
    yticklabel={\pgfmathparse{\tick*100}\pgfmathprintnumber{\pgfmathresult}\%},
    xmin=0, xmax=1000,                          % Set x-axis range
    font=\scriptsize,             % Sets the base size for everything
    title style={yshift=-5pt, font=\small},    % Optional: Make title slightly bigger than the rest
    label style={font=\scriptsize}, % Font for xlabel and ylabel
    tick label style={font=\scriptsize}, % Font for the numbers on the axes
      xshift=5.9cm,
      width=6.5cm,
      height=5.5cm,
]

\addplot[color=blue, thick] coordinates {(50,0.7545364668479843) (100,0.7255458162063846) (150,0.7119201995335014) (200,0.7033429812767781) (250,0.696524071161093) (300,0.6916367464120456) (350,0.6884894463427089) (400,0.6855218289916332) (450,0.6834237878735445) (500,0.6813894001504142) (550,0.67951864487514) (600,0.6782205725951692) (650,0.6766564862096361) (700,0.6755847039022301) (750,0.6745882831724183) (800,0.6740555022443864) (850,0.6732813897423021) (900,0.6727000496701429) (950,0.6722167921964767) (1000,0.6714457175976168)};
\addlegendentry{MultiPPI}

\addplot[color=red, thick, densely dashed] coordinates {(50,0.7546919891124703) (100,0.7264575235638037) (150,0.7138411916595542) (200,0.705318049356559) (250,0.6994838122370611) (300,0.6947826491856386) (350,0.6907713549920115) (400,0.6878615793518422) (450,0.6857378326776609) (500,0.6836570874070872) (550,0.6816466010392656) (600,0.6800658785909818) (650,0.6787083486202546) (700,0.6776799697953044) (750,0.6765784277590318) (800,0.6754752258388018) (850,0.6751655396684646) (900,0.6743394913152534) (950,0.6734953468258082) (1000,0.6727259078236244)};
\addlegendentry{MultiPPI (Restricted)}

\legend{};
\end{axis}
\begin{axis}[
    title={Mean-squared Error (frac. of classical)},
    xlabel={Budget},
    xlabel style={yshift=5pt},   % Moves xlabel closer to the axis
    grid=both,
    outer sep=0pt,
    enlarge x limits=false,    minor grid style={dashed, gray!50},
    major grid style={solid, gray!80},
    width=\textwidth,
    yticklabel style={/pgf/number format/fixed}, % No scientific notation
    yticklabel={\pgfmathparse{\tick*100}\pgfmathprintnumber{\pgfmathresult}\%},
    xmin=0, xmax=1000,                          % Set x-axis range
    font=\scriptsize,             % Sets the base size for everything
    title style={yshift=-5pt, font=\small},    % Optional: Make title slightly bigger than the rest
    label style={font=\scriptsize}, % Font for xlabel and ylabel
    tick label style={font=\scriptsize}, % Font for the numbers on the axes
      xshift=11.7cm,
      width=6.5cm,
      height=5.5cm,
            legend pos=north east,
            legend cell align={left},
            legend style={font=\tiny,inner sep=1pt,row sep=-2pt},
]

\addplot[color=blue, thick] coordinates {(50,0.7949081400250686) (100,0.7786140849376558) (150,0.7727034268397065) (200,0.7688334978778444) (250,0.7656795057417095) (300,0.7631636433328535) (350,0.7660535033967734) (400,0.7622304664272552) (450,0.7635189532789504) (500,0.7603842564882731) (550,0.7598959203567006) (600,0.7586542968122403) (650,0.7606197542972088) (700,0.7583829198571969) (750,0.7593437537641989) (800,0.7594086697945458) (850,0.7576758823212617) (900,0.7585728076381071) (950,0.7587632228894517) (1000,0.7582761252461898)};
\addlegendentry{MultiPPI}

\addplot[color=red, thick, densely dashed] coordinates {(50,0.7973455913189056) (100,0.7778161281512482) (150,0.7675868115691713) (200,0.7647830323119169) (250,0.7656927768168506) (300,0.7650503047923587) (350,0.7625271652819368) (400,0.763126888177619) (450,0.7590195808850055) (500,0.7593930233868563) (550,0.7584588800310204) (600,0.7605764124389637) (650,0.7571096526959882) (700,0.7578682686786946) (750,0.7601932925579366) (800,0.7585011213429403) (850,0.7590207310747368) (900,0.7587973215673303) (950,0.757816564819311) (1000,0.7579175443679933)};
\addlegendentry{MultiPPI (Restricted)}

\end{axis}
\end{tikzpicture}
}
    \caption{Comparison of MultiPPI for $\mathcal{I}=2^{\{1,\dots,k\}}$ (default settings) with MultiPPI (Restricted), as defined in \Cref{sec:restricted}.}
    \label{fig:restricted}
\end{figure}

\subsection{Autorater accuracy scaling}

\begin{figure}[!h]
    \centering
    \includegraphics[width=0.8\linewidth]{figures/download__35_.png}
    \caption{Performance at determination of process error vs. word budget. This is calculated via the procedure described in \autoref{sec:experimental-details}. The majority of the improvement observed due to thinking occurs once 500 words of thought is reached, and plateaus around 1,000 words of thought.\looseness=-1}
    \label{fig:scaling-budget}
\end{figure}
\begin{figure}[!h]
    \centering
    \includegraphics[width=0.5\linewidth]{figures/scaling_accuracy_processbench.png}
    \caption{Performance at determination of process error vs. word budget. This is calculated via the procedure described in \autoref{sec:experimental-details}. The majority of the improvement observed due to thinking occurs once 500 words of thought is reached, and plateaus around 1,000 words of thought.\looseness=-1}
    \label{fig:scaling-budget}
\end{figure}
\begin{figure}[!h]
    \centering
    \includegraphics[width=\linewidth]{figures/scaling_accuracy_biographies.png}
    \caption{Performance at factuality evaluation with increasing number of agents and rounds of debate. Soft accuracy awards half a point to reporting an uncertain answer, while hard accuracy awards nothing. }
    \label{fig:scaling-accuracy-biographies}
\end{figure}
\begin{figure}[!h]
    \centering
    \includegraphics[width=0.5\linewidth]{figures/scaling_uncertainty_biographies.png}
    \caption{Proportion of uncertain predictions by number of agents and rounds of debate. An increased number of agents leads to fewer uncertain predictions, and almost all predictions are certain by the end of the second round of debate.}
    \label{fig:scaling-uncertainty-biographies}
\end{figure}

\begin{figure}[!h]
    \centering
    \includegraphics[width=\linewidth]{figures/Screenshot_2026-02-28_11.08.01_AM.png}
    \caption{Different schemes for evaluation with autoraters on the ProcessBench dataset. Gray: classical sampling---no autoraters. Orange: pure autoraters in increasing order of thinking budget (from top to bottom, word budget is 50, 500, 1000, 1500, 2000 and 3000) ---note that the bias is increasingly pronounced with thinking budget. Green: various schemes for debiasing autoraters (from top to bottom, MultiPPI, vector PPI++, followed by scalar PPI++ with 3000, 2000, 1500, and 1000 word thinking budget, respectively).} %\adam{I don't quite understand this figure. The takeaway seems to be that, counter-intuitively, simply making models smarter won't necessarily solve the bias problem: debiasing methods like multi-ppi are necessary. What are the different levels of green?} \charlie{That's right for the takeaway. The different levels of green are different debiasing methods (e.g. our baselines: PPI++ with different models) with MultiPPI at the top. I've updated the figure to clarify this, but I think maybe it's not a very good figure, and we can remove it.}}
    \label{fig:need-to-debias}
\end{figure}

\clearpage
\section{Additional theoretical results}

%\adam{Can you add a brief paragraph here explaining what these next theoretical results address, and why they are important / interesting?}

In this section, we provide additional theoretical guarantees on the MultiPPI estimator. First, we establish finite-sample bounds on the estimation risk, under bounded, sub-Gaussian, and auto-regressive data distributions, and extend these results to the Ledoit-Wolf shrinkage estimator. Next, we characterize the behavior of the estimator at the boundaries of the budget spectrum. We prove that in the extreme low-budget limit, the allocation deterministically isolates the single model with the best correlation-to-cost ratio; conversely, in the large-budget regime, we show that the continuous relaxation of our discrete allocation problem is asymptotically optimal. Finally, we formally explain the empirically observed phenomenon of coverage decay, which occurs when the unlabeled budget grows while the labeled sample size remains fixed.

\subsection{Finite-sample bounds}
\label{app:finite-sample-bounds}
We consider the setting of \autoref{app:generalized}, in which we may have several budget constraints. For the time being, we fix $a=(1,0,\dots,0)$ as in all experiments. Let $I^0\in\mathcal{I}$ contain $1$. A procedure which is similar to classical sampling is the following: Consider the choice $\ul{n}^0,\ul{\lambda}^0$ defined such that $n^0_I=0$ if $I\neq I^0$, and let $n^0_{I^0}$ be the maximal choice afforded by the budget (i.e. $n^0_{I^0} = \max_{1\leq \ell\leq m}\left \lfloor B^{(\ell)}/c_{I^0}^{(\ell)}\right\rfloor$). Then setting $\lambda^0_I=0$ if $I\neq I^0$, and $\lambda^0_{I^0}$ to be $a$ restricted to $I^0$, we recover the classical estimator
\[
\frac{1}{n^0_{I^0}} \sum_{j=1}^{n^0_{I^0}} X_1^{(j)}
\]
which has MSE $\sigma^2_1/n^0_{I^0}$, where $\sigma_1^2=\Sigma_{11}$. We let $\sigma^2_{\text{classical}}:=\sigma^2_1/n^0_{I^0}$ denote this quantity.

We will compare $\hat{\theta}_{\text{MultiPPI}}$ to this in finite samples. Let $\wh{\Sigma}_N$ denote the empirical covariance matrix constructed from $N$ i.i.d. samples from $P$, and let $\wh{\underline{n}},\wh{\underline{\lambda}}$ denote the solution to MultiAllocate$(\wh{\Sigma}_N)$, i.e. the minimizer of
\[
\wh{R}_N(\ul{n},\ul{\lambda}) = \sum_{I\in\mathcal{I}:n_I>0}\frac{1}{n_I}\lambda_I^\top\wh{\Sigma}_N\lambda_I
\]
such that $\mathbf{U}$ and $\mathbf{B}$ hold. On the other hand, let ${\ul{n}}^*,{\ul{\lambda}}^*$ denote the solution to MultiAllocate$(\Sigma)$, i.e. the minimizer of
\[
R(\ul{n},\ul{\lambda}) = \sum_{I\in\mathcal{I}:n_I>0}\frac{1}{n_I}\lambda_I^\top\Sigma\lambda_I
\]
such that $\mathbf{U}$ and $\mathbf{B}$ hold. In this section, we bound
\[
R(\wh{\ul{n}},\wh{\ul{\lambda}}) - R({\ul{n}}^*,{\ul{\lambda}}^*).
\]

\begin{theorem}\label{thm:finite-sample-v2-1}
    Let $\gamma_{\min}$ denote the minimal eigenvalue of $\Sigma$, and $\delta=\|\Sigma-\wh{\Sigma}_N\|_{op}$. Then for all $\delta\leq\gamma_{\min}/2$,
    \[
    R(\wh{\ul{n}},\wh{\ul{\lambda}}) \leq R({\ul{n}}^*,{\ul{\lambda}}^*) + 4\frac{\delta}{\gamma_{\min}}\cdot \sigma^2_{\classical}
    \]
\end{theorem}

\begin{corollary}\label{cor:bounded}
    Suppose that $X_i\in[0,1]$ almost surely. Then with high probability,
    \[
    R(\wh{\ul{n}},\wh{\ul{\lambda}}) \leq R({\ul{n}}^*,{\ul{\lambda}}^*) + c\left(\frac{\gamma_{\max}^{1/2}}{\gamma_{\min}} \sqrt{\frac{k\log k}{N}} + \frac{1}{\gamma_{\min}}\frac{k\log k}{N}\right) \sigma^2_{\classical}
    \]
    for a universal constant $c$, and so
    \[
    \mathbb{E}R(\wh{\ul{n}},\wh{\ul{\lambda}}) \leq R({\ul{n}}^*,{\ul{\lambda}}^*) + c'\left(\frac{\gamma_{\max}^{1/2}}{\gamma_{\min}} \sqrt{\frac{k}{N}} + \frac{1}{\gamma_{\min}} \frac{k}{N}\right)\sigma^2_{\classical}
    \]
    for another constant $c'$, where the expectation is taken over the $N$ labeled samples used to construct $\wh{\Sigma}_N$.
\end{corollary}

\begin{corollary}\label{cor:subg}
    Suppose that $X$ is a subgaussian with variance proxy $K$. Then
    \[
    \Ex R(\wh{\ul{n}},\wh{\ul{\lambda}})\leq R(\ul{n}^*,\ul{\lambda}^*)+c' K^2\left(\sqrt{\frac{k}{N}}+\frac{k}{N}\right)\sigma^2_{\classical}
    \]
\end{corollary}

In the AR(1) model, and with bounded observations, choosing $N\gg k$ in the limit $k,N\to\infty$ is enough that $\mathbb{E}R(\wh{\ul{n}},\wh{\ul{\lambda}}) \to R({\ul{n}}^*,{\ul{\lambda}}^*)$. This follows as a special case of the following result.

\begin{corollary}\label{cor:gersch}
Suppose, in addition to the conditions of \Cref{thm:finite-sample-v2-1}, that $X_1,X_2,\dots$ is a stochastic process such that $\Var X_t>c$ for all $t$, and $\Corr(X_t,X_s)\leq (1-\rho)\rho^{|t-s|}$ for some $0<c,\rho<1$. Then we have
\[
\Ex R(\wh{\ul{n}},\wh{\ul{\lambda}}) = R(\ul{n},\ul{\lambda}) +o(1)
\]
whenever $k/N = o(1)$.
\end{corollary}

Lastly, below we include a finite-sample bound for the performance of the estimator using Ledoit-Wolf shrinkage.
\begin{theorem}[Finite-sample bounds specialized to Ledoit-Wolf shrinkage]\label{thm:ledoit}
    Let $\widehat{\Sigma}_N^
    \text{LW}$ denote the Ledoit-Wolf shrinkage estimator of $\Sigma$ based on $N$ samples. Let $\gamma_{\min}$ denote the minimum eigenvalue of $\Sigma$, and suppose that $X\in \R^k$ is sub-Gaussian with proxy $K$. Lastly, suppose that $\Sigma$ is not a multiple of the identity. Then for absolute constants $c_1,c_2$, we have
    \[
    \Ex \left[\left(\hat{\theta}_{\multiPPI(\wh{\Sigma})}-\theta^*\right)^2\right] \leq \mathcal{V}_B+\frac{4\sigma^2_{\text{classical}}}{\gamma_{\min}} \frac{1}{\sqrt{N}}\sqrt{c_1K^4\gamma_{\max}^2 k^2 + c_2 K^8 \gamma_{\max} k^3/a^2}
    \]
    where $a^2:=\frac{1}{k}\left\|\Sigma - I\cdot \frac{\tr(\Sigma)}{k}\right\|_F^2$.
\end{theorem}
For proofs, see \Cref{sec:finite_sample_proofs}.


% \subsection{Finite-sample bounds}\label{sec:finite}
% \begin{theorem}[Stability of Multi-PPI]\label{thm:stability} Let $\Sigma,\Sigma'\succ0$, and $\epsilon:=\|\Sigma-\Sigma'\|_{\text{op}}/\lambda_{\min}(\Sigma)$. Then
% \[
% \Ex_\Sigma\left[\left(\hat{\theta}_{\multiPPI(\Sigma')}-\theta^*\right)^2\right] \leq (1+2\epsilon)^2 \Ex_\Sigma\left[\left(\hat{\theta}_{\multiPPI(\Sigma)}-\theta^*\right)^2\right]
% \]
% where the expectations are taken with respect to any distribution with covariance $\Sigma$.
% \end{theorem}
% This immediately yields a bound on the mean-squared error of Multi-PPI in the case where a small number of fully-labeled samples $m$ is reserved for estimating $\widehat{\Sigma}$. Lastly, set $\gamma=\lambda_{\max}(\Sigma)^{1/2}/\lambda_{\min}(\Sigma)$.
% \begin{corollary}\label{cor:finite} Suppose that $X_i\in[0,1]$ almost surely. Then for an absolute constant $c'$,
% \[
% \Ex_\Sigma\left[\left(\hat{\theta}_{\multiPPI(\widehat{\Sigma})}-\theta^*\right)^2\right] \leq \left(1+c'\gamma k\sqrt{\frac{k}{m}} + c'\gamma^2\frac{k^2}{m}\right) \Ex_\Sigma\left[\left(\hat{\theta}_{\multiPPI(\Sigma)}-\theta^*\right)^2\right]
% \]
% Thus any choice of burn-in such that $k^2\max\{\gamma,\gamma^2\}\ll m\ll B$ is asymptotically optimal, in that such a procedure achieves the asymptotic variance $V$ specified in \autoref{thm:asymptotic}.
% \end{corollary}
% For example, in the AR(1) model, it is enough to take $k^2\ll m\ll B$; if features are equicorrelated, then one may take $k^{5/2}\ll m \ll B$.
% \begin{proof}[Proof of \cref{thm:stability}]
%     Recall that the solution to
% \[
% \text{M}(\Sigma):=\underset{\substack{\underline{n},\underline{\lambda} \\ \mathbf{B}(\underline{n}),\mathbf{U}(\underline{n},\underline{\lambda})}}{\argmin} \sum_{I\in\mathcal{I},n_I>0} \frac{1}{n_I}\lambda_I^\top\Sigma_I\lambda_I
% \]
% satisfies
% \[
% n_I=\left(\frac{B}{c_I}\right) \frac{\sqrt{c_I\lambda_I^\top\Sigma_I\lambda_I}}{\sum_{J\in\mathcal{I}} \sqrt{c_J\lambda_J^\top\Sigma_J\lambda_J}}
% \]
% by the Cauchy-Schwartz inequality. Now let $\underline{n}',\underline{\lambda}'$ denote the corresponding solutions for an arbitrary other SPD matrix $\Sigma'$. Let us assume that $c_I>0$ for all $I$, so that $n_I>0$ for all $I$. We proceed to bound
% \[
% \sum_{I\in\mathcal{I}} \frac{1}{n'_I}\lambda_I'^\top\Sigma_I\lambda_I'
% \]
% Rearranging, we have
% \begin{align*}
%     B\sum_{I\in\mathcal{I}} \frac{1}{n'_I}\lambda_I'^\top\Sigma_I\lambda_I' &= \left(\sum_I \sqrt{c_I}\frac{\|\lambda_I'\|_{\Sigma_I}^2}{\|\lambda'_I\|_{\Sigma'_I}}\right)\left(\sum_J \sqrt{c_J} \|\lambda'_J\|_{\Sigma'_J}\right) \\
%     &= \left(\sum_I \sqrt{c_I} \|\lambda'_I\|_{\Sigma_I'}\cdot\left(\frac{\|\lambda_I'\|_{\Sigma_I}}{\|\lambda'_I\|_{\Sigma'_I}}\right)^2\right)\left(\sum_J \sqrt{c_J} \|\lambda'_J\|_{\Sigma'_J}\right) \\
%     &\leq \left(\sum_I \sqrt{c_I} \|\lambda'_I\|_{\Sigma_I'}\right)^2 \cdot\sup_I\left(\frac{\|\lambda_I'\|_{\Sigma_I}}{\|\lambda'_I\|_{\Sigma'_I}}\right)^2
% \end{align*}
% by H\"{o}lder's inequality. The term on the right is bounded by
% \[
% \frac{\lambda'^\top\Sigma\lambda'}{\lambda'^\top \Sigma'\lambda'} = 1 + \frac{\lambda'^\top(\Sigma- \Sigma')\lambda'}{\lambda'^\top\Sigma'\lambda'}\leq 1+\frac{\|\Sigma-\Sigma'\|}{\lambda_{\min}(\Sigma')}
% \]
% where the operator norm is intended.

% % The term on the left is exactly $a^\top S'a$, where $S'=\left(\sum_I c_I^{-1}P_I^\top\Sigma_I^{-1}P_I\right)^\dagger$, as a consequence of the following lemma.

% % \begin{lemma}
% %     \[
% %     \min_{\underline{\lambda}\,:\,\sum_I P_I^\top\lambda_I = a} \left(\sum_I \sqrt{c_I} \|\lambda_I\|_{\Sigma_I}\right)^2 = a^\top Sa,\quad\quad S = \left(\sum_I c_I^{-1} P_I^\top\Sigma_I^{-1}P_I\right)^\dagger
% %     \]
% % \end{lemma}
% % \begin{proof}
% %     We first prove the lower bound. We must show that for all such $\underline{\lambda}$, we have
% %     \[
% %     \sum_I\sqrt{c_I\lambda_I^\top\Sigma_I\lambda_I} \geq \sqrt{a^\top Sa}
% %     \]
% %     Suppose not, and consider setting $n_I$ as defined above. The variance of the multi-PPI estimator with this particular choice is
% %     \[
% %     \sum_I \frac{1}{n_I}\lambda_I^\top\Sigma_I\lambda_I = \frac{1}{B}\left(\sum_I \sqrt{c_I}\|\lambda_I\|_{\Sigma_I}\right)^2 \geq \min_{\underline{n}\,:\,\mathbf{B}(\underline{n})} a^\top S_{\underline{n}}a
% %     \]
% %     We proved in lemma 4 that
% %     \[
% %     \min_{\underline{\lambda}:\mathbf{U}} \sum_{I:n_I>0} \frac{1}{n_I} \lambda_I^\top\Sigma_I \lambda_I = a^\top\left(\sum_I n_I P_I^\top \Sigma_I^{-1} P_I\right)^\dagger a
% %     \]
% %     for every fixed $\underline{n}$. With the choice of $n_I$ given above, we then have
% %     \[
% %     \frac{1}{B}\left(\sum_I \sqrt{c_I}\|\lambda_I\|_{\Sigma_I}\right)^2  = \frac{1}{B}a^\top \left(\sum_I c_I^{-1/2} \|\lambda_I\|_{\Sigma_I} P_I^\top \Sigma_I^{-1} P_I\right)^\dagger\left(\sum_J \sqrt{c_J}\|\lambda_J\|_{\Sigma_J}\right)a
% %     \]
% % \end{proof}

% We now bound the term on the left.

% \begin{lemma}[The dual formulation]\label{lemma:dual}
%     We have
%     \[
%     v(\Sigma) := \inf_{\underline{\lambda}:\mathbf{U}} \sum_I \sqrt{c_I}\|\lambda_I\|_{\Sigma_I} = \sup_{x} a^\top x
%     \]
%     where the supremum is taken over all $x\in\R^k$ such that $x_I\Sigma_I x_I\leq c_I^{-1}$ for all $I\in\mathcal{I}$.
% \end{lemma}
% \begin{proof}
%     See \autoref{sec:dual}.
% \end{proof}

% \begin{lemma}[Stability of $v$]
%     Let $\epsilon := \|\Sigma-\Sigma'\|/\lambda_{\min}(\Sigma)$. Then we have
%     \[
% 1+\epsilon\geq    \left(\frac{v(\Sigma)}{v(\Sigma')}\right)^2 \geq 1-\epsilon
%     \]
% \end{lemma}
% \begin{proof}
%     Putting $\|\Sigma-\Sigma'\|=\delta,$ we have
%     \[
%     (1-\epsilon)\Sigma\prec \Sigma'\prec (1+\epsilon)\Sigma
%     \]
%     for $\epsilon=\delta/\lambda_{\min}$. To see this, consider that $E=\Sigma^{-1/2}(\Sigma'-\Sigma)\Sigma^{-1/2}$ satisfies $\|E\|\leq \frac{\delta}{\lambda_{\min}}=\epsilon$. Then $-\epsilon I\prec E \prec \epsilon I$. Thus
%     \[
%     (1-\epsilon)x_I\Sigma_I x_I\leq x_I\Sigma_I'x_I \leq (1+\epsilon) x_I\Sigma_I x_I
%     \]
%     Now suppose that $x_I\Sigma_I'x_I\leq c_I^{-1}$ for all $I$. We then have that $(1-\epsilon)x_I\Sigma_Ix_I\leq c_I^{-1}$ for all $I$, and so
%     \[
%     v(\Sigma) \geq (1-\epsilon)^{1/2}a^\top x
%     \]
%     But this holds for all $x$ lying in the feasible set for $\Sigma'$, and so taking a supremum, we see that
%     \[
% \sqrt{1+\epsilon}v(\Sigma')\geq   v(\Sigma) \geq \sqrt{1-\epsilon} v(\Sigma')
%     \]
%     or in other words
%     \[
%  1+\epsilon\geq    \left(\frac{v(\Sigma)}{v(\Sigma')}\right)^2 \geq 1-\epsilon
%     \]
%     as desired.
% \end{proof}
% \begin{remark}
%     Reversing $\Sigma$ and $\Sigma'$ above, if $\epsilon'=\|\Sigma-\Sigma'\|/\lambda_{\min}(\Sigma')$ we have similarly that $v(\Sigma')^2\leq (1+\epsilon')v(\Sigma)$. Lastly note that $\lambda_{\min}(\Sigma')\geq \lambda_{\min}(\Sigma)-\delta$, and so for all $\delta\leq \lambda_{\min}(\Sigma)/2$ we have $\epsilon'\leq 2\epsilon$
% \end{remark}
%     Taken together, we have bounded
%     \[
%     \sum_{I\in\mathcal{I}}\frac{1}{n'_I}\lambda_I'^\top \Sigma_I \lambda'_I \leq (1+2\epsilon)^2\cdot v(\Sigma)/B 
%     \]
%     where $v(\Sigma)/B$ is the mean-squared error of $\hat{\theta}_{\multiPPI(\Sigma)}$ under a distribution of covariance $\Sigma$.
% \end{proof}

% \subsection{Finite-sample bounds, V2}
% We consider the setting of \autoref{sec:generalized}, in which we may have several budget constraints. For the time being, we fix $a=(1,0,\dots,0)$ as in all experiments. Let $I^0\in\mathcal{I}$ contain $1$. A procedure which is similar to classical sampling is the following: Consider the choice $\ul{n}^0,\ul{\lambda}^0$ defined such that $n^0_I=0$ if $I\neq I^0$, and let $n^0_{I^0}$ be the maximal choice afforded by the budget (i.e. $n^0_{I^0} = \max_{1\leq \ell\leq m}\left \lfloor B^{(\ell)}/c_{I^0}^{(\ell)}\right\rfloor$). Then setting $\lambda^0_I=0$ if $I\neq I^0$, and $\lambda^0_{I^0}$ to be $a$ restricted to $I^0$, we recover the classical estimator
% \[
% \frac{1}{n^0_{I^0}} \sum_{j=1}^{n^0_{I^0}} X_1^{(j)}
% \]
% which has MSE $\sigma^2_1/n^0_{I^0}$, where $\sigma_1^2=\Sigma_{11}$. We let $\sigma^2_{\text{classical}}:=\sigma^2_1/n^0_{I^0}$ denote this quantity.

% We will compare $\hat{\theta}_{\text{MultiPPI}}$ to this in finite samples. Let $\wh{\Sigma}_N$ denote the empirical covariance matrix constructed from $N$ i.i.d. samples from $P$, and let $\wh{\underline{n}},\wh{\underline{\lambda}}$ denote the solution to MultiAllocate$(\wh{\Sigma}_N)$, i.e. the minimizer of
% \[
% \wh{R}_N(\ul{n},\ul{\lambda}) = \sum_{I\in\mathcal{I}:n_I>0}\frac{1}{n_I}\lambda_I^\top\wh{\Sigma}_N\lambda_I
% \]
% such that $\mathbf{U}$ and $\mathbf{B}$ hold. On the other hand, let ${\ul{n}}^*,{\ul{\lambda}}^*$ denote the solution to MultiAllocate$(\Sigma)$, i.e. the minimizer of
% \[
% R(\ul{n},\ul{\lambda}) = \sum_{I\in\mathcal{I}:n_I>0}\frac{1}{n_I}\lambda_I^\top\Sigma\lambda_I
% \]
% such that $\mathbf{U}$ and $\mathbf{B}$ hold. In this section, we bound
% \[
% R(\wh{\ul{n}},\wh{\ul{\lambda}}) - R({\ul{n}}^*,{\ul{\lambda}}^*).
% \]

% \begin{theorem}\label{thm:finite-sample-v2-1}
%     Let $\gamma_{\min}$ denote the minimal eigenvalue of $\Sigma$, and $\delta=\|\Sigma-\wh{\Sigma}_N\|_{op}$. Then for all $\delta\leq\gamma_{\min}/2$,
%     \[
%     R(\wh{\ul{n}},\wh{\ul{\lambda}}) \leq R({\ul{n}}^*,{\ul{\lambda}}^*) + 4\frac{\delta}{\gamma_{\min}}\cdot \sigma^2_{\classical}
%     \]
% \end{theorem}
% \charlie{So it seems to me that this result is roughly the same as/implied by \cref{thm:stability}, since the MSE written on the right-hand side is at most $\sigma^2_{\classical}$. Not sure if I'm doing something wrong, just wanted to check.}
% \begin{theorem}[Slight refinement]\label{thm:finite-sample-v2-2}
%     Let $\gamma_{\min,I}$ denote the minimal eigenvalue of $\Sigma_I$, and $\delta_I=\|\Sigma_I-\wh{\Sigma}_{N,I}\|_{op}$. Then whenever $\delta_I\leq\gamma_{\min,I}/2$ for all $I\in\mathcal{I}$,
%     \[
%     R(\wh{\ul{n}},\wh{\ul{\lambda}}) \leq R({\ul{n}}^*,{\ul{\lambda}}^*) + 4\max_{I\in\mathcal{I}}\frac{\delta_I}{\gamma_{\min,I}}\cdot \sigma^2_{\classical}
%     \]
% \end{theorem}
% \charlie{It would be great to prove something which is stronger when $|\mathcal{I}|$ is smaller. I haven't really been able to do that: I could only refine \cref{thm:finite-sample-v2-1} to the maximum written above, but this is not helpful when $I=\{1,\dots,k\}$ (i.e. the set of all indices) is included in $\mathcal{I}$, as is usually the case.}



% \begin{proof}[Proof of \cref{thm:finite-sample-v2-2}]
%     We can see that we can refine \autoref{eqn:Delta} to
%     \[
%     \Delta_N(\ul{n},\ul{\lambda}) \leq \sum_{I\in\mathcal{I}:n_I>0}\|\Sigma_I-\wh{\Sigma}_N\| \frac{1}{n_I} \|\lambda_I\|^2_2
%     \]
%     and \autoref{eqn:bound-norm} to
%     \[
%     \sigma^2_1/n_{I^0}^0 \geq \sum_{I\in\mathcal{I}:n_I>0} \gamma_{\min}(\Sigma_I)\frac{1}{n_I} \|\lambda_I\|^2_2
%     \]
%     from which the result follows.
% \end{proof}


\subsection{Behavior of the estimator in the limiting regimes}\label{sec:limiting}

In this section, we explain a certain limiting behavior of the estimator in the regime of very low budget. Let $X=(X_1,\dots,X_k)$ be a random vector of bounded second moment. We take $a=(1,0,\dots,0)$, so that our target is $\mathbb{E}[X_1]$. We consider the setting (as is the case in all experiments) in which $\mathcal{I}=\{1,\dots,k\}\cup \mathcal{I}_{\text{models}}$, where for each $I\in\mathcal{I}_{\text{models}}$ we have $1\not\in I$.

As in the experiments, we consider the budget model in which we have a fixed number of 

For $I\in\mathcal{I}_{\text{models}}$, $\rho_I$ denote the multiple correlation coefficient of $X_I$ with $X_1$; that is, let $\rho_I = \text{Cov}_I^\top \Sigma_I^{-1} \text{Cov}_I$, where we define $\text{Cov}_I:=(\text{Cov}(X_i,X_1))_{i\in I}$. The following result shows that, in the low-budget regime, MultiAllocate$(\Sigma)$ returns $n_I$ such that the only $I\in\mathcal{I}_{\text{models}}$ for which $n_I\neq0$ is the one which minimizes the correlation/cost ratio $\rho_I/c_I$.

\begin{theorem}\label{thm:limiting}
    Fix $B>0$ and consider the limit as $n_{[k]}\to\infty$. For each $I\in\mathcal{I}$, let $\alpha_I:=\rho_I/c_I$. Suppose that $I^*$ uniquely minimizes $\alpha_I$ over $I\in\mathcal{I}_{\text{models}}$. Then the solution to MultiAllocate$(\Sigma)$ satisfies 
    \[
    n_I\longrightarrow \frac{B}{c_{I}}\cdot \begin{cases}
        1 & I=I^* \\
        0 & I\neq I^*
    \end{cases} 
    \]
\end{theorem}

% \subsection{A note on the PPI++ Pareto frontier}\label{sec:pareto}
% Here we briefly discuss the Pareto frontier of models with which to perform PPI++, described in \autoref{sec:results}. 
% \[
% \hat{\theta}
% \]

\subsection{Rounding in the large budget regime}\label{sec:rounding}
In this section, we consider the suboptimality of the rounding scheme in the large budget regime. We consider the general setup in which we optimize
\[
V_B(\underline{n}) = a^\top \left(\sum_{I} n_I P_I^\top \Sigma_I^{-1} P_I\right)^\dagger a\quad\text{s.t.}\quad n_I\geq0,\, \sum_I c_In_I\leq B,\, \supp(a)\subseteq\bigcup\{I:n_I>0\}
\]
We let $\underline{n}^*_\text{frac}$ denote the solution to this problem over all $\underline{n}\in\mathbb{R}^{|\mathcal{I}|}_{\geq0}$, and $\underline{n}^*_{\text{int}}$ denote the solution over all $\underline{n}\in\mathbb{Z}^{|\mathcal{I}|}_{\geq0}$. Let $\underline{n}_{\text{round}}$ denote the component-wise floor of $\underline{n}^*_{\text{frac}}$. Here we show that
\[
\lim_{B\to\infty} \frac{V_B(\underline{n}_{\text{frac}})}{V_B(\underline{n}^*_{\text{int}})}=1
\]
This follows from the fact that
\[
V_B(\underline{n}^*_{\text{frac}}) \leq V_B(\underline{n}^*_{\text{int}}) \leq V_B(\underline{n}_{\text{round}})
\]
and the limit $V_B(\underline{n}^*_{\text{frac}}) / V_B(\underline{n}_{\text{round}}) \to1$, to be proven next. Consider the difference vector $\underline{\delta} = \underline{n}^*_{\text{frac}}-\underline{n}_{\text{round}}\in[0,1]^{|\mathcal{I}|}$. Now observe that there is some $\underline{\nu}^*\in\mathbb{R}_{\geq0}^{|\mathcal{I}|}$ such that
\[
BV_B(\underline{n}_{\text{frac}}^*)= V_1(\underline{\nu}^*)
\]
for all $B$, and equality holds if we take $\underline{n}_{\text{frac}}^*=B\underline{\nu}^*$. In particular, since we must have $\bigcup\{I:n^*_{\text{frac},I}>0\}\supseteq \supp(a)$, we may take the same to hold for $\underline{\nu}^*$. We therefore have
\begin{align*}
    BV_B(\underline{n}_{\text{round}})&=Ba^\top\left(B\sum_I \nu_I P_I^\top\Sigma_I^{-1}P_I + \sum_I \delta_I P_I^\top \Sigma_I^{-1}P_I\right)^\dagger a \\
    &= a^\top \left(\sum_I \nu_I P_I^\top\Sigma_I^{-1}P_I + \frac{1}{B}\sum_I \delta_I P_I^\top \Sigma_I^{-1}P_I\right)^\dagger a 
\end{align*}
Now since $\bigcup\{I:{\nu}^*_I>0\}\supseteq\supp(a)$, we may apply continuity of the inverse to conclude that
\[
\lim_{B\to\infty} BV_B(\underline{n}_{\text{round}}) = a^\top \left(\sum_I \nu_I^* P_I^\top\Sigma_I^{-1}P_I \right)^\dagger a =V_1(\underline{\nu}^*)
\]
and the limit is proven.

\subsection{The continuous problem and its properties}\label{sec:continuous_problem}
In this section, we study several properties of the the following continuous version of \Cref{eq:variance}, stated below:
\begin{equation}\label{eqn:continous-problem}
    \min_{\substack{\ul{\nu}\geq0,\, \sum_I \nu_I c_I\leq B \\ \supp(a) \subseteq \{I:\nu_I>0\}}} a^\top S(\ul{\nu})a,\quad\quad S(\ul{\nu}) = \left(\sum_{I\in\mathcal{I}} n_I\Sigma_I^\dagger\right)^\dagger
\end{equation}
where $\ul{\nu}$ may vary over $\R^{|\mathcal{I}|}$. Here, as is the case throughout, the inequalities $\ul{\nu}\geq0,\, \sum_I \nu_I c_I\leq B $ should be interpreted in the vector-valued sense (see \Cref{app:generalized}).

We are interested in the (set-valued) minimizers of \Cref{eqn:continous-problem}. More specifically, we are interested in the continuity of these minimizers in $\Sigma$, and how this relates to the discrete optimization problem of \Cref{eq:variance}.
To address these topics, we introduce the notation, following \citet{boyd}, that $\mathbf{S}^k_{++}$ denotes the set of symmetric positive definite $k\times k$ matrices. We are interested in the set-valued solution to the minimization problem
\begin{equation}\label{eqn:nu}
        \underline{\nu}^*(A) = \underset{\substack{\ul{\nu} :\ul{\nu}\geq0 \\ \sum_I \nu_Ic_I\leq B_0 \\ \supp(a)\subseteq \bigcup\{I:\nu_I>0\}}}{\argmin} a^\top \left(\sum_I \nu_I A_I^\dagger \right)^\dagger a
\end{equation}
for $A\in \mathbf{S}^k_{++}$. More formally, recalling the set of feasible allocations $K=\{\ul{\nu}:\ul{\nu}\geq0,\sum_I\nu_Ic_I\leq B\}\subseteq \R^{|\mathcal{I}|}$, we have
\begin{equation}\label{eqn:nustar}
\begin{split}
    \ul{\nu}^*:\mathbf{S}^k_{++}&\rightrightarrows K \\
    A&\mapsto \argmin_{\ul{\nu}\in K} F(\ul{\nu},A)
\end{split}    
\end{equation}
where
\begin{equation}\label{eqn:F}
F(\ul{\nu},A)=\begin{cases}
    a^\top \left(\sum_I \nu_I A_I^\dagger \right)^\dagger a & a \in \text{range}\left(\sum_I\nu_I A_I^\dagger\right) \\
    \infty & \text{otherwise.}    
\end{cases}
\end{equation}

%\adam{Can we move the supporting lemmas to the proofs section, and only state Theorem E.9 here?}\charlie{Yes, good call, done thanks!}

We now connect the continuous problem to the discrete problem. Below we recall the discrete problem: define
\begin{equation}\label{eqn:n}
        \underline{n}^*_t(A) = \underset{\substack{\ul{n}\in\mathbb{Z}^{|\mathcal{I}|} :\ul{n}\geq0 \\ \sum_I n_Ic_I\leq tB_0 \\ \supp(a)\subseteq \bigcup\{I:n_I>0\}}}{\argmin} a^\top \left(\sum_I n_I A_I^\dagger \right)^\dagger a
\end{equation}
Then we have the following:
\begin{lemma}\label{lemma:discrete-convergence}
    Suppose that $\ul{\nu}^*(A^*)=\{\ul{{\nu}}^0\}$ is a singleton and that $A^{(N)}\to A^*$. Then for every $\epsilon>0$, there is some $N_1$ so that
    \[
    \left\| \frac{\ul{n}}{t}-\ul{\nu}^0 \right\| <\epsilon
    \]
    whenever $\ul{n}\in\ul{n}^*_t(A^{(N)})$ and $N,t\geq N_1$. 
\end{lemma}

For the proofs, see \Cref{app:proofs}.


\subsection{Decay of coverage in the large budget regime}\label{sec:coverage-decay}

In this section, we discuss the phenomenon of decaying coverage as $B\to\infty$. Note that this is not unique to MultiPPI: it can be seen occuring to all baselines we compare to, and is especially pronounced for PPI++ vector. After discussing the phenomenon, we describe one way to avoid it.

Since, to the best of our knowledge, this phenomenon has not been observed in other works concerning PPI++, we focus our discussion on the PPI++ estimator and explain why it happens in that setting. Recall from \Cref{eqn:ppi++} the PPI++ estimator
\[
\hat{\theta}_{\text{PPI++}} = \frac{1}{n}\sum_{i=1}^n\left(Y_i-\widehat{\lambda}X_i\right) + \frac{1}{N}\sum_{j=1}^N \widehat{\lambda}\tilde{X}_j
\]
where $\{(X_i,Y_i)\}_{i\leq n}$ are i.i.d. according to some joint distribution $\p$, and $\{\tilde{X}_j\}_{j\leq N}$ are i.i.d. $\p_X$.

\citet{angelopoulos2023ppi++} (as well as many works before, in the context of control variates) propose a choice of $\widehat{\lambda}$ which depends on $\{(X_i,Y_i)\}_{i\leq n}$; namely, they let
\[
\widehat{\lambda} = \frac{N}{n+N}\frac{\widehat{\Cov}(X_{1:n},Y_{1:n})}{\widehat{\Var}(X_{1:n})}
\]
where $\widehat{\Cov}(X_{1:n},Y_{1:n})$ and $\widehat{\Var}(X_{1:n})$ are the relevant empirical covariance and variance computed from $\{(X_i,Y_i)\}_{i\leq n}$. This choice introduces bias in finite samples, and MultiPPI exhibits a similar behavior, as discussed in \Cref{sec:methods}. In the limit theorems provided in this work, c.f. \Cref{thm:asymptotic_revised}, and in \citet{angelopoulos2023ppi++}, it is assumed that the number of labeled samples (here, denoted $n$) tends to infinity. But this is not the situation presented in our experimental results.

Here we consider the bias of $\hat{\theta}_{\text{PPI++}}$ for fixed $n$ as $N\to\infty$. This bias is exactly
\[
\text{bias}(\hat{\theta}_{\text{PPI++}}):=\left|\Ex[\hat{\theta}_{\text{PPI++}}]-\Ex[Y]\right| = \left|\Ex[\widehat{\lambda}(X_1-\tilde{X}_1)]\right|=\frac{N}{n+N}\left|\Cov\left(X_1,\frac{\widehat{\Cov}(X_{1:n},Y_{1:n})}{\widehat{\Var}(X_{1:n})}\right)\right|
\]
by independence of $\widehat{\lambda}$ and $\tilde{X}_1$. Now for fixed $n$, and $N\to\infty$, the right-hand side converges upward precisely to the covariance of $X_1$ with the sample regression slope of $Y$ onto $X$, which is not in general zero. Therefore, the bias will increase but stay bounded as $N\to\infty$, as observed.

Note that this analysis does not apply to the setting in which the ratio $N/n$ is bounded. We find, accordingly, that this decay is unobserved in our experiments in which the number of labeled samples is in constant proportion with the budget.


\clearpage
\section{Proofs}
\label{app:proofs}
Unless explicitly stated otherwise, we prove results for the generalized setup outlined in \Cref{app:generalized}.

\subsection{Proof of \Cref{thm:minimax}}\label{app:minimax_proof}
For $\Sigma\in\R^{k\times k}$ symmetric positive-definite, let $\mathcal{P}_\Sigma$ denote the set of distributions on $\R^k$ with covariance $\Sigma$. For a fixed collection of index subsets $\mathcal{I}$ with associated costs $c_I$, let $\Theta_B$ denote the set of budget satisfying estimators $\hat{\theta}$, i.e. the estimators $\hat{\theta}$ which are measurable functions of $n_I$ independent copies of $X_I=(X_i)_{i\in I}$, for each $I\in\mathcal{I}$, such that $\mathbf{B}(\underline{n})$ holds. Note that here we allow for multiple budget inequalities, as described in \Cref{app:generalized}, and so $\mathbf{B}(\underline{n})$ denotes the proposition that all budget constraints are satisfied simultaneously. We emphasize that we make no explicit restriction to linear or unbiased estimators.

%\adam{Make an explicit note that this solves the more general case}

\begin{theorem}[Minimax optimality for general budget constraints]\label{thm:1}
    We have
    \[
    \inf_{\hat{\theta}\in\Theta_B}\sup_{P\in\mathcal{P}_\Sigma} \Ex\left[(\hat{\theta}-\theta^*)^2\right] = \Var 
    \left(\hat{\theta}_{\text{Multi-allocate$(\Sigma)$}}\right) = \mathcal{V}_B
    \]
    where the variance is with respect to any distribution $P\in\mathcal{P}_\Sigma$.
\end{theorem}
    
\begin{proof}
    We first reduce to the case of known and fixed $\underline{n}$.
    \begin{lemma}\label{lemma:1}
        Let $\Theta^{(\underline{n})}$ denote the set of measurable functions $\hat{\theta}$ which are functions of $n_I$ independent copies of $X_I$, for each $I\in\mathcal{I}$. Then if $\supp(a)\subseteq\bigcup\{I:n_I>0\}$,
        \[
        \inf_{\hat{\theta}\in\Theta^{(\underline{n})}}\sup_{P\in\mathcal{P}_\Sigma} \Ex\left[(\hat{\theta}-\theta^*)^2\right] = \min_{\substack{\underline{\lambda}
\,:\,\mathbf{U}(\underline{n},\underline{\lambda})}} \sum_{I:n_I>0} \frac{1}{n_I}\lambda_I^\top \Sigma_I\lambda_I;
        \]
        otherwise, $\sup_{P\in\mathcal{P}_\Sigma} \Ex\left[(\hat{\theta}-\theta^*)^2\right]$ is unbounded for all $\hat{\theta}\in\Theta_B$.
    \end{lemma}

    We now reduce the conjecture to this lemma. Observe that
    \[
    \Theta_B = \bigcup_{\underline{n}\,:\,\mathbf{B}(\underline{n})} \Theta^{(n)}
    \]
    and so the left hand-side of the conjecture is equal to
    \[
    \inf_{\underline{n}\,:\,\mathbf{B}(\underline{n})} \inf_{\hat{\theta}\in\Theta^{(\underline{n})}}\sup_{P\in\mathcal{P}_\Sigma} \Ex\left[(\hat{\theta}-\theta^*)^2\right] = \inf_{\underline{n}\,:\,\mathbf{B}(\underline{n})} \min_{\underline{\lambda}\,:\,\mathbf{U}(\underline{n},\underline{\lambda})} \sum_{I:n_I>0} \frac{1}{n_I}\lambda_I^\top \Sigma_I\lambda_I =: \Var(\hat{\theta}_{\text{Multi-allocate$(\Sigma)$}})
    \]
    since $\mathbf{U}(\underline{n},\underline{\lambda})$ is feasible for $\underline{\lambda}$ if and only if $\supp(a)\subseteq \bigcup\{I:n_I>0\}$.
    
    It now suffices to prove the lemma.
\end{proof}

\begin{proof}[Proof of Lemma~\ref{lemma:1}] The claim that  $\sup_{P\in\mathcal{P}_\Sigma} \Ex\left[(\hat{\theta}-\theta^*)^2\right]$ is unbounded for all $\hat{\theta}\in\Theta_B$ if $\supp(a)\not\subseteq\bigcup\{I:n_I>0\}$ follows from the observation that if $i\in \supp(a)\setminus \bigcup\{I:n_I>0\}$, there exist distributions $P\in\mathcal{P}_\Sigma$ such that $\theta^*_i=\Ex[X_i]$ may be made arbitrary large, while $\hat{\theta}$ cannot depend on such $X_i$.

    Therefore, in what follows, we assume $\supp(a)\subseteq \bigcup\{I:n_I>0\}$. The upper bound is clear from that fact that
    \[
    \{\hat{\theta}_{\underline{n},\underline{\lambda}} : \,\mathbf{U}(\underline{n},\underline{\lambda})\}\subseteq \Theta^{(\underline{n})}
    \]
    i.e., the set of unbiased linear estimators depending on $\underline{n}$ samples is a subset of the set of all estimators depending on $\underline{n}$ samples; and from the fact that $\Var (\hat{\theta}_{\underline{n},\underline{\lambda}}) = \sum_{I:n_I>0}\frac{1}{n_I}\lambda_I^\top\Sigma_I\lambda_I$ for every $P\in\mathcal{P}_\Sigma$, hence the minimal MSE of such estimators is precisely the right-hand side.

    We now prove the lower bound. Since the Bayes risk for any prior $\mu$ lower bounds the minimax risk, it suffices to construct a sequence of priors $\mu$ for which the risk of the Bayes estimator tends upward to our claimed lower bound. Let us choose the distribution $X\sim \mathcal{N}(\mu,\Sigma)$, and supply the prior $\mu\sim \mathcal{N}(0,\tau^2\text{Id}_k)$ for $\tau>0$ arbitrary; we will later take $\tau\to\infty$. Note that we then have $X_I=P_IX\sim\mathcal{N}(P_I\mu,P_I\Sigma P_I^\top)$.
    
    By construction, any estimator $\hat{\theta}\in\Theta^{(\underline{n})}$ depends on the independent set $
    \bigcup_{I\in\mathcal{I}}\{X_I^{(j)}\}_{1\leq j\leq n_I}$ where each $X_I^{(j)}$ is distributed according to $\mathcal{N}(\mu_I,\Sigma_I)$. The posterior\footnote{Morally, we are done at this point: the posterior mean is linear in $(\overline{X}_I)_I$, and the Multi-PPI estimator is the best such linear estimator. However, this does not yet directly imply the result. See next page for calculation of the posterior.} is then
    
    \begin{align*}
        \mu\,\,\bigg|\,\,\bigcup_{I\in\mathcal{I}}\{X_I^{(j)}\}_{1\leq j\leq n_I} \,&\sim\, \mathcal{N}(m_\tau,S_\tau) \\
        S_\tau &= \left(\frac{1}{\tau^2}\text{Id}_k + \sum_{I\in\mathcal{I}}n_I P_I^\top \Sigma_I^{-1}P_I \right)^{-1} \\
        m_\tau &= S_\tau\left(\sum_{I}n_I P_I^\top \Sigma_I^{-1}\overline{X}_I\right)
    \end{align*}
    The Bayes risk of estimating $\theta=a^\top \mu$ is then $a^\top S_\tau a$. Letting $\tau\to\infty$, we have shown that the minimax risk is at least\footnote{Here we use the assumption that $\text{supp}(a)\subseteq \bigcup\{I:n_I>0\}$, and thus $a$ lies in the range of $\sum_I n_IP_I^\top \Sigma_I^{-1}P_I$.}
    \[
    a^\top Sa,\quad S = \left(\sum_{I\in\mathcal{I}}n_I P_I^\top \Sigma_I^{-1}P_I \right)^\dagger.
    \]
    It remains to show that this risk is achievable by the $\hat{\theta}_{\underline{n},\underline{\lambda}}$ for some choice of $\underline{\lambda}$ satisfying $\mathbf{U}(\underline{n},\underline{\lambda})$. We quickly verify this below:

    Putting\footnote{To find this choice organically, one may solve an infimal norm convolution with Lagrange multipliers.}
    \[
    \lambda_I = \left(n_I \Sigma_I^{-1}P_I\right)S a
    \]
    we see that indeed $\mathbf{U}(\underline{n},\underline{\lambda})$ holds. Moreover, we calculate
    \[
    \Var(\hat{\theta}_{\underline{n},\underline{\lambda}})=\sum_{I} n_Ia^\top SP_I^\top \Sigma_I^{-1}\Sigma_I\Sigma_I^{-1}P_ISa = a^\top S \left(\sum_{I:n_I>0}n_I P_I^\top \Sigma_I^{-1}P_I\right)S a = a^\top Sa
    \]
    as desired. This concludes the proof.
\end{proof}

\subsection{Proof of \Cref{thm:asymptotic_revised}}\label{sec:asymptotic_proof}
We prove a generalization of \Cref{thm:asymptotic_revised} in which we allow for multiple budget inequalities.

Fix a vector $B_0\in \R^m_{>0}$. We consider the limit in which our budget is $B=t\cdot B_0$ and let $t\to\infty$. Suppose that $\widehat{\Sigma}\overset{p}{\to} \Sigma$ in the operator norm, potentially dependent on the variables sampled $X_I$.

We assume the following condition: Suppose that the following problem has a unique minimizer $\ul{\nu}$:
\begin{equation}\label{eqn:condition}
\begin{split}
    \ul{\nu}^\star = &\argmin_{\ul{\nu}} V(\ul{\nu}) :=a^\top\left(\sum_I \nu_I P_I^\top \Sigma_I^{-1} P_I\right)^\dagger a \\
    & \text{s.t.}\quad \ul{\nu}\geq0,\quad \sum_I \nu_I c_I\leq B_0,\quad \supp(a)\subseteq\bigcup \{I:\nu_I>0\}
\end{split}
\end{equation}


\begin{theorem}[Asymptotic normality in the general setting]\label{thm:asymptotic2} Suppose $X\in\R^k$ has finite second moment, and suppose that $\Sigma=\Cov(X)$ satisfies \Cref{eqn:condition}. Suppose that $\wh{\Sigma}\overset{p}{\to}\Sigma$ in the operator norm as $t\to\infty$. Let $\hat{\theta}_{\multiPPI(\widehat{\Sigma})}$ denote the MultiPPI estimator with budget $tB_0$. Then for $\hat{\theta}_{\multiPPI(\widehat{\Sigma})}$ arbitrarily dependent on any potential samples used to estimate $\widehat{\Sigma}$, we have
    \[
    \sqrt{t}\left(\hat{\theta}_{\multiPPI(\widehat{\Sigma})}-\theta^*\right)\overset{d}{\to}\mathcal{N}(0,\mathcal{V}^*)
    \]
    as $B\to\infty$, where $\mathcal{V}^*=\lim_{B\to\infty} B\mathcal{V}_{B}$, and $\mathcal{V}_{B}$ is defined in \Cref{eq:variance}.
\end{theorem}
While $\hat{\theta}_{\multiPPI(\Sigma)}$ is minimax optimal in the setting of fixed and known covariance $\Sigma$, it is in general not efficient, and the variance $\mathcal{V}$ can in general be improved by slowly concatenating onto $X$ nonlinear functions of its components. %It may be that such a version of $\hat{\theta}_{\multiPPI(\Sigma)}$, in which $k$ is increased slowly by adding appropriate nonlinear transformations of the components of $X$, is semiparametrically efficient if this is done at such a rate that $k\ll B^{1/2},N^{1/2}$.

\begin{proof}[Proof of \Cref{thm:asymptotic2}]
% Let $\hat{\theta} = \hat{\theta}_{\text{MultiPPI}(\widehat{\Sigma})}$. Note that \ref{eqn:condition} is simply a rounded version of the optimization problem which is solved by $\hat{\theta}_{\text{MultiPPI}(\widehat{\Sigma})}$. Let $\wh{\ul{\nu}}$ denote the solution to \ref{eqn:condition} with $\Sigma$ replaced by $\wh{\Sigma}$.

% We first show that, as a result of the assumed condition, we have $\hat{\nu}_I\to \nu_I^*$ whenever $\wh{\Sigma}\to\Sigma$; that rounded solutions are optimal in the limit $t\to\infty$ is justified by \autoref{sec:rounding}. Since $a$ lies in the range of $\sum_I \nu^*_I P_I^\top \Sigma_I^{-1} P_I$, the objective function is continuous in $(\nu,\Sigma)$ at $\nu^*$. 

% The allocation $\underline{\hat{n}}$ and weights $\underline{\hat{\lambda}}$ are chosen to minimize the variance under $\widehat{\Sigma}$ subject to the budget $B = tB_0$. Let $\hat{\nu}_I = \hat{n}_I/t$. As $t \to \infty$, the optimal proportions $\underline{\hat{\nu}}$ converge to the solution $\underline{\nu}^*$ of the continuous optimization problem \ref{eqn:condition}. The convergence $\underline{\hat{\nu}} \xrightarrow{p} \underline{\nu}^*$ follows from $\widehat{\Sigma} \xrightarrow{p} \Sigma$ and Berge's Maximum Theorem, as the objective function is continuous and the feasible set is compact. By the continuous mapping theorem, we similarly have $\hat{\lambda}_I\overset{p}{\to} \lambda^*_I$.

% EDIT: 

For fixed $t$, we let allocation $\underline{\hat{n}}$ and weights $\underline{\hat{\lambda}}$ be chosen to minimize the variance of  $\hat{\theta} = \hat{\theta}_{\text{MultiPPI}(\widehat{\Sigma})}$ under $\widehat{\Sigma}$ subject to the budget $B = tB_0$, as in the MultiPPI procedure. Since $\wh{\Sigma}\overset{p}{\to}\Sigma$, \Cref{lemma:discrete-convergence} ensures that $\ul{\hat{n}}/t\overset{p}{\to}\ul{\nu}^*$, as defined in \Cref{eqn:condition}.

We can write $\sqrt{t}(\hat{\theta} - \theta^*) = \sum_{I \in \mathcal{I}} \hat{\lambda}_I^\top \sqrt{\frac{t}{\hat{n}_I}} W_{I, \hat{n}_I}$, where $W_{I, \hat{n}_I} = \frac{1}{\sqrt{\hat{n}_I}} \sum_{j=1}^{\hat{n}_I} (X_I^{(I,j)} - \mu_I)$. For indices $I$ with $\nu_I^* > 0$, we have $\hat{n}_I \xrightarrow{p} \infty$. Define $n_I^* = \lfloor t\nu_I^*\rfloor $, and let 
\[
W^*_I:=\frac{1}{\sqrt{n_I^*}}\sum_{i=1}^{n_I^*} (X_I^{(I,j)}-\mu_I)
\]
It is now enough to show that $W_{I,\hat{n}_I}-W^*_I\overset{p}{\to}0$, and this will follow from Kolmogorov's inequality. To simplify notation, let us focus on a single subset $I$, and define $Y_j = X_I^{(I,j)}-\mu_I$. Let us also define $S_m=\sum_{j=1}^m Y_j$. We must show that
\[
\frac{S_{\hat{n}}}{\sqrt{\hat{n}}}-\frac{S_{n^*}}{\sqrt{n^*}}\overset{p}{\to}0
\]
where we have dropped dependence on $I$ for convenience. We decompose
\[
\frac{S_{\hat{n}}}{\sqrt{\hat{n}}}-\frac{S_{n^*}}{\sqrt{n^*}} = \underbrace{\frac{S_{\hat{n}}-S_{n^*}}{\sqrt{n^*}}}_A + \underbrace{\frac{S_{\hat{n}}}{\sqrt{\hat{n}}}\left(1-\sqrt{\hat{n}/n^*}\right)}_B
\]
Fix $0<\delta<1$. We work on the event $E_\delta(t)=\{|\hat{n}-n^*|\leq \delta t\}$, which holds with high probability.

We first control $A$. On $E_\delta(t)$, $\sqrt{n^*}|A|$ is a sum of at most $\delta t+1$ i.i.d. copies of $Y_j$. Kolmogorov's inequality then yields
\[
\p\left(A>\epsilon \right) \leq \frac{\delta t+1}{\epsilon^2 n^*}\leq 4\frac{\delta}{\epsilon^2}
\]
because $n^*=\lfloor t\nu^*\rfloor$. Taking $\delta\to0$ yields that $A\overset{p}{\to}0$.

We next control $B$. Working again on $E_\delta(t)$, we have
\[
\frac{S_{\hat{n}}}{\sqrt{\hat{n}}} \leq \frac{1}{\sqrt{1-\delta}}\left(\underbrace{\frac{S_{n^*}}{\sqrt{n^*}}}_{O_p(1)}  + \underbrace{\frac{S_{\hat{n}}-S_{n^*}}{\sqrt{n^*}}}_A\right) 
\]
Recognizing the second term as $A\overset{p}{\to}0$, and the first term as tight by the central limit theorem, we conclude that $S_{\hat{n}}/\sqrt{\hat{n}}$ is tight. Now we conclude that $B\overset{p}{\to}0$ because $\hat{n}/n^*\overset{p}{\to}1$.

Having proven $W_{I,\hat{n}_I}-W^*_I\overset{p}{\to}0$, we conclude that
\[
\sqrt{t}(\hat{\theta}-\theta^*) = \sum_{I:\nu_I^*>0} \frac{1}{n^*_I} \sum_{j=1}^{n^*_I} (\lambda^*_I)^\top (X_I^{(I,j)}-\mu_I) + o_p(1)
\]
But this is precisely the desired result, since this is the solution to the continuous optimization problem, and we are done.
\end{proof}

\subsection{Proof of \Cref{thm:stab}}
\begin{proof} 
The proof is immediate from \Cref{thm:finite-sample-v2-1}, proven in \Cref{sec:finite_sample_proofs}.
\end{proof}

\subsection{Proofs for \Cref{app:finite-sample-bounds}}\label{sec:finite_sample_proofs}

\subsubsection{Proof of \Cref{thm:finite-sample-v2-1}}
\begin{proof}
    We have
    \begin{align}\label{eqn:R}
    \begin{split}
        R(\wh{\ul{n}},\wh{\ul{\lambda}}) - R({\ul{n}}^*,{\ul{\lambda}}^*) &= R(\wh{\ul{n}},\wh{\ul{\lambda}}) - \wh{R}_N(\wh{\ul{n}},\wh{\ul{\lambda}}) \\
        &+ \underbrace{\wh{R}_N(\wh{\ul{n}},\wh{\ul{\lambda}}) - \wh{R}_N({\ul{n}}^*,{\ul{\lambda}}^*)}_{\leq0} \\
        &+ \wh{R}_N({\ul{n}}^*,{\ul{\lambda}}^*) - R({\ul{n}}^*,{\ul{\lambda}}^*)
    \end{split}
    \end{align}
and so it suffices to bound $|R(\wh{\ul{n}},\wh{\ul{\lambda}}) - \wh{R}_N(\wh{\ul{n}},\wh{\ul{\lambda}})|$ and $|\wh{R}_N({\ul{n}}^*,{\ul{\lambda}}^*) - R({\ul{n}}^*,{\ul{\lambda}}^*)|$. Define
\begin{align}\label{eqn:Delta}
\begin{split}
    \Delta_N(\ul{n},\ul{\lambda}) &= |R({\ul{n}},{\ul{\lambda}}) - \wh{R}_N({\ul{n}},{\ul{\lambda}})| \\
    &=\left|\sum_{I\in\mathcal{I}:n_I>0}\frac{1}{n_I}\lambda_I^\top(\Sigma-\wh{\Sigma}_N)\lambda_I\right|\\
    &\leq \|\Sigma-\wh{\Sigma}_N\|\sum_{I\in\mathcal{I}:n_I>0}\frac{1}{n_I}\|\lambda_I\|_2^2
\end{split}
\end{align}
Now since $\ul{n}^0,\ul{\lambda}^0$ satisfies $\mathbf{U}$ and $\mathbf{B}$, we have
\[
R({\ul{n}}^*,{\ul{\lambda}}^*)\leq R(\ul{n}^0,\ul{\lambda}^0),\quad\quad \wh{R}_N(\wh{\ul{n}},\wh{\ul{\lambda}})\leq \wh{R}_N(\ul{n}^0,\ul{\lambda}^0)
\]
from which it follows that
\begin{equation}\label{eqn:bound-norm}
    \sigma^2_1/n^0_{I^0} \geq \sum_{I\in\mathcal{I}:n_I^*>0}\frac{1}{n_I^*}(\lambda_I^*)^\top\Sigma(\lambda_I^*) \geq \gamma_{\min}(\Sigma)\sum_{I\in\mathcal{I}:n_I^*>0} \frac{1}{n_I^*} \|\lambda_I^*\|_2^2
\end{equation}
and similarly
\[
\wh{\sigma^2_1}/n^0_{I^0} \geq \sum_{I\in\mathcal{I}:\wh{n}_I>0}\frac{1}{\wh{n}_I}\wh{\lambda}_I^\top\wh{\Sigma}_N\wh{\lambda_I} \geq \gamma_{\min}(\wh{\Sigma}_N)\sum_{I\in\mathcal{I}:\wh{n}_I>0} \frac{1}{\wh{n}_I} \|\wh{\lambda}_I\|_2^2,
\]
where $\gamma_{\min}(A)$ denotes the minimum eigenvalue of the matrix $A$. We deduce that
\begin{align*}
    \sum_{I\in\mathcal{I}:n_I^*>0} \frac{1}{n_I^*} \|\lambda_I^*\|_2^2 &\leq \frac{\Sigma_{11}}{n^0_{I^0}\gamma_{\min}(\Sigma)} \\
    \sum_{I\in\mathcal{I}:\wh{n}_I>0} \frac{1}{\wh{n}_I} \|\wh{\lambda}_I\|_2^2 &\leq \frac{\wh{\Sigma}_{N,11}}{n^0_{I^0}(\gamma_{\min}(\Sigma)-\delta)} \leq \frac{{\Sigma}_{11}+\delta}{n^0_{I^0}(\gamma_{\min}(\Sigma)-\delta)}
\end{align*}
by Weyl's inequality, where we let $\delta=\|\Sigma-\wh{\Sigma}_N\|$. Coupled with \Cref{eqn:Delta}, we have
\begin{align*}
\Delta_N(\ul{n}^*,\ul{\lambda}^*) &\leq \delta \frac{\Sigma_{11}}{n^0_{I^0}\gamma_{\min}(\Sigma)} \\
\Delta_N(\wh{\ul{n}},\wh{\ul{\lambda}}) &\leq \delta \frac{{\Sigma}_{11}+\delta}{n^0_{I^0}(\gamma_{\min}(\Sigma)-\delta)}
\end{align*}
Taken together with \Cref{eqn:R} and the definition of $\Delta_N$, we conclude that
\[
R(\wh{\ul{n}},\wh{\ul{\lambda}}) \leq R({\ul{n}}^*,{\ul{\lambda}}^*) + 4\frac{\delta}{\gamma_{\min}(\Sigma)}\cdot \frac{\sigma^2_1}{n^0_{I^0}}
\]
for all $\delta \leq \gamma_{\min}(\Sigma)/2$.
\end{proof}

\subsubsection{Proof of \Cref{cor:bounded}}
\begin{proof}
    This follows immediately from \Cref{thm:finite-sample-v2-1} and Corollary 6.20 of \citet{wainwright2019high}.
\end{proof}

\subsubsection{Proof of \Cref{cor:subg}}
\begin{proof}
    This follows immediately from \Cref{thm:finite-sample-v2-1} and Theorem 4.7.1 of \citet{vershynin2018high}.
\end{proof}

\subsubsection{Proof of \Cref{cor:gersch}}
\begin{proof}
    This follows immediately from the Gershgorin circle theorem, as $\sum_{t\neq s} \Cov(X_t,X_s) \leq \sqrt{\Var(X_t)\Var(X_s)}<c$, and so $\lambda_{\min}(\Sigma)$ is bounded below for all $k$. On the other hand, $\lambda_{\max}(\Sigma)$ is bounded above on account of the same argument and the assumption that $X_i$ are bounded.
\end{proof}

\subsubsection{Proof of \Cref{thm:ledoit}}
\begin{proof}
The result follows immediate from \Cref{thm:finite-sample-v2-1} after the following lemma.
\end{proof}
\begin{lemma}
    Suppose that $\Sigma$ is not a multiple of the identity, and that $X\in\R^k$ is sub-Gaussian with proxy $K$. Let $\gamma_{\max}$ denote the maximum eigenvalue of $\Sigma$. Then the Ledoit-Wolf shrinkage estimator $\widehat{\Sigma}_N^\text{LW}$ satisfies the bound
    \[
    \Ex \|\wh{\Sigma}_N^\text{LW}-\Sigma\|_{op} \leq \frac{1}{\sqrt{N}}\sqrt{c_1K^4\gamma_{\max}^2 k^2 + c_2 K^8 \gamma_{\max} k^3/a^2}
    \]
    where $a^2:=\frac{1}{k}\left\|\Sigma - I\cdot \frac{\tr(\Sigma)}{k}\right\|_F^2$.
\end{lemma}
\begin{proof}
    Let $\wh{\Sigma}_N$ denote the empirical covariance matrix. Recall that by definition
    \[
    \wh{\Sigma}^{LW}_N = (1-\hat{\delta})\wh{\Sigma}_N +\hat{\delta} \hat{m}I
    \]
    where $\hat{m}=\tr(\wh{\Sigma}_N)/k$, and $\hat{\delta}=\hat{b}^2/\hat{d}^2$; we have $b^2= \Ex\|\wh{\Sigma}_N-\Sigma\|_F^2/k$ and $d^2=a^2+b^2$, and $\hat{b}$ and $\hat{d}$ are such that $\hat{b}\to b$ and $\hat{d}\to d$ in quartic mean. Our strategy will be to employ the observation that
    \begin{align*}
        \|\wh{\Sigma}_N^{LW}-\Sigma\|_F^2 &= \|(1-\hat{\delta})(\wh{\Sigma}_N-\Sigma)+\hat{\delta}(\Sigma-mI)\|_F^2 \\
        &\leq \left(|1-\hat{\delta}|\|\wh{\Sigma}_N-\Sigma\|_F + |\hat{\delta}|\|\Sigma-mI\|_F\right)^2 \\
        &\leq 6\|\wh{\Sigma}_N-\Sigma\|_F^2 + 4\hat{\delta}^2 \|\Sigma-mI\|_F^2
    \end{align*}
    using the coarse bounds that $|1-\hat{\delta}|\leq1,|\hat{\delta}|\leq1$ and $(u+v)^2\leq2u^2+2v^2$. It therefore suffices to bound $\Ex \|\wh{\Sigma}_N-\Sigma\|_F^2$ and $\Ex\hat{\delta}^2$.
    
    Since $X$ is sub-Gaussian with proxy $K$, $\wh{\Sigma}_N$ satisfies
    \[
    \Ex \|\wh{\Sigma}_N-\Sigma\|_F^2 \lesssim \frac{K^4}{N}\gamma_{\max}(k^2+k)
    \]
    by \citet{wainwright2019high}. This provides a bound on $b^2$; the estimator $\hat{b}$ is (after truncation) a average of $N$ i.i.d. quartic functionals of $X$ of the form $\|XX^\top-\wh{\Sigma}_N\|_F^2/k$, each of which have finite second-moment bounded by $cK^8\gamma_{\max}^4 k^2$ by the sub-Gaussian assumption.
    
    We conclude that we may bound
    \[
    \Ex \hat{b}^2 \lesssim \frac{K^4}{N}\gamma_{\max}k.
    \]
    We proceed by cases to bound $\Ex \hat{\delta}^2$. On the event $\{\hat{d}^2>a^2/2\}$, we have $\hat{\delta}\leq 2\hat{b}^2/a^2$, so it will suffices to bound the probability that $\{\hat{d}^2\leq a^2/2\}$. Since $\hat{d}^2$ is again an average of $N$ i.i.d. quartics in $X$, each of which have second moment bounded by $cK^8\gamma_{\max}^4k^2$, we have
    \[
    \Ex(\hat{d}^2-d^2)^2 \lesssim \frac{K^8}{N}\gamma_{\max}^4p^2
    \]
    We conclude that by Chebyshev's inequality, we have
    \[
    \p(\hat{d}^2\leq a^2/2) \leq c''\frac{K^8}{a^4N}\gamma_{\max}^4p^2
    \]
    Lastly, since $0\leq \hat{\delta}\leq 1$ (since $\hat{b}$ is truncated by $\hat{d}$), we conclude that in all cases
    \[
    \hat{\delta}^2\leq \hat{\delta} \leq \frac{2\hat{b}^2}{a^2} + \I_{\{\hat{d^2}\leq a^2/2\}}
    \]
    and so
    \[
    \Ex \hat{\delta}^2 \leq \frac{2}{a^2}\Ex \hat{b}^2 + \p(\hat{d}^2\leq a^2/2) \leq \frac{1}{N}\left(c'''K^4\gamma_{\max}^2\frac{k}{a^2} + c''''K^8\gamma_{\max}^4\frac{k^2}{a^4}\right)
    \]
    Taken together, we have shown that
    \[
    \Ex\|\wh{\Sigma}_N^{LW}-\Sigma\|_F^2 \leq \frac{1}{N}\left[c_1 K^4\gamma_{\max}^2k^2+c_2K^8\gamma_{\max}^4\frac{k^3}{a^2}\right]
    \]
    as desired.
\end{proof}


\subsection{Proofs for \Cref{sec:limiting}}

\subsubsection{Proof of \Cref{thm:limiting}}
\textbf{Note:} For the purpose of this proof only, we slightly change notation, letting $m$ denote the number of labeled samples rather than $n$. This just has the purpose of clarifying the potential conflict with the notation $n_I$.

\begin{proof}
Let us introduce the notation that $P_I$ is the orthogonal projection onto coordinates $I$, and thus $P_I^\top \lambda_I$ shares its values with $\lambda_I$ on coordinates $I$, and is 0 elsewhere. As a result, note that we have required
\[
\sum_{I:n_I>0} P_I^\top \lambda_I = \mu.
\]
Now we aim to minimize
\[
\frac{1}{m}\left(\sigma^2_Y - 2\mu^\top \Cov + \mu^\top \Sigma\mu\right) + \sum_{I:n_I>0} \frac{1}{n_I} \lambda_I^\top\Sigma_I\lambda_I
\]
or, expanding,
\[
V(n,\lambda):=\frac{1}{m}\left(\sigma^2_Y - 2\sum_{I:n_I>0} \lambda_I^\top \Cov_I + \sum_{I,J:n_I,n_J>0}\lambda_I^\top \Sigma_{IJ}\lambda_J\right) + \sum_{I:n_I>0} \frac{1}{n_I} \lambda_I^\top\Sigma_I\lambda_I
\]
We are interested in minimizing $V(n,\lambda)$ over all $\lambda$ (by which we mean $(\lambda_I)_{I\in\mathcal{I}}$) and $n$ satisfying the budget constraint $\sum_I c_In_I\leq C$. We will first minimize over $\lambda$ for fixed $n$: define $U(n):=\min_\lambda V(n,\lambda)$. But
\[
V(n,\lambda) = \lambda^\top \begin{pmatrix}
    \left(\frac{1}{m}+\frac{1}{n_{I_1}}\right)\Sigma_{I_1}\I_{n_I>0} & \dots & \frac{1}{m}\Sigma_{I_1I_k}\I_{n_{I_1},n_{I_k}>0} \\
    \vdots & \ddots & \vdots\\
    \frac{1}{m}\Sigma_{I_kI_1}\I_{n_{I_k},n_{I_1}>0} & \dots & \left(\frac{1}{m}+\frac{1}{n_{I_1}}\right)\Sigma_{I_1}\I_{n_I>0}
\end{pmatrix}\lambda  -2 \lambda^\top\begin{pmatrix}
    \frac{1}{m}\Cov_{I_1}\I_{n_{I_1}>0} \\
    \vdots \\
    \frac{1}{m}\Cov_{I_k}\I_{n_{I_k}>0}
\end{pmatrix} +\frac{\sigma^2_Y}{m}
\]
is a quadratic form in $\lambda$, where we define $\Sigma_{IJ}=(\Sigma_{ij})_{i\in I, j\in J} = P_I\Sigma P_J^\top$. This is of the form
\[
V(n,\lambda) = \lambda^\top \left(\frac{1}{m}S_1 + S_2\right)\lambda - 2\frac{1}{m}\lambda^\top T + d
\]
where
\[
S_1 = \begin{pmatrix}
    \Sigma_{IJ} \I_{n_I,n_J>0}
\end{pmatrix}_{I,J\in\mathcal{I}},\quad\quad S_2 = \texttt{block\_diag}\left(\frac{1}{n_I}\Sigma_I \I_{n_I>0}\right)_{I\in\mathcal{I}},\quad\quad T = \begin{pmatrix}
    \Cov_{I}\I_{n_I>0}
\end{pmatrix}_{I\in\mathcal{I}}
\]
and $d$ is constant in $n,\lambda$. It is known that the minimum value of such a quadratic form is
\[
U(n) = \min_\lambda V(n,\lambda) = -\frac{1}{m^2}T^\top \left(\frac{1}{m}S_1+S_2\right)^+T.
\]
This is because $T$ lies in the range of $\frac{1}{m}S_1+S_2$. To see this, let us introduce the notation that $\mathcal{I}^+=\{I\in\mathcal{I}:n_I>0\}$ and let $\mathcal{I}^0$ be its complement. Reorder $\mathcal{I}$ if necessary so that $\mathcal{I}^+$ strictly precedes $\mathcal{I}^0$. Then $\frac{1}{m} S_1+S_2$ takes the block form
\[
\frac{1}{m}\begin{pmatrix}
    (\Sigma_{IJ})_{I,J\in\mathcal{I}^+} & 0 \\
    0 & 0
\end{pmatrix} + \begin{pmatrix}
    \texttt{block\_diag}(\Sigma_I/n_I)_{I\in\mathcal{I}^+} & 0 \\
    0 & 0
\end{pmatrix}.
\]
Now, both $(\Sigma_{IJ})_{I,J\in\mathcal{I}^+}$ and $\texttt{block\_diag}(\Sigma_I/n_I)_{I\in\mathcal{I}^+}$ are symmetric positive-definite, hence invertible, on the coordinates $\mathcal{I}^+$, and $T$ has support in the span of the coordinates $\mathcal{I}^+$.

Given the block form shown above, we see that
\[
\left(\frac{1}{m}S_1 + S_2\right)^+ = \begin{pmatrix}
    \left(\frac{1}{m}(\Sigma_{IJ})_{I,J\in\mathcal{I}^+}+\texttt{block\_diag}(\Sigma_I/n_I)_{I\in\mathcal{I}^+}\right)^{-1} & 0 \\
    0 & 0
\end{pmatrix}
\]
again in the coordinates in which $\mathcal{I}^+$ precedes $\mathcal{I}^0$.

Continuity of the inverse is now enough to conclude that
\[
\lim_{m\to0} m^2U(n) = -T^\top\texttt{block\_diag}\left(n_I\Sigma_I^{-1}\right)_{I\in\mathcal{I}}  T = -\sum_{I\in\mathcal{I}} n_I \Cov_I^\top\Sigma_I^{-1}\Cov_I =: L(n)
\]
But now this is a linear function $L(n)$ in $n$. Consider minimizing this in $n$, subject to the (simplex) budget constraint $n_I\geq0$, $\sum_I c_In_I\leq C$. The minimum is achieved on a vertex of the simplex, and the minimizer is unique except in the unlikely situation that
\[
\frac{\Cov_I^\top\Sigma_I^{-1}\Cov_I}{c_I} = \text{constant in $I$}
\]
assuming that $\Cov_I\neq0$ for some $I$.

Now we claim that $m^2 U(n)\rightarrow L(n)$ uniformly in $n$ subject to the budget constraint. For this, it suffices to show that
\[
\left(\frac{1}{m}(\Sigma_{IJ})_{I,J\in\mathcal{I}^+}+\texttt{block\_diag}(\Sigma_I/n_I)_{I\in\mathcal{I}^+}\right)^{-1} \to \texttt{block\_diag}(\Sigma_I/n_I)_{I\in\mathcal{I}^+}^{-1}
\]
in the operator norm, uniformly in $n$. The Woodbury matrix identity implies that the difference is exactly
\[
\texttt{block\_diag}(n_I\Sigma_I^{-1})_{I\in\mathcal{I}^+} (I+m\texttt{block\_diag}(\Sigma_I^{-1}/n_I)_{\mathcal{I}^+}(\Sigma_{IJ})_{I,J\in\mathcal{I}^+})^{-1}
\]
Now, we have $0<n_I\leq C/c_I$ for all $I\in\mathcal{I}^+$ by the constraint. The operator norm is sub-multiplicative, and the first factor is bounded in norm by a constant multiple of $1/\min_I c_I$. Similarly, we have
\[
I+m\texttt{block\_diag}(\Sigma_I^{-1}/n_I)_{\mathcal{I}^+}(\Sigma_{IJ})_{I,J\in\mathcal{I}^+} \succ I+\frac{mC}{\min_I c_I}\texttt{block\_diag}(\Sigma_I^{-1})_{\mathcal{I}^+}(\Sigma_{IJ})_{I,J\in\mathcal{I}^+}
\]
The operator norm of the right-hand side goes to $\infty$ uniformly in $n$, so the operator norm of its inverse goes to $0$ uniformly as well. In conclusion, we have uniform convergence. Therefore, we have
\[
n^*(m):=\argmin_n \min_\lambda V(n,\lambda)\xrightarrow[m\to\infty]{} n^*
\]
\end{proof}

\subsection{Proofs for \Cref{sec:continuous_problem}}


% \begin{lemma}\label{lemma:cont}
%     Define the set-valued function
    
%     for $A$. Let $A^*\in\mathbf{S}^k_{++}$ be such that the above minimization problem has a unique solution; i.e., $\nu$ 
%     be unique as a minimizer. Then the map $A\mapsto \ul{\nu}^*(A)$ is continuous in the operator norm for $A\in \mathbf{S}^k_{++}$.
% \end{lemma}
% \begin{proof}[Proof of \Cref{lemma:cont}]
% We may alternatively write
% \[
% \ul{\nu}^*(A) = \argmin_{\ul{\nu}\in K} F(\ul{\nu},A)
% \]
% where $K=\{\ul{\nu}:\ul{\nu}\geq0,\sum_I\nu_Ic_I\leq B_0\}\subseteq \R^{|\mathcal{I}|}$ is compact, and
% \[
% F(\ul{\nu},A)=\begin{cases}
%     a^\top \left(\sum_I \nu_I A_I^\dagger \right)^\dagger a & a \in \text{range}\left(\sum_I\nu_I A_I^\dagger\right) \\
%     \infty & \text{otherwise.}
% \end{cases}
% \]
% This is because, since $A\succ0$, the range of $A_I^\dagger$ is precisely $\R^I$, so for each $i\in[k]$, the condition $\exists I\ni i:\nu_I>0$ is equivalent to the condition that $e_i\in \text{range}\left(\sum_I\nu_I A_I^\dagger\right)$. This will be formalized in \Cref{lemma:l1} below. We are essentially now ready to apply Berge's maximum theorem, after compactifying the domain. Since the usual statement of Berge's theorem applies to objective functions taking values in $\R$ rather than $\overline{\R}=\R\cup\{\infty\}$, let $\tilde{F}(\ul{\nu},A)$ represent the composition of a homeomorphism $\R\longrightarrow \overline{\R}$ with $F$, e.g. $\tilde{F}(\ul{\nu},A)=1-\exp\{- F(\ul{\nu},A)\}$. We may then alternatively write
% \[
% \ul{\nu}^*(A)=\argmin_{\ul{\nu}\in K} \tilde{F}(\ul{\nu},A)
% \]
% where, importantly, $\tilde{F}(\ul{\nu},A)$ is now real-valued.

% It remains to show that $\tilde{F}$ is jointly continuous in $(\ul{\nu},A)$; Berge's theorem will then imply the desired result. This is equivalent to showing that $F$ is continuous when $F<\infty$, and lower semi-continuous (l.s.c.) otherwise.

% To show that $F$ is continuous when $F<\infty$, we need the following technical lemma related to the pseudoinverse:

% \begin{lemma}\label{lemma:PI1}
%     Let $M\succeq0$ by symmetric and take the block form
%     \[
%     M = \begin{pmatrix}
%         M_{JJ}  & 0 \\
%         0 & 0
%     \end{pmatrix}
%     \]
%     for some index set $J\subseteq[k]$, where $M_{JJ}\succ0$. Let $a\in\R^k$ have support in $J$. Then for any matrix $N$, we have
%     \[
%     a^\top(M+\epsilon N)^\dagger a = a^\top M_{JJ}^{-1} a+ O(\epsilon).
%     \]
% \end{lemma}

% \begin{lemma}\label{lemma:PI2}
%     Fix a symmetric $M\succeq0$ such that, for an index set $I$, $M$ is positive definite on $I$ and zero on $I^c$, and let $a\in\range(M)$. Let $N_1,\dots,N_\ell$ be symmetric matrices such that, for index sets $I_1,\dots,I_\ell$, $N_i$ is positive definite on $I_i$ and zero on $I_i^c$. Then we have
%     \[
%     a^\top \left(M+\sum_i N_i\right)^\dagger a = a^\top M^\dagger a + O(\max_i\|N_i\|)
%     \]
% \end{lemma}
% \begin{proof}[Proof of \Cref{lemma:PI2}]
% Let $I^*=I\cup I_1\cup\dots\cup I_\ell$. We first show that we may assume WLOG that $I^*=[k]$ and that $M+\sum_i N_i$ is invertible.

% Since $\supp(a)\subseteq I\subseteq I^*$, by reordering coordinates if necessary we may write
% \[
% M+\sum_i N_i = \begin{pmatrix}
%     M_{I^*,I^*}+\sum_i (N_i)_{I^*,I^*} & 0 \\
%     0 & 0
% \end{pmatrix},\quad\quad\quad a = \begin{pmatrix}
%     a_{I^*} \\ 0
% \end{pmatrix}
% \]
% Now we claim that $M_{I^*,I^*}+\sum_i (N_i)_{I^*,I^*}$ is positive definite. To see this, note that we have
% \[
% x^\top (M_{I^*,I^*}+\sum_i (N_i)_{I^*,I^*})x = x^\top M_{{I^*,I^*}} x + \sum_i x^\top (N_i)_{I^*,I^*} x = x^\top_I M_{I,I} x_I + \sum_i x_{I_i}^\top (N_i)_{I_i,I_i} x_{I_i}
% \]
% Now, if the right-hand side vanishes, we must have $x_I=0,x_{I_i}=0$ for all $i$. But then $x_{I^*}=0$, and so we are done. With this said, we have
% \[
% a^\top (M+\sum_i N_i)^\dagger a = a_{I^*}^\top (M_{I^*,I^*}+\sum_i (N_i)_{I^*,I^*})^{-1}a_{I^*}
% \]
% From this point onward, we therefore assume that $I^*=[k]$ and that $M+\sum_i N_i$ is invertible. Let us abbreviate $S=\sum_iN_i$. We then have the block form
% \[
% M+S = \begin{pmatrix}
%     M_{I,I}+S_{I,I} & S_{I,I^c} \\
%     S_{I^c,I} & S_{I^c,I^c}
% \end{pmatrix}
% \]
% We have seen that $M+S$ is invertible, and moreover $S_{I^c,I^c}$ is invertible because $I^c\subseteq  I_1\cup\dots,I_\ell$, therefore if $\supp(x)\subseteq I^c$, then the fact that $N_i$ is positive definite on $I_i$ is enough to conclude.

% Therefore we may apply the block inverse formula based on the Schur complement: the complement of $M+S$ is
% \[
% M_{I,I}+S_{I,I}-S_{I,I^c}S_{I^c,I^c}^{-1}S_{I,I^c}
% \]
% acting on $\R^I$. Now since $\supp(a)\subseteq I$, we again see from the block form that
% \[
% a^\top(M+S)^{-1}a = a_I^\top (M_{I,I}+S_{I,I}-S_{I,I^c}S_{I^c,I^c}^{-1}S_{I,I^c})^{-1}a_I
% \]
% Now $M_{I,I}$ is invertible by assumption, 
% \end{proof}
% \begin{lemma}
%     The function $F$
%     Let $A^*\succ0$ be a symmetric matrix, and set $A = A^*+\Delta$. Let $\ul{\nu}^*\in K$, and set $\ul{\nu}=\ul{\nu}^*+\ul{\delta}$. Suppose that $\|\Delta\|,\|\ul{\delta}\|_\infty<\epsilon$. Then if $a\in\range\left(\sum_I\nu_I^* (A^*_I)^\dagger\right)$, as $\epsilon\to0$ we have
%     \[
%     F(\ul{\nu},A)\to F(\ul{\nu}^*,A^*)<\infty
%     \]
%     and if $a\not\in\range\left(\sum_I\nu_I^* (A^*_I)^\dagger\right)$, as $\epsilon\to0$ we have
%     \[
%     F(\ul{\nu},A)\to\infty
%     \]
% \end{lemma}

In this section, we prove \Cref{lemma:discrete-convergence}. The following suffices for most of the proof:
\begin{lemma}\label{lemma:Fcont}
    The map $F:K\times \mathbf{S}^k_{++}\longrightarrow\R\cup\{\infty\}$ is continuous.
\end{lemma}
From this, we have the following:
\begin{lemma}\label{lemma:nucont}
    The correspondence $\ul{\nu}^*$ is upper hemi-continuous. Consequently, if $\ul{\nu}^*(A^*)=\{\ul{\nu}^0\}$ is a singleton and $A^{(N)}\to A^*$, then for every $\epsilon>0$ there is some $N_0$ so that $\|\ul{\nu}-\ul{\nu}^0\|<\epsilon$ whenever $\ul{\nu}\in\ul{\nu}^*(A^{(N)})$ and $N\geq N_0$.
\end{lemma}

\begin{proof}[Proof of \Cref{lemma:Fcont}]
We need three warm-up lemmas.

\begin{lemma}\label{lemma:l1}
    Suppose that, for $I\in\mathcal{J}$, $M_I\in\R^{k\times k}$ is a symmetric PSD matrix which is SPD on the subspace $\R^I$, and zero off of the coordinates. Then $\range(\sum_{I\in\mathcal{I}} M_I)=\R^{\bigcup \mathcal{J}}$.
\end{lemma}
\begin{proof}
    It is clear that $\sum_I M_I$ is SPD on $\bigcup \mathcal{I}$; since it is symmetric, it is invertible on $\bigcup\mathcal{I}$.
\end{proof}

\begin{lemma}\label{lemma:l2}
        For all $\epsilon$ small enough, $\range(\sum_I \nu_I^*(A_I^*)^\dagger)\subseteq \range(\sum_I \nu_IA_I^\dagger)$.
\end{lemma}
\begin{proof}
    Let us take $\epsilon$ small enough that $\nu_I\geq \nu_I^*/2>0$ whenever $\nu_I^*>0$. Let us also take $\epsilon$ small enough that $A\succ0$ by Weyl's inequality. Then each $A_I^\dagger$ and $(A_I^*)^\dagger$ is symmetric PSD and SPD on the subspace $\R^I$, and zero off the coordinates.
    
    It now follows from \Cref{lemma:l1} that $\range(\sum_I \nu_I^*(A^*_I)^\dagger) = \text{span}(e_i:i\in\bigcup_{I:\nu_I^*>0}I)\subseteq \text{span}(e_i:i\in\bigcup_{I:\nu_I>0} I)=\range(\sum_I \nu_I A_I^\dagger)$.
\end{proof}

It follows from the assumption and this lemma that $a\in\range(\sum_I \nu_I A_I^\dagger)$.

%Since $a\in\range(\sum_I \nu_I^*(A_I^*)^\dagger)$, we have $\supp(a)\subseteq I^*_+ := \bigcup_{I:\nu_I^*>0} I$.

We have two final lemmas:

\begin{lemma}\label{lemma:l3}
    Suppose that $M$ is symmetric PSD and that $\range(M)=\text{span}(J)$. Then we have
    \[
    x^\top M^\dagger x = x_{(J)}^\top M_{(J)}^{-1} x_{(J)}
    \]
    where $x_{(J)}$ denotes the restriction of $x$ to the coordinates $J$, and $M_{(J)}$ denotes the restriction of $M$ to the coordinates $J\times J$.
\end{lemma}
\begin{proof}
    Since $My\in \text{span}(J)$ for all $y\in\R^k$, we have $(My)_i=0$ whenever $i\in J^c$. Thus every row of $M$ indexed by an element of $J^c$ is zero. By symmetry, the same is true for every column. Thus $M$ takes the block form
    \[
    M = \begin{pmatrix}
        M_{(J)} &0 \\
        0 & 0
    \end{pmatrix}
    \]
    in the ordered coordinates such that $J$ precedes $J^c$. Now the fact that $M$ is symmetric means that it is invertible on its range, and so $M_{(J)}$ is invertible, and we are done. 
\end{proof}

Our last lemma is this:

\begin{lemma}\label{lemma:l4}
    Assume that $M,\Xi \in \R^{\ell\times\ell}$ are symmetric PSD, and suppose that $\range(M)=\R^J$. Assume that $\|\Xi\|\leq \lambda_{\min}(M_{(J)})/2$ and $M+\Xi$ is invertible. Then if $x\in \range(M)$, we have
    \[
    \left|x^\top (M+\Xi)^{-1}x -  x_{(J)}^\top M_{(J)}^{-1}x_{(J)}\right| \leq 2 \|M_{(J)}^{-1}\|^2 \|\Xi\| \|x\|^2.
    \]
    On the other hand, if $x\not\in \range(M)$, we have
    \[
    \left|x^\top (M+\Xi)^{-1}x\right| \geq \frac{1}{2\|\Xi\|} \|x_{(J^c)}\|^2.
    \]
\end{lemma}
% \begin{lemma}\label{lemma:l4}
%     Assume that $M,\Xi \in \R^{\ell\times\ell}$ are symmetric PSD, and let $\widetilde{M}$ denote the restriction of $M$ to its range. Similarly, for a vector $x\in\R^{\ell}$, let $\widetilde{x}$ denote the projection of $x$ onto $\range(M)$, so that $x-\widetilde{x}$ is the projection of $x$ onto the kernel of $M$. Suppose that $\|\Xi\|\leq \|M_{(J)}^{-1}\|/2$ and $M+\Xi$ is invertible. Then if $x\in \range(M)$, we have
%     \[
%     \left|x^\top (M+\Xi)^{-1}x -  \widetilde{x}^\top \widetilde{M}^{-1}\widetilde{x}\right| \leq 2 \|\widetilde{M}^{-1}\|^2 \|\Xi\| \|x\|^2.
%     \]
%     On the other hand, if $x\not\in \range(M)$, we have
%     \[
%     \left|x^\top (M+\Xi)^{-1}x\right| \geq \frac{1}{2\|\Xi\|} \|x-\widetilde{x}\|^2.
%     \]
% \end{lemma}
% \begin{proof}
% Since $M$ is symmetric, we may choose a basis in which 
    
% \end{proof}

\begin{proof}
    In block form, again in ordered coordinates such that $J$ precedes $J^c$, we have
    \[
    x^\top (M+\Xi)^{-1}x = \begin{pmatrix}
        x_{(J)} \\
        0
    \end{pmatrix}^\top\begin{pmatrix}
        M_{(J,J)}+\Xi_{(J,J)} & \Xi_{(J,J^c)} \\
        \Xi_{J^c,J} & \Xi_{J^c,J^c}
    \end{pmatrix}^{-1} \begin{pmatrix}
        x_{(J)} \\ 0
    \end{pmatrix} = x_{(J)}^\top \left[M_{(J)} + S\right]^{-1}x_{(J)}
    \]
    where $S=\Xi / \Xi_{(J^c,J^c)}$ denotes the Schur complement. The norm of the Schur complement is bounded $\|S\|\leq \|\Xi\|$, as a result of the forthcoming lemma. By submultiplicativity, we have
    \[
    \left\|(M_{(J)}+S)^{-1}-M_{(J)}^{-1}\right\| \leq \left\| M_{(J)}^{-1}\right\| \|S\|\left\|(M_{(J)}+S)^{-1}\right\|
    \]
    Now the minimum eigenvalue of $M_{(J)}+S$ is at least $\lambda_{\min}(M_{(J)}) - \|S\| \geq \lambda_{\min}(M_{(J)})-\|\Xi\|\geq \lambda_{\min}(M_{(J)})/2=\|M_{(J)}^{-1}\|^{-1}/2$ by Weyl's inequality, and so we conclude that
    \[
    \left\|(M_{(J)}+S)^{-1}-M_{(J)}^{-1}\right\|\leq 2 \|M_{(J)}^{-1}\|^2 \|\Xi\|
    \]
    as desired, and we are done.

    On the other hand, by standard properties of the Schur complement, we have
    \[
    x^\top(M+\Xi)^{-1}x = \begin{pmatrix}
        x_{(J)} \\
        x_{(J^c)}
    \end{pmatrix}^\top\begin{pmatrix}
        M_{(J,J)}+\Xi_{(J,J)} & \Xi_{(J,J^c)} \\
        \Xi_{J^c,J} & \Xi_{J^c,J^c}
    \end{pmatrix}^{-1} \begin{pmatrix}
        x_{(J)} \\ x_{(J^c)}
    \end{pmatrix} \geq x_{(J^c)}^\top (S')^{-1}x_{(J^c)}
    \]
    where
    \[
    S'=(M+\Xi)/(M+\Xi)_{(J^c,J^c)} = \Xi_{(J^c,J^c)} - \Xi_{(J,J^c)}(M_{(J,J)}+\Xi_{(J,J)})^{-1}\Xi_{(J^c,J)} 
    \]
    denotes the Schur complement. We proceed to bound $\|S'\|$:
    \[
    \|S'\|\leq \|\Xi\| + \|\Xi\|^2\|(M_{(J,J)}+\Xi_{(J,J)})^{-1}\|\leq \|\Xi\|\left(1+\|\Xi\|\cdot 2\|M_{(J))}^{-1}\|\right)\leq 2\|\Xi\|
    \]
    since we assume $\|\Xi\|\leq \lambda_{\min}(M_{(J)})/2$. We conclude that $\lambda_{\min}((S')^{-1})\geq 1/(2\|\Xi\|)$, establishing the desired result.
\end{proof}

Below, we summarize the properties of the Schur complement used above:

\begin{lemma}[Properties of Schur complement]\label{lemma:schur}
Let $N$ be a symmetric SPD matrix, and $S$ denote one of its Schur complements. Then we have
\[
x^\top S x \leq \begin{pmatrix}
    x \\
    y
\end{pmatrix}^\top N\begin{pmatrix}
    x \\ y
\end{pmatrix}
\]
and, if $N$ is invertible,
\[
x^\top S^{-1} x \leq \begin{pmatrix}
    x \\
    y
\end{pmatrix}^\top N^{-1}\begin{pmatrix}
    x \\ y
\end{pmatrix}
\]
for all $x,y$, whence $\|S\|\leq \|N\|$.
\end{lemma}
\begin{proof}
    The proof is immediate from equation 7.7.5. of \citet{horn2012matrix}, using the fact that the invertibility of $N$ implies that of $S$ by Schur's formula.
\end{proof}

Now we conclude the proof. Let $I^+_*$ denote $\bigcup_{I:\nu_I^*>0} I$ and $I^+$ denote $\bigcup_{I:\nu_I>0}I$. By Lemma~\ref{lemma:l1} and Lemma~\ref{lemma:l2} we have $\range(\sum_I\nu_I^*(A^*_I)^\dagger)=\R^{I^+_*} \subseteq \R^{I^+}=\range(\sum_I\nu_I(A_I)^\dagger)$. By assumption, we have $a\in \R^{I^+_*}$.

Applying Lemma~\ref{lemma:l3} to $F$, we have
\[
F(\ul{\nu},A) = a^\top\left(\sum_I\nu_I A_I^\dagger\right)^\dagger a = a_{(I^+)}^\top\left(\left(\sum_I\nu_I A_I^\dagger\right)_{(I^+)}\right)^{-1} a_{(I^+)}
\]
Now we write
\[
\left(\sum_I\nu_I A_I^\dagger\right) = \left(\sum_I\nu_I^* (A_I^*)^\dagger\right) + \Xi
\]
whence
\[
\left(\sum_I\nu_I A_I^\dagger\right)_{(I^+)} = \left(\sum_I\nu_I^* (A_I^*)^\dagger\right)_{(I^+)} + \Xi_{(I^+)}
\]
with $\|\Xi_{(I^+)}\|\leq\|\Xi\|$. Now, the first term is symmetric PSD with range $\R^{I^+_*}$, since $I^+_*\subseteq I^+$. Thus by Lemma~\ref{lemma:l4} we conclude
\[
|F(\ul{\nu},A)-F(\ul{\nu}^*,A^*)|\leq 2\left\|\left(\sum_I \nu^*_I(A_I^*)^\dagger\right)_{(I^+_*)}^{-1}\right\|^2\|\Xi\|\|a\|^2
\]
and we are done.
\end{proof}

\subsubsection{Proof of \Cref{lemma:nucont} }
\begin{proof}
    The result follows immediately from Berge's theorem after an appropriate compactification of the domain. Let $\tilde{F}(\ul{\nu},A)$ represent the composition of a homeomorphism $\R_{\geq0}\longrightarrow \R_{\geq0}\cup\{\infty\}$ with $F$, e.g. $\tilde{F}(\ul{\nu},A)=1-\exp\{- F(\ul{\nu},A)\}$. We may then alternatively write
\[
\ul{\nu}^*(A)=\argmin_{\ul{\nu}\in K} \tilde{F}(\ul{\nu},A)
\]
where, importantly, $\tilde{F}(\ul{\nu},A)$ is now real-valued.

It remains to show that $\tilde{F}$ is jointly continuous in $(\ul{\nu},A)$; Berge's theorem will then imply the desired result. This is equivalent to showing that $F$ is continuous, and so \Cref{lemma:Fcont} suffices.
\end{proof}

Finally, we prove the required lemma regarding convergence of the discrete problem to the continuous one.
\begin{proof}[Proof of Lemma~\ref{lemma:discrete-convergence}]
    Suppose to the contrary that there is some sequence $N_{\alpha},t_{\alpha}$ of arbitrarily large $N,t$ so that
    \[
    \left\| \frac{\ul{n}}{t_\alpha}-\ul{\nu}^0 \right\| \geq \epsilon
    \]
    for some $\ul{n}\in \ul{n}^*_{t_\alpha}(A^{(N_\alpha)})$, for each $\alpha$. By the preceding lemma, we know that there is some $N_0$ so that we may assume $\|\ul{\nu}-\ul{\nu}^0\|<\epsilon/2$ whenever $N\geq N_0$, for all $\ul{\nu}\in\ul{\nu}^*(A^{(N)})$.

    For large enough $\alpha$, therefore, we have $\|\ul{\nu}-\ul{\nu}^0\|<\epsilon/2$ for all $\ul{\nu}\in\ul{\nu}^*(A^{(N_\alpha)})$. By the reverse triangle inequality, then, we must have
    \[
    \left\| \frac{\ul{n}}{t_\alpha}-\ul{\nu} \right\| \geq \epsilon/2
    \]
    for all large $\alpha$, for each $\ul{n}\in \ul{n}_{t_\alpha}^*(A^{(N_\alpha)})$ and $\ul{\nu}\in \ul{\nu}^*(A^{(N_\alpha)})$.

    Since every such $\ul{n}/t_\alpha$ is contained in the compact set $K$, we may find a limit point $\ul{\nu}'\in K$ so that for all $\delta>0$, there is some $\alpha$ so that $\|\ul{n}/t_\alpha-\ul{\nu}'\|<\delta$ for some $\ul{n}\in\ul{n}^*_{t_\alpha}(A^{(N_\alpha)})$. For this $x'$ we must then also have
    \[
    \|\ul{\nu}'-\ul{\nu}\|\geq\epsilon/2
    \]
    for each $\ul{\nu}\in\ul{\nu}^*(A^{(N_\alpha)})$. 

    We now argue that in fact $\ul{\nu}'\in \ul{\nu}^*(A^*)$, which will serve to contradict the assumption that $\ul{\nu}^*(A^*)$ is a singleton. Recall that every $\ul{n} \in \ul{n}^*_{t_\alpha}(A^{(N_\alpha)})$ is optimal, and therefore minimizes $F(\ul{n},A^{(N_\alpha)})$ over the set $K_t\cap \mathbb{Z}^{|\mathcal{I}|}$, where
    \[
    K_t:=\left\{\ul{\nu}:\ul{\nu}\geq0,\sum_I\nu_Ic_I\leq tB_0\right\}
    \]
    Since we have $F(\kappa \ul{\nu},A)=\frac{1}{\kappa}F(\ul{\nu},A)$ for all $\kappa>0$, it is not hard to see that if $\ul{\nu}$ is a minimizer of $F(\cdot,A)$ over $K_t$, then $\ul{\nu}/t$ is a minimizer of $F(\cdot,A)$ over $K$. Let us define
    \[
    \ul{n}^*_{t,\text{round}}(A)=\text{round}(t\ul{\nu}^*(A))
    \]
    in the set-valued sense. It follows that $\ul{n}^*_{t,\text{round}}(A)\subseteq  K_t\cap \mathbb{Z}^{|\mathcal{I}|}$ for all $A$, and so
    \[
    F(\ul{n},A^{(N_\alpha)})\leq F(\ul{n}_{\text{round}},A^{(N_\alpha)})
    \]
    for all $\ul{n}\in\ul{n}^*_{t_\alpha}(A^{(N_\alpha)})$ and $\ul{n}_{\text{round}}\in \ul{n}^*_{t,\text{round}}(A^{(N_\alpha)})$. By definition, it follows that
    \[
    \ul{n}^*_{t_\alpha,\text{round}}(A^{(N_\alpha)})/t_\alpha\to \ul{\nu^0}
    \]
    in the set valued sense, and so by continuity of $F$ it follows that
    \[
    t_\alpha F(\ul{n}_{\text{round}},A^{(N_\alpha)}) = F(\ul{n}_{\text{round}}/t_\alpha,A^{(N_\alpha)}) \to F(\ul{\nu}^0,A^*).
    \]
    We conclude that
    \[
    t_\alpha F(\ul{n},A^{(N_\alpha)}) = F(\ul{n}/t_\alpha ,A^{(N_\alpha)}) \leq F(\ul{\nu}^0,A^*)
    \]
    and so taking the limit along the sequence of $\alpha$ previously specified we conclude that
    \[
    F(\ul{\nu}',A^*)\leq F(\ul{\nu}^0,A^*)
    \]
    by continuity of $F$ again. Thus $\ul{\nu}'\in \ul{\nu}^*(A^*)$, yet $\ul{\nu}'$ is separated from $\ul{\nu}^0$ by at least $\epsilon/2$, contradicting the assumption that $\ul{\nu}^*(A^*)$ is a singleton. This concludes the proof.
    
    
    % By the preceding section, we have that
    % \[
    % \lim_{t\to\infty} tV_t(\ul{n}^*_{\text{int}})=V_1(\ul{x}^*)
    % \]
    % For $t$ large enough, therefore, and any fixed $N$, we have $tF(\ul{n},A^{(N)}) \geq F(\ul{\nu}^*(A^{(N)}),A^{(N)})-\epsilon/2$ for each $\ul{n} \in \ul{n}^*_{t}(A^{(N)})$.

    % By Berge's theorem, the value function is continuous, and so 

    % From scratch: 
    
    % Now each $\ul{n}$ considered is optimal, so in particular, by the analysis of the preceding section, 

    
    
    % We now compare $\ul{\nu}^*(A^{(N)})$ to $\ul{n}^*_t(A^{(N)})$. Suppose to the contrary that for arbitrarily large $t$ we may choose $\ul{n}^{(t)}\in \ul{n}^*_t(A^{(N)})$
    % it suffices to show that there is some $N_1$ so that 
\end{proof}
\section{Computational considerations}\label{sec:computation}
% \begin{center}
% \fbox{%
%   \begin{minipage}{1.0\textwidth}\label{alg_multiallocate}
%     \paragraph{Subroutine 1:} 
%     Multi-Allocate$(\Sigma)$
%     \paragraph{Input:} a symmetric positive-define matrix $\Sigma\in\R^{k\times k}$. Implicitly given $a$, $\mathcal{I}$, $(c_I)_I$, $B$.
%     \paragraph{Output:} the solution $(n_I,\lambda_I)_{I\in\mathcal{I}}$ to the problem
%     \begin{align*}
%         \min_{(n_I,\lambda_I)_{I\in\mathcal{I}}} &\sum_{I\in\mathcal{I}:n_I>0} \frac{1}{n_I}\lambda_I^\top \Sigma_I \lambda_I \\
%         \text{s.t.} & \text{ \textbf{U} and \textbf{B} hold} 
%     \end{align*}
%     where $\Sigma_I=(\Sigma_{ij})_{i,j\in I}$.
%   \end{minipage}
% }\label{fact}
% \end{center}

Here we show that the Multi-Allocate procedure reduces to a SOCP in the case of a single budget constraint, and to an SDP in the general case. The proof of \Cref{prop:og-primal} shows that the minimization problem over $\underline{n},\underline{\lambda}$ may be reduced to one only over $\underline{\lambda}$ via the Cauchy-Schwartz inequality. This minimization over $\underline{\lambda}$ is the dual of an SOCP, as shown by \Cref{prop:primal-dual}, and the KKT conditions hold. This is
\[
\sup_x a^\top x
\]
where the supremum is taken over all $x\in\R^k$ such that $x_I^\top\Sigma_Ix_I\leq c_I^{-1}$ for all $I\in\mathcal{I}$. This SOCP is simple to implement in the Python package cvxpy.

In the general case, \Cref{thm:minimax} shows that the optimal choice of $\underline{n}$ is
\[
\argmin_{\underline{n}\,:\,\mathbf{B}(\underline{n})} a^\top\left(\sum_{I\in\mathcal{I}} n_I P_I^\top \Sigma_I^{-1} P_I\right)^\dagger a
\]
Let us denote
\[
M(\underline{n}) = \sum_{I\in\mathcal{I}} n_IP_I^\top \Sigma_I^{-1}P_I
\]
so that our goal is to solve
\[
\min t
\]
subject to the constraints that
\[
a^\top M(\underline{n})^\dagger a \geq t
\]
and $\mathbf{B}(\underline{n})$, which denotes a set of linear constraints on $\underline{n}$. But this is equivalent to the SDP
\[
\min t
\]
subject to the constraint that
\[
\begin{pmatrix}
    M(\underline{n}) & a \\
    a^\top & t 
\end{pmatrix}\succeq0
\]\
and linear constraints on $\underline{n}$. Once again, this is straightforward to implement in cvxpy.
\section{The dual problem}\label{sec:dual}
We briefly recall the setup. Let $\Sigma\in\R^{k\times k}$ be SPD, let $\mathcal{I}$ denote a collection of index subsets $I\subseteq \{1,\dots,k\}$, and let $c_I$ be a positive scalar defined for every $I\in\mathcal{I}$. It will be convenient to define, for every $I\in\mathcal{I}$, a vector $\lambda_I\in\R^{|I|}$. We denote the concatenation of such vectors by $\underline{\lambda}\in\Lambda = \prod_{I\in\mathcal{I}} \R^{|I|}$. We further recall that $P_I:\R^k\longrightarrow\R^{|I|}$ is the orthogonal projection onto the coordinates indexed by $I$, and set $\Sigma_I=P_I\Sigma P_I^\top$. We define the norm $\|v\|_{\Sigma_I}=\sqrt{v^\top\Sigma_I v}$ on $\R^{|I|}$; this induces the seminorms $\|y\|_{\Sigma_I}=\|P_Iy\|_{\Sigma_I}$ on $\R^k$, and $\|\underline{\lambda}\|_{\Sigma_I}=\|\lambda_I\|_{\Sigma_I}$ on $\Lambda$. Lastly, we employ
\[
A:\Lambda\rightarrow\R^k,\quad A(\underline{\lambda}) = \sum_{I\in\mathcal{I}} P_I^\top \lambda_I
\]
to enforce the linear (unbiasedness) constraint $A(\underline{\lambda})=a$, for some fixed $a\neq0\in\R^k$.

Our first step will be to show how to alleviate the budget constraint. To do so, we first briefly recall this constraint. To describe the budget, recall that we define $\underline{n}=(n_I)_{I\in\mathcal{I}}\in \mathbb{Z}_{\geq0}^{|\mathcal{I}|}$, and employ a budget constraint of the form $\sum_{I\in\mathcal{I}} n_Ic_I \leq B$ for a fixed $B>0$. Denoting $\underline{c}=(c_I)_{I\in\mathcal{I}}\in\R_{>0}^{|\mathcal{I}|}$, our budget constraint may be written $\underline{c}^\top\underline{n}\leq B$. With all of this said, recall that our original problem of interest is
\begin{equation}\label{eqn:og}
    V(a) = \min_{\underline{n},\underline{\lambda}} \sum_{I\in\mathcal{I}:n_I>0} \frac{1}{n_I}\lambda_I^\top\Sigma_I\lambda_I\quad\quad\text{s.t.}\quad 
    \sum_{I\in\mathcal{I}:n_I>0}P_I^\top \lambda_I=a,\quad     \underline{c}^\top\underline{n}\leq B
\end{equation}
We begin by deriving tractable methods to solve \Cref{eqn:og}. Let us assume for the moment that $\underline{n}\in\R_{\geq0}^{|\mathcal{I}|}$; we will later construct the final budget allocation by rounding. Our first step is to remove the dependence on $\underline{n}$: we show that the above problem is equivalent to the following:
\begin{equation}\label{eqn:primal}
    U(a)=\min_{\underline{\lambda}\in\Lambda} \sum_{I\in\mathcal{I}} \sqrt{c_I}\|\lambda_I\|_{\Sigma_I}\quad\quad\text{s.t.}\quad A\underline{\lambda} = a
\end{equation}
We next show that this is equivalent to the dual problem
\begin{equation}\label{eqn:dual}
    U(a) = \sup_{y\in \R^k} a^\top y \quad\quad\text{s.t.}\quad \bigwedge_{I\in\mathcal{I}} \big\{\|y\|_{\Sigma_I}^2\leq c_I\big\}
\end{equation}
Finally, this is a second order cone program, and can be solved with off-the-shelf tools. After we have shown these things, we describe how to convert solutions to \Cref{eqn:dual} into solutions to \Cref{eqn:og}.

\begin{proposition}\label{prop:og-primal}
    The problems described in \Cref{eqn:og} and \Cref{eqn:primal} yield the same optimum $V=U^2/B$.
\end{proposition}
\begin{proposition}\label{prop:primal-dual}
    The problems described in \Cref{eqn:primal} and \Cref{eqn:dual} yield the same optimum $U$.
\end{proposition}
\begin{proof}[Proof of \cref{prop:og-primal}] We now begin the proof.
    \paragraph{$(2)\leq  (3)$:} Let $A\underline{\lambda}=a$. Define $\underline{n}$ by\footnote{This is defined as long as $\lambda_J\neq0$ for some $J$; if this fails then $\underline{\lambda}=0$ and $A\underline{\lambda}=0$ yields a contradiction.}
    \[
    n_I:=\left(\frac{B}{c_I}\right) \frac{\sqrt{c_I}\|\lambda_I\|_{\Sigma_I}}{\sum_{J\in\mathcal{I}} \sqrt{c_J} \|\lambda_J\|_{\Sigma_J}}
    \]
    It is clear that $\underline{c}^\top\underline{n}=B$ by construction, and we have
    \[
    BV(a) \leq \sum_{I:n_I>0} \frac{B}{n_I} \lambda_I^\top\Sigma_I\lambda_I = \sum_{I:\lambda_I\neq0} \sqrt{c_I}\|\lambda_I\|_{\Sigma_I} \sum_{J}\sqrt{c_J}\|\lambda_J\|_{\Sigma_J} = \left(\sum_{I\in\mathcal{I}}\sqrt{c_I}\|\lambda_I\|_{\Sigma_I}\right)^2
    \]
    \paragraph{$(3)\leq(2)$:} Let $\underline{n},\underline{\lambda}$ satisfy the constraints of \Cref{eqn:og}. Consider the vectors $\underline{c}^{1/2}\odot \underline{n}^{1/2}=(\sqrt{c_In_I})_{I\in\mathcal{I}}$ and $\left(\I_{n_I>0}n_I^{-1/2} \|\lambda_I\|_{\Sigma_I}\right)_{I\in\mathcal{I}}$ in $\R^{|\mathcal{I}|}$. The Cauchy-Schwartz inequality yields that the product of their squared norms is
    \[
    \left(\sum_I c_In_I\right)\left(\sum_{I:n_I>0} \frac{1}{n_I}\|\lambda_I\|_{\Sigma_I}^2\right) \geq \left(\sum_{I:n_I>0} \sqrt{c_I}\|\lambda_I\|_{\Sigma_I}\right)^2
    \]
    Now let us define $\underline{\tilde{\lambda}}$ by $\tilde{\lambda}_I=\lambda_I$ if $n_I>0$, and $\tilde{\lambda_I}=0$ otherwise. Then we have
    \[
    A\underline{\tilde{\lambda}} = \sum_I P_I^\top\tilde{\lambda_I} = \sum_{I:n_I>0} P_I^\top\lambda_I=a
    \]
    by assumption, and
    \[
    U(a)^2\leq \left(\sum_I \sqrt{c_I}\|\tilde{\lambda_I}\|_{\Sigma_I}\right)^2 = \left(\sum_{I:n_I>0} \sqrt{c_I} \|\lambda_I\|_{\Sigma_I}\right)^2 \leq BV(a)
    \]
    and we are done.
\end{proof}
\begin{remark}
    Note that in general, many $n_I$ will be zero.
\end{remark}

\begin{proof}[Proof of \cref{prop:primal-dual}] Let $\iota_{\{a\}}$ denote the indicator $b\mapsto \begin{cases}
    0 & a= b\\
    \infty & a\neq b
\end{cases}$.

Then \Cref{eqn:primal} is alternatively written
\[
V(a) = \min_{\underline{\lambda}\in\Lambda} g(\underline{\lambda}) + \iota_{\{a\}}(A\underline{\lambda})
\]
where $g(\underline{\lambda})=\sum_I g_I(\lambda_I)$ and $g_I(\lambda_I)=\sqrt{c_I}\|\lambda_I\|_{\Sigma_I}$. We now apply the Fenchel duality theorem. Note that $\iota_{\{a\}}^*(y)=a^\top y$, and $g^*(A^\top y) = \sum_I g_I^*(P_I^\top y) = \sum_I \iota_{\{\|y_I\|_{\Sigma_I}\leq c_I\}} = \iota_{\bigwedge_I \|y_I\|_{\Sigma_I}^2\leq c_I}$.
\end{proof}
% \subsection{Reconstructing $\underline{n}^\star,\underline{\lambda}^\star$ from the solution $y^\star$ to the dual problem}
% By KKT, we have the multipliers $\alpha_I^*\geq0$ such that
% \[
% \lambda_I^\star = 2\alpha_I^\star \Sigma_I^{-1} y_I^\star\quad\quad\forall I\in\mathcal{I}
% \]
% Finally, as in the proof of \cref{prop:og-primal}, we set
% \[
% n_I^\star = \left(\frac{B}{c_U}\right) \frac{\sqrt{c_I}\|\lambda_I^\star\|_{\Sigma_I}}{\sum_J \sqrt{c_J} \|\lambda^\star_J\|_{\Sigma_J}}
% \]
% and this is optimal. Round down if needed.

% \subsection{Attempt 2, re: Alekh's suggestion}
% Recall that we define $g_I(\lambda_I):=\sqrt{c_I}\|\lambda_I\|_{\Sigma_I}$. Since $\Sigma_I\succ0$ for each $I$, $g_I$ is a norm on $\mathbb{R}^I$, and $g(\underline{\lambda})=\sum_I g_I(\lambda_I)$ is a norm on $\Lambda=\prod_I \mathbb{R}^I$. We form the Lagrangian
% \[
% \mathcal{L}(\underline{\lambda},y) = \sum_{I\in\mathcal{I}} g_I(\lambda_I) + y^\top(a-A\underline{\lambda}) = a^\top y+\sum_{I\in\mathcal{I}} (g_I(\lambda_I)-y_I^\top\lambda_I)
% \]
% Minimizing over $\underline{\lambda}\in\Lambda$, we find
% \[
% \inf_{\underline{\lambda}\in\Lambda} \mathcal{L}(\underline{\lambda},y) = a^\top y + \sum_{I\in\mathcal{I}} \inf_{\lambda_I}(g_I(\lambda_I)-y_I^\top\lambda_I) = a^\top y + \sum_{I\in\mathcal{I}} g_I^*(y_I)
% \]
% where $g_I^*$ is the convex conjugate of $g_I$. Because $g_I$ is a norm, its convex conjugate is the indicator of the dual norm ball
% \[
% g_I^*(z) = \begin{cases}
%     0 & z^\top\Sigma_I^{-1}z\leq c_I \\
%     \infty & \text{else}
% \end{cases}
% \]
% We therefore have the dual problem
% \[
% \max_y a^\top y \quad\text{s.t.}\quad y_I^\top\Sigma_I^{-1}y_I\leq c_I \quad\forall I\in\mathcal{I}
% \]
\section{Experimental details}\label{sec:experimental-details}

Here we detail the experimental setup used. We do so in two parts: first, we explain the details for generating the model predictions $(X_2,\dots,X_k)$ in each experiment; second, we explain the details for constructing the proposed estimator, $\hat{\theta}_{\multiPPI}$, and the baselines from such predictions.

\subsection{Generating model predictions}

\subsubsection{Experiment 1: Chatbot Arena}
We follow the implementation of \citet{angelopoulos2025costoptimalactiveaimodel} to request autoratings from Gemini 2.5 Pro and Gemini 2.5 Flash. See section E of \citet{angelopoulos2025costoptimalactiveaimodel} for implementation details.

\begin{table}[]
    \centering
    \begin{tabular}{c|c}
        Model collection & Cost \\
        \hline
        Gemini 2.5 Pro & \$1.25 \\
        Gemini 2.5 Flash & \$0.30 \\
        Both & \$1.55
    \end{tabular}
    \caption{Cost structure for experiment 1.}
    \label{tab:cost-exp1}
\end{table}

\subsubsection{Experiment 2: Processbench}
We evaluate our method on 500 samples from the OlympiadBench subset of the ProcessBench dataset \citep{zheng2024processbench}. Binary labels are determined according to whether or not a process error occurred in the given (problem, attempted solution) pair.

To generate autoratings, we use Gemini 2.5 Pro and truncate its reasoning process at various checkpoints. Specifically, using the prompt shown in \autoref{fig:processbench-prompt}, we instruct the model to think for up to 3,000 tokens but interrupt it and demand an answer after $B$ words of thought have been produced, for $B\in\{125,250,375,500\}$, as described in \S\ref{sec:exp-setup}. To elicit a definite judgement at each checkpoint, we provide ``So, the answer is:'' as the assistant and attempt to extract an answer from the subsequent 20 tokens of output with our template.

\begin{figure}
    \centering
    \includegraphics[width=\linewidth]{figures/Processbench_prompt.pdf}
    \caption{Prompt used to generate autoratings for Experiment 2.}
    \label{fig:processbench-prompt}
\end{figure}

\subsubsection{Experiment 3: Biography factuality}

We consider evaluating the factuality of a set of biographies generated by Gemini 2.5 Pro. We replicate the setting of \citet{du2023improving}: Gemini 2.5 Pro is asked to generate biographies for 524 computer scientists, and we evaluate the factual consistency of such biographies with a set of grounding facts collected by \citet{du2023improving}.

More specifically, for every person $p\in\mathcal{P}$, we associate a Gemini-generated biography $b^p$ and a set of collected grounding facts $\mathcal{F}^p=\{f_1^p,\dots,f_{m_p}^p\}$ about the person. Following \citet{du2023improving}, we estimate the proportion of \textit{factually consistent pairs} $(b^p,f^p_i)$ of generated biographies $b^p$ with each of the collected grounding facts $f^p_i$. Concretely, given the set of all pairs
\[
\mathcal{S}=\{(b^p,f^p): p\in\mathcal{P},f^p\in\mathcal{F}^p\}
\]
we \textit{target} the proportion of factually consistent pairs
\[
\frac{\#\{(b,f)\in\mathcal{S}:(b,f)\text{ is factually consistent}\}}{\# \mathcal{S}}
\]
We determine the factual consistency, or lack thereof, of a pair $(b,f)$ by majority voting over 5 independent judgments from Gemini 2.5 Pro with thinking. \citet{du2023improving} found that judgments by ChatGPT achieved over 95\% agreement with human labelers on a set of 100 samples. This level of agreement is evidently not achieved by certain cheaper models, as we proceed to demonstrate experimentally. In \autoref{fig:scaling-accuracy-biographies}, we explore using Gemini 2.0 Flash Lite as an autorater for evaluating the factuality consistency of pairs $(b,f)\in\mathcal{S}$.

To elicit better autoratings from queries to Gemini 2.0 Flash Lite, we bootstrap performance via multi-round debate. For a fixed number of agents $A\in\{1,\dots,5\}$, and a fixed number of maximum rounds $R\in\{1,2\}$, we perform the following procedure:
\begin{outline}[enumerate]
    \1 $A$ instances of Flash Lite are independently prompted to consider the factual consistency of pairs $(b,f)\in\mathcal{S}$, and provide an explanation for their reasoning.
    \1 A ``pooler'' instance of Flash Lite is then asked to review the pair $(b,f)$ and the responses generated by each of the $A$ other instances, and output a judgment in the form of a single word: yes, no, or uncertain.
    \2 If the pooler outputs ``yes'' or ``no,'' the judgment is final.
    \2 If the pooler outputs ``uncertain'' and the number of maximum rounds $R$ has not yet been reached, the $A$ instances of Flash Lite are independently shown their prior responses, and the prior responses of each other, and prompted to continue reasoning given this additional information. This procedure continues until either the pooler no longer reports ``uncertain,'' or the maximum number of rounds $R$ has been reached.
    \2 If the pooler outputs ``uncertain'' and the maximum number of rounds $R$ has been reached, a fair coin is flipped and ``yes'' or ``no'' are reported with equal probability.
\end{outline}
Since the dataset is balanced, the outcome described in (c) is fair insofar as it is as good as random guessing. We impose the maximum round restriction to encapsulate our budget constraint. To reduce randomness, we generate all autoratings twice, so that the resulting dataset has an effective size of 1048.

\begin{description}
    \item[Target:] Proportion of factually-consistent pairs, $\#\{(b,f)\in\mathcal{S}:(b,f)\text{ is factually consistent}\}/\#\mathcal{S}$
    \item[Model family:] $\{\text{The output of the above procedure given $(A,R)$}: A\in\{1,\dots,5\},R\in\{1,2\}\}$
    \item[Cost structure:] For a given $(A,R)$, the cost is $
A\cdot R$. For collections of models, the cost is additive.%\footnote{This we can obviously change. Just a simplification for draft purposes.}
\end{description}

\begin{figure}
    \centering
    \includegraphics[width=\linewidth]{figures/Factuality_task2.png}
    \caption{Depiction of biography-fact pairs $(b,f)$ as in Experiment 3. Judgements about factual consistency of $(b,f)$ are made by a language model.}
    \label{fig:biography-prompt}
\end{figure}

\subsection{Constructing the MultiPPI estimator}
For the results shown in \S\ref{sec:results}, we draw 250 fully-labeled samples from each dataset above. We then follow the procedure described in \S\ref{sec:algorithm_unknown} for $N=250$, using the empirical distribution over each dataset as our source of randomness. In \autoref{sec:experiments-num-labeled}, we replicate the study over a range of values of $N$.


\end{document}
